\hypertarget{ux3054ux3042ux3044ux3055ux3064ux304bux3064ux3066ux306eux79c1ux3078ux306eux8d08ux308aux7269ux3042ux308bux3044ux306fsnsux6642ux4ee3ux306eux4e88ux9632ux7dda}{%
\chapter{ごあいさつ：かつての私への贈り物、あるいはSNS時代の予防線}\label{ux3054ux3042ux3044ux3055ux3064ux304bux3064ux3066ux306eux79c1ux3078ux306eux8d08ux308aux7269ux3042ux308bux3044ux306fsnsux6642ux4ee3ux306eux4e88ux9632ux7dda}}

この本は、大学生時代の筆者自身の疑問を解決するための本として、かつての筆者の前提知識を想定して書きました。

同時に、インターネット上に数学めいた文章を公開する以上、いくつかの予防線も必要でした。本書には、より進んだ立場の教科書なら最初に抽象化してしまうところを、あえて行列表示や有限セルの計算から始める箇所があります。

これはより進んだ立場を否定するためではありません。むしろ、いずれこの先の世界に立ち向かう際に恐れることなく頼りにできる具体的な道具、直観、そして「厳密さ」への嗅覚を養うことを目的としています。

詳しい問題設定は、次の「はじめに」で述べます。

さて、本題です。

みなさんは、\(dx\) とは何かを説明できるでしょうか。

\begin{quote}
\textbf{注}（注釈について）
本書には注釈が多くあります。読者によっては、本文の流れを止める不親切なものに見えるかもしれません。

しかし、本書では注釈に三つの役割を持たせています。

第一に、筆者自身のための予防線です。どこで省略し、どこで同一視し、この先にどのような厳密化があるのかを、本文の脇に少しだけ残しておくためです。

第二に、読者の嗅覚を養うためです。いまは完全には分からなくても、「ここには何かあるらしい」「この言い方には限界があるらしい」と感じておくことは、先へ進むときの助けになります。

第三に、ミステリ小説の回想のような効果です。最初は断片的にしか見えなかった注釈が、読み進めるにつれて少しずつ意味を持ち、あとから「あれはこのためだったのか」と分かる。そういう読み味も、少しだけ狙っています。

注釈は、すべてをその場で理解するためのものではありません。読めるところだけ拾って、分からないところは通り過ぎてもらって構いません。
\end{quote}

\begin{quote}
\textbf{注}（歴史順と分かる順）
現代の読者は、行列、線型代数、微分形式、ベクトル解析を、すでに整備された道具として学ぶことができます。少なくとも筆者の受けた教育では、行列はまず、高校生時代に手で計算できる代数的な道具として現れました。しかし、実際の歴史では、これらの道具は互いに影響し合いながら、行きつ戻りつ発展してきたように見えます。そのため、教育現場にも歴史的経緯による混乱が残っているように筆者には思えます。本書は物理数学の歴史を再現する本ではありません。むしろ、現代の読者がすでに持っている行列代数の言葉を使って、多変数微積分とベクトル解析を別の順序でほどき直す試みという側面があります。
\end{quote}

\begin{quote}
\textbf{注}（本書には怒りがある）
それは、ある特定の人物への怒りではなく、見えているはずの問題が、私の学生時代にも、そして現代に至るまで、十分には解決されていないことへの怒りです。ベクトル解析は、計量・向き・双対性・外微分を暗黙に含んでいるにもかかわらず、それらを隠したまま公式として教えられる。一方で、微分形式は多様体論の高みに置かれ、初学者がベクトル解析を理解するための道具としては降ろされてこなかった。本書は、その断絶に対する攻撃であり、同時に対案でもあります。
\end{quote}

\newpage

\hypertarget{ux306fux3058ux3081ux306bdxux3068ux306fux4f55ux304bux30caux30d6ux30e9ux3068ux306fux4f55ux304b}{%
\chapter{\texorpdfstring{はじめに：\(dx\)とは何か・ナブラとは何か}{はじめに：dxとは何か・ナブラとは何か}}\label{ux306fux3058ux3081ux306bdxux3068ux306fux4f55ux304bux30caux30d6ux30e9ux3068ux306fux4f55ux304b}}

\hypertarget{ux524dux63d0ux77e5ux8b58}{%
\section{前提知識}\label{ux524dux63d0ux77e5ux8b58}}

本書は、以下の基本事項に馴染みがあれば、十分に読み進めることができます。

\begin{itemize}
\item
  \textbf{微積分}：一変数関数の微積分ができること。置換積分に抵抗がないと、なお望ましいです。
\item
  \textbf{線形代数}：行列の積の計算に抵抗がないこと。また、「矢印としてのベクトル」と「数ベクトルとしてのベクトル」を行き来するイメージを持っていること。
\end{itemize}

行列式を既に知っている必要はありません。途中で一緒に発見しましょう。

高校生でも読みこなせる、と言うと頑張ってもらう必要がありますが、少なくとも20歳の頃の筆者であれば通読できる、つもりで書きました。

\hypertarget{dxux3068ux306fux4f55ux304b}{%
\section{\texorpdfstring{\(dx\)とは何か}{dxとは何か}}\label{dxux3068ux306fux4f55ux304b}}

積分記号の最後には、いつも \(dx\)
が付いています。高校数学では、これは「\(x\)
で積分する」という印のように扱われます。 また、 \(\frac{dy}{dx}\)
は割り算ではないと習うものの、あたかも割り算のように置換積分を行うことが、事実上許容されています。

大学一年生の教養の微積分では、全微分 \(dx\)
として登場しますが、実際のところこれが何なのかは深入りしないことが多いように思います。

さらに、特に物理学や工学系の専門課程では、当然のように微小量 \(dx\) や
\(\delta x\) という実体が操作されます。

筆者の個人的な経験では、大学院一年生（M1）の頃に、数学系の友人に聞いたところ、「\(dx\)
は接ベクトル空間に対する線形汎関数、あるいは余接空間の元」と、分かるような分からないような説明をされた覚えがあります。

さて、結局 \(dx\)
とは何でしょうか。進んだ専門課程でないと、深入りしないか、微小量として理解するしかないのでしょうか。

もちろん、数学的により進んだ説明をすれば、先の大学院時代の友人の説明が圧倒的に正しく、そのように理解できたほうがよいです。しかし、難しい。

そこで本書は、本当に高校生にその理解へたどり着いてもらうにはどうすればいいかを考えてみました。そのために、少なくとも筆者が高校時代にすでに訓練されていた行列代数を足がかりとして使います。

すなわち、本書ではまず、\(dx\) を「変位ベクトルから \(x\)
成分を取り出す測定器」として読みます。行列表現で、象徴的には、

\[
dx=\begin{pmatrix}1&0&0\end{pmatrix}
\]

です。

知りたいのは、積分記号の尻尾にある \(dx\)
が、実際に何をしているのかです。本書では、その働きをまず見える形にします。

\(dx\) は変位を食べて \(x\) 成分を返す。\(dx\wedge dy\)
は、二本の変位が張る向き付き面積を測る。\(dx\wedge dy\wedge dz\)
は、三本の変位が張る向き付き体積を測る。

高校数学の \(dx\) と微分形式の \(dx\)
のあいだに、いきなり抽象語を置くのではなく、まず行列と測定器を置いてみる。これが本書の出発点です。

\hypertarget{ux30caux30d6ux30e9ux3068ux306fux4f55ux304b}{%
\section{ナブラとは何か}\label{ux30caux30d6ux30e9ux3068ux306fux4f55ux304b}}

もう一つ、本書が取り上げたい記号があります。\(\nabla\) です。

ベクトル解析を学んでいると、あるところから急にベクトル解析の体系が一連の公式集のように見えてくることがあります。少なくとも筆者の経験ではそうでした。デカルト座標のうちはよくても、円柱座標や球座標に入ると、\(r\)
や \(\sin\theta\)
が現れ、どこから来たのか分からない係数を覚えることになります。

たとえば、二次元でいえば

\[
\frac{\partial F_y}{\partial x}-\frac{\partial F_x}{\partial y}
\qquad=\qquad
\frac{1}{r}\frac{\partial (r F_\theta)}{\partial r}-\frac{1}{r}\frac{\partial F_r}{\partial \theta}
\]

が現れます。左辺がデカルト座標、右辺が同じものを極座標で書いたものです。これのどこが「回転」なのか、明日の中間テストを控えた状態で納得するのはなかなか難しいでしょう。

もちろん、この点については物理数学の副読本や、近年の丁寧な書籍でも解説されています。しかし、教科書の体系のなかで、これを発見的に組み立てているものは、筆者の知る限り、和書ではまだ珍しいように思います。

さて、本書では、いきなりナブラから始めません。せっかく \(dx\)
の行列表現を考えるのですから、これを使います。

その後 \(dx\)
から、面積とは何か、体積とは何かを考え直します。それから、積分、変数変換、外微分、ホッジスターへ進みます。そのあとで、ようやくナブラに戻ってきます。

明日の中間テストには間に合いませんが、期末テストなら間に合うのではないでしょうか。

\newpage

\hypertarget{ux30ddux30fcux30bfux30ebux30b5ux30a4ux30c8ux3068ux3061ux3087ux3063ux3068ux3057ux305fux8a66ux307f}{%
\chapter{ポータルサイトとちょっとした試み}\label{ux30ddux30fcux30bfux30ebux30b5ux30a4ux30c8ux3068ux3061ux3087ux3063ux3068ux3057ux305fux8a66ux307f}}

本書の PDF 版、正誤表、改訂履歴、GitHub
リポジトリ、および周辺メモへの導線は、以下のポータルサイトにまとめています。

\hypertarget{ux30ddux30fcux30bfux30ebux30b5ux30a4ux30c8}{%
\subsection{ポータルサイト}\label{ux30ddux30fcux30bfux30ebux30b5ux30a4ux30c8}}

\url{https://covectorspace.xyz/jp/}

本書は無料で公開しています。明らかな誤植レベルの指摘は GitHub Issues
でも受け付けます。

\begin{quote}
\textbf{注}（名称について）\\
ちなみに、\emph{covector space}、あるいは \emph{Co-Vector Space}
という語はあまり一般的ではありません。数学では \emph{dual space}
などと言う方が普通です。
\end{quote}

\hypertarget{discord-ux30b5ux30fcux30d0ux30fc}{%
\subsection{Discord サーバー}\label{discord-ux30b5ux30fcux30d0ux30fc}}

公開 SNS
に投げにくい未整理のアイデアや初歩的な質問を、少しだけクローズドに持ち寄れる場を目指して、試験的に
Discord サーバーを開けてみました。

これから整備していく想定で、現状は本当に「空き部屋」に近い状態です。理工系に関心のある人たちが、本書の内容に限らず、ほどよい距離感で場所を目指しています。

興味のある方は、ポータルサイトから覗いてみてください。
招待リンクはこちらです。

\url{https://discord.gg/NCffyR9gj}

\newpage

\hypertarget{ux7b2c1ux7ae0dx-ux3068ux306fux4f55ux304b-ux30d9ux30afux30c8ux30ebux3092ux98dfux3079ux308bux6e2cux5b9aux5668ux3042ux308bux3044ux306fux6a2aux30d9ux30afux30c8ux30eb}{%
\chapter{\texorpdfstring{第1章：\(dx\) とは何か ------
ベクトルを食べる測定器、あるいは横ベクトル}{第1章：dx とは何か ------ ベクトルを食べる測定器、あるいは横ベクトル}}\label{ux7b2c1ux7ae0dx-ux3068ux306fux4f55ux304b-ux30d9ux30afux30c8ux30ebux3092ux98dfux3079ux308bux6e2cux5b9aux5668ux3042ux308bux3044ux306fux6a2aux30d9ux30afux30c8ux30eb}}

\hypertarget{ux6570ux5b66ux8005ux306e1ux6b21ux5143ux7269ux7406ux5b66ux8005ux306e1ux6b21ux5143}{%
\subsection{§1.0
数学者の1次元、物理学者の1次元}\label{ux6570ux5b66ux8005ux306e1ux6b21ux5143ux7269ux7406ux5b66ux8005ux306e1ux6b21ux5143}}

高校で \(\int f(x) dx\)
という1次元の積分を習ったとき、教科書の記述では「\(x\)
軸という1本の線」が置かれる。数学の教科書としてはそれで妥当だ。\textbf{数学者}は1次元空間
\(\mathbb{R}^1\)
という自己完結した抽象世界を仮定し、その上で論理を積み上げる。点は実数
\(x\)、変位は実数
\(\Delta x\)、積分は「関数×微小幅」の極限として定義される。すべてが完結している。

\textbf{しかし、我々は物理学者である。}

\begin{quote}
\textbf{注}（筆者は物理学者ではない）\\
もちろんこれはレトリック（修辞法）である。筆者の学位は物理学ではなく化学であり、現在アカデミアの中にいるわけでもない。

ここで言いたいのは、本書の見方が、情報系の読者にも、機械系や電子系の読者にも、その他の理工系の読者にも、あるいは単にベクトル解析や微積分の記法で迷子になった読者にも、役に立つ可能性があるということだ。

「物理学者」は所属ではなく、物理数学的なものの見方のことである。数学者や数学書を下げる意図もない。

それでも少し大げさに書いたのは、単に言ってみたかったからである。
\end{quote}

我々が扱う\textbf{現実の物理空間は、常に3次元だ。}質点が直線上を運動しているように見えても、それは「架空の1次元空間」にいるのではなく、「3次元空間
\(\mathbb{R}^3\) の中で、たまたま \(y\) 方向と \(z\)
方向への変位が観測されない（あるいは無視できる）断面」を見ているに過ぎない。

例えば、摩擦のない直線レール上を滑る台車を考えよう。我々は「これは \(x\)
軸方向の1次元運動だ」と言うが、実際には：

\begin{itemize}
\item
  レールは3次元空間内に設置されている。
\item
  台車は微小に上下に揺れているかもしれない。
\item
  空気抵抗による横方向の微小変動もある。
\end{itemize}

物理学者の「1次元運動」とは、\textbf{3次元空間内で特定の方向への変位が支配的であり、他の方向の変位が無視できるほど小さいか、系の対称性によって厳密にゼロに拘束されている状況}を指す。

\begin{quote}
\textbf{注}（埋め込みという見方）\\
ここでいう「物理学者の1次元」とは、1次元のパラメータ空間を、3次元空間の中の曲線として写し込んで見る、という意味である。数学では、このような写像をまず曲線の\textbf{パラメータ表示}として扱い、自己交差がなく、速度が消えないなどの条件がよい場合には、\textbf{埋め込み}と呼ぶ。

本書は、抽象的な1次元空間そのものを否定しているわけではない。ここではむしろ、物理数学でしばしば現れる「空間内の曲線に沿った運動」を、まずは3次元空間への写し込みとして捉えている。
\end{quote}

したがって、物理学者が「1次元の微小変位」と呼ぶものの真の姿は、単なるスカラー
\(\Delta x\)
ではなく、次のような\textbf{縦ベクトル（列ベクトル）}で表すのがふさわしい：

\[
\mathbf{v} :=
\begin{pmatrix}
\Delta x \\
0 \\
0
\end{pmatrix}
\]

\begin{quote}
\textbf{注}（標準的でない呼び方）\\
本書全体を通して\textbf{縦ベクトル・横ベクトル}という言い方をする（縦が列、横が行、と対応させて読んでほしい）。筆者は行と列という語の区別が苦手だからだ。
\end{quote}

我々はこのように、3次元空間内の一点からの変位を3成分で表したものを\textbf{変位ベクトル}と呼ぶ。成分の並びは慣習に従い、\textbf{第1成分が}\(x\)\textbf{軸方向、第2成分が}\(y\)\textbf{軸方向、第3成分が}\(z\)\textbf{軸方向}の変位に対応するものとする。

この第2成分・第3成分のゼロは、「無視されている」あるいは「存在しない」という消極的な意味ではない。\textbf{我々が今、}\(x\)\textbf{軸方向の運動だけに注目し、他の成分を測定対象から意図的に除外している}という積極的な選択の結果なのである。

この「物理学者としての3次元」という見方が、\(dx\)
が積分記号の末尾につくおまけではない、という視点を与えてくれる。次節では、この視点から
\(dx\)\textbf{を「微小量」から「行列」あるいは「演算子」へと読み直していく。}

\begin{quote}
\textbf{注}（本書の立場）\\
本書はあくまで、初等微分積分学とベクトル解析をほどき直すための本である。一般次元の微分形式を体系的に展開する本ではない。

したがって、最後まで基本的には3次元デカルト座標 \((x,y,z)\)
に固定し、線形代数学として最も素直な表現------行列表現------を採用する。本書で用いる
\(dx,dy,dz\)
は、この座標系において成分を取り出す具体的な行列として定義する。

まずは「3次元・デカルト・行列」という素直な道具立ての中で、具体的なシンボルとその操作として、また微小量の誤魔化しなくベクトル解析を身体に染み込ませることが狙いである。
\end{quote}

\begin{quote}
\textbf{注}（デカルト座標の呼称について）\\
本書の「デカルト座標」とは、\textbf{正規直交座標系}（原点が一致し、各軸が互いに直交し、間隔が等間隔である座標系）を指す。数学や他書籍では「ユークリッド空間上の直交直線座標」「デカルト座標系」などとも呼ばれるが、本書では単に「デカルト座標」と呼ぶ。
\end{quote}

\begin{quote}
\textbf{注}（第III部での拡張）\\
ただし第III部では、必要に応じて曲線座標や高次元への拡張にも触れる。
\end{quote}

\begin{center}\rule{0.5\linewidth}{0.5pt}\end{center}

\hypertarget{ux30eaux30fcux30deux30f3ux548cux3068ux884cux5217ux306eux7a4d}{%
\subsection{§1.1
リーマン和と行列の積}\label{ux30eaux30fcux30deux30f3ux548cux3068ux884cux5217ux306eux7a4d}}

前節で、物理的な微小変位を縦ベクトル \(\mathbf{v}\)
として捉え直した。§1.0 と同じく
\[\mathbf{v} = \begin{pmatrix} \Delta x \\ 0 \\ 0 \end{pmatrix}\]
である。この視点を持って、今度は高校以来慣れ親しんだリーマン積分のプロセスを書き直し、「行列としての
\(dx\)」がどのような見方を与えるかを見てみよう。

\begin{quote}
\textbf{注}（転置の記法について）\\
数学や他文献では、紙面の都合で縦ベクトルを横に寝かせ、右上に \({}^T\)
を添えて表すことがある。たとえば \((\Delta x,\Delta y,\Delta z)^T\)
のような書き方である。本書では、この紙面都合の転置記号をできるだけ避ける。縦ベクトルは縦に、横ベクトルは横に書く。もし例外的に
\({}^T\)
が出てきたときは、何かしらの意図がある転置操作だと読んでほしい。
\end{quote}

\hypertarget{ux30eaux30fcux30deux30f3ux548cux306eux6a19ux6e96ux7684ux69cbux6210ux5fa9ux7fd2}{%
\subsubsection{1.1.1
リーマン和の標準的構成（復習）}\label{ux30eaux30fcux30deux30f3ux548cux306eux6a19ux6e96ux7684ux69cbux6210ux5fa9ux7fd2}}

関数 \(f(x)\) の区間 \([a, b]\)
における定積分は、次のように定義される（高校数学で言うと区分求積法だ）。

\begin{quote}
\textbf{注} （閉区間）記号 \([a,b]\) は、端点 \(a\) と \(b\)
を両方含む\textbf{閉区間}（\(a \le x \le b\) の実数 \(x\)
の集合）を表す。
\end{quote}

\begin{enumerate}
\def\labelenumi{\arabic{enumi}.}
\tightlist
\item
  区間を \(n\) 個の小区間に分割： \(a = x_0 < x_1 < \cdots < x_n = b\)
\item
  小区間の幅を \(\Delta x_i = x_i - x_{i-1}\)
\item
  各小区間 \([x_{i-1}, x_i]\) の\textbf{中から}代表点 \(\xi_i\)
  を1つ選ぶ
\item
  リーマン和 \(R_n := \sum_{i=1}^n f(\xi_i) \Delta x_i\) を作る（この
  \(n\) 分割に対する和を \(R_n\) と書く）
\item
  分割を細かくする極限をとる（厳密には、小区間の幅の\textbf{最大値}が
  \(0\) に近づく。分割数 \(n\)
  だけを大きくしても幅が残る取り方がありうるので、より進んだ数学ではこちらを重視する。以下の
  \(\lim_{n\to\infty} R_n\)
  は、この条件を満たす分割の取り方を前提にしている）：
  \(\int_a^b f(x)\,dx = \lim_{n\to\infty} R_n\)
\end{enumerate}

この教科書的な説明では、「微小な幅
\(\Delta x_i\)」がそのまま小さくなったものが\(dx\)と読める。しかし、我々はもう一歩踏み込めるはずだ。

\hypertarget{ux5c0fux533aux9593ux3054ux3068ux306eux5909ux4f4d}{%
\subsubsection{1.1.2
小区間ごとの変位}\label{ux5c0fux533aux9593ux3054ux3068ux306eux5909ux4f4d}}

第 \(i\) 小区間 \([x_{i-1}, x_i]\) における変位は、§1.0
の\textbf{変位ベクトル}を、その区間の幅に合わせて
\[\mathbf{v}_i = \begin{pmatrix} \Delta x_i \\ 0 \\ 0 \end{pmatrix}\]
と書く。添字 \(i\) は「第 \(i\)
小区間の分」という意味だけで、\(\mathbf{v}_i\)
もまた変位ベクトルである。今は直線運動を考えているので \(y,z\)
成分はゼロだが、それは条件ではなく結果としてゼロになっているだけだ。

\clearpage

\hypertarget{dx-ux306eux767bux5834-ux672cux66f8ux6700ux5927ux306eux7279ux5fb4ux3068ux3057ux3066ux306eux65adux8a00}{%
\subsubsection{\texorpdfstring{1.1.3 \(dx\) の登場 ---
本書最大の特徴としての「断言」}{1.1.3 dx の登場 --- 本書最大の特徴としての「断言」}}\label{dx-ux306eux767bux5834-ux672cux66f8ux6700ux5927ux306eux7279ux5fb4ux3068ux3057ux3066ux306eux65adux8a00}}

さて、本書では \(dx\)
という記号に次のような意味を与える。\textbf{大胆で奇妙に思えるかもしれないが、これが本書最大の特徴でもある}------ここでは
\(dx\) を次の \(1\times 3\) 行列だと\textbf{断言する}。

すなわち、\textbf{横ベクトル}（\(1\times 3\)
\textbf{行列}として成分を横に並べて書く）として

\begin{center}\scalebox{1.8}{$\displaystyle 
dx := \begin{pmatrix} 1 & 0 & 0 \end{pmatrix}
$}\end{center}

と\textbf{定義する}。

そして、実はこれは標準的な考え方から外れた独自記法ではない。標準的な線形代数・テンソル解析・多様体論で現れる対象を、\(\mathbb R^3\)
の標準座標に固定して、最初から行列として書き下したものである。詳しくは、すぐ下の注釈で示しておく。

より進んだ数学に詳しい読者向けに釈明しておくと、これは入力された縦ベクトルから
\(x\)
成分だけを抜き出して返す線形演算子の行列表現であり、\(\mathbb R^3\)
の標準座標における \(dx\) の表示である。

\begin{quote}
\textbf{注}（\(dx = (1\ 0\ 0)\) と教科書との対応）\\
この表記は通常の教科書ではあまり前面に出てこないが、実は暗に含まれている。

\textbf{ベクトル解析を学んだ読者へ}

スカラー値関数 \(f:\mathbb R^3\to\mathbb R\) の勾配は
\(\nabla f = \begin{pmatrix}\frac{\partial f}{\partial x} \\ \frac{\partial f}{\partial y} \\ \frac{\partial f}{\partial z}\end{pmatrix}\)
である。この転置
\((\nabla f)^T = (\frac{\partial f}{\partial x}\ \frac{\partial f}{\partial y}\ \frac{\partial f}{\partial z})\)
は \(1\times3\) の横ベクトル（行ベクトル）であり、\(f=x\) と置けば
\((\nabla x)^T = (1\ 0\ 0)\)
となる。デカルト座標においては、実はこの勾配の転置こそが \(dx\) である。

\textbf{テンソル解析を学んだ読者へ}

\(dx^i\) は座標基底 \(\frac{\partial}{\partial x^j}\)
の双対基底なので、\(dx^i\left(\frac{\partial}{\partial x^j}\right)=\delta^i_j\)
である。したがって、デカルト座標で横ベクトル（行ベクトル）表示すれば
\(dx^1=(1\ 0\ 0)\), \(dx^2=(0\ 1\ 0)\), \(dx^3=(0\ 0\ 1)\) である。

\textbf{多様体論に詳しい読者へ}

\(df_p:T_p\mathbb R^3\to\mathbb R\), \(df_p(v)=v(f)\),
\(dx^i_p(v)=v(x^i)\) である。標準座標で
\(v=v^i\frac{\partial}{\partial x^i} \longmapsto \begin{pmatrix}v^1\\v^2\\v^3\end{pmatrix}\)
と表せば \(dx^1_p(v)=v^1\), \(dx^2_p(v)=v^2\), \(dx^3_p(v)=v^3\)
となる。すなわち、標準座標表示では \(dx^1_p\longmapsto(1\ 0\ 0)\),
\(dx^2_p\longmapsto(0\ 1\ 0)\), \(dx^3_p\longmapsto(0\ 0\ 1)\) である。
\end{quote}

変位ベクトル \(\mathbf{v}_i\) に左から掛けると

\[
dx\,\mathbf{v}_i
=
\begin{pmatrix} 1 & 0 & 0 \end{pmatrix}
\begin{pmatrix} \Delta x_i \\ 0 \\ 0 \end{pmatrix}
=
\Delta x_i
\]

となる。ここで「\(dx\) が変位ベクトル \(\mathbf{v}_i\) を食べてスカラー
\(\Delta x_i\)
を吐く」というイメージを、\textbf{関数の値のように書く}ことにする：

\[
dx(\mathbf{v}_i)
:=
dx\,\mathbf{v}_i
=
\Delta x_i
\]

本書では、しばしばこの \(dx(\mathbf{v})\) の形で「横ベクトル \(dx\)
が縦ベクトル \(\mathbf{v}\)
に作用する」と強調する。何度繰り返しても強調しすぎることはない。

\begin{quote}
\textbf{注}（演算子・作用素・関数・写像の呼び分け）\\
この \(dx(\mathbf{v})\)
を\textbf{演算子}と呼ぶか\textbf{作用素}と呼ぶか\textbf{関数}と呼ぶか\textbf{写像}と呼ぶかは、文献によってさまざまな流儀がある。筆者は\textbf{演算子}と呼ぶことを好むので、この先も演算子と呼ぶことが多いだろう。
\end{quote}

リーマン和に現れる \(\Delta x_i\)
は、「1次元の謎の微小量」ではない。\textbf{3次元空間の中での変位ベクトル}\(\mathbf{v}_i\)\textbf{に対し、上の行列としての}\(dx\)\textbf{が}\(x\)\textbf{成分だけを抜き出した結果}である。

つまり

\[
\Delta x_i = dx(\mathbf{v}_i)
\]

である。

\textbf{これが決定的な瞬間だ}：

\begin{itemize}
\item
  左辺：\(dx(\mathbf{v}_i)\) は、行列 \(dx\)
  が\textbf{変位ベクトル}\(\mathbf{v}_i\)\textbf{に}作用する代数的操作
\item
  右辺：従来のリーマン和に現れる「小区間の幅 \(\Delta x_i\)」
\end{itemize}

つまり \(\Delta x_i\)
は、\textbf{1次元の謎の微小変位ではなく、3次元の微小変位}\(\mathbf{v}_i\)\textbf{から}\(x\)\textbf{方向の成分だけを取り出した}\(dx(\mathbf{v}_i)\)\textbf{である}。

\begin{quote}
\textbf{注}（成分並びの約束）\\
横ベクトル（行）や行成分の略記は \((1\ 0\ 0)\)
のように\textbf{コンマを打たず}、成分の間だけを区切って書く。\((1, 0, 0)\)
のように\textbf{コンマのあとにスペース}を入れる書き方は\textbf{座標}（点の位置など）を強調するときに用い、\textbf{行列・横ベクトルの略記には使わない}。
\end{quote}

この一行が分かれば、積分記号の末尾にある \(dx\)
の見え方が変わる。次に、通常のリーマン和をこの記法で書き直してみよう。

\hypertarget{ux30eaux30fcux30deux30f3ux548cux306eux30d9ux30afux30c8ux30ebux518dux89e3ux91c8ux3068ux7a4dux5206ux8a18ux53f7}{%
\subsubsection{1.1.4
リーマン和のベクトル再解釈と積分記号}\label{ux30eaux30fcux30deux30f3ux548cux306eux30d9ux30afux30c8ux30ebux518dux89e3ux91c8ux3068ux7a4dux5206ux8a18ux53f7}}

この視点でリーマン和を書き直すと、

\[
R_n
=
\sum_{i=1}^n f(\xi_i)\,dx(\mathbf{v}_i)
\]

となる。ここで

\[
dx(\mathbf{v}_i)=\Delta x_i
\]

だから、これは通常のリーマン和

\[
R_n
=
\sum_{i=1}^n f(\xi_i)\,\Delta x_i
\]

と同じものである。

しかし、この書き方にはもう一つの読み方がある。各小区間で、関数値
\(f(\xi_i)\) を測定器 \(dx\) に掛けた横ベクトル

\[
\bigl(f(\xi_i)\,dx\bigr)
:=
f(\xi_i)\,dx
=
\begin{pmatrix}
f(\xi_i) & 0 & 0
\end{pmatrix}
\]

を考える。この横ベクトルを変位ベクトル \(\mathbf{v}_i\) に作用させると、

\[
\bigl(f(\xi_i)\,dx\bigr)(\mathbf{v}_i)
=
f(\xi_i)\,dx(\mathbf{v}_i)
=
f(\xi_i)\,\Delta x_i
\]

となる。つまり、リーマン和の各項そのものが得られる。

したがって、リーマン和は

\[
R_n
=
\sum_{i=1}^n
\bigl(f(\xi_i)\,dx\bigr)(\mathbf{v}_i)
\]

とも読める。

分割を無限に細かくする極限では、つまり幅の最大値が \(0\)
に近づく極限では、リーマン積分の定義より

\[
\int_a^b f(x)\,dx
=
\lim_{n\to\infty} R_n
=
\lim_{n\to\infty}
\sum_{i=1}^n
\bigl(f(\xi_i)\,dx\bigr)(\mathbf{v}_i)
\]

となる。

左辺の \(\int_a^b f(x)\,dx\)
は、高校以来おなじみの積分記号そのものである。すなわち、上の手順で書いた
\(\lim_{n\to\infty}R_n\) と同じ量だ。

右辺は、\textbf{各小区間で「測定器}\(f(\xi_i)\,dx\)\textbf{を変位ベクトル}\(\mathbf{v}_i\)\textbf{に作用させた値」を足し上げたものの極限}という、同じ積分の別顔だ。

つまり、積分記号の末尾にある \(dx\)
は、単なる飾りではない。少なくともこの読み方では、各小区間の変位ベクトルから
\(x\) 方向の幅を取り出す測定器として働いている。そして \(f(x)\,dx\)
は、その測定器に関数の値を掛けた、新しい横ベクトルとして読める。

\begin{quote}
\textbf{注}（仕事の積分）\\
力学で現れる \(W=\int F(x)\,dx\)
も、同じ読み方ができる。各小区間では、横ベクトル \(F(\xi_i)\,dx\)
が変位ベクトル \(\mathbf v_i\) を食べて \(F(\xi_i)\Delta x_i\)
を返す。その和の極限が仕事である。
\end{quote}

\hypertarget{ux4e00ux6b21ux5f62ux5f0f1-form}{%
\subsubsection{\texorpdfstring{1.1.5
一次形式（\(1\)-form）}{1.1.5 一次形式（1-form）}}\label{ux4e00ux6b21ux5f62ux5f0f1-form}}

ここまで、我々は \(dx\)
を「変位ベクトルに作用して特定の成分を抽出し、スカラー（実数）を返す横ベクトル」として定義してきた。このような、ベクトルを食べてスカラーを吐き出す\textbf{線形}な測定器のことを、専門的には\textbf{一次形式（}\(1\)\textbf{-form）}と呼ぶ。本稿で「行列としての
\(dx\)」と呼んできたものは、まさにこの一次形式に他ならない。

\begin{quote}
\textbf{注} （余ベクトルという用語）数学者は
\textbf{一次形式（}\(1\)\textbf{-form）} の別名として
\textbf{余ベクトル（covector）}と呼ぶ流儀をとることもよくある。他書を読むときの辞書として、頭の片隅に置いておけばよい。
\end{quote}

命名の由来は単純だ：

\begin{itemize}
\item
  \textbf{「一次」}：変位ベクトルを\textbf{一次（線形）}に処理するから
\item
  \textbf{「形式」}：「測定の形式」「作用の形式」を意味する
\end{itemize}

これ以降、この測定器のことを「行列」とも「一次形式（\(1\)-form）」とも、演算子とも呼ぶことにする。いずれの呼び方も、同じものを指している。

\begin{quote}
\textbf{注}（行列・一次形式・測定器という呼び分け）\\
繰り返すが、本書の立場では、これらは同じ配列で表される。したがって本書を通して必要に応じて「行列」「一次形式」「測定器」と呼び分ける。
\end{quote}

\hypertarget{ux672cux66f8ux306eux8a18ux53f7ux5951ux7d04}{%
\subsubsection{1.1.6
本書の記号契約}\label{ux672cux66f8ux306eux8a18ux53f7ux5951ux7d04}}

本書では、\(dx\)\textbf{を単独で「微小な変位」や「微小変化」の意味には用いない}------変位の幅と測定器（横ベクトル）を混同しないために、次のように約束する。\(x\)
方向の（微小）変位の大きさは \(\Delta x\) を使うか、変位ベクトル
\(\mathbf{v}\) を明示して \(dx(\mathbf{v})\) と書く。

一方、\textbf{単独で現れる}\(dx\)
は、\(x\)\textbf{成分を抜き出す横ベクトル（一次形式）としての演算子}を指す。式
\(df = f'(x)\,dx\) のように並ぶ \(dx\)
も、\textbf{常に演算子として読む}。\textbf{この区別を徹底することが、以降のすべての理解の土台となる。}

以降、積分記号 \(\int_a^b f(x)\,dx\) の末尾の \(dx\)
は、高校以来の\textbf{慣用記法}として残すが、積分記号の外では必ず
\(dx(\mathbf{v})\)
の形で示し、\textbf{変位や微小幅を語る本文では}\(\Delta x\)\textbf{か}\(dx(\mathbf{v})\)\textbf{に揃え、裸の}\(dx\)\textbf{には幅や変化量のイメージを載せない}。

\begin{quote}
\textbf{注} （文献における \(dx\) の用法）物理学者は変位の成分そのものを
\(dx\)
と書く慣習もよく使う。しかし、本書の記法に慣れると、\textbf{測定器としての}\(dx\)\textbf{と変位のスカラーを切り分けた}読みがしやすくなる------より進んだ本では、その区別がそのまま効いてくる。
\end{quote}

次節では、この視点を関数の微分へと拡張し、全微分 \(df\)
を行列として定義していく。

\begin{center}\rule{0.5\linewidth}{0.5pt}\end{center}

\begin{quote}
\textbf{【ここまでのチェックポイント】}

\begin{itemize}
\item
  微小変位は縦ベクトル \(\mathbf{v}\)。\(\Delta x\)
  はスカラー幅であり、単独の \(dx\)
  は横ベクトル（一次形式）としての演算子である。
\item
  デカルトの \(dy\), \(dz\) は\textbf{次節 §1.2.3}で置く（§1.1.6 の
  \(dx\) と\textbf{同じ考え方の}一次形式）。\(\Delta y\), \(\Delta z\)
  や \(dy(\mathbf{v})\), \(dz(\mathbf{v})\) との区別は §1.1.6
  の\textbf{記号契約}に沿う。
\item
  リーマン和の各項は \((f\,dx)(\mathbf{v}_i)\)
  の形で、積分はその極限として理解できる。
\item
  積分記号 \(\int_a^b f(x)\,dx\) の末尾の \(dx\)
  は慣用記法であり、変位そのものは \(\Delta x\) や \(dx(\mathbf{v})\)
  で書く（§1.1.6 の契約）。
\end{itemize}
\end{quote}

\begin{center}\rule{0.5\linewidth}{0.5pt}\end{center}

\hypertarget{ux5168ux5faeux5206-df-ux5909ux5316ux7387ux3092ux884cux5217ux306bux307eux3068ux3081ux305fux6f14ux7b97ux5b50}{%
\subsection{\texorpdfstring{§1.2 全微分 \(df\) ---
変化率を行列にまとめた演算子}{§1.2 全微分 df --- 変化率を行列にまとめた演算子}}\label{ux5168ux5faeux5206-df-ux5909ux5316ux7387ux3092ux884cux5217ux306bux307eux3068ux3081ux305fux6f14ux7b97ux5b50}}

前節では、リーマン積分を「行列 \(f(x)\,dx\)
の変位ベクトルへの作用の和の極限」として解体した。この節では、その考え方を関数自体の微分へと拡張し、\textbf{全微分}\(df\)\textbf{を横ベクトル（行列）として定義する}。「微分」は、単なる数値の変化率ではなく、\textbf{変位ベクトルに掛けて初めて「変化量の一次部分」が出てくる演算子}として扱う。
\#\#\#\# 1.2.1 微分可能性の行列表現 関数 \(f(x)\) が点 \(x\)
で微分可能であるとは、次の一次近似が成り立つことである：
\[\Delta f = f(x + \Delta x) - f(x) = f'(x) \Delta x + o(|\Delta x|) \quad (|\Delta x| \to 0)\]

\begin{quote}
\textbf{注} （ランダウのオーダー記法）\(o(|\Delta x|)\)
は\textbf{ランダウのオーダー記法}である。\(\Delta x \to 0\)
のとき、\(o(|\Delta x|)\) でまとめて書いた量は
\(|\Delta x|\)\textbf{よりずっと速くゼロに近づく余り}を意味し、主項
\(f'(x)\Delta x\)
に比べれば\textbf{無視してよいほど小さい}、という約束である。理工系の専門書ではよく出るが、初めて見る読者もいるかもしれない。
\end{quote}

ここで \(\Delta x\)
はスカラーだが、我々はこれを前節同様、3次元の変位ベクトルとして捉え直す：
\[\mathbf{v} = \begin{pmatrix} \Delta x \\ 0 \\ 0 \end{pmatrix}\]

すると、関数の変化 \(\Delta f\) は、この \(\mathbf{v}\)
によって引き起こされる効果と見なせる。

\hypertarget{df-ux306eux884cux5217ux3068ux3057ux3066ux306eux5b9aux7fa9ux3068-dfmathbfv}{%
\subsubsection{\texorpdfstring{1.2.2 \(df\) の行列としての定義と
\(df(\mathbf{v})\)}{1.2.2 df の行列としての定義と df(\textbackslash mathbf\{v\})}}\label{df-ux306eux884cux5217ux3068ux3057ux3066ux306eux5b9aux7fa9ux3068-dfmathbfv}}

点 \(x\) における関数 \(f\) の\textbf{全微分} \(df\) を、次の
\(1 \times 3\) 横ベクトルとして定義する：
\[df := f'(x) \, dx = \begin{pmatrix} f'(x) & 0 & 0 \end{pmatrix}\]

この \(df\) を変位 \(\mathbf{v}\) に作用させた結果を、前節と同様に
\(df(\mathbf{v})\) と書く：
\[df(\mathbf{v}) = \begin{pmatrix} f'(x) & 0 & 0 \end{pmatrix} \begin{pmatrix} \Delta x \\ 0 \\ 0 \end{pmatrix} = f'(x)\,\Delta x\]

\textbf{重要な認識の転換}（\textbf{常に次のように読む}）：

\begin{itemize}
\item
  式 \(df = f'(x)\,dx\)
  は、微積分でおなじみの\textbf{全微分の記法}だが、本書では
  \(df\)\textbf{も}\(dx\)\textbf{も演算子としての横ベクトル}であり、\textbf{単体では微小変化量ではない}、と読む。
\item
  \(df(\mathbf{v})\)
  と書いたときは、\textbf{有限の（ただし十分小さい）変位}\(\mathbf{v}\)\textbf{に対する一次近似としての変化量}を意味する。
\end{itemize}

\(df\)
それ自体は変化量ではない。\textbf{変化量の一次部分を測定する「測定器」}なのである。

\begin{quote}
\textbf{注} （\(df\) は演算子、\(df(\mathbf{v})\) が変化量）§1.1.6
の記号契約と同じ区別を、ここでもう一度述べる。くどいと感じるかもしれないが、\textbf{ここを誤読すると以降すべてがずれる}ので、繰り返しに価値がある。
\end{quote}

\hypertarget{y-ux65b9ux5411ux3068-z-ux65b9ux5411ux306eux6e2cux5b9aux5668-dy-dz}{%
\subsubsection{\texorpdfstring{1.2.3 \(y\) 方向と \(z\) 方向の測定器
\(dy\),
\(dz\)}{1.2.3 y 方向と z 方向の測定器 dy, dz}}\label{y-ux65b9ux5411ux3068-z-ux65b9ux5411ux306eux6e2cux5b9aux5668-dy-dz}}

この視点の強みを、§1.1.5
の力学の例に戻せば、\textbf{力が}\(y\)\textbf{方向にも成分を持つ一般的な場合に自然に拡張できる}点だ。
\textbf{いま定める}\(dy\)を使えば、\(F_x dx + F_y dy\)
という形で2次元の仕事を統一的に扱える。 \(dx\) が \(x\)
成分を抜き出すのと同様に、実空間では
\(dy\)\textbf{は}\(y\)\textbf{成分}、\(dz\)\textbf{は}\(z\)\textbf{成分}を抜き出す横ベクトル（一次形式）と定める。行列表現は
\[dy := \begin{pmatrix} 0 & 1 & 0 \end{pmatrix}, \qquad dz := \begin{pmatrix} 0 & 0 & 1 \end{pmatrix}\]
である。変位
\(\mathbf{v} = \begin{pmatrix} \Delta x \\ \Delta y \\ \Delta z \end{pmatrix}\)
に対して \(dy(\mathbf{v}) = \Delta y\), \(dz(\mathbf{v}) = \Delta z\)
となる。

\begin{quote}
\textbf{注} （\(dy\), \(dz\) の契約）単独の \(dy\), \(dz\)
は演算子であり、\(y\) や \(z\) の微小幅を語るときは \(\Delta y\),
\(\Delta z\) または \(dy(\mathbf{v})\), \(dz(\mathbf{v})\)
と書く。積分記号の末尾に並ぶ \(dy\), \(dz\) も、高校以来の慣用記法として
\(dx\)\textbf{と同じ仕方で読めばよい}。細部の約束は、§1.1.6 で \(dx\)
についてまとめた\textbf{記号契約}と\textbf{同じ考え方}を適用する。
\end{quote}

3次元の空間には、\textbf{三つの測定器}がそろった。本章では説明の主線として
\(x\) 方向の断面に寄せてきたが、座標 \(y,z\) と測定子 \(dy,dz\)
は最初からそろっている、と考えてほしい。 いま、一変数の \(df=f'(x)\,dx\)
に \(dy\), \(dz\) を同じ型の測定器として足しそろえた。以下の具体例と
§1.2.6 での三次元への拡張で、この三つがどう効いてくるかを追う。

\hypertarget{ux5177ux4f53ux4f8b-fx-x2-ux306eux5834ux5408}{%
\subsubsection{\texorpdfstring{1.2.4 具体例： \(f(x) = x^2\)
の場合}{1.2.4 具体例： f(x) = x\^{}2 の場合}}\label{ux5177ux4f53ux4f8b-fx-x2-ux306eux5834ux5408}}

\(f(x) = x^2\) とする。\(f'(x) = 2x\) である。

たとえば
\(x=3\)\textbf{において}全微分を具体的に書き下ろす。横ベクトルと縦ベクトルを並べて、
\[df = 6\,dx = \begin{pmatrix} 6 & 0 & 0 \end{pmatrix}, \qquad \mathbf{v} = \begin{pmatrix} 0.1 \\ 0 \\ 0 \end{pmatrix}\]
であるから、行列の積は
\[df(\mathbf{v}) = \begin{pmatrix} 6 & 0 & 0 \end{pmatrix} \begin{pmatrix} 0.1 \\ 0 \\ 0 \end{pmatrix} = 6\cdot 0.1 + 0\cdot 0 + 0\cdot 0 = 0.6\]

実際、\(f(3.1) - f(3) = 9.61 - 9 = 0.61\) であり、一次近似 \(0.6\)
はよく一致している。

\hypertarget{ux7f6eux63dbux7a4dux5206ux3068ux6e2cux5b9aux5668ux306eux4f5cux308aux66ffux3048}{%
\subsubsection{1.2.5
置換積分と測定器の作り替え}\label{ux7f6eux63dbux7a4dux5206ux3068ux6e2cux5b9aux5668ux306eux4f5cux308aux66ffux3048}}

ここまで、\(dx\) を「変位を測る測定器」として読んできた。

すなわち、\(dx\) は単独の微小量ではない。変位ベクトルを食べて、その
\(x\) 成分を返す横ベクトルである。本書ではこの読み方を基本にする。

この読み方を持つと、高校以来おなじみの置換積分も、少し違って見える。

変数 \(x\) が、別の変数 \(t\) によって

\[
x=\gamma(t)
\]

と表されているとする。このとき、置換積分ではしばしば

\[
dx=\gamma'(t)\,dt
\]

のように書く。

通常の微積分では、この式は合成関数の微分から説明される。たとえば、\(F'(x)=f(x)\)
として、\(F(\gamma(t))\) を \(t\) で微分すると、

\[
\frac{d}{dt}F(\gamma(t))
=
f(\gamma(t))\gamma'(t)
\]

となる。したがって、

\[
\int_{t_0}^{t_1} f(\gamma(t))\gamma'(t)\,dt
=
F(\gamma(t_1))-F(\gamma(t_0))
\]

である。端点を

\[
x_0=\gamma(t_0),\qquad x_1=\gamma(t_1)
\]

と書けば、右辺は

\[
F(x_1)-F(x_0)
=
\int_{x_0}^{x_1} f(x)\,dx
\]

である。よって、

\[
\int_{x_0}^{x_1} f(x)\,dx
=
\int_{t_0}^{t_1} f(\gamma(t))\gamma'(t)\,dt
\]

が得られる。

ここまでは、標準的な置換積分の説明である。

しかし本書では、この式を少し違う目で見る。ここで現れた

\[
dx=\gamma'(t)\,dt
\]

は、単なる計算上の記号ではない。\(x\) 側の測定器 \(dx\) を、\(t\)
側の測定器 \(dt\) に作り替えた姿として読める。

\(t\) が少し変わったとき、\(x=\gamma(t)\) はその \(\gamma'(t)\)
倍だけ変わる。だから、\(t\) 側の小区間を食べたときに、\(x\)
側の変位と同じ値を返す測定器を作るには、\(dt\) に \(\gamma'(t)\)
を掛ける必要がある。

である。

この見方を本格的に扱うのが、第4章の引き戻しである。

第4章では、いきなりこの式を天下りに使うのではなく、有限区間を写し、その像を測るところから、同じ係数
\(\gamma'(t)\)
を発見し直す。さらに同じ構造を、有限マス目、有限箱へ拡張する。

ここでは、置換積分の標準的な計算を一度だけ確認した。だが本書の本筋では、これを後に「測定器のつじつま合わせ」として読み直していく。

\begin{quote}
\textbf{注}（\(dx=\gamma'(t)\,dt\)）\\
通常の微積分では、置換 \(x=\gamma(t)\) に対して、この関係を
\(dx=\gamma'(t)\,dt\) と書く。

ただし本書の立場では、これは「\(x\) 側の測定器 \(dx\) そのものが \(t\)
側へ移動した」という意味ではなく、写像 \(\gamma:I\to\mathbb{R}\)
に沿って、値域側の一次形式 \(dx\)
を定義域側へ引き戻した、という意味に読む。

より正確には、第4章で述べる\textbf{引き戻し}の記法を使って、\(\gamma^\ast(dx)=\gamma'(t)\,dt\)
と書くほうが安全である。つまり、左辺の \(dx\) はもともと \(x\)
側の変位を測る道具であり、右辺の \(\gamma'(t)\,dt\) は、\(t\)
側の変位を測って同じ一次変化を返すように作り直された道具である。

少なくとも本書の立場から見ると、この区別が隠れていることが、置換積分における
\(dx\) の説明を分かりにくくしている。
\end{quote}

\hypertarget{ux4e09ux6b21ux5143ux3078ux306eux62e1ux5f35}{%
\subsubsection{1.2.6
三次元への拡張}\label{ux4e09ux6b21ux5143ux3078ux306eux62e1ux5f35}}

空間の各点 \((x,y,z)\) にスカラー（実数）を対応させる関数 \(f\)
は、実空間上の\textbf{実数値関数}として扱えばよい。

\begin{quote}
\textbf{注}
（スカラー場）物理学ではこのような関数を\textbf{スカラー場}と呼ぶことが多い。
\end{quote}

その \(f(x,y,z)\) が点 \((x,y,z)\)
で\textbf{（全）微分可能である}とは、変位
\(\mathbf{v}=\begin{pmatrix}\Delta x\\\Delta y\\\Delta z\end{pmatrix}\)
に対して \[
\begin{aligned}
\Delta f
&= f(x+\Delta x,\,y+\Delta y,\,z+\Delta z) - f(x,y,z) \\
&= \frac{\partial f}{\partial x}\,\Delta x
{}+ \frac{\partial f}{\partial y}\,\Delta y
{}+ \frac{\partial f}{\partial z}\,\Delta z
{}+ o(\|\mathbf{v}\|) \quad (\|\mathbf{v}\|\to 0)
\end{aligned}
\] が成り立つことである。§1.2.1
の1変数の定義と形式がそろっている。\(o(\|\mathbf{v}\|)\)
は余り項の略記（\(o(|\Delta x|)\)
と同趣旨）であり、ここでは深追いしないが、理解に支障はないであろう。

§1.2.3 の測定器 \(dx,dy,dz\) を使い、§1.2.2 と同じ型の演算子として
\[df := \frac{\partial f}{\partial x}\,dx + \frac{\partial f}{\partial y}\,dy + \frac{\partial f}{\partial z}\,dz = \begin{pmatrix} \frac{\partial f}{\partial x} & \frac{\partial f}{\partial y} & \frac{\partial f}{\partial z} \end{pmatrix}\]
と定義する。いちばん初めの式にあった \(o(\|\mathbf{v}\|)\)
は、ここから数行では省略し、\textbf{記号の骨格だけ}を追う。\(\mathbf{v}=\begin{pmatrix}\Delta x\\\Delta y\\\Delta z\end{pmatrix}\)
のとき、微分可能性の定義より
\(\Delta f = df(\mathbf{v}) + o(\|\mathbf{v}\|)\) である------すなわち
\(df(\mathbf{v})\)\textbf{は厳密な変化量}\(\Delta f\)\textbf{そのものではなく、剰余を除いた一次（主）項}である。以下ではその主項を行列の積で書き下ろす：
\[df(\mathbf{v}) = \frac{\partial f}{\partial x}\begin{pmatrix}1&0&0\end{pmatrix}\begin{pmatrix}\Delta x\\\Delta y\\\Delta z\end{pmatrix} + \frac{\partial f}{\partial y}\begin{pmatrix}0&1&0\end{pmatrix}\begin{pmatrix}\Delta x\\\Delta y\\\Delta z\end{pmatrix} + \frac{\partial f}{\partial z}\begin{pmatrix}0&0&1\end{pmatrix}\begin{pmatrix}\Delta x\\\Delta y\\\Delta z\end{pmatrix}\]
\[= \frac{\partial f}{\partial x}\,\Delta x + \frac{\partial f}{\partial y}\,\Delta y + \frac{\partial f}{\partial z}\,\Delta z.\]
第一行は、§1.2.3 の測定器 \(dx,dy,dz\) が行ベクトルとして \(\mathbf{v}\)
から \(x,y,z\)
成分を抜き出し、偏導関数が係数としてかかる様子を\textbf{行列の積}で見せたものである。第二行はその評価結果であり、先の
\(\Delta f\)
の定義式における\textbf{線形主項}と同じ形である。\textbf{横ベクトル（演算子）× 縦ベクトル（変位）→ スカラー}という
§1.2.2 の読みが、三次元でも変わらない。

座標が二つまでしか現れない \(f(x,y)\) も、\(z\) に依らなければ
\(\partial f/\partial z=0\)
として上の枠に含めればよい------\textbf{三次元に載せた枠組みが、素直に埋め込まれる}。

見比べてみよう。 1変数では：
\(df := f'(x) dx = \begin{pmatrix} f'(x) & 0 & 0 \end{pmatrix}\)
2変数関数 \(f(x, y)\) では（§1.2.3 の \(dy\) を用いる）：
\(df := \frac{\partial f}{\partial x} dx + \frac{\partial f}{\partial y} dy = \begin{pmatrix} \frac{\partial f}{\partial x} & \frac{\partial f}{\partial y} & 0 \end{pmatrix}\)
3変数 \(f(x,y,z)\) では、§1.2.6
の式のとおり、横ベクトルの成分が三つそろうだけである。

このように、\textbf{形式は同じで、ただ成分の個数が増える}。これが\(3\)次元の行列表示の利点である。\(df\)
を行列と見なすことで、微分は「変位に対する線形近似を与える演算子」として統一的に扱える。

\hypertarget{ux7ddaux7a4dux5206ux306eux5148ux53d6ux308a}{%
\subsubsection{1.2.7
線積分の先取り}\label{ux7ddaux7a4dux5206ux306eux5148ux53d6ux308a}}

空間内の曲線をパラメータ \(t\) で
\(\mathbf{r}(t)=\begin{pmatrix}x(t)\\y(t)\\z(t)\end{pmatrix}\)
と表し、細かいステップの変位を \(\Delta\mathbf{r}\)
と書くとき、各ステップで \(\bigl(df\bigr)(\Delta\mathbf{r})\)
を足し上げる操作の極限を、記号では \[\int_\gamma df\]
の形で表す（\(\gamma\) は曲線）。§1.2.5 で見た \(\int f'(x)\,dx\)
と同じく、\textbf{一次形式}\(df\)\textbf{が各ステップの変位に作用するリーマン和の極限}という骨格である。閉曲線・向き・パラメータの取り換えなど、詳細な定式化は\textbf{後の章に譲る}。

\begin{center}\rule{0.5\linewidth}{0.5pt}\end{center}

\begin{quote}
\textbf{【ここまでのチェックポイント】}

\begin{itemize}
\item
  微分可能性は \(\Delta f = f'(x)\Delta x + o(|\Delta x|)\)（\(o\)
  はランダウのオーダー記法）で表される一次近似として捉えられる。
\item
  \(df = f'(x)\,dx\) は横ベクトル；数値の変化量は \(df(\mathbf{v})\)
  として読む。
\item
  置換積分の標準的な入口を一度だけ見るが、本書では「測定器のつじつま合わせ」として読み直す（§1.2.5）。
\end{itemize}
\end{quote}

\begin{center}\rule{0.5\linewidth}{0.5pt}\end{center}

\hypertarget{ux30e9ux30a4ux30d7ux30cbux30c3ux30c4ux306eux8a18ux6cd5ux3068ux4ee3ux6570ux7684ux76f4ux611f}{%
\subsection{§1.3
ライプニッツの記法と代数的直感}\label{ux30e9ux30a4ux30d7ux30cbux30c3ux30c4ux306eux8a18ux6cd5ux3068ux4ee3ux6570ux7684ux76f4ux611f}}

ライプニッツが \(dx\), \(dy\)
という記号を導入したとき、彼はこれらを「無限小」として直感的に扱った。

\begin{quote}
\textbf{注} （ライプニッツと記号）\textbf{ライプニッツ}（Gottfried
Wilhelm Leibniz, 1646--1716）は、\(dx\), \(dy\)
など微積分の記号体系を整えた数学者である。以下では歴史の経緯には立ち入らず、記号を導入した人物として名前だけに触れる（記号
\(dx\)
自体は慣れている読者も多いが、人名に心当たりがなければ、この一行で足りる）。
\end{quote}

現代の我々は、その直感を\textbf{線形代数の言葉}で再配置したと言える。
ライプニッツの記法 \(df = f'(x)dx\) は、単なる形式的等式ではない：

\begin{itemize}
\item
  \(dx\)：ライプニッツの無限小の直感 → 本書では \textbf{演算子}としての
  \(x\) 成分抽出（微小変位そのものを語るときは \(\Delta x\) や
  \(dx(\mathbf{v})\)）。\(dy\), \(dz\) も §1.2.3 で \(dx\)
  と同様に定義する。
\item
  \(df\)：ライプニッツの無限小変化の直感 → 本書では
  \textbf{演算子}としての全微分（数値の変化量は \(df(\mathbf{v})\)
  などで）
\end{itemize}

\textbf{ライプニッツの天才は、微分・積分が本質的に代数的操作であることを見抜いていた}。我々は、彼の直感に行列という具体的な骨格を与えたに過ぎない。

次節では、この枠組みが依存している暗黙の前提------デカルト座標系の「使いやすさ」------を明らかにし、より一般的な座標への接続を示唆する。

\begin{center}\rule{0.5\linewidth}{0.5pt}\end{center}

\begin{quote}
\textbf{【ここまでのチェックポイント】}

\begin{itemize}
\item
  ライプニッツの \(dx\)、\(df\)
  の直感は、本書ではいずれも「変位に作用する横ベクトル（演算子）」として再配置されている。
\item
  微小変化量を語るときは \(\Delta x\) や \(df(\mathbf{v})\)
  とし、記号の役割を混同しない。
\end{itemize}
\end{quote}

\begin{center}\rule{0.5\linewidth}{0.5pt}\end{center}

\hypertarget{ux5ea7ux6a19ux5909ux63dbux6e2cux5b9aux5668ux3092ux5225ux306eux76eeux76dbux308aux306bux4f5cux308aux66ffux3048ux308b}{%
\subsection{§1.4
座標変換------測定器を別の目盛りに作り替える}\label{ux5ea7ux6a19ux5909ux63dbux6e2cux5b9aux5668ux3092ux5225ux306eux76eeux76dbux308aux306bux4f5cux308aux66ffux3048ux308b}}

\begin{quote}
\textbf{注}（この節の読み方）

この節は、多少分からなくても問題ない。ここで本当に理解してほしいのは、円柱座標の公式そのものではなく、「\(dx\)
を行列として見ると、座標を替えたときにも何か計算できそうだ」という雰囲気である。

繰り返すが、ここは雰囲気でよい。第4章にて、同じ考え方をもっと紙面を割いて扱う。そこでは、有限区間や有限マス目を実際に写し、その像を測るところから出発することで、「十分小さい」という仮定を必要とせず、測定器の作り替えをより整理して説明する。以下は、パラメータ空間での変化が十分に小さいとして、一次の項だけを取り出す。これは、従来の微積分や物理数学でよく使われる近似的な見方に寄せた説明である。
\end{quote}

本章を通じて、我々は \(dx\) の\textbf{行列表現}を

\[
dx:=\begin{pmatrix}1&0&0\end{pmatrix}
\]

と置いて議論してきた。これは、実空間をデカルト座標 \((x,y,z)\)
で見ている限り、\(x\) 成分を抜き出す測定器である。

しかし物理学の現実は、いつでも四角いとは限らない。円形パイプ内の流れや、直線電流の周りの磁場のような問題では、点を
\((r,\theta,z)\)
という数の組で指定する\textbf{円柱座標}を使ったほうが扱いやすい。

では、実空間の測定器 \(dx\)
は、円柱座標の目盛りではどう見えるのだろうか。

ここで大事なのは、「座標変換の公式を暗記する」ことではない。実空間の測定器を、別の変数の世界で同じ一次の値を返す測定器へ作り替えることである。

\begin{quote}
\textbf{注}（円柱座標の局所性）\\
ここでは、\(\theta\) の周期性や \(r=0\)
の特異性は扱わず、局所的なパラメータ表示として計算する。

実際、円柱座標や球座標を用いる物理数学の計算では、積分範囲や対象領域を適切に取ることで、座標の特異点を避けて計算できる場合が多い。

ただし、原点や軸を含む領域、あるいは角度の周期性が本質的に効く問題では、この局所的な表示だけでは不十分になる。そのような場合には、座標パッチや境界条件をより慎重に扱う必要がある。
\end{quote}

\hypertarget{ux5186ux67f1ux5ea7ux6a19ux304bux3089ux5b9fux7a7aux9593ux3078ux306eux5909ux63db}{%
\subsubsection{1.4.1
円柱座標から実空間への変換}\label{ux5186ux67f1ux5ea7ux6a19ux304bux3089ux5b9fux7a7aux9593ux3078ux306eux5909ux63db}}

円柱座標のパラメータ空間から、実空間への変換を

\[
\Phi(r,\theta,z)
=
\begin{pmatrix}
r\cos\theta\\
r\sin\theta\\
z
\end{pmatrix}
\]

と書く。

これは、パラメータ空間の点 \((r,\theta,z)\) を、実空間の点 \((x,y,z)\)
に送る変換である。つまり、

\[
x=r\cos\theta,\qquad
y=r\sin\theta,\qquad
z=z
\]

である。

いま知りたいのは、実空間で \(x\) 成分を測る測定器 \(dx\)
が、パラメータ空間側ではどのような測定器として見えるかである。

\hypertarget{ux5c0fux3055ux306aux4e00ux6b69ux3092ux5199ux3057ux3066-dx-ux3067ux6e2cux308b}{%
\subsubsection{\texorpdfstring{1.4.2 小さな一歩を写して \(dx\)
で測る}{1.4.2 小さな一歩を写して dx で測る}}\label{ux5c0fux3055ux306aux4e00ux6b69ux3092ux5199ux3057ux3066-dx-ux3067ux6e2cux308b}}

パラメータ空間の点

\[
p=(r,\theta,z)
\]

から、小さな一歩

\[
\mathbf{h}
=
\begin{pmatrix}
\Delta r\\
\Delta\theta\\
\Delta z
\end{pmatrix}
\]

だけ動くことを考える。

このとき、移動後の点は

\[
p+\mathbf{h}
=
(r+\Delta r,\theta+\Delta\theta,z+\Delta z)
\]

である。

ただし、この \(\mathbf h\)
は実空間の変位ベクトルではない。これは、円柱座標のパラメータ空間での変化量を並べたものである。

この小さな一歩を、変換 \(\Phi\) によって実空間へ送る。実空間での変位は

\[
\Phi(p+\mathbf h)-\Phi(p)
\]

である。

成分で書けば、

\[
\Phi(r+\Delta r,\theta+\Delta\theta,z+\Delta z)
-
\Phi(r,\theta,z)
=
\begin{pmatrix}
(r+\Delta r)\cos(\theta+\Delta\theta)-r\cos\theta\\
(r+\Delta r)\sin(\theta+\Delta\theta)-r\sin\theta\\
\Delta z
\end{pmatrix}
\]

である。

このベクトルは、実空間における本当の変位である。

実空間の測定器 \(dx\) は、実空間の変位ベクトルから \(x\)
成分だけを抜き出す。したがって、いまの変位に \(dx\) を作用させると、

\[
dx\bigl(\Phi(p+\mathbf h)-\Phi(p)\bigr)
=
(r+\Delta r)\cos(\theta+\Delta\theta)-r\cos\theta
\]

となる。

これは、「パラメータ空間で \(\mathbf h\) だけ動いたとき、実空間の \(x\)
成分がどれだけ変化したか」を測った値である。

\hypertarget{ux4e00ux6b21ux9805ux304bux3089ux6e2cux5b9aux5668ux3092ux4f5cux308b}{%
\subsubsection{1.4.3
一次項から測定器を作る}\label{ux4e00ux6b21ux9805ux304bux3089ux6e2cux5b9aux5668ux3092ux4f5cux308b}}

ここで、\(\Delta r,\Delta\theta,\Delta z\)
が小さいとして、一次の項だけを見る。

まず、

\[
\Delta x
=
(r+\Delta r)\cos(\theta+\Delta\theta)-r\cos\theta
\]

である。

一次近似を使うと、

\[
\cos(\theta+\Delta\theta)
=
\cos\theta-\sin\theta\,\Delta\theta
+
\text{高次の項}
\]

だから、

\[
\begin{aligned}
\Delta x
&=
(r+\Delta r)\cos(\theta+\Delta\theta)-r\cos\theta\\
&=
(r+\Delta r)(\cos\theta-\sin\theta\,\Delta\theta+\text{高次の項})-r\cos\theta\\
&=
r\cos\theta
+\cos\theta\,\Delta r
-r\sin\theta\,\Delta\theta
+\text{高次の項}
-r\cos\theta\\
&=
\cos\theta\,\Delta r
-r\sin\theta\,\Delta\theta
+\text{高次の項}
\end{aligned}
\]

となる。

\(\Delta z\) は \(x\) 成分には影響しないので、一次の項としては

\[
\Delta x
=
\cos\theta\,\Delta r
-r\sin\theta\,\Delta\theta
+0\cdot\Delta z
+
\text{高次の項}
\]

である。

いま得られた一次項は

\[
\cos\theta\,\Delta r
-r\sin\theta\,\Delta\theta
+0\cdot\Delta z
\]

である。

これは、パラメータ空間の小さな一歩

\[
\mathbf h
=
\begin{pmatrix}
\Delta r\\
\Delta\theta\\
\Delta z
\end{pmatrix}
\]

に対して、次の横ベクトルを作用させた値に等しい。

\[
\begin{pmatrix}
\cos\theta & -r\sin\theta & 0
\end{pmatrix}
\begin{pmatrix}
\Delta r\\
\Delta\theta\\
\Delta z
\end{pmatrix}
=
\cos\theta\,\Delta r
-r\sin\theta\,\Delta\theta
+0\cdot\Delta z
\]

したがって、実空間の測定器 \(dx\)
を、円柱座標のパラメータ空間側で同じ一次の値を返す測定器として書き直すと、

\[
\begin{pmatrix}
\cos\theta & -r\sin\theta & 0
\end{pmatrix}
\]

になる。

これが、円柱座標の目盛りで見た \(dx\) の姿である。

\begin{quote}
\textbf{注}（大きな幾何を小さな幾何へ分ける）

多くの教科書では、実空間に \((r,\theta)\)
の「曲がったグリッド」を描き込み、「\(\theta\) 方向の弧の長さは約
\(r\Delta\theta\)」と幾何学的に説明する。この直感自体は正しい。

ただし本書では、まずパラメータ空間の小さな一歩を実空間へ写し、その像を測る、という代数的な手順を優先する。

係数 \(-r\sin\theta\) は「\(\theta\) 方向の一歩が実空間の \(x\)
方向にどれだけ効くか」という情報を持っている。同様に、\(\cos\theta\) は
\(r\) 方向が \(x\) 方向にどれだけ効くかを、\(0\) は \(z\) 方向が \(x\)
方向に効かないことを示している。

つまり \(\begin{pmatrix}\cos\theta & -r\sin\theta & 0\end{pmatrix}\)
という一行は、\(x\)
成分の変化に関する幾何を、三つの小さな成分へ分解した姿なのである。
\end{quote}

\hypertarget{ux5f15ux304dux623bux3055ux308cux305fux6e2cux5b9aux5668ux3068ux3057ux3066ux8aadux3080}{%
\subsubsection{1.4.4
引き戻された測定器として読む}\label{ux5f15ux304dux623bux3055ux308cux305fux6e2cux5b9aux5668ux3068ux3057ux3066ux8aadux3080}}

ここまでの計算は、次のように読める。

実空間の測定器 \(dx\) は、実空間の変位を食べて \(x\)
成分を返す。ところが、円柱座標で計算したいときには、入力されるのは実空間の変位そのものではなく、パラメータ空間の小さな一歩

\[
\begin{pmatrix}
\Delta r\\
\Delta\theta\\
\Delta z
\end{pmatrix}
\]

である。

そこで、パラメータ空間の小さな一歩を食べたときに、実空間で \(dx\)
が測った値と同じ一次の値を返す測定器を作る。

その測定器が

\[
\begin{pmatrix}
\cos\theta & -r\sin\theta & 0
\end{pmatrix}
\]

であった。

測定器の数学では、この操作を\textbf{引き戻し}と呼ぶ。記号で書けば、

\[
\Phi^*(dx)
=
\cos\theta\,dr-r\sin\theta\,d\theta
\]

である。

ここで \(dr,d\theta,dz\) は、パラメータ空間の変位

\[
\begin{pmatrix}
\Delta r\\
\Delta\theta\\
\Delta z
\end{pmatrix}
\]

から、それぞれ第1成分、第2成分、第3成分を取り出す測定器である。したがって、

\[
\cos\theta\,dr-r\sin\theta\,d\theta
\]

は、パラメータ空間側の横ベクトル

\[
\begin{pmatrix}
\cos\theta & -r\sin\theta & 0
\end{pmatrix}
\]

を、測定器の記法で書いたものである。

つまり、

\[
\Phi^*(dx)
=
\begin{pmatrix}
\cos\theta & -r\sin\theta & 0
\end{pmatrix}
\]

と読んでよい。ただし右辺は、円柱座標のパラメータ空間の変位

\[
\begin{pmatrix}
\Delta r\\
\Delta\theta\\
\Delta z
\end{pmatrix}
\]

を食べる横ベクトルである。

\begin{quote}
\textbf{注}（ここでは有限の一歩から作った）

ここでは、既知の公式として
\(dx=\frac{\partial x}{\partial r}dr+\frac{\partial x}{\partial\theta}d\theta+\frac{\partial x}{\partial z}dz\)
をいきなり使ったのではない。

まず、パラメータ空間の小さな一歩を実空間へ写し、その像を実空間の \(dx\)
で測った。その一次の値と同じ値を返すように、パラメータ空間側の測定器を作ったのである。

第4章では、これと同じ考え方を有限区間、有限マス目、有限箱へ拡張する。そこで、区間の長さ、面積、体積を保つために測定器がどう作り替わるかを見る。
\end{quote}

\hypertarget{ux4e8cux3064ux306eux30c7ux30abux30ebux30c8ux5ea7ux6a19}{%
\subsubsection{1.4.5
二つのデカルト座標}\label{ux4e8cux3064ux306eux30c7ux30abux30ebux30c8ux5ea7ux6a19}}

ここで、いま扱っている二つの数の組------実空間の \((x,y,z)\)
と円柱座標のパラメータ空間
\((r,\theta,z)\)------の関係をはっきりさせておこう。

実空間の \((x,y,z)\)
は、通常のデカルト座標である。一方、\((r,\theta,z)\)
も、計算用紙の上では三つの数を並べた座標である。\(\theta\)
は物理的には角度だが、パラメータ空間では一つの数直線として扱う。

つまり、実空間もパラメータ空間も、計算の上ではそれぞれ「まっすぐな成分」を持つ。ただし、その二つの空間を結ぶ変換
\(\Phi\) が曲がっている。その曲がり方が、\(\cos\theta\) や
\(-r\sin\theta\) という係数として現れる。

実空間の中に曲がった座標軸が生えている、と考えることもできる。しかし本書では、なるべくそう考えない。むしろ、

\textbf{実空間とパラメータ空間という二つの場所があり、そのあいだに翻訳規則}\(\Phi\)\textbf{がある。}

と考える。

この見方を取ると、測定器の作り替えはかなり単純になる。実空間の測定器を、変換
\(\Phi\)
を通してパラメータ空間側へ引き戻す。すると、実空間で測った値と同じ一次の値を、パラメータ空間の変位から直接読めるようになる。

\begin{quote}
\textbf{注}（\(dr\), \(d\theta\), \(dz\) の行列表現）

パラメータ空間もまた、計算上は三つの成分を持つ空間である。したがって、\(dr\)
はパラメータ空間の第1成分を抜く測定器、\(d\theta\)
は第2成分を抜く測定器、\(dz\) は第3成分を抜く測定器である。

行列表現で書けば、\(dr=\begin{pmatrix}1&0&0\end{pmatrix}\)、\(d\theta=\begin{pmatrix}0&1&0\end{pmatrix}\)、\(dz=\begin{pmatrix}0&0&1\end{pmatrix}\)
である。ただし、これはパラメータ空間の変位に作用する測定器であり、実空間の
\((x,y,z)\) に作用する \(dx,dy,dz\) とは、載っている空間が違う。
\end{quote}

\hypertarget{ux7b2c4ux7ae0ux3078ux306eux4f0fux7dda}{%
\subsubsection{1.4.6
第4章への伏線}\label{ux7b2c4ux7ae0ux3078ux306eux4f0fux7dda}}

このように、元の世界の測定器を、別の目盛りで同じ値を返すように作り替える操作が、引き戻しである。

第1章では、\(dx\)
という最も小さい測定器について、その一例を見たに過ぎない。しかし第4章では、同じ発想を面積と体積へ拡張する。

そこでは、

\[
dx\wedge dy
\]

のような面積測定器や、

\[
dx\wedge dy\wedge dz
\]

のような体積測定器を、別の変数の世界へ引き戻す。

そのとき現れる係数が、円柱座標でおなじみの \(r\)
や、一般の変数変換で現れるヤコビアン \(J\) である。

つまり、第4章で出てくる

\[
\Phi^*(dx\wedge dy)=r\,dr\wedge d\theta
\]

や

\[
\Phi^*(dx\wedge dy\wedge dz)=J\,du\wedge dv\wedge dw
\]

は、いま見た

\[
\Phi^*(dx)
=
\cos\theta\,dr-r\sin\theta\,d\theta
\]

と同じ構造を持っている。

\begin{quote}
\textbf{注}（計量遅延の哲学）

ここで出てきた \(\cos\theta\) や \(-r\sin\theta\) は、のちに第6章で計量
\(g\) と関係して整理される。ただし、現時点では「変換 \(\Phi\)
によって測定器の成分が作り替わった」と読めば十分である。

実空間に曲がった図形を描いて、弧の長さを \(r\Delta\theta\)
と見積もる発想も大切である。しかし本書では、まずパラメータ空間の小さな一歩を写し、その像を測る、という代数的な手順を優先する。この手順が、後の面積・体積の計算で効いてくる。
\end{quote}

\begin{center}\rule{0.5\linewidth}{0.5pt}\end{center}

\begin{quote}
\textbf{【ここまでのチェックポイント】}

\begin{itemize}
\item
  \(dx\) という測定器自体は実空間の \(x\) 成分を抜き出す測定器である。
\item
  円柱座標で計算したいときは、実空間の測定器 \(dx\)
  を、パラメータ空間側の測定器へ作り替える。
\item
  パラメータ空間の小さな一歩を実空間へ写し、その像を \(dx\)
  で測ることで、パラメータ空間側の測定器の成分が分かる。
\item
  円柱座標では、\(\Phi^*(dx)=\cos\theta\,dr-r\sin\theta\,d\theta\)
  である。
\item
  第4章では、同じ考え方を面積測定器と体積測定器に拡張する。
\end{itemize}
\end{quote}

\begin{center}\rule{0.5\linewidth}{0.5pt}\end{center}

\hypertarget{ux5225ux306eux5ea7ux6a19ux3067ux66f8ux304fux6f14ux7fd2}{%
\subsection{§1.5
別の座標で書く------演習}\label{ux5225ux306eux5ea7ux6a19ux3067ux66f8ux304fux6f14ux7fd2}}

本章の終盤で、読者自身の手で「別の座標で書くこと」を一度、計算として体験してもらう。ここではまだ\textbf{座標系の取り替えを一般論として定義したわけではない}が、§1.4.3
で見たように、\textbf{同じ}\(dx\)\textbf{でも座標の取り方で行列表現が変わる}ことは既に述べた。この演習は、その感覚を手で確かめるためのものである。\textbf{いま座標系の取り替えの厳密な言語を全部そろえるより先に、なぜ計算をしておきたいのか}------それは、後の章で式が一気に忙しくなるときに、「成分が変わるのは当たり前」という腹落ちを持っておくためだ。

\hypertarget{ux6f14ux7fd2ux554fux984cux5186ux67f1ux5ea7ux6a19ux3067ux306eux6e2cux5b9a}{%
\subsubsection{1.5.1
演習問題：円柱座標での測定}\label{ux6f14ux7fd2ux554fux984cux5186ux67f1ux5ea7ux6a19ux3067ux306eux6e2cux5b9a}}

\textbf{問題設定}\\
前節で見た通り、円柱座標 \((r, \theta, z)\) における \(dx\)
の行列表示は以下のようになる：

\[
dx = \begin{pmatrix} \cos\theta & -r\sin\theta & 0 \end{pmatrix}
\]

（上の行は §1.4.3
で作った、パラメータ空間の変位に作用する横ベクトルであり、§1.1.3
のデカルト表示 \(\begin{pmatrix}1&0&0\end{pmatrix}\)
とは\textbf{右に立てる縦ベクトルの意味が違う}。）

このとき、次の問いに答えよ。

\textbf{問1}\\
円柱座標空間内の点 \(P(r=2, \theta=\pi/6, z=0)\)
に質点があるとする。この質点が、\(r\) 方向に \(0.1\)
だけ微小変位した。\textbf{この節では記号}\(\mathbf{v}\)\textbf{を、}\(r,\theta,z\)\textbf{のパラメータの変化量を並べた}\(\begin{pmatrix}\Delta r\\\Delta\theta\\\Delta z\end{pmatrix}\)\textbf{とする}（§1.0--§1.2
のデカルト列
\(\begin{pmatrix}\Delta x\\\Delta y\\\Delta z\end{pmatrix}\)
とは別の並べ方である）。この変位ベクトル \(\mathbf{v}\)
を、その成分で表せ。

\textbf{問2}\\
問1の変位ベクトル \(\mathbf{v}\) に対して、点 \(P\) における行列 \(dx\)
を作用させよ。（\(\cos(\pi/6) = \sqrt{3}/2\) を用いよ）

\textbf{問3}\\
問2で計算した結果 \(dx(\mathbf{v})\)
は、物理的に何を意味しているか説明せよ。

\hypertarget{ux89e3ux7b54ux3068ux89e3ux8aac}{%
\subsubsection{1.5.2 解答と解説}\label{ux89e3ux7b54ux3068ux89e3ux8aac}}

\textbf{問1の解答}\\
変位は \(r\) 方向に \(0.1\)、\(\theta\) 方向と \(z\)
方向にはゼロなので、\textbf{円柱座標のパラメータの変化量}を縦に並べたベクトルとして次のように書ける：

\[
\mathbf{v} = \begin{pmatrix} 0.1 \\ 0 \\ 0 \end{pmatrix}
\]

\textbf{問2の解答}\\
点 \(P\) における行列 \(dx\) は、\(\theta = \pi/6,\ r=2\)
を代入して得られる \(1\times 3\) 行ベクトルとして

\[
dx
=
\begin{pmatrix}
\cos(\frac{\pi}{6}) & -2\sin(\frac{\pi}{6}) & 0
\end{pmatrix}
=
\begin{pmatrix}
\frac{\sqrt{3}}{2} & -2 \times \frac{1}{2} & 0
\end{pmatrix}
=
\begin{pmatrix}
\frac{\sqrt{3}}{2} & -1 & 0
\end{pmatrix}
\]

と書く（円柱座標では \(dx\) の成分は \(r,\theta\)
に依存するので、ここでは点 \(P\)
の座標で評価したあとを指す。\(\mathbf{v}\) への作用の記号
\(dx(\mathbf{v})\) 自体は §1.1.3 以来と同じである）。

重要なのは、第2成分は一般に \(-r\sin\theta\) という係数であり、\(r\)
のおかげで\textbf{長さの次元を持つ}ことだ。\(r=2,\;\theta=\pi/6\)
と代入するとその係数の値として \(-1\)
が見えるが、\textbf{数値の}\(-1\)\textbf{そのものに次元が付いているのではない}。行列の成分としての
\(-r\sin\theta\)
が長さの次元を運んでいる、と読むのが正確である。これが次元解析の鍵となる。

これを変位ベクトル \(\mathbf{v}\) に作用させると：

\[
dx(\mathbf{v})
=
\begin{pmatrix}
\frac{\sqrt{3}}{2} & -1 & 0
\end{pmatrix}
\begin{pmatrix}
0.1\\
0\\
0
\end{pmatrix}
=
0.1 \times \frac{\sqrt{3}}{2}
\approx 0.0866
\]

\textbf{問3の解答（物理的意味）}\\
この結果は、\textbf{「円柱座標において外側（}\(r\)\textbf{方向）へ}\(0.1\)\textbf{だけ進むという運動は、実空間の}\(x\)\textbf{軸方向から見ると約}\(0.0866\)\textbf{の前進に相当する」}という物理的事実を表している。

\begin{quote}
\textbf{注} （「約」と厳密一致）\\
この変位では \(\Delta\theta=\Delta z=0\) なので、\(x=r\cos\theta\) より
\(\Delta x=(\Delta r)\cos\theta\)
が\textbf{厳密に}成り立つ。ここで「約」と書いたのは、小数表示 \(0.0866\)
に丸めたためである。
\end{quote}

\hypertarget{ux89d2ux5ea6ux65b9ux5411ux306bux52d5ux304bux3057ux305fux5834ux5408}{%
\subsubsection{1.5.3
角度方向に動かした場合}\label{ux89d2ux5ea6ux65b9ux5411ux306bux52d5ux304bux3057ux305fux5834ux5408}}

同じ点 \(P(r=2,\theta=\pi/6,z=0)\) で、今度は \(\theta\) 方向に \(0.1\)
だけ動く場合も見ておこう。

このとき、円柱座標のパラメータ空間での変位は

\[
\mathbf{v}
=
\begin{pmatrix}
0\\
0.1\\
0
\end{pmatrix}
\]

である。点 \(P\) における \(dx\) は先ほどと同じく

\[
dx
=
\begin{pmatrix}
\frac{\sqrt{3}}{2} & -1 & 0
\end{pmatrix}
\]

であるから、

\[
dx(\mathbf{v})
=
\begin{pmatrix}
\frac{\sqrt{3}}{2} & -1 & 0
\end{pmatrix}
\begin{pmatrix}
0\\
0.1\\
0
\end{pmatrix}
=
-0.1
\]

となる。

この場合、\textbf{入力ベクトルの第2成分} \(0.1\)
は角度（無次元）だが、\textbf{行列の第2列に対応する係数}
\(-r\sin\theta\) が長さの次元を運ぶため、積としての出力 \(-0.1\)
は長さの次元を持つ。

ここに\textbf{一次形式}の重要な性質がある。\(dx\) の行列成分自体が \(r\)
を含む関数であり、入力が座標成分（次元がバラバラでも）であっても、出力は常に「\(x\)
方向の長さ」という正しい物理的次元を持つように自動調整される。\textbf{一次形式は、座標系の歪みを吸収し、物理的に意味のある測定値を出力する測定器}なのである。

\hypertarget{ux672cux7ae0ux306eux307eux3068ux3081ux3068ux6b21ux7ae0ux3078ux306eux5c55ux671b}{%
\subsection{§1.6
本章のまとめと次章への展望}\label{ux672cux7ae0ux306eux307eux3068ux3081ux3068ux6b21ux7ae0ux3078ux306eux5c55ux671b}}

\begin{quote}
\textbf{【ここまでのチェックポイント — 第1章全体】}

\begin{itemize}
\item
  第1章の主役は「\(dx\) を行列・一次形式として読み、\(\int f\,dx\)
  を作用の極限として読む」ことである。説明の主線は \(x\)
  方向に寄せたが、デカルトの測定器 \(dy\), \(dz\) は §1.2.3 で \(dx\)
  と同様に定義済みである。
\item
  \(df\)
  も横ベクトルとして統一し、多次元への拡張は成分の増加として素直に繋がる。
\item
  次章以降は、これらの測定器の\textbf{ウェッジ積}、外微分、ホッジスターへと拡張する。
\end{itemize}
\end{quote}

本章では\textbf{積分と}\(df\)\textbf{の具体例}の主線として \(x\)
方向に寄せ、\(y,z\)
は多くの場面で「断面」として抑えてきた。とはいえ、実空間での一次形式は
\(dx\)\textbf{,}\(dy\)\textbf{,}\(dz\)\textbf{の三つがそろっている}（§1.2.3）。
次章では、この三つを組み合わせて\textbf{面積計・体積計（}\(2\)\textbf{-form,}\(3\)\textbf{-form）}を作る\textbf{ウェッジ積（外積}\(\wedge\)\textbf{）}を導入し、ベクトルに作用してスカラーを返す型にとどまらない\textbf{高次の微分形式}へと進む。曲線・曲面・領域に沿った集計（線積分・面積分・体積分）は、それらの測定器を揃えた後の章で扱う。
その後の部では、\textbf{微分演算子}\(\mathrm{d}\)、\textbf{ホッジスター演算子}\(\ast\)、ベクトル解析の
\(\mathrm{grad}\)（\(\nabla\)）、\(\mathrm{rot}\)（\(\nabla\times\)、\(\mathrm{curl}\)）、\(\mathrm{div}\)（\(\nabla\cdot\)）、ストークスの定理、マクスウェル方程式や流体力学の基礎方程式へと進む。3次元での微分形式とベクトル場の対応は、ホッジスターで整理できる。

3次元ユークリッド空間の真の姿を解き明かす旅へ、いざ出発しよう。

\newpage

\hypertarget{ux7b2c2ux7ae0ux9762ux7a4dux3068ux306fux4f55ux304b-ux5e73ux884cux591aux9762ux4f53ux306bux6f5cux3080ux7b26ux53f7ux306eux30ebux30fcux30eb}{%
\chapter{第2章：面積とは何か ------
平行多面体に潜む、符号のルール}\label{ux7b2c2ux7ae0ux9762ux7a4dux3068ux306fux4f55ux304b-ux5e73ux884cux591aux9762ux4f53ux306bux6f5cux3080ux7b26ux53f7ux306eux30ebux30fcux30eb}}

\hypertarget{ux6e2cux5b9aux5668ux3068ux9762ux7a4dux4f53ux7a4dux305dux3057ux3066ux9577ux3055}{%
\subsection{§2.0
測定器と面積・体積、そして長さ}\label{ux6e2cux5b9aux5668ux3068ux9762ux7a4dux4f53ux7a4dux305dux3057ux3066ux9577ux3055}}

　第1章では、\(dx, dy, dz\)
を「ベクトルを1つ食べてスカラーを返す横ベクトル（\(1\)-form）」として定義し、\(\int f\,dx\)
を行列作用の極限として読み直した。本章ではこれを拡張し、\textbf{2つのベクトルを食べる面積測定器}（\(2\)-form）と\textbf{3つのベクトルを食べる体積測定器}（\(3\)-form）を、同じ部品
\(dx, dy, dz\) から発見する。

\hypertarget{ux5c0fux5b66ux6821ux306eux9762ux7a4dux306eux9650ux754c}{%
\subsection{§2.1
小学校の面積の限界}\label{ux5c0fux5b66ux6821ux306eux9762ux7a4dux306eux9650ux754c}}

多くの読者は「積分とはつまり面積を求める計算だ」ということを知っているだろう。
ところで、面積とはなんだろうか。改めて考えてみると、小学校以来、改めて定式化したことがないのではないだろうか。

ではその小学校での定式化とは何だったか。それは\textbf{「}\(1 \times 1\)\textbf{の正方形（タイル）が、その図形の中にいくつ敷き詰められるか」}という、極めて素朴な直感に基づいていただろう。

2次元の紙の上（\(xy\)平面）に描かれた図形であれば、この「正方形の敷き詰め」はうまく機能する。例えば、2つの2次元ベクトル
\(\mathbf{v}_1 = (x_1, y_1)\) と \(\mathbf{v}_2 = (x_2, y_2)\)
が張る平行四辺形の面積 \(S\) は次のように書ける：

\[S = \sqrt{(x_1 y_2 - x_2 y_1)^2}\]

しかし、我々が生きているのは3次元空間である。3次元空間の中で「斜めに傾いた平面」の面積を計算しようとしたとき、事態は少し大変である。

多くの読者は、高校数学で3次元空間のベクトル
\(\mathbf{v}_1 = (x_1, y_1, z_1)\) と \(\mathbf{v}_2 = (x_2, y_2, z_2)\)
が張る平行四辺形の面積を求めさせられた経験があるだろう。高校の記法を使って、次のように書ける：
\[
S = \sqrt{|\mathbf{v}_1|^2 |\mathbf{v}_2|^2 - (\mathbf{v}_1 \cdot \mathbf{v}_2)^2}
\]

成分表示を考えると、ノートの幅を埋め尽くすような計算の末に、ようやく次の式にたどり着く。

\[S = \sqrt{(y_1z_2 - z_1y_2)^2 + (z_1x_2 - x_1z_2)^2 + (x_1y_2 - x_2y_1)^2}\]

その過程はなかなか長大であり、人生で何度もやりたくないほどには大変である。
なぜそんなに大変になってしまうのか？
それは、我々が\textbf{面積というものを、はじめに「2次元世界の中」で定義してきたから}だ。\textbf{3次元空間では}、平面の「傾き」と、そこに載る2本のベクトルの関係とが絡み、\textbf{角度の取り扱い}まで含めて式が膨らんでしまう。

\begin{quote}
\textbf{注}（未定義の語の先取り）ここでは「角度」や「内積」という言葉を使っているが、この段階では小学校以来の直感的な意味で十分である。形式的な内積の定義は第6章で与えるが、ここでの議論の理解には不要だ。
\end{quote}

ここで我々は、がらりと考え方を変えよう。\textbf{小学校以来の「面積の定義」を一旦棚に上げ}、\textbf{「面積を出力する代数的な測定装置」}をゼロから設計し直すのである。

\begin{center}\rule{0.5\linewidth}{0.5pt}\end{center}

\hypertarget{ux9762ux7a4dux6e2cux5b9aux5668ux304cux6e80ux305fux3059ux3079ux304d3ux3064ux306eux30ebux30fcux30eb}{%
\subsection{§2.2
面積測定器が満たすべき「3つのルール」}\label{ux9762ux7a4dux6e2cux5b9aux5668ux304cux6e80ux305fux3059ux3079ux304d3ux3064ux306eux30ebux30fcux30eb}}

設計の第一歩として、「\(1\times 1\)
の正方形だけ」を足掛かりにするのではなく、\textbf{「2つのベクトルが張る平行四辺形」}を基本に据えることにする。

\begin{quote}
\textbf{注}（平行四辺形から始める理由）いままで「面積＝正方形の敷き詰め」と数えていた読者には、\textbf{なぜいきなり平行四辺形なのか}という動機づけが見えにくいかもしれない。平行四辺形を基本に据えると、辺が直交する、という条件が緩和されることになり、これから組み立てるルールがすっきりする。正方形はその\textbf{特別な場合}として自然に含まれる。また\textbf{積分を高次元へと拡張する}ときにも、こちらのほうが扱いやすい。
\end{quote}

上記を言い換えると、我々が欲しいのは、2つのベクトル \(\mathbf{v}_1\) と
\(\mathbf{v}_2\)
を入力（インプット）として食べると、それらが張る平行四辺形の面積 \(S\)
を出力（アウトプット）として吐き出す、魔法のブラックボックス
\(S(\mathbf{v}_1, \mathbf{v}_2)\) である。

やや天下り的であるが、このブラックボックスが「面積測定器」としてまともに機能するためには、次の3つのルールを満たせばよい。原点から伸びる2本の矢印を隣り合う辺とする平行四辺形を思い浮かべながら確認してほしい。

\hypertarget{ux30ebux30fcux30eb1ux7247ux65b9ux304cux4f38ux3073ux308cux3070ux9762ux7a4dux3082ux4f38ux3073ux308bux6bd4ux4f8bux95a2ux4fc2}{%
\subsubsection{2.2.1
ルール1：片方が伸びれば、面積も伸びる（比例関係）}\label{ux30ebux30fcux30eb1ux7247ux65b9ux304cux4f38ux3073ux308cux3070ux9762ux7a4dux3082ux4f38ux3073ux308bux6bd4ux4f8bux95a2ux4fc2}}

もしベクトル \(\mathbf{v}_1\)
の長さが2倍になれば、作られる平行四辺形の面積も当然2倍になるはずだ。
\[S(2\mathbf{v}_1, \mathbf{v}_2) = 2 S(\mathbf{v}_1, \mathbf{v}_2)\]
同様に、2つのベクトルを足し合わせたものを入力すれば、面積も足し算に分解できる：
\[S(\mathbf{v}_1 + \mathbf{u}, \mathbf{v}_2) = S(\mathbf{v}_1, \mathbf{v}_2) + S(\mathbf{u}, \mathbf{v}_2)\]
この2つの性質を合わせて\textbf{線形性}と呼ぶ。第2引数 \(\mathbf{v}_2\)
についても、同様に「スカラー倍」と「和に関する分配」が成り立つ------すなわち
\(S\) は\textbf{両方の引数について線形}（\textbf{双線形}）である。

\hypertarget{ux30ebux30fcux30eb2ux91cdux306aux308cux3070ux30daux30c1ux30e3ux30f3ux30b3ux3067ux30bcux30edux4ea4ux4ee3ux6027}{%
\subsubsection{2.2.2
ルール2：重なれば、ペチャンコでゼロ（交代性）}\label{ux30ebux30fcux30eb2ux91cdux306aux308cux3070ux30daux30c1ux30e3ux30f3ux30b3ux3067ux30bcux30edux4ea4ux4ee3ux6027}}

これが最も重要だ。もし \(\mathbf{v}_1\) と \(\mathbf{v}_2\)
が全く同じベクトルだったらどうなるか。平行四辺形は潰れてただの「線」になってしまい、面積はゼロになる。
\[S(\mathbf{v}_1, \mathbf{v}_1) = 0\]

実は、この「同じならゼロ」というルールから、非常に奇妙な性質が導かれる。ルール1（線形性）を使って、\(\mathbf{v}_1 + \mathbf{v}_2\)
を両方の引数に入れてみよう。「同じならゼロ」だから：

\[S(\mathbf{v}_1 + \mathbf{v}_2,\; \mathbf{v}_1 + \mathbf{v}_2) = 0\]

左辺を線形性で展開すると：

\[S(\mathbf{v}_1, \mathbf{v}_1) + S(\mathbf{v}_1, \mathbf{v}_2) + S(\mathbf{v}_2, \mathbf{v}_1) + S(\mathbf{v}_2, \mathbf{v}_2) = 0\]

第1項と第4項はルール2で消えるから、残るのは
\(S(\mathbf{v}_1, \mathbf{v}_2) + S(\mathbf{v}_2, \mathbf{v}_1) = 0\)。すなわち：

\[S(\mathbf{v}_1, \mathbf{v}_2) = - S(\mathbf{v}_2, \mathbf{v}_1)\]

入力するベクトルの順番を入れ替えると、\textbf{面積の符号がマイナスに反転してしまう}のだ。

\begin{quote}
\textbf{注}（面積の符号1）「面積がマイナスになる？
それは困る」と驚くかもしれない。しかし、物理学者は自然界の「向き」を無視しない。電流の向き、磁場の向き------すべてベクトルの「向き」が本質だ。だから「面積にも向き（向き付き面積）がある」という発想は、我々物理学者にはむしろ自然なのだ。この概念は後で曲面の「裏表」を区別したり、空間を貫く磁束や流体の流れを計算する際に絶対に不可欠なものとなる。今のところはこだわりすぎずに進んでほしい。
\end{quote}

\begin{quote}
\textbf{注}（面積の符号2）1変数でも、\(f(x)<0\) の区間では
\(\displaystyle\int_a^b f(x)\,dx<0\)
となり得る------符号付き「面積」はそこで既に現れている。向き付き面積は、その考え方を2次元に持ち上げたものだ。
\#\#\#\# 2.2.3 ルール3：基準の決定（規格化）
最後に、「どのくらいの大きさを1とするか」を決めなければならない。ここでは素直に、\(x\)方向の\textbf{基底ベクトル}（標準基底の一つ）\(\hat{e}_x\)
と \(y\)方向の基底ベクトル \(\hat{e}_y\) が作る正方形の面積を \(1\)
としよう。本書の慣例どおり\textbf{縦ベクトル}で書けば
\[\hat{e}_x = \begin{pmatrix} 1 \\ 0 \\ 0 \end{pmatrix}, \quad \hat{e}_y = \begin{pmatrix} 0 \\ 1 \\ 0 \end{pmatrix}\]
である。
\end{quote}

\[S(\hat{e}_x, \hat{e}_y) = 1\]

\begin{quote}
\textbf{注}（基底ベクトル（標準基底））デカルト座標の各軸方向に長さ
\(1\)
の足を伸ばしたものを、基底ベクトル（標準基底）と呼ぶ、といって良いだろう。より進んだ定義は線形代数に任せる。
\end{quote}

\begin{quote}
\textbf{【ここまでのチェックポイント】}

\begin{itemize}
\item
  面積測定器は「2つのベクトルを食べてスカラーを返す」装置。3つのルール------線形性・交代性・規格化------を満たさなければならない。
\item
  交代性（同じベクトルを2つ入れるとゼロ）から
  \(S(\mathbf{v}_1,\mathbf{v}_2) = -S(\mathbf{v}_2,\mathbf{v}_1)\)
  が導かれ、面積は\textbf{符号付き}になる。
\item
  次節では、この3ルールを満たす代数的な装置を具体的に探す。
\end{itemize}
\end{quote}

\begin{center}\rule{0.5\linewidth}{0.5pt}\end{center}

\hypertarget{ux9762ux7a4dux6e2cux5b9aux5668ux306fux53cdux5bfeux79f0ux884cux5217ux3067ux3042ux308b}{%
\subsection{§2.3
面積測定器は「反対称行列」である}\label{ux9762ux7a4dux6e2cux5b9aux5668ux306fux53cdux5bfeux79f0ux884cux5217ux3067ux3042ux308b}}

我々は今、この「平行四辺形」という\textbf{図形のイメージ}を頭に描きながら3つのルールを設定した。しかし、一度ルールを決めてしまえば、もはや図形を思い浮かべる必要はない。\textbf{「この3つのルールを満たす代数的な計算式」を探せば、それが面積の定義になるからだ。}

簡単のために、まずは \(z\)成分が常にゼロのベクトル、つまり
\(xy\)平面上にある2つのベクトルで考えよう。

\[
\mathbf{v}_1 = \begin{pmatrix} x_1 \\ y_1 \\ 0 \end{pmatrix}, \quad \mathbf{v}_2 = \begin{pmatrix} x_2 \\ y_2 \\ 0 \end{pmatrix}
\]

結論から言うと、この
\(xy\)\textbf{平面上のベクトルという制約のもとで}、先ほどの3ルールを完璧に満たす量を出力する測定器は、線形代数の世界ではただ一つの自然な形に落ち着く。

それは、真ん中に特定の \(3 \times 3\)\textbf{行列}
を挟んだ、次のような「\textbf{横ベクトル} \(\times\) \textbf{行列}
\(\times\) \textbf{縦ベクトル}」の計算である：

\[
\begin{pmatrix} x_1 & y_1 & 0 \end{pmatrix}
\begin{pmatrix}
  0 & 1 & 0 \\
  -1 & 0 & 0 \\
  0 & 0 & 0
\end{pmatrix}
\begin{pmatrix} x_2 \\ y_2 \\ 0 \end{pmatrix}
\]

実際に \(\mathbf{v}_1\) と \(\mathbf{v}_2\) を代入して計算してみよう：

\[
\begin{pmatrix} x_1 & y_1 & 0 \end{pmatrix}
\begin{pmatrix}
  0 & 1 & 0 \\
  -1 & 0 & 0 \\
  0 & 0 & 0
\end{pmatrix}
\begin{pmatrix} x_2 \\ y_2 \\ 0 \end{pmatrix}
= \begin{pmatrix} -y_1 & x_1 & 0 \end{pmatrix}
\begin{pmatrix} x_2 \\ y_2 \\ 0 \end{pmatrix}
= x_1 y_2 - x_2 y_1
\]

ここで得られた \(x_1 y_2 - x_2 y_1\)
は、「2次元の面積公式」の中身そのものである。

\begin{quote}
\textbf{注}（反対称行列と計量）我々は今、\(\mathbf{v}_1^T M \mathbf{v}_2\)
（横ベクトル \(\times\) 行列 \(\times\)
縦ベクトル）という書き方をした。より進んだ立場、すなわち一般のテンソル解析や一般の多様体論の立場から見ると、計量を明示せずにベクトルを横に寝かせてしまうこの書き方は、少々お行儀が悪い。しかし本書は当面、標準デカルト座標に固定して進む。この範囲では問題ないので、ここではこれで押し進めよう。第5章以降で計量を導入した後に、改めてこの点を振り返る。
\end{quote}

\hypertarget{ux6e2cux5b9aux5668ux306eux9ad8ux6b21ux5143ux5316}{%
\subsubsection{2.3.1
測定器の高次元化}\label{ux6e2cux5b9aux5668ux306eux9ad8ux6b21ux5143ux5316}}

繰り返しになるが、我々はここで、次の立場をはっきりさせる。
小学校で数えていた\textbf{タイルの枚数}に相当する素朴な直感を先に置くのではなく、\textbf{いま定めたこの}\(3 \times 3\)\textbf{反対称行列に2本のベクトルを食わせて得られる量}を、\textbf{この測定器（}\(xy\)\textbf{平面専用の面積測定器）が返す面積の成分}------\(dx \wedge dy\)
に沿った読み取り値------だと定義するのである。三次元空間における向き付き面積の\textbf{すべての情報}を一つの数で尽くすわけではなく、あくまで\textbf{一つの眼鏡が切り出したスカラー}にすぎない。全体像は
§2.4.6 で3種の測定器をそろえて述べる。

前章の
\(1\)-form（一次形式）は、1つのベクトルを捕まえる\textbf{「}\(1 \times 3\)\textbf{行列」}だった。
今目の前にある行列は、\textbf{2つのベクトルを同時に受け取り面積を計算}する\textbf{「}\(3 \times 3\)\textbf{反対称行列」}へと進化したのである。

この行列（あるいは演算子）を、\(2\)-formと呼ぶことにする。

\begin{quote}
\textbf{注}（二次形式と \(2\)-form）線形代数の「二次形式（quadratic
form）」は、多くの場合「二次の多項式」を指す。ここでの \(2\)-form
は微分形式の用語であり、上の quadratic form とは別物である。
\end{quote}

さて、読者はこう思っているはずだ。
「なるほど、行列で面積が出るのは分かった。でも、この行列は第3行・第3列がゼロだから、\(z\)成分を完全に無視しているじゃないか。これではあの『3次元の大変な計算』の解決になっていないぞ」と。

\textbf{まったくその通りである。} 今我々が手にした
\(\begin{pmatrix} 0 & 1 & 0 \\ -1 & 0 & 0 \\ 0 & 0 & 0 \end{pmatrix}\)
という行列は、3次元空間に3種類ある基本測定器の一つ------\(xy\)\textbf{平面専用の面積測定器}------に過ぎない。
言い換えれば、これは\(xy\)\textbf{平面「だけ」を見る「偏った眼鏡」}のようなものだ。空間には、\(yz\)平面を見る眼鏡や、\(zx\)平面を見る眼鏡も存在する。

\begin{center}\rule{0.5\linewidth}{0.5pt}\end{center}

\begin{quote}
\textbf{【ここまでのチェックポイント】}

\begin{itemize}
\item
  面積測定器（\(2\)-form）とは「2つのベクトルの絡み合い」を測る代数的な装置であり、測定対象である平行四辺形そのものではない。
\item
  面積測定器の正体は「反対称行列」であり、物理的ルールが行列の形を完全に決める。
\item
  3次元空間には3種類の「偏った眼鏡」（基底\(2\)-form）が存在する。
\end{itemize}
\end{quote}

\begin{center}\rule{0.5\linewidth}{0.5pt}\end{center}

\hypertarget{ux9762ux7a4dux6e2cux5b9aux5668ux306eux5185ux90e8ux69cbux9020}{%
\subsection{§2.4
面積測定器の内部構造}\label{ux9762ux7a4dux6e2cux5b9aux5668ux306eux5185ux90e8ux69cbux9020}}

\hypertarget{ux3064ux306eux5165ux529bux3092ux3055ux3070ux304fux88c5ux7f6eux306eux5f62ux3092ux6c7aux3081ux308b}{%
\subsubsection{2.4.1
「2つの入力をさばく装置」の形を決める}\label{ux3064ux306eux5165ux529bux3092ux3055ux3070ux304fux88c5ux7f6eux306eux5f62ux3092ux6c7aux3081ux308b}}

§2.3
で、面積測定器（\(2\)-form）の正体は真ん中に挟まった「\(3 \times 3\)
の反対称行列」であると述べた。しかし、あの行列は筆者から天下り的に与えられたものだった。

我々の次なる目標はこうだ：
\textbf{「第1章で導入したシンプルな物差し}\(dx, dy, dz\)\textbf{を部品として組み合わせて、あの複雑な行列をシステマティックに作り出せないだろうか？」}
こうすれば、いまは \(xy\) 平面専用の面積測定器しかないが、\(yz\)
平面用・\(zx\)
平面用の面積測定器も同じ部品から同じ手順で量産できるようになる。

これを実現するために、まずは「2つのベクトルを食べてスカラーを吐き出す装置」の外枠がどうなるべきかを考えよう。

第1章で、ベクトルを1つ食べる装置（\(1\)-form）は「横ベクトル」で表現できることを見た（横
\(\times\) 縦 \(=\) スカラー）。 では、\(\mathbf{v}_1\) と
\(\mathbf{v}_2\) という2つのベクトルを同時に食べて、しかも §2.2
のルール1（比例関係・線形性）を両方に対して保つ装置はどう書けるだろうか？

それは、2つのベクトルの間に真四角の表（行列）を挟み込む形、すなわち
\textbf{「横}\(\times\)\textbf{行列}\(\times\)\textbf{縦（}\(\mathbf{v}_1^T M \mathbf{v}_2\)\textbf{）」}
の形が自然に導かれる。これは線形代数の要請であるとして認めよう。我々が探すべきは、この真ん中に挟まる行列
\(M\) の中身だ。

\hypertarget{ux7e2eux7d04ux884cux5217ux3092ux6d88ux3057ux53bbux308bux64cdux4f5c}{%
\subsubsection{2.4.2
縮約------行列を消し去る操作}\label{ux7e2eux7d04ux884cux5217ux3092ux6d88ux3057ux53bbux308bux64cdux4f5c}}

ここで、先に導いた「横 \(\times\) 行列 \(\times\)
縦」という形に、名前を与えておこう。

\(\mathbf{v}_1^T M \mathbf{v}_2\) を成分で書き下すと、

\[
\mathbf{v}_1^T M \mathbf{v}_2 = \sum_{i=1}^3 \sum_{j=1}^3 v_{1i}\, M_{ij}\, v_{2j}
\]

である。左辺には行列 \(M\)（\(3 \times 3 = 9\)
個の数字）と2本のベクトル（\(3+3=6\)
個の数字）があったが、右辺はただ1つのスカラーだ。\(i\) と \(j\)
という添字が和によって「消し去られ」、15個の数字が1個の数字に縮まった。これは、\textbf{縮約}（contraction）と呼ばれる操作の一例である。

\textbf{狭い意味での縮約}は、「1組の共通添字について和をとる」操作である。もっとも単純な例は第1章の
\(1\)-form だ。\(dx(\mathbf{v}) = \sum_i (dx)_i v_i\) は、添字 \(i\)
について1回だけ縮約してスカラーを得ている。行列とベクトルの積
\(\sum_j M_{ij} v_j\) も同様に、添字 \(j\)
について1回縮約してベクトル（添字 \(i\) が残る）を得る操作だ。

\textbf{広い意味での縮約}は、これを複数回重ねたものである。いま見た
\(\mathbf{v}_1^T M \mathbf{v}_2 = \sum_i \sum_j v_{1i} M_{ij} v_{2j}\)
は、添字 \(i\) と \(j\)
の\textbf{両方}で縮約することで、行列と2本のベクトルからスカラーを作り出す------これは\textbf{2回の縮約}である。本章の後半（§2.5）では、3本のベクトルとレヴィ=チヴィタ記号
\(\epsilon_{ijk}\)
の縮約として行列式が現れる。付録Aではその全過程を追う。

\begin{quote}
\textbf{注}（縮約の感覚）添字が消える------これが縮約の感覚である。\(3 \times 3\)
行列は2つの添字（行と列）を持ち、ベクトルは1つの添字を持つ。縮約では、ベクトルの添字と行列の添字を「同じ文字」にして和をとることで、その添字が式から姿を消す。添字が2つ消えれば、残る添字はゼロ------すなわちスカラーになる。この「添字を消す」操作が、本書のあらゆる計算の根底にある。
\end{quote}

\hypertarget{ux6b32ux3057ux3044ux7d50ux679cux304bux3089ux30deux30b9ux76eeux3092ux57cbux3081ux308b}{%
\subsubsection{2.4.3
欲しい結果からマス目を埋める}\label{ux6b32ux3057ux3044ux7d50ux679cux304bux3089ux30deux30b9ux76eeux3092ux57cbux3081ux308b}}

闇雲に探す必要はない。我々はすでに、出力されるべき\textbf{「答え」}を知っている。
\(xy\) 平面の面積公式だ。

\[\text{欲しい出力} = x_1 y_2 - x_2 y_1\]

この出力を吐き出すためには、行列 \(M\)
のマス目にどんな数字を配置すればよいか？
読者も一緒に、\(\mathbf{v}_1^T M \mathbf{v}_2\)
の展開式を想像しながらマス目を埋めるパズルをしてほしい。

\[
\begin{pmatrix} x_1 & y_1 & z_1 \end{pmatrix}
\begin{pmatrix}
  ? & ? & ? \\
  ? & ? & ? \\
  ? & ? & ?
\end{pmatrix}
\begin{pmatrix} x_2 \\ y_2 \\ z_2 \end{pmatrix}
\]

\begin{enumerate}
\def\labelenumi{\arabic{enumi}.}
\tightlist
\item
  \(x_1 y_2\)\textbf{を作る：} \(\mathbf{v}_1\) の \(x\)
  成分（\(x_1\)）と \(\mathbf{v}_2\) の \(y\)
  成分（\(y_2\)）が掛け合わされる場所、すなわち「1行2列目」には \(+1\)
  を置けばよい。
\item
  \(-x_2 y_1\)\textbf{を作る：} \(\mathbf{v}_1\) の \(y\)
  成分（\(y_1\)）と \(\mathbf{v}_2\) の \(x\)
  成分（\(x_2\)）が掛け合わされる場所、すなわち「2行1列目」には \(-1\)
  を置けばよい。
\item
  \textbf{それ以外：} 面積の式には \(x_1 x_2\) や \(z\)
  成分などは一切登場しない。したがって、残りのマス目はすべて \(0\)
  である。
\end{enumerate}

こうして§2.3
で天下り的に登場したあの反対称行列が以下の形に一意に定まる：

\[
M = \begin{pmatrix}
  0 & 1 & 0 \\
  -1 & 0 & 0 \\
  0 & 0 & 0
\end{pmatrix}
\]

\hypertarget{ux30c6ux30f3ux30bdux30ebux7a4dux306eux5c0eux5165}{%
\subsubsection{2.4.4
テンソル積の導入}\label{ux30c6ux30f3ux30bdux30ebux7a4dux306eux5c0eux5165}}

さて、この面積測定器の行列 \(M\)
をよく見ると、2つの部品の「引き算」でできていることがわかる。

\[
\begin{pmatrix} 0 & 1 & 0 \\ -1 & 0 & 0 \\ 0 & 0 & 0 \end{pmatrix}
= \begin{pmatrix} 0 & 1 & 0 \\ 0 & 0 & 0 \\ 0 & 0 & 0 \end{pmatrix} - 
 \begin{pmatrix} 0 & 0 & 0 \\ 1 & 0 & 0 \\ 0 & 0 & 0 \end{pmatrix}
\]

右辺の第1項は「1行2列目だけが \(1\) の行列」であり、これは
\(\mathbf{v}_1\) から \(x\) 成分を、\(\mathbf{v}_2\) から \(y\)
成分を抽出して掛ける（\(x_1 y_2\)）という演算を表している。

ここで第1章の物差しを思い出そう。\(x\) 成分を抽出する物差しは
\(dx\)、\(y\) 成分を抽出する物差しは \(dy\) だった。
この2つの物差しを組み合わせて「\(\mathbf{v}_1\) は \(dx\)
で測り、\(\mathbf{v}_2\) は \(dy\)
で測って掛ける」という新しい演算子を作り、これを \textbf{テンソル積}
と呼び、\(dx \otimes dy\) と書くことにしよう。

\[(dx \otimes dy)(\mathbf{v}_1, \mathbf{v}_2) := dx(\mathbf{v}_1)\,dy(\mathbf{v}_2) = x_1 y_2\]

テンソル積 \(dx \otimes dy\)
というといかにも難しそうな響きだが、恐れることはない。行列表現の世界では、これは単に\textbf{「1行2列目のマス目だけを}\(1\)\textbf{にする」}スイッチのオンオフのような操作に過ぎない。

\begin{quote}
\textbf{注}（テンソル積の行列表現）数式が好きな読者のために書いておくと、\(dx = \begin{pmatrix} 1 & 0 & 0 \end{pmatrix}\)、\(dy = \begin{pmatrix} 0 & 1 & 0 \end{pmatrix}\)
としたとき、「前のベクトルを転置して縦にし、後ろの横ベクトルと隣り合わせに置き（\(dx^T dy\)）、縦と横の成分を総当たりで掛け合わせて行列のマス目に配置する」という操作で\(dx \otimes dy\)の行列表現が機械的に得られる。おそらく手元の線形代数の教科書にも書いてあるだろう。
\#\#\#\# 2.4.5 ウェッジ積の完成
\end{quote}

これで準備は整った。 我々が導き出した \(xy\)
平面の面積測定器行列は、「\(x_1 y_2\) のスイッチ」から「\(y_1 x_2\)
のスイッチ」を引き算したものだった。

これを物差し（\(1\)-form）の言葉で翻訳すれば、新しい演算子
\textbf{ウェッジ積（}\(dx \wedge dy\)\textbf{）}
の定義が自然に立ち上がる。

\[
dx \wedge dy := dx \otimes dy - dy \otimes dx
\]

記号 \(\wedge\) は英語で
``wedge''、くさびを意味する。この定義をベクトルに作用させるとこうなる：

\[
(dx \wedge dy)(\mathbf{v}_1, \mathbf{v}_2) = dx(\mathbf{v}_1)\,dy(\mathbf{v}_2) - dy(\mathbf{v}_1)\,dx(\mathbf{v}_2) = x_1 y_2 - x_2 y_1
\]

見事に、第1章の物差し（\(1\)-form）の組み合わせから、面積測定器（\(2\)-form）をシステマティックに構成することができた。
テンソル積の引き算で構成したことにより、\textbf{引数を入れ替えると符号が反転する（交代性）}という
§2.2 のルール2が、自動的に「仕込まれている」ことにも注目してほしい。

\hypertarget{ux3064ux306eux57faux5e952-form}{%
\subsubsection{\texorpdfstring{2.4.6
3つの基底\(2\)-form}{2.4.6 3つの基底2-form}}\label{ux3064ux306eux57faux5e952-form}}

同じ方法で、残る2つの「偏った眼鏡」も構成できる。

\(yz\)\textbf{平面を見る眼鏡：}\(dy \wedge dz\)

\[
dy \otimes dz - dz \otimes dy = \begin{pmatrix} 0 & 0 & 0 \\ 0 & 0 & 1 \\ 0 & 0 & 0 \end{pmatrix} - \begin{pmatrix} 0 & 0 & 0 \\ 0 & 0 & 0 \\ 0 & 1 & 0 \end{pmatrix} = \begin{pmatrix} 0 & 0 & 0 \\ 0 & 0 & 1 \\ 0 & -1 & 0 \end{pmatrix}
\]

\(zx\)\textbf{平面を見る眼鏡：}\(dz \wedge dx\)

\[
dz \otimes dx - dx \otimes dz = \begin{pmatrix} 0 & 0 & 0 \\ 0 & 0 & 0 \\ 1 & 0 & 0 \end{pmatrix} - \begin{pmatrix} 0 & 0 & 1 \\ 0 & 0 & 0 \\ 0 & 0 & 0 \end{pmatrix} = \begin{pmatrix} 0 & 0 & -1 \\ 0 & 0 & 0 \\ 1 & 0 & 0 \end{pmatrix}
\]

\textbf{これで全部だ！}

では、3つの眼鏡を一つずつ一般のベクトルにかけて、それぞれ何を測るか見てみよう。

一般の3次元ベクトル
\(\mathbf{v}_1 = \begin{pmatrix} x_1 \\ y_1 \\ z_1 \end{pmatrix}\),
\(\mathbf{v}_2 = \begin{pmatrix} x_2 \\ y_2 \\ z_2 \end{pmatrix}\)
に対し、§2.4.5 の定義どおりに（第1引数に \(\mathbf{v}_1\)、第2引数に
\(\mathbf{v}_2\) を入れる）計算すればよい：

\[
(dy \wedge dz)(\mathbf{v}_1, \mathbf{v}_2) = y_1 z_2 - z_1 y_2
\]

\[
(dz \wedge dx)(\mathbf{v}_1, \mathbf{v}_2) = z_1 x_2 - x_1 z_2
\]

\[
(dx \wedge dy)(\mathbf{v}_1, \mathbf{v}_2) = x_1 y_2 - x_2 y_1
\]

各計算は引き算1回。ベクトルの向きがどうであれ、ウェッジ積の定義が自動的に正しい値を返す。

\hypertarget{ux3064ux306eux6570ux5024ux306eux610fux5473ux5404ux5ea7ux6a19ux9762ux3078ux306eux6b63ux5c04ux5f71}{%
\subsubsection{2.4.7
3つの数値の意味------各座標面への正射影}\label{ux3064ux306eux6570ux5024ux306eux610fux5473ux5404ux5ea7ux6a19ux9762ux3078ux306eux6b63ux5c04ux5f71}}

さて、読者は当然こう期待していただろう------「3つの眼鏡の計算結果を足し合わせて2つのベクトルを食わせればそれが小学校で習った面積になるのだろう？」と。

すなわち、こうだ。

\[
\text{小学校で習った面積} \stackrel{?}{=} (dy \wedge dz)(\mathbf{v}_1, \mathbf{v}_2) + (dz \wedge dx)(\mathbf{v}_1, \mathbf{v}_2) + (dx \wedge dy)(\mathbf{v}_1, \mathbf{v}_2)
\]

しかしながら、\textbf{実は、そう単純ではない}。

§2.2
で面積測定器を設計したとき、\(xy\)平面に張り付いた話では確かに\textbf{一つの数}で足りた。しかし3次元空間の一般の平行四辺形では、向きは
\(+/-\) の2択では済まない。3つの基底 \(2\)-form（眼鏡）を通した測定値
\[A_{yz} = (dy \wedge dz)(\mathbf{v}_1,\mathbf{v}_2),\quad A_{zx} = (dz \wedge dx)(\mathbf{v}_1,\mathbf{v}_2),\quad A_{xy} = (dx \wedge dy)(\mathbf{v}_1,\mathbf{v}_2)\]
は、空間の平行四辺形そのものの面積ではなく、各座標面へ正射影したときの符号付き面積を表している。

\textbf{各座標面への正射影}------これが3つの測定値の幾何学的な意味だ。\(A_{yz}\)
は平行四辺形を \(yz\) 平面に正射影したときの符号付き面積、\(A_{zx}\) は
\(zx\) 平面、\(A_{xy}\) は \(xy\) 平面への射影面積である。

\begin{quote}
\textbf{注}（正射影）正射影とは、図形をある平面に垂直に押し潰したときにできる「影」のことである。光源が平面の真上から照らす状況を思い浮かべるとよい。たとえば
\(xy\) 平面への正射影では、図形の \(z\) 座標を無視して \(x,y\)
座標だけを取り出した像が得られる。第3章以降、本書では正射影のことを単に\textbf{影}と呼ぶことがある。
\end{quote}

では、これらの射影面積から平行四辺形そのもののスカラー面積 \(S\)
はどう取り出せるか？　答えは、ベクトルの長さと同じ構造をしている。

読者は、ベクトル \(\mathbf{v} = x\hat{e}_x + y\hat{e}_y + z\hat{e}_z\)
から「長さ（スカラー）」を取り出すとき、各成分の二乗和の平方根
\(\sqrt{x^2+y^2+z^2}\)
を計算すればよいことを知っているだろう。3つの射影面積から平行四辺形のスカラー面積を取り出す操作も、これとまったく同じ形である。

\[S = \sqrt{A_{yz}^2 + A_{zx}^2 + A_{xy}^2} = \sqrt{(y_1 z_2 - z_1 y_2)^2 + (z_1 x_2 - x_1 z_2)^2 + (x_1 y_2 - x_2 y_1)^2}\]

これはまさしく §2.1
で「大変な計算」と呼んだ高校数学の面積の式と完全に一致する。我々はついに、代数的な測定器の設計から出発し、小学校の「タイルの枚数」にまで無事に帰還したのである。

\begin{quote}
\textbf{注}（一般の\(2\)-form）3つの測定値の単純な和
\((dy \wedge dz + dz \wedge dx + dx \wedge dy)(\mathbf{v}_1,\mathbf{v}_2)\)
も、立派に一つのスカラーを返す\(2\)-formである。しかしこの値はユークリッド面積（小学校で習った面積）とは一致しない。これを得るには「二乗和の平方根」という非線形な操作が必要であり、この非線形性こそが第6章で導入する\textbf{計量（内積）}の出番を予告している。ベクトルの長さ
\(\sqrt{x^2+y^2+z^2}\) が内積 \(\mathbf{v}\cdot\mathbf{v}\)
の平方根であったことと、完全にパラレルな構造だ。
\end{quote}

\hypertarget{ux6b21ux5143ux306eux9762ux7a4dux9577ux3055ux30680ux6b21ux5143ux306eux9762ux7a4dux70b9}{%
\subsubsection{2.4.8
1次元の「面積」（長さ）と0次元の「面積」（点）}\label{ux6b21ux5143ux306eux9762ux7a4dux9577ux3055ux30680ux6b21ux5143ux306eux9762ux7a4dux70b9}}

ここで少し視点を下げてみよう。「成分の二乗和の平方根」という操作に見覚えがあると言ったが、実は我々は\textbf{1次元の図形（線分）}に対しても、全く同じ「はかり」と「図形」の構造をすでに知っている。

第1章で扱った
\(1\)-form（\(dx, dy, dz\)）は、「ベクトルを1つだけ食べる」測定器だった。
はかられる図形は、ベクトルそのもの（\(\mathbf{v} = x\hat{e}_x + y\hat{e}_y + z\hat{e}_z\)）であり、これは「1次元の向き付き面積（向き付きの長さ）」のデータに他ならない。ここから「スカラーの長さ」を取り出す操作が
\(\sqrt{x^2+y^2+z^2}\) だった。

では、\(0\)\textbf{-form} とは何か？ 「ベクトルを \(k\) 個食べる測定器が
\(k\)-form」なのだから、\(0\)-form
とは\textbf{「ベクトルを1つも食べない測定器」}である。入力する方向（矢印）すら必要なく、その場に置くだけで1つのスカラー（温度や密度など）を返すもの------すなわち、ただの\textbf{「スカラー場（関数）」}のことだ。
\(0\)-form
が測る「0次元の図形」とは単なる「点」であり、成分は1つしかなく、大きさはその絶対値をとるだけだ。標準的な数学の言葉では
\(0\)-form
は関数そのものだが、本書で「点を測る」と言うのは、入力方向を必要とせずその点に置くだけで値を返す、という性質を次元の系列に並べるための比喩である。

点（0次元）、線分（1次元）、平行四辺形（2次元）。すべては「ベクトルを
\(k\) 個食べてスカラーを返す」という同じ構造の反復だったのである。

§2.5
ではこれをさらに一段上げて、3次元の平行六面体（体積）に同じ構造を適用する。まとめると：

\begin{itemize}
\item
  \(k=0\)：\(0\)-form は点を測る。成分1つ、大きさは絶対値
\item
  \(k=1\)：\(1\)-form
  はベクトル（\textbf{向き付きの長さ}）を測る。成分3つ、大きさは
  \(\sqrt{x^2+y^2+z^2}\)
\item
  \(k=2\)：\(2\)-form
  は2本のベクトルが張る平行四辺形（\textbf{向き付きの面積}）を測る。成分3つ、大きさは
  \(\sqrt{A_{yz}^2+A_{zx}^2+A_{xy}^2}\)
\item
  \(k=3\)：\(3\)-form
  は3本のベクトルが張る平行六面体（\textbf{向き付きの体積}）を測る。成分1つ、大きさは絶対値
\end{itemize}

本書の舞台は3次元空間だから \(k=0,1,2,3\) で打ち止めだが、「\(k\)-form
が \(k\) 本のベクトルを食べてスカラーを返す」という原理は \(k\)
によらない。このパターンが第3章の積分、第5章の外微分へと受け継がれていく。

\hypertarget{ux30afux30edux30b9ux7a4dux3068ux306eux7b26ux5408}{%
\subsubsection{2.4.9
クロス積との符合}\label{ux30afux30edux30b9ux7a4dux3068ux306eux7b26ux5408}}

大学の線形代数で「ベクトルの外積（クロス積）」を学んだ読者なら、次の3成分に見覚えがあるだろう。

\[\mathbf{v}_1 \times \mathbf{v}_2 = (y_1 z_2 - z_1 y_2)\,\hat{e}_x + (z_1 x_2 - x_1 z_2)\,\hat{e}_y + (x_1 y_2 - x_2 y_1)\,\hat{e}_z = \begin{pmatrix} y_1 z_2 - z_1 y_2 \\ z_1 x_2 - x_1 z_2 \\ x_1 y_2 - x_2 y_1 \end{pmatrix}\]

これは、いまの \(dy \wedge dz\), \(dz \wedge dx\), \(dx \wedge dy\)
の\textbf{値}を、この順に並べたものにほかならない。

\begin{quote}
\textbf{注}（不安な読者へ）見覚えがなくとも、本書ではウェッジ積をクロス積よりも先に発見した、というだけでこの先の理解に支障はない。
\end{quote}

\begin{quote}
\textbf{注}（\(2\)-formとクロス積）\(2\)-formとしての面積と、3成分ベクトルとしてのクロス積のあいだには、後の章で\textbf{ホッジ双対}と名付ける関係が成立する。デカルト座標ではこれらが同じ成分を持つことが必然的であることは第6章で改めて明らかにするとして、ここでは「向き付き面積の情報の持ち方の違い」程度に留めておこう。
\end{quote}

\begin{center}\rule{0.5\linewidth}{0.5pt}\end{center}

\begin{quote}
\textbf{【ここまでのチェックポイント】}

\begin{itemize}
\item
  ウェッジ積 \(dx \wedge dy := dx \otimes dy - dy \otimes dx\)
  は、テンソル積の引き算で反対称行列を自動生成する。\(dy \wedge dz\),
  \(dz \wedge dx\) も同様に構成され、3つの基底 \(2\)-form
  の線形結合で任意の面積測定器を組み立てられる。
\item
  3つの基底 \(2\)-form に2本のベクトル \(\mathbf{v}_1,\mathbf{v}_2\)
  を食わせると、各座標平面（\(yz\), \(zx\),
  \(xy\)）への正射影の符号付き面積 \(A_{yz},A_{zx},A_{xy}\) が得られる。
\item
  スカラー面積（タイルの枚数に相当する正の大きさ）は、射影面積の二乗和の平方根
  \(\sqrt{A_{yz}^2 + A_{zx}^2 + A_{xy}^2}\) で §2.1
  の式に戻る（ラグランジュの恒等式）。ベクトルの長さの公式と同じ形の操作だ。この非線形性は第6章で導入する計量を予告している。
\end{itemize}
\end{quote}

\begin{center}\rule{0.5\linewidth}{0.5pt}\end{center}

\hypertarget{ux4f53ux7a4dux6e2cux5b9aux5668ux3068ux884cux5217ux5f0f}{%
\subsection{§2.5
体積測定器と行列式}\label{ux4f53ux7a4dux6e2cux5b9aux5668ux3068ux884cux5217ux5f0f}}

前節で我々は、2つの一次形式（物差し）\(dx, dy\)
をウェッジ積で掛け合わせ、面積測定器（\(2\)-form）\(dx \wedge dy\)
を組み立てた。面積測定器は「2つのベクトルを食べて、それらが張る平行四辺形の符号付き面積を返す装置」だった。

いよいよ\textbf{体積}だが、ここでも立場は同じである。まず\textbf{標準基底}\(\hat{e}_x,\hat{e}_y,\hat{e}_z\)\textbf{を並べた単位立方体}に体積
\(1\)
を対応させ（ルール3）、交代性・線形性で測定器を定める。小学校の「単位立方体がいくつ入るか」という\textbf{正の体積}に戻すには、やはり\textbf{成分の二乗和の平方根}（あるいは絶対値）が入る------面積のときと同じ構造だ。

自然な問いが湧く。
\textbf{「3つのベクトルが張る平行六面体の体積を測る装置——体積測定器（}\(3\)\textbf{-form）——も、同じ方法で組み立てられるだろうか？」}

\hypertarget{ux4f53ux7a4dux6e2cux5b9aux5668ux304cux6e80ux305fux3059ux3079ux304dux30ebux30fcux30eb}{%
\subsubsection{2.5.1
体積測定器が満たすべきルール}\label{ux4f53ux7a4dux6e2cux5b9aux5668ux304cux6e80ux305fux3059ux3079ux304dux30ebux30fcux30eb}}

§2.2
で面積測定器を設計したときとまったく同じ戦略をとろう。まず「体積測定器
\(V\)
が満たすべきルール」を決め、その後でルールから装置の形を一意に導く。

体積測定器 \(V(\mathbf{v}_1, \mathbf{v}_2, \mathbf{v}_3)\)
は3つのベクトルを食べて1つのスカラーを返す関数である。これに課すルールは、面積測定器の3ルールを\textbf{3引数へそのまま拡張}したものだ：

\textbf{ルール1（各引数について線形）}
片方のベクトルが2倍になれば体積も2倍。ベクトルの足し算にも分配する。
\[V(a\mathbf{v}_1 + b\mathbf{u},\; \mathbf{v}_2,\; \mathbf{v}_3) = a\,V(\mathbf{v}_1, \mathbf{v}_2, \mathbf{v}_3) + b\,V(\mathbf{u}, \mathbf{v}_2, \mathbf{v}_3)\]
第2・第3引数についても同様。

\textbf{ルール2（交代性 — どの2つを入れ替えても符号反転）}
\[V(\mathbf{v}_2, \mathbf{v}_1, \mathbf{v}_3) = -V(\mathbf{v}_1, \mathbf{v}_2, \mathbf{v}_3)\]
他のペアについても同様。§2.2
と同じ論法で、\textbf{同じベクトルを2つのスロットに入れると必ずゼロ}になる------\(\mathbf{v}_1 = \mathbf{v}_2\)
なら
\(V(\mathbf{v}_1,\mathbf{v}_1,\mathbf{v}_3) = -V(\mathbf{v}_1,\mathbf{v}_1,\mathbf{v}_3)\)
だから \(V=0\)
である。図形としては平行六面体が潰れるが、\textbf{代数では交代性だけで「潰れてゼロ」が強制される}。これが2次元の面積測定器のルール2と同じ構造だ。

\textbf{ルール3（規格化 — 単位立方体の体積を1とする）}
\[V(\hat{e}_x, \hat{e}_y, \hat{e}_z) = 1\]

\hypertarget{ux30ebux30fcux30ebux304bux3089ux5f62ux3092ux6c7aux3081ux308bux57faux5e95ux5c55ux958bux3068ux6d88ux53bb}{%
\subsubsection{2.5.2
ルールから形を決める------基底展開と消去}\label{ux30ebux30fcux30ebux304bux3089ux5f62ux3092ux6c7aux3081ux308bux57faux5e95ux5c55ux958bux3068ux6d88ux53bb}}

§2.3
で面積測定器の行列を導いたときと同じ手順を踏もう。3つのベクトルを基底で展開する：
\[\mathbf{v}_1 = x_1 \hat{e}_x + y_1 \hat{e}_y + z_1 \hat{e}_z, \quad \mathbf{v}_2 = x_2 \hat{e}_x + y_2 \hat{e}_y + z_2 \hat{e}_z, \quad \mathbf{v}_3 = x_3 \hat{e}_x + y_3 \hat{e}_y + z_3 \hat{e}_z\]

ルール1（多重線形性）により、\(V(\mathbf{v}_1, \mathbf{v}_2, \mathbf{v}_3)\)
は \(3 \times 3 \times 3 = 27\) 個の項に展開される：
\[V(\mathbf{v}_1, \mathbf{v}_2, \mathbf{v}_3) = \sum_i\sum_j\sum_k v_{1i}\,v_{2j}\,v_{3k}\; V(\hat{e}_i, \hat{e}_j, \hat{e}_k)\]

27項は多いが、ルール2（交代性）が大量の項を消してくれる：

\begin{itemize}
\item
  添字に\textbf{重複がある項はすべてゼロ}（例：\(V(\hat{e}_x, \hat{e}_x, \hat{e}_z) = 0\)）
\item
  重複のない組み合わせは \(i, j, k\) がすべて異なる場合だけ、すなわち
  \((x,y,z)\) の\textbf{置換}
\end{itemize}

\((x, y, z)\) の置換は全部で \(3! = 6\) 通り。交代性により、偶置換は
\(V(\hat{e}_x, \hat{e}_y, \hat{e}_z)\)
と同符号、奇置換は反対符号になる。ルール3で
\(V(\hat{e}_x, \hat{e}_y, \hat{e}_z) = 1\) と決めたから：

\begin{longtable}[]{@{}lll@{}}
\toprule\noalign{}
置換 & 偶/奇 & 値 \\
\midrule\noalign{}
\endhead
\bottomrule\noalign{}
\endlastfoot
\((x,y,z)\) & 偶（0回交換） & \(+1\) \\
\((y,z,x)\) & 偶（2回交換） & \(+1\) \\
\((z,x,y)\) & 偶（2回交換） & \(+1\) \\
\((y,x,z)\) & 奇（1回交換） & \(-1\) \\
\((x,z,y)\) & 奇（1回交換） & \(-1\) \\
\((z,y,x)\) & 奇（1回交換） & \(-1\) \\
\end{longtable}

したがって：
\[V(\mathbf{v}_1, \mathbf{v}_2, \mathbf{v}_3) = x_1 y_2 z_3 + y_1 z_2 x_3 + z_1 x_2 y_3 - y_1 x_2 z_3 - x_1 z_2 y_3 - z_1 y_2 x_3\]

線形代数を学んだ読者なら、3つのベクトルを列に並べた行列
\[\begin{pmatrix} x_1 & x_2 & x_3 \\ y_1 & y_2 & y_3 \\ z_1 & z_2 & z_3 \end{pmatrix}\]
の\textbf{行列式（determinant）}に見えるはずだ。知らない読者も、\textbf{いまこの6項の和を「行列式」と呼ぶ}と決めて先に進んでかまわない------上のルールからこの形に落ちた、という事実のほうが本質である。

ここで押さえたいのは次の一点だけだ。行列式は、単なる「成分を掛けて引く手続き」にとどまらず、\textbf{3本の列ベクトルをこの順に食わせたときの「符号付き体積」という1つの数}を返す演算子としても読める。第1章の
\(1\)-form は「ベクトル1つ → スカラー」、§2.3 の \(2\)-form
は「ベクトル2つ → スカラー」------行列式は「ベクトル3つ →
スカラー」で、この系列をつなぐ。

面積測定器のときは反対称行列、体積測定器のときは行列式が立ち上がったが、\textbf{いずれも、先に課したルールが代数的な形を決めた}のだ。

\begin{quote}
\textbf{注}（行列式の6項の順序）教科書によっては、いきなり「行列式の定義はこの6項」と置く。本書の流れでは、\textbf{測定器のルール → 6項}の順である。
\end{quote}

\hypertarget{form-ux306eux884cux5217ux5f0fux3068ux3057ux3066ux306eux518dux89e3ux91c8}{%
\subsubsection{\texorpdfstring{2.5.3 \(2\)-form
の行列式としての再解釈}{2.5.3 2-form の行列式としての再解釈}}\label{form-ux306eux884cux5217ux5f0fux3068ux3057ux3066ux306eux518dux89e3ux91c8}}

実は、§2.4 で定義した \(2\)-form
のウェッジ積にも、行列式が隠れていた。定義を振り返ろう：
\[(dx \wedge dy)(\mathbf{v}_1, \mathbf{v}_2) = dx(\mathbf{v}_1)\,dy(\mathbf{v}_2) - dy(\mathbf{v}_1)\,dx(\mathbf{v}_2)\]

右辺は、次の \(2 \times 2\) 行列式に他ならない：
\[(dx \wedge dy)(\mathbf{v}_1, \mathbf{v}_2) = \det\begin{pmatrix} dx(\mathbf{v}_1) & dx(\mathbf{v}_2) \\ dy(\mathbf{v}_1) & dy(\mathbf{v}_2) \end{pmatrix}\]

つまり、横に「どの物差しか」、縦に「どのベクトルか」を並べた行列の行列式が、ウェッジ積の値だったのだ。

\hypertarget{dx-wedge-dy-wedge-dz-ux306eux5b9aux7fa9ux884cux5217ux5f0fux30d1ux30bfux30fcux30f3ux306eux5b8cux6210}{%
\subsubsection{\texorpdfstring{2.5.4 \(dx \wedge dy \wedge dz\)
の定義------行列式パターンの完成}{2.5.4 dx \textbackslash wedge dy \textbackslash wedge dz の定義------行列式パターンの完成}}\label{dx-wedge-dy-wedge-dz-ux306eux5b9aux7fa9ux884cux5217ux5f0fux30d1ux30bfux30fcux30f3ux306eux5b8cux6210}}

この「行列式パターン」を素直に \(3 \times 3\)
に拡張すれば、体積測定器が手に入る。

\textbf{定義：} 3つのベクトル
\(\mathbf{v}_1, \mathbf{v}_2, \mathbf{v}_3\) に対し、
\[(dx \wedge dy \wedge dz)(\mathbf{v}_1, \mathbf{v}_2, \mathbf{v}_3) := \det\begin{pmatrix} dx(\mathbf{v}_1) & dx(\mathbf{v}_2) & dx(\mathbf{v}_3) \\ dy(\mathbf{v}_1) & dy(\mathbf{v}_2) & dy(\mathbf{v}_3) \\ dz(\mathbf{v}_1) & dz(\mathbf{v}_2) & dz(\mathbf{v}_3) \end{pmatrix}\]

右辺の各成分は「一次形式 \(dx, dy, dz\) がベクトル \(\mathbf{v}_i\)
から抽出した座標成分」だから、これはまさに：
\[= \det\begin{pmatrix} x_1 & x_2 & x_3 \\ y_1 & y_2 & y_3 \\ z_1 & z_2 & z_3 \end{pmatrix}\]

§2.5.2 で導いた \(V\)
と完全に一致する。\textbf{体積測定器は、3つの物差し}\(dx, dy, dz\)\textbf{のウェッジ積として構成できた。}

線形代数で行列式を学んだ読者なら、この定義が §2.5.1
の3ルールを満たすことが容易に分かるだろう：

\begin{itemize}
\item
  \textbf{多重線形性}：行列式は各列について線形だから、自動的に満たされる。
\item
  \textbf{交代性}：行列式の2つの列を入れ替えると符号が反転する。
\item
  \textbf{規格化}：\(\hat{e}_x, \hat{e}_y, \hat{e}_z\)
  を入れると行列は単位行列になり、\(\det(I) = 1\)。
\end{itemize}

\hypertarget{ux4f53ux7a4dux6e2cux5b9aux5668ux306eux30c6ux30f3ux30bdux30ebux7a4dux8868ux73fe3-times-3-ux30d6ux30edux30c3ux30afux3092ux6a2aux306bux4e26ux3079ux305fux914dux5217}{%
\subsubsection{\texorpdfstring{2.5.5
体積測定器のテンソル積表現------\(3 \times 3\)
ブロックを横に並べた配列}{2.5.5 体積測定器のテンソル積表現------3 \textbackslash times 3 ブロックを横に並べた配列}}\label{ux4f53ux7a4dux6e2cux5b9aux5668ux306eux30c6ux30f3ux30bdux30ebux7a4dux8868ux73fe3-times-3-ux30d6ux30edux30c3ux30afux3092ux6a2aux306bux4e26ux3079ux305fux914dux5217}}

§2.4.4--2.4.5 で我々は、面積測定器 \(dx \wedge dy\) がテンソル積の引き算
\(dx \otimes dy - dy \otimes dx\) で構成でき、その結果が \(3 \times 3\)
の反対称行列であることを見た。体積測定器 \(dx \wedge dy \wedge dz\)
にも同じ手が使えるだろうか？

なぜここで敢えてテンソル配列を持ち出すのか------それは、\textbf{ベクトル3本を食べて行列式を返す操作が、面積測定器のときとまったく同じ「縮約」の考え方で書ける}ことを示すためである。§2.4.2
で定義したように、縮約とは共通添字について和をとって添字を消し去る操作だ。\(1\)-form・\(2\)-form・\(3\)-form
はすべて、このたった一つの原理で動いている。

答えはイエスだ。ただし、2つの物差しのテンソル積は \(3 \times 3\)
行列に収まったが、3つの物差しでは添字が1つ増えて
\(3 \times 3 \times 3 = 27\)
成分の\textbf{3次元配列}になる。添字を3つ持つこのような多次元配列は数学でも物理学でも\textbf{テンソル}（tensor）と呼ばれる。行列（添字2つ）はテンソルの特殊例であり、今回登場するのは添字を3つ持つ\textbf{3階のテンソル}だ。

ところで、3階以上のテンソルは、全成分を書き並べると紙面が爆発するので、さすがの物理学者もめったに配列表示はしない。代わりに成分
\(\epsilon_{ijk}\)
のように\textbf{添字つきの代表}で書いて、あとの計算は縮約に任せるのが通常のやり方だ。本節でもその流儀に倣い、必要に応じて成分だけを取り出して話を進める。

\begin{quote}
\textbf{注}（テンソルの記法）テンソル代数では、しばしば
\(\epsilon_{ijk}\)
という\textbf{書き方そのもの}が「3階テンソル全体」を指す。本書では紛らわしいので、\textbf{3階テンソル全体}を
\(\widehat{\epsilon}\) と書き、\textbf{1成分}を \(\epsilon_{ijk}\)
と書き分ける。本書では行列表示（§2.3 の
\(3\times3\)）と「成分1つ」を表記上区別する、ということだ。他書を読む際には気を付けてほしい。
\end{quote}

\begin{quote}
\textbf{注}（上付き・下付き添字）本書をのぞき見している詳しい読者のために言い訳しておく。より進んだ微分幾何やテンソル解析では、ベクトルと余ベクトル、基底と双対基底を区別するために\textbf{上付き・下付きの位置}で添字を書き分ける流儀がよく用いられる。本書は\textbf{デカルト直交座標}に固定して議論しており、この座標系では余ベクトル成分とベクトル成分の区別が表面化しないため、\(\epsilon_{ijk}\)
のように\textbf{すべて下付き}で通す。
\end{quote}

さて、この27成分の配列を反対称化（6つの置換を符号つきで足す）すると、27成分のうち6箇所だけが
\(\pm 1\)、残りはゼロの配列が立ち上がる。§2.5.2
で導いた置換の表そのものだ------成分の取り方は次のとおりだ（\(\widehat{\epsilon}\)
の \((i,j,k)\) 成分）：

\[
\epsilon_{ijk} = \begin{cases} +1 & (i,j,k) \text{ が } (x,y,z) \text{ の偶置換} \\ -1 & (i,j,k) \text{ が } (x,y,z) \text{ の奇置換} \\ 0 & \text{添字に重複あり} \end{cases}
\]

27個の数をどう並べて見せるか------良い方法がある。\(3 \times 3\)\textbf{行列を3枚、横一列に並べた}と考えるとよい：

\[
\widehat{\epsilon} \;{=}\;
\begin{pmatrix}
\begin{pmatrix} 0 & 0 & 0 \\ 0 & 0 & 1 \\ 0 & -1 & 0 \end{pmatrix} &
\begin{pmatrix} 0 & 0 & -1 \\ 0 & 0 & 0 \\ 1 & 0 & 0 \end{pmatrix} &
\begin{pmatrix} 0 & 1 & 0 \\ -1 & 0 & 0 \\ 0 & 0 & 0 \end{pmatrix}
\end{pmatrix}
\]

第1添字 \(i\) が「左から何枚目の断面か」を指定し、各断面の中身は \(j,k\)
で添字づけされた \(3 \times 3\) ブロックだ。

\textbf{ベクトル3本を食べてスカラーに潰す——3重の縮約}

式で書けば一発である：

\[
V(\mathbf{v}_1, \mathbf{v}_2, \mathbf{v}_3) = \sum_i\sum_j\sum_k \epsilon_{ijk}\, v_{1i}\, v_{2j}\, v_{3k}
\]

これは\textbf{3回の縮約}だ。添字 \(i, j, k\)
の3つすべてについて和をとり、27成分の配列と3本のベクトルからただ1つのスカラーを作り出す。§2.4.2
で見た
\(\mathbf{v}_1^T M \mathbf{v}_2\)（2回の縮約）の自然な拡張である。

この3重の縮約を2段階に分解して眺めると、構造がさらによく見える：

\begin{itemize}
\tightlist
\item
  \textbf{ステップ1（添字 $i$ について縮約 → 行列）：} \(\mathbf{v}_1\)
  と \(\widehat{\epsilon}\) を添字 \(i\)
  で縮約する。\(\displaystyle M = \sum_{i=1}^3 v_{1i}\,\epsilon_{i,\cdot,\cdot}\)（\(\epsilon_{i,\cdot,\cdot}\)
  は \(\widehat{\epsilon}\) の左から \(i\) 枚目の \(3 \times 3\)
  ブロック）。結果は1つの \(3 \times 3\) 行列 \(M\) に降りる。この \(M\)
  の正体は、§2.4.5--2.4.6 で構成した面積測定器の行列そのものだ：
\end{itemize}

\begin{longtable}[]{@{}cc@{}}
\toprule\noalign{}
左から \(i\) 枚目 & 正体 \\
\midrule\noalign{}
\endhead
\bottomrule\noalign{}
\endlastfoot
\(i=1\) & \(dy \wedge dz\)（\(yz\) 平面の面積測定器） \\
\(i=2\) & \(dz \wedge dx\)（\(zx\) 平面の面積測定器） \\
\(i=3\) & \(dx \wedge dy\)（\(xy\) 平面の面積測定器） \\
\end{longtable}

\begin{itemize}
\tightlist
\item
  \textbf{ステップ2（添字 $j,k$ について縮約 → スカラー）：}
  得られた行列 \(M\) を \(\mathbf{v}_2, \mathbf{v}_3\)
  でサンドイッチして縮約する------
  \(V = \mathbf{v}_2^T M \mathbf{v}_3\)。§2.4 で慣れ親しんだ操作だ。
\end{itemize}

この2段階の縮約で §2.5.2
の行列式の6項が一つ残らず再現されること、および数値検算は、\textbf{付録A}で全成分を書き下して確認する。

\hypertarget{ux5177ux4f53ux4f8bux5e73ux884cux516dux9762ux4f53ux306eux7b26ux53f7ux4ed8ux304dux4f53ux7a4d}{%
\subsubsection{2.5.6
具体例：平行六面体の符号付き体積}\label{ux5177ux4f53ux4f8bux5e73ux884cux516dux9762ux4f53ux306eux7b26ux53f7ux4ed8ux304dux4f53ux7a4d}}

手を動かして確認しよう。§2.4.6
で面積を測った2つのベクトルに、3つ目を加えてみる：

\[\mathbf{v}_1 = \begin{pmatrix}1\\0\\1\end{pmatrix}, \quad \mathbf{v}_2 = \begin{pmatrix}0\\1\\1\end{pmatrix}, \quad \mathbf{v}_3 = \begin{pmatrix}1\\1\\0\end{pmatrix}\]

体積測定器に投入する：
\[(dx \wedge dy \wedge dz)(\mathbf{v}_1, \mathbf{v}_2, \mathbf{v}_3) = \det\begin{pmatrix}1 & 0 & 1 \\ 0 & 1 & 1 \\ 1 & 1 & 0\end{pmatrix}\]

これを定義に従って計算すると：
\[= 1\cdot1\cdot0 + 0\cdot1\cdot1 + 1\cdot0\cdot1 - 1\cdot1\cdot1 - 0\cdot0\cdot0 - 1\cdot1\cdot1 = 0 + 0 + 0 - 1 - 0 - 1 = -2\]

\begin{quote}
\textbf{注}（体積が負？）§2.2
の向き付き面積と同じで、\textbf{順序}によってはマイナスがあり得る。
\end{quote}

試しに \(\mathbf{v}_1\) と \(\mathbf{v}_2\) を入れ替えてみよう：
\[(dx \wedge dy \wedge dz)(\mathbf{v}_2, \mathbf{v}_1, \mathbf{v}_3) = -(-2) = +2\]

\hypertarget{ux30b9ux30abux30e9ux30fcux4e09ux91cdux7a4dux3068ux306eux7b26ux5408}{%
\subsubsection{2.5.7
スカラー三重積との符合}\label{ux30b9ux30abux30e9ux30fcux4e09ux91cdux7a4dux3068ux306eux7b26ux5408}}

ベクトル解析で「スカラー三重積」を学んだ読者は、次の恒等式を知っているかもしれない：
\[\mathbf{v}_1 \cdot (\mathbf{v}_2 \times \mathbf{v}_3) = \det\begin{pmatrix}x_1 & x_2 & x_3 \\ y_1 & y_2 & y_3 \\ z_1 & z_2 & z_3\end{pmatrix}\]

右辺は、我々が定義した
\((dx \wedge dy \wedge dz)(\mathbf{v}_1, \mathbf{v}_2, \mathbf{v}_3)\)
そのものだ。

\begin{quote}
\textbf{注}（内積記号の先取り）左辺の「\(\cdot\)」は\textbf{内積}の記号であり、ここでは小学校以来おなじみの「成分ごとの積の和」という意味で使っている。第6章で計量
\(g\) を使った形式的な定義を与える。
\end{quote}

\hypertarget{ux4f53ux7a4dux6e2cux5b9aux5668ux306eux51faux529bux72ecux7acbux6210ux5206ux306f1ux3064}{%
\subsubsection{2.5.8
体積測定器の出力------独立成分は1つ}\label{ux4f53ux7a4dux6e2cux5b9aux5668ux306eux51faux529bux72ecux7acbux6210ux5206ux306f1ux3064}}

我々が組み立てた \(dx \wedge dy \wedge dz\)
は、3つのベクトルを待ち構える\textbf{体積測定器（}\(3\)\textbf{-form）}である。2本のベクトルを食わせた面積測定器と異なり、体積測定器に3本のベクトルを食わせた結果は、ただ1つの符号付きスカラー
\(V\) に収まる（§2.5.9 で見るように、\(3\)-form
の独立成分は1つしかないからだ）。

そして「小学校の正の体積（スカラー体積）」が欲しければ、成分の二乗和の平方根（\(\sqrt{V^2} = |V|\)）をとればよい。独立成分が1つしかないので、ただの絶対値になる。面積のときと同じく、この非線形な操作には計量が必要であり、第6章で改めて扱う。

\hypertarget{form-ux306eux30b5ux30a4ux30baux306aux305cux72ecux7acbux6210ux5206ux306f1ux3064ux3060ux3051ux304b}{%
\subsubsection{\texorpdfstring{2.5.9 \(3\)-form
の「サイズ」------なぜ独立成分は1つだけか}{2.5.9 3-form の「サイズ」------なぜ独立成分は1つだけか}}\label{form-ux306eux30b5ux30a4ux30baux306aux305cux72ecux7acbux6210ux5206ux306f1ux3064ux3060ux3051ux304b}}

面積測定器（\(2\)-form）は \(3 \times 3\)
反対称行列で表現でき、独立な成分は3つだった（\(dy \wedge dz\),
\(dz \wedge dx\), \(dx \wedge dy\) の係数）。

では体積測定器（\(3\)-form）の独立成分はいくつか？ 基底 \(3\)-form
として作れる候補を考えてみよう。3つの物差し \(dx, dy, dz\)
から\textbf{重複なく}3つを選ぶ方法は\ldots\ldots1通りしかない。\(dx \wedge dy \wedge dz\)
だけだ（順序を変えると符号が変わるだけで、新しい基底にはならない）。

\begin{quote}
\textbf{注}（ホッジ双対の伏線）面積のとき、3つの基底 \(2\)-form
に2本のベクトルを食わせると3つの数値が得られ、スカラー面積を得るには非線形な操作（二乗和の平方根）が必要だった。一方、体積測定器
\(dx \wedge dy \wedge dz\) の出力は最初から1つのスカラー \(V\)
に収まる。そしてこれは偶然ではない。\(0\)-form（スカラー場）の独立成分も1つ、\(1\)-form
の独立成分は3つ、\(2\)-form の独立成分も3つ------つまり \(1, 3, 3, 1\)
と、両端から鏡写しになっている。この並びは、\textbf{ホッジ双対}という対称性の前触れだ。詳しくは後の章で。
\end{quote}

\hypertarget{ux6e2cux5b9aux5668ux306eux968eux5c64ux7b2c2ux7ae0ux3067ux624bux306bux5165ux308cux305fux6b66ux5668ux306eux6574ux7406}{%
\subsubsection{2.5.10
測定器の階層------第2章で手に入れた武器の整理}\label{ux6e2cux5b9aux5668ux306eux968eux5c64ux7b2c2ux7ae0ux3067ux624bux306bux5165ux308cux305fux6b66ux5668ux306eux6574ux7406}}

ここで、我々が第1章から積み上げてきた「測定器」の全体像を俯瞰しよう。

§2.4.8 以降、とくに §2.5.8 と
§2.5.9（独立成分の個数）までを通じて、\textbf{次元 $k$ ごとに「$k$-form が $k$ 本のベクトルを食べてスカラーを返す」という同じ型が繰り返される}ことが見えてきた。データの成分の個数と、「小学校の面積
\({=}\)
スカラー大きさ」の取り出し方を一枚の表にまとめると、次のとおりだ。

\begin{longtable}[]{@{}
  >{\raggedright\arraybackslash}p{(\columnwidth - 8\tabcolsep) * \real{0.2000}}
  >{\raggedright\arraybackslash}p{(\columnwidth - 8\tabcolsep) * \real{0.2000}}
  >{\raggedright\arraybackslash}p{(\columnwidth - 8\tabcolsep) * \real{0.2000}}
  >{\raggedright\arraybackslash}p{(\columnwidth - 8\tabcolsep) * \real{0.2000}}
  >{\raggedright\arraybackslash}p{(\columnwidth - 8\tabcolsep) * \real{0.2000}}@{}}
\toprule\noalign{}
\begin{minipage}[b]{\linewidth}\raggedright
次元 (\(k\))
\end{minipage} & \begin{minipage}[b]{\linewidth}\raggedright
測定器（\(k\)-form）
\end{minipage} & \begin{minipage}[b]{\linewidth}\raggedright
測る対象
\end{minipage} & \begin{minipage}[b]{\linewidth}\raggedright
成分の個数
\end{minipage} & \begin{minipage}[b]{\linewidth}\raggedright
小学校の面積の取り出し方
\end{minipage} \\
\midrule\noalign{}
\endhead
\bottomrule\noalign{}
\endlastfoot
\textbf{0次元} & \(0\)-form（スカラー場 \(f\)） & 点 & 1成分（\(f\)） &
\(\sqrt{f^2}\) \\
\textbf{1次元} & \(1\)-form（\(dx, dy, dz\)） & ベクトル（線分） &
3成分（\(x, y, z\)） & \(\sqrt{x^2 + y^2 + z^2}\) \\
\textbf{2次元} & \(2\)-form（\(dy \wedge dz\) 等） & 平行四辺形 & 3成分
& 係数の二乗和平方根 \\
\textbf{3次元} & \(3\)-form（\(dx \wedge dy \wedge dz\)） &
3本のベクトルが張る平行六面体 & 1成分（\(V\)） & \(\sqrt{V^2}\) \\
\end{longtable}

\begin{quote}
\textbf{【ここまでのチェックポイント】}

\begin{itemize}
\item
  体積測定器
  \(dx \wedge dy \wedge dz\)（\(3\)-form）は、3つの物差しのウェッジ積であり、その値は
  \(3 \times 3\) 行列式に等しい。
\item
  \(0\)-form（1成分）→ \(1\)-form（3成分）→ \(2\)-form（3成分）→
  \(3\)-form（1成分）と、\(1, 3, 3, 1\) の対称性がある。
\item
  次元が変わっても、\(k\)-form が \(k\)
  本のベクトルを食べてスカラーを返すという原理は同じ。スカラーの大きさは成分の二乗和の平方根で取り出す。
\end{itemize}
\end{quote}

\hypertarget{ux672cux7ae0ux306eux307eux3068ux3081ux3068ux7b2c3ux7ae0ux3078ux306eux5c55ux671b}{%
\subsection{§2.6
本章のまとめと第3章への展望}\label{ux672cux7ae0ux306eux307eux3068ux3081ux3068ux7b2c3ux7ae0ux3078ux306eux5c55ux671b}}

\begin{quote}
\textbf{【ここまでのチェックポイント — 第2章全体】}

\begin{itemize}
\item
  面積とは「2つのベクトルの絡み合い」を測る反対称行列（\(2\)-form）であり、物理的ルールが行列の形を完全に決める。
\item
  ウェッジ積 \(\wedge\)
  は「テンソル積の引き算」で反対称行列を構成する演算。3つの基底
  \(2\)-form から場の係数を載せて一般の \(2\)-form
  を組み立てる（係数の意味は後章）。
\item
  体積測定器
  \(dx \wedge dy \wedge dz\)（\(3\)-form）は、3つの物差しのウェッジ積であり、その値は
  \(3 \times 3\) 行列式に等しい。反対称化した
  \(\widehat{\epsilon}\)（成分
  \(\epsilon_{ijk}\)、レヴィ=チヴィタ記号）と縮約でも表せる（付録Aで全成分を確認）。
\end{itemize}
\end{quote}

本章では、第1章で揃えた物差し \(dx, dy, dz\)
から出発し、\textbf{面積（}\(2\)\textbf{-form）と体積（}\(3\)\textbf{-form）を代数的に組み立てる}ところまで来た。ウェッジ積は「幾何を直接描く」のではなく、\textbf{測定器のルールを満たす式を機械的に生成する}ための道具である。

次章（第3章）では、これらの形式を\textbf{曲線・曲面・領域に沿って集計する}------すなわち線積分・面積分・体積分を、\(1\)-form
/ \(2\)-form / \(3\)-form
の言葉で統一的に書き直す。さらに変数変換と\textbf{引き戻し（pullback）}が登場し、ヤコビアンが「影の計算」として自然に現れる。

\begin{center}\rule{0.5\linewidth}{0.5pt}\end{center}

\hypertarget{ux4ed8ux9332aux4f53ux7a4dux6e2cux5b9aux5668ux306eux30c6ux30f3ux30bdux30ebux7a4dux8868ux73fe-ux5168ux6210ux5206ux306eux8a08ux7b97}{%
\section{付録A：体積測定器のテンソル積表現 ---
全成分の計算}\label{ux4ed8ux9332aux4f53ux7a4dux6e2cux5b9aux5668ux306eux30c6ux30f3ux30bdux30ebux7a4dux8868ux73fe-ux5168ux6210ux5206ux306eux8a08ux7b97}}

§2.5.5 で、\(\widehat{\epsilon}\)（成分
\(\epsilon_{ijk}\)）とベクトル3本との縮約で行列式に一致すると述べた。ここではその全過程を手計算で追い、一つ残らず確認する。

\hypertarget{a.1-ux30c6ux30f3ux30bdux30ebux7a4dux306e3ux5f15ux6570ux3078ux306eux62e1ux5f35}{%
\subsubsection{A.1
テンソル積の3引数への拡張}\label{a.1-ux30c6ux30f3ux30bdux30ebux7a4dux306e3ux5f15ux6570ux3078ux306eux62e1ux5f35}}

復習しよう。\(dx \otimes dy\) の \((i,j)\) 成分は、\(dx\) の第 \(i\)
成分と \(dy\) の第 \(j\) 成分の積だった：
\[(dx \otimes dy)_{ij} = (dx)_i\,(dy)_j\]

3つの物差しへの拡張は自然だ------添字を1つ増やすだけ：
\[(dx \otimes dy \otimes dz)_{ijk} = (dx)_i\,(dy)_j\,(dz)_k\]

\(dx = (1\ 0\ 0)\), \(dy = (0\ 1\ 0)\), \(dz = (0\ 0\ 1)\)
を代入すると、\(3^3 = 27\) マスのうち \((i,j,k) = (1,2,3)\)
の1箇所だけが \(1\)、残り26マスはすべて \(0\)
になる。他の5つの並べ替え（\(dx \otimes dz \otimes dy\)
など）もそれぞれ1箇所だけ非ゼロになる------\(dy \otimes dx \otimes dz\)
なら位置 \((2,1,3)\) に、\(dz \otimes dy \otimes dx\) なら \((3,2,1)\)
に \(1\) が立つ。

\hypertarget{a.2-ux53cdux5bfeux79f0ux53166ux3064ux306eux914dux5217ux3092ux7b26ux53f7ux3064ux304dux3067ux91cdux306dux308b}{%
\subsubsection{A.2
反対称化------6つの配列を符号つきで重ねる}\label{a.2-ux53cdux5bfeux79f0ux53166ux3064ux306eux914dux5217ux3092ux7b26ux53f7ux3064ux304dux3067ux91cdux306dux308b}}

§2.4.5 で面積のとき
\(dx \wedge dy = dx \otimes dy - dy \otimes dx\)（2項の符号つき和）だった。3つの物差しは
デカルト座標における \(dx, dy, dz\) である。その並べ替えは \(3! = 6\)
通りだから、置換 \(\sigma\) ごとに \((dx,dy,dz)\)
を並べ替えた順にテンソル積 \(\otimes\) でつなぎ、符号
\(\mathrm{sgn}(\sigma)\) を付けて \(S_3\)
上で和をとる、という意味である------すなわち次の6行の和を、置換の言葉で一行にまとめたものだ。

\begin{quote}
\textbf{注}（対称群と置換の符号）\(S_3\)
は3次対称群（置換の集合）。\(\mathrm{sgn}(\sigma)\)
は\textbf{偶置換なら}\(+1\)\textbf{、奇置換なら}\(-1\) と定める符号。
\end{quote}

書き下すと：

\[\underbrace{dx \otimes dy \otimes dz}_{(1,2,3)\text{ に }+1} - \underbrace{dx \otimes dz \otimes dy}_{(1,3,2)\text{ に }-1} + \underbrace{dy \otimes dz \otimes dx}_{(2,3,1)\text{ に }+1}\]
\[- \underbrace{dy \otimes dx \otimes dz}_{(2,1,3)\text{ に }-1} + \underbrace{dz \otimes dx \otimes dy}_{(3,1,2)\text{ に }+1} - \underbrace{dz \otimes dy \otimes dx}_{(3,2,1)\text{ に }-1}\]

6つの「1箇所だけ非ゼロの配列」を符号つきで重ねると、27成分のうち6箇所だけが
\(\pm 1\) で、残り21箇所はゼロになる。§2.5.2
の置換の表と見比べてほしい------これが \(\widehat{\epsilon}\) の各成分
\(\epsilon_{ijk}\) を与える構成原理だ。

\hypertarget{a.3-ux5404ux6210ux5206ux306eux884cux5217ux3092ux66f8ux304dux4e0bux3059}{%
\subsubsection{A.3
各成分の行列を書き下す}\label{a.3-ux5404ux6210ux5206ux306eux884cux5217ux3092ux66f8ux304dux4e0bux3059}}

§2.5.5 で示した3つの成分行列を、ここでは1マスずつ確認しよう。

\textbf{第1成分（}\(i = 1\)\textbf{,}\(x\)\textbf{）：} 6項のうち
\(i = 1\) を持つのは \((1,2,3)\) に \(+1\) と \((1,3,2)\) に \(-1\)。

\[\epsilon_{1,\cdot,\cdot} = \begin{pmatrix} 0 & 0 & 0 \\ 0 & 0 & 1 \\ 0 & -1 & 0 \end{pmatrix}\]

\textbf{第2成分（}\(i = 2\)\textbf{,}\(y\)\textbf{）：} \((2,3,1)\) に
\(+1\) と \((2,1,3)\) に \(-1\)。

\[\epsilon_{2,\cdot,\cdot} = \begin{pmatrix} 0 & 0 & -1 \\ 0 & 0 & 0 \\ 1 & 0 & 0 \end{pmatrix}\]

\textbf{第3成分（}\(i = 3\)\textbf{,}\(z\)\textbf{）：} \((3,1,2)\) に
\(+1\) と \((3,2,1)\) に \(-1\)。

\[\epsilon_{3,\cdot,\cdot} = \begin{pmatrix} 0 & 1 & 0 \\ -1 & 0 & 0 \\ 0 & 0 & 0 \end{pmatrix}\]

繰り返しになるが、これらはそれぞれ
\(dy \wedge dz\)、\(dz \wedge dx\)、\(dx \wedge dy\)
の行列そのものである。

\hypertarget{a.4-ux30d9ux30afux30c8ux30eb3ux672cux3092ux98dfux308fux305bux308bux7e2eux7d04ux306eux5168ux904eux7a0b}{%
\subsubsection{A.4
ベクトル3本を食わせる------縮約の全過程}\label{a.4-ux30d9ux30afux30c8ux30eb3ux672cux3092ux98dfux308fux305bux308bux7e2eux7d04ux306eux5168ux904eux7a0b}}

\textbf{ステップ1：添字}\(i\)\textbf{について}\(\mathbf{v}_1\)\textbf{の成分で加重和（添字}\(i\)\textbf{を潰す）}

\(\mathbf{v}_1\)
の各成分を重みにして3つの面積測定器行列を足し合わせる。3次元配列がベクトル1本を食べて
\(3 \times 3\) 行列に降りる：

\[M := \sum_i v_{1i}\, \epsilon_{i,\cdot,\cdot} = x_1 \begin{pmatrix} 0 & 0 & 0 \\ 0 & 0 & 1 \\ 0 & -1 & 0 \end{pmatrix} + y_1 \begin{pmatrix} 0 & 0 & -1 \\ 0 & 0 & 0 \\ 1 & 0 & 0 \end{pmatrix} + z_1 \begin{pmatrix} 0 & 1 & 0 \\ -1 & 0 & 0 \\ 0 & 0 & 0 \end{pmatrix}\]

各成分を足し合わせると：

\[M = \begin{pmatrix} 0 & z_1 & -y_1 \\ -z_1 & 0 & x_1 \\ y_1 & -x_1 & 0 \end{pmatrix}\]

面積測定器3つを \(\mathbf{v}_1\) の成分で構成した、「\(\mathbf{v}_1\)
専用の面積測定器」が出現した。反対称行列であることに注目------\(\widehat{\epsilon}\)
の添字に対する反対称性が受け継がれている。

\textbf{ステップ2：残り2本で行列を挟む（添字}\(j, k\)\textbf{を潰す）}

あとは §2.4 で馴染んだ「横 \(\times\) 行列 \(\times\) 縦」だ：

\[V = \mathbf{v}_2^T\, M\, \mathbf{v}_3 = \begin{pmatrix} x_2 & y_2 & z_2 \end{pmatrix} \begin{pmatrix} 0 & z_1 & -y_1 \\ -z_1 & 0 & x_1 \\ y_1 & -x_1 & 0 \end{pmatrix} \begin{pmatrix} x_3 \\ y_3 \\ z_3 \end{pmatrix}\]

まず \(\mathbf{v}_2^T \times M\)（横 \(\times\) 行列 → 横）を計算する：

\[= \begin{pmatrix} y_1 z_2 - z_1 y_2, & z_1 x_2 - x_1 z_2, & x_1 y_2 - x_2 y_1 \end{pmatrix}\]

最後に、得た横ベクトルと \(\mathbf{v}_3\) の対応成分の積の和（横
\(\times\) 縦 → スカラー）：

\[V = (y_1 z_2 - z_1 y_2)\,x_3 + (z_1 x_2 - x_1 z_2)\,y_3 + (x_1 y_2 - x_2 y_1)\,z_3\]

展開して整理すると：

\[= x_1 y_2 z_3 + y_1 z_2 x_3 + z_1 x_2 y_3 - y_1 x_2 z_3 - x_1 z_2 y_3 - z_1 y_2 x_3\]

となり、\textbf{§2.5.2 の行列式と、6つの項が1つ残らず一致する。}

\newpage

\hypertarget{ux7b2c3ux7ae0ux7a4dux5206ux3059ux308bux3068ux306fux4f55ux304b-ux6709ux9650ux306eux30deux30b9ux3092ux6570ux3048ux6700ux5f8cux306bux6975ux9650ux3092ux53d6ux308b}{%
\chapter{第3章：積分するとは何か ------
有限のマスを数え、最後に極限を取る}\label{ux7b2c3ux7ae0ux7a4dux5206ux3059ux308bux3068ux306fux4f55ux304b-ux6709ux9650ux306eux30deux30b9ux3092ux6570ux3048ux6700ux5f8cux306bux6975ux9650ux3092ux53d6ux308b}}

\hypertarget{ux66f2ux304cux3063ux305fux3082ux306eux3092ux6e2cux308bux5c0fux5b66ux6821ux4ee5ux6765ux306eux501fux308aux3092ux8fd4ux3059}{%
\subsection{§3.0
曲がったものを測る------小学校以来の借りを返す}\label{ux66f2ux304cux3063ux305fux3082ux306eux3092ux6e2cux308bux5c0fux5b66ux6821ux4ee5ux6765ux306eux501fux308aux3092ux8fd4ux3059}}

第2章では、平行四辺形の面積や平行六面体の体積といった\textbf{平らな図形}の測り方を定義した。

\(dx \wedge dy\)
という面積測定器に二本のベクトルを食わせれば符号付き面積が返り、\(dx \wedge dy \wedge dz\)
という体積測定器に三本のベクトルを食わせれば符号付き体積が返る。小学校の「\(1 \times 1\)
の正方形がいくつ入るか」という直感を、代数的な測定器として再設計したのである。

しかし、小学校以来の借りはまだ残っている。

我々は円周 \(2\pi r\)、球の表面積 \(4\pi r^2\)、球の体積
\(\frac{4}{3}\pi r^3\)
を、長いあいだ「そういう公式」として使ってきた。だが、それらは本来、どこから来るのか。

答えは一つである。

\textbf{有限の小片に刻み、それぞれを測り、足し上げ、最後に刻みを無限に細かくする。}

これが積分である。

本章の目的は、球の公式を覚え直すことではない。むしろ、見慣れた球という題材を使って、「積分する」とは何をしているのかを、リーマン和のレベルから確認することにある。

そのため、ここではあえて極座標や球面座標にすぐ逃げない。\(x,y,z\)
のまま小片を数え、どの小片が対象に含まれるかを判定し、和を整理する。泥臭いが、この泥臭さこそが積分の本体である。

さて、いま我々の手元には、第1章と第2章で作り上げた測定器------\(0\)-form、\(1\)-form、\(2\)-form、\(3\)-form------がそろっている。本章の仕事は、これらを\textbf{曲線・曲面・領域に沿って適用し、集計する}ことである。

原理はどの次数でも同じだ。まずは係数1の体積測定器がそのまま領域の体積を返してくれるケースから始め、次元を降りながら「形式は何を測り、何を測らないのか」を確かめていこう。

\begin{quote}
\textbf{注} （お約束）
お約束として断っておく。以下では、曲線や曲面は十分滑らかで、小片に刻んで足し上げる操作が素直に極限を持つ場合を扱う。向きの反転や自己交差などの細部は「物理でよく使う状況では問題にならない」レベルに留める。厳密な可積分論は本書の射程外である。
\end{quote}

\begin{center}\rule{0.5\linewidth}{0.5pt}\end{center}

\hypertarget{ux4f53ux7a4d3ux6b21ux5143ux4fc2ux65701}{%
\subsection{§3.1
体積------3次元、係数1}\label{ux4f53ux7a4d3ux6b21ux5143ux4fc2ux65701}}

\hypertarget{ux76f4ux65b9ux4f53ux304bux3089ux66f2ux304cux3063ux305fux9818ux57dfux3078ux30eaux30fcux30deux30f3ux548cux3067ux5b9aux7fa9ux3059ux308b}{%
\subsubsection{3.1.1
直方体から曲がった領域へ------リーマン和で定義する}\label{ux76f4ux65b9ux4f53ux304bux3089ux66f2ux304cux3063ux305fux9818ux57dfux3078ux30eaux30fcux30deux30f3ux548cux3067ux5b9aux7fa9ux3059ux308b}}

第2章 §2.5 で、我々は \(dx \wedge dy \wedge dz\)
という体積測定器を組み立てた。これに三本のベクトル
\(\mathbf{v}_1, \mathbf{v}_2, \mathbf{v}_3\)
を食わせると、それらが張る平行六面体の符号付き体積が返ってくる。

いま、空間内の領域 \(V\) の体積を測りたい。\(V\)
の形がどんなに曲がっていても、十分細かく刻めば各小片はほぼ直方体と見なせる。そこで
\(V\) を \(x\) 方向・\(y\) 方向・\(z\) 方向に刻み、\(V\)
の内部にある小直方体に番号 \(i\) を振る。第 \(i\) 小直方体の三辺を

\[\Delta x_i\,\hat{e}_x,\quad \Delta y_i\,\hat{e}_y,\quad \Delta z_i\,\hat{e}_z\]

とする。これら三本のベクトルを体積測定器 \(dx \wedge dy \wedge dz\)
に食わせると：

\[(dx \wedge dy \wedge dz)(\Delta x_i\,\hat{e}_x,\; \Delta y_i\,\hat{e}_y,\; \Delta z_i\,\hat{e}_z) = \det\begin{pmatrix}\Delta x_i & 0 & 0 \\ 0 & \Delta y_i & 0 \\ 0 & 0 & \Delta z_i\end{pmatrix} = \Delta x_i\,\Delta y_i\,\Delta z_i\]

\begin{quote}
\textbf{注} （測定器と図形の縮約）ここで起きていることは、第2章 §2.5.10
で見たパターンそのものだ。\textbf{$3$-form（体積測定器）}が、三本の変位ベクトルが張る\textbf{体積素（図形）}を食べてスカラーを吐き出す。この縮約が各小直方体で行われる。ここでは、この三本の変位ベクトルが張る向き付きの小直方体、すなわち微小な体積図形を「体積素」と呼ぶ。
\end{quote}

各小直方体から得られたスカラー \(\Delta x_i\,\Delta y_i\,\Delta z_i\)
を、\(V\) の内部にあるすべての小直方体について足し上げる：

\[\sum_{\text{小直方体 }i \subset V} \Delta x_i\,\Delta y_i\,\Delta z_i\]

そして刻みを無限に細かくする極限をとる。この極限を

\[\iiint_V dx \wedge dy \wedge dz\]

と\textbf{書くことにする}。これが \(3\)-form
の積分------体積分------の定義である。

\begin{quote}
\textbf{注} （\(\iiint\) という記号について）本書では \(\iiint_V\)
を「3次元領域 \(V\) における \(3\)-form
の積分」を表す記号として\textbf{ここで初めて定義する}。高校数学の体積分を前提にした記号ではなく、上記のリーマン和の極限に与えた名前である。
\end{quote}

さて、定義はした。ではこの和を実際にどう計算するか。\(V\) が
\(x^2+y^2+z^2 \le R^2\) のような曲がった領域のとき、どの小直方体が \(V\)
の内部にあるかを判定しながら和をとるのは面倒だ。しかし和のとり方を工夫すれば、見通しがよくなる。次節で実際にやってみよう。

\hypertarget{ux7403ux306eux4f53ux7a4dux6ce5ux81edux304fux611aux76f4ux306bux8db3ux3059}{%
\subsubsection{3.1.2
球の体積------泥臭く、愚直に足す}\label{ux7403ux306eux4f53ux7a4dux6ce5ux81edux304fux611aux76f4ux306bux8db3ux3059}}

半径 \(R\) の球 \(x^2 + y^2 + z^2 \le R^2\)
の体積を、本当に泥臭く計算してみよう。これから示したいことはただ一つ------特別な座標系に逃げずとも、\(x, y, z\)
のまま愚直にリーマン和を整理するだけで、見慣れた \(\frac{4}{3}\pi R^3\)
に到達できることである。特別な座標系は使わない。\(x, y, z\)
のまま、小直方体 \(\Delta x\,\Delta y\,\Delta z\)
を愚直に積み重ねるだけだ。

球を、関数

\[F(x,y,z)=x^2+y^2+z^2\]

の値で判定する領域として見る。球の内部とは \(F\le R^2\)
を満たす点全体であり、境界の球面は \(F=R^2\) である。

境界の向きや、球面に沿う変位を読むには、第1章で定義した全微分 \(dF\)
が効いてくる。しかし体積を数える段階では、まだ \(dF\)
は使わない。ここではまず、\(F\le R^2\)
という条件をそのままほどいて、小直方体を数える。体積を測る主役は、あくまで
\(dx\wedge dy\wedge dz\) である。

§3.1.1 のリーマン和をどう整理すれば計算できるか。\(z\)
をある値に固定すると、\(F\le R^2\) は

\[x^2+y^2\le R^2-z^2\]

になる。つまり、その高さで許される \((x,y)\) は半径 \(\sqrt{R^2-z^2}\)
の円板である。さらに \(y\) も固定すれば、

\[x^2\le R^2-z^2-y^2\]

だから、\(x\) の動く範囲は \(-\sqrt{R^2 - z^2 - y^2}\) から
\(+\sqrt{R^2 - z^2 - y^2}\) までである。同様に、\(y\) の範囲は
\(-\sqrt{R^2 - z^2}\) から \(+\sqrt{R^2 - z^2}\) まで、\(z\) の範囲は
\(-R\) から \(+R\) まで。

これは結局、\(F\le R^2\) という判定を、\(z\)、\(y\)、\(x\)
の順にほどいているだけである。この順に和をとる------第1章 §1.1
でやった1次元のリーマン和を三回重ねる------と考えれば、極限では：

\[\iiint_V dx \wedge dy \wedge dz = \int_{-R}^{R} \Biggl( \int_{-\sqrt{R^2 - z^2}}^{\sqrt{R^2 - z^2}} \Biggl( \int_{-\sqrt{R^2 - z^2 - y^2}}^{\sqrt{R^2 - z^2 - y^2}} dx \Biggr) dy \Biggr) dz\]

右辺の \(dx, dy, dz\) は、左辺の \(dx \wedge dy \wedge dz\)
を順に和の整理に使った結果として並んでいる。\(\wedge\)
を省略しているのではない。カッコの内側から順に、\(x\)
についての和（\(\int dx\)）→ \(y\) についての和（\(\int dy\)）→ \(z\)
についての和（\(\int dz\)）を実行する、という意味である。

一番内側の \(x\) についての和（極限で \(x\) 積分）はすぐに実行できる：

\[\int_{-\sqrt{R^2 - z^2 - y^2}}^{\sqrt{R^2 - z^2 - y^2}} dx = 2\sqrt{R^2 - z^2 - y^2}\]

次に \(y\) 積分。\(a = \sqrt{R^2 - z^2}\) とおけば：

\[\int_{-a}^{a} 2\sqrt{a^2 - y^2}\,dy\]

ここで置換積分（高校数学）を使う。\(y = a\sin t\)
とおくと、\(dy = a\cos t\,dt\)。\(y\) が \(-a \to a\) のとき \(t\) は
\(-\pi/2 \to \pi/2\)。また
\(\sqrt{a^2 - y^2} = \sqrt{a^2 - a^2\sin^2 t} = a\cos t\)（\(t \in [-\pi/2,\pi/2]\)
では \(\cos t \ge 0\) だから絶対値が外れる）。したがって：

\[\begin{aligned}
\int_{-a}^{a} 2\sqrt{a^2 - y^2}\,dy
&= \int_{-\pi/2}^{\pi/2} 2\cdot a\cos t \cdot a\cos t\,dt \\
&= 2a^2\!\int_{-\pi/2}^{\pi/2} \cos^2 t\,dt
\end{aligned}\]

\(\cos^2 t = (1 + \cos 2t)/2\) だから：

\[\begin{aligned}
2a^2\!\int_{-\pi/2}^{\pi/2} \frac{1 + \cos 2t}{2}\,dt
&= a^2\int_{-\pi/2}^{\pi/2} (1 + \cos 2t)\,dt \\[4pt]
&= a^2\Bigl[t + \frac{\sin 2t}{2}\Bigr]_{-\pi/2}^{\pi/2} \\[4pt]
&= a^2\Bigl[\Bigl(\frac{\pi}{2} + 0\Bigr) - \Bigl(-\frac{\pi}{2} + 0\Bigr)\Bigr] \\[4pt]
&= \pi a^2 = \pi(R^2 - z^2)
\end{aligned}\]

最後に \(z\) 積分：

\[\int_{-R}^{R} \pi(R^2 - z^2)\,dz = 2\pi\Bigl[R^2 z - \frac{z^3}{3}\Bigr]_{0}^{R} = 2\pi\cdot\frac{2R^3}{3} = \frac{4}{3}\pi R^3\]

\begin{quote}
\textbf{注} （この計算の意味）ここで我々がやったことは、§3.1.1
のリーマン和の整理にすぎない。\(x,y,z\)
のまま、どの小直方体が球の内部にあるかを判定しながら和をとるかわりに、和の順序を
\(x \to y \to z\)
と整理した。各段階は第1章でやった1変数のリーマン和であり、極限で見慣れた1変数積分の計算になる。それだけである。測定器と図形の縮約を愚直に集計すれば
\(\frac{4}{3}\pi R^3\) に到達する。
\end{quote}

こうして、\(dx \wedge dy \wedge dz\)
という\textbf{一つの測定器}の積分が、曲がった領域の体積
\(\frac{4}{3}\pi R^3\)
を正しく与えることが確認できた。3次元では、係数1の体積測定器がそのまま体積を直接測れるのである。

\begin{quote}
\textbf{【ここまでのチェックポイント】}

\begin{itemize}
\item
  \(dx \wedge dy \wedge dz\)
  の積分は、領域を小直方体に刻み、各小片の体積素を体積測定器に食わせた値を足し上げるリーマン和の極限である。\(\iiint_V\)
  は本書がこの極限に与えた記号である。
\item
  係数1でそのまま体積が測れる。球の体積 \(\frac{4}{3}\pi R^3\)
  が、リーマン和を \(x \to y \to z\)
  の順に整理することで得られる。特別な座標系は不要。
\item
  次節では次元を一つ下げ、曲面の面積を測る。
\end{itemize}
\end{quote}

\begin{center}\rule{0.5\linewidth}{0.5pt}\end{center}

\hypertarget{ux8868ux9762ux7a4d2ux6b21ux5143ux4fc2ux65701ux306eux9650ux754c}{%
\subsection{§3.2
表面積------2次元、係数1の限界}\label{ux8868ux9762ux7a4d2ux6b21ux5143ux4fc2ux65701ux306eux9650ux754c}}

\hypertarget{ux5e73ux884cux56dbux8fbaux5f62ux304bux3089ux66f2ux9762ux3078}{%
\subsubsection{3.2.1
平行四辺形から曲面へ}\label{ux5e73ux884cux56dbux8fbaux5f62ux304bux3089ux66f2ux9762ux3078}}

第2章 §2.4 で、我々は \(dx \wedge dy\)
という面積測定器を組み立てた。これに二本のベクトル
\(\mathbf{v}_1, \mathbf{v}_2\) を食わせると、それらが \(xy\)
平面に落とす影の符号付き面積が返ってくる。同様に \(dy \wedge dz\) は
\(yz\) 平面への影、\(dz \wedge dx\) は \(zx\) 平面への影を測る。

では曲面の面積を測るにはどうすればよいか。まず平らな場合で確認しておこう。\(xy\)
平面に乗った単位正方形を、細かな小正方形に刻む。一つの小正方形の二辺は

\[\Delta\mathbf{x}=\begin{pmatrix}\Delta x\\0\\0\end{pmatrix},\qquad
\Delta\mathbf{y}=\begin{pmatrix}0\\\Delta y\\0\end{pmatrix}\]

である。これを \(dx \wedge dy\) に食わせると、

\[(dx \wedge dy)(\Delta\mathbf{x},\Delta\mathbf{y})=\Delta x\,\Delta y\]

となる。全小正方形で集計すれば面積は
\(1\)。これは第2章の平らな平行四辺形の拡張にすぎない。ここで確認しているのは、\(dx\wedge dy\)
という \(2\)-form が、\(xy\)
平面上の小片に対してはその符号付き面積をそのまま返す、ということである。

では曲面が曲がっていたらどうなるか。曲面を十分細かい小片に刻む。各小片の上で、隣り合う二方向の変位ベクトルを
\(\Delta\mathbf{a},\Delta\mathbf{b}\)
と書く。この二本の変位ベクトル、またはそれらが張る向き付きの微小平行四辺形を、この章では曲面の「面素」と呼ぶ。曲面を細かく刻めば刻むほど、この面素の集計は曲面そのものの集計に近づく。

この曲面小片を張る二本の変位ベクトル
\(\Delta\mathbf{a},\Delta\mathbf{b}\) を、面積測定器 \(dx \wedge dy\)
に食わせる。第2章 §2.4.4 の定義より、もし

\[\Delta\mathbf{a}=\begin{pmatrix}\Delta a_x\\\Delta a_y\\\Delta a_z\end{pmatrix},\qquad
\Delta\mathbf{b}=\begin{pmatrix}\Delta b_x\\\Delta b_y\\\Delta b_z\end{pmatrix}\]

なら、

\[(dx \wedge dy)(\Delta\mathbf{a},\Delta\mathbf{b})
=\Delta a_x\Delta b_y-\Delta a_y\Delta b_x\]

である。これは、その小片が \(xy\) 平面に落とす影の符号付き面積だ。

刻みを無限に細かくする極限を

\[\iint_S dx \wedge dy\]

と\textbf{書くことにする}。これが \(2\)-form
の積分------面積分------の定義である。構成は §3.1
の体積分と同じ型で、食わせるベクトルが2本になっただけだ。

\hypertarget{dx-wedge-dy-ux304cux6e2cux308bux3082ux306eux7403ux9762ux3067ux8a66ux3059}{%
\subsubsection{\texorpdfstring{3.2.2 \(dx \wedge dy\)
が測るもの------球面で試す}{3.2.2 dx \textbackslash wedge dy が測るもの------球面で試す}}\label{dx-wedge-dy-ux304cux6e2cux308bux3082ux306eux7403ux9762ux3067ux8a66ux3059}}

実際に球面で計算してみよう。半径 \(R\)
の球面の上半分（\(z \ge 0\)）を考える。球面上の点 \((x,y,z)\) は

\[x^2+y^2+z^2=R^2,\qquad z=\sqrt{R^2-x^2-y^2}\]

を満たす。ここで、第1章で定義した全微分を使う。関数

\[F(x,y,z)=x^2+y^2+z^2\]

の全微分は

\[dF=2x\,dx+2y\,dy+2z\,dz\]

である。第1章の約束どおり、これは変位ベクトルを食べて \(F\)
の変化を読み取る横ベクトルである。球面に接する変位ベクトルなら、一次の変化としては
\(F=x^2+y^2+z^2\) の値を変えない。したがって、その変位ベクトル
\(\mathbf{v}\) は \(dF(\mathbf{v})=0\)
を満たす。球面小片を張る二本の変位ベクトルも、この条件を満たすように選ぶ。

たとえば、\(y\) を止めて \(x\) 方向へ微小幅 \(\Delta x\)
だけ進む変位ベクトルを

\[\Delta\mathbf{x}_S=\begin{pmatrix}\Delta x\\[2pt]0\\[2pt]\Delta z\end{pmatrix}\]

とおく。これを \(dF\) に食わせると、

\[dF(\Delta\mathbf{x}_S)=2x\,\Delta x+2z\,\Delta z\]

である。球面に沿う方向として \(dF(\Delta\mathbf{x}_S)=0\)
となるように選ぶと、\(\Delta z=-(x/z)\Delta x\) となる。同様に、\(x\)
を止めて \(y\) 方向へ微小幅 \(\Delta y\)
だけ進む変位ベクトルについては、\(\Delta z=-(y/z)\Delta y\) となる。

したがって、球面小片を張る二本の変位ベクトルは

\[\Delta\mathbf{x}_S=\begin{pmatrix}\Delta x\\[2pt]0\\[2pt]-\dfrac{x}{z}\Delta x\end{pmatrix},\qquad
\Delta\mathbf{y}_S=\begin{pmatrix}0\\[2pt]\Delta y\\[2pt]-\dfrac{y}{z}\Delta y\end{pmatrix}\]

と書ける。曲面を十分細かく刻めば、この二本が張る平行四辺形の集計は、球面そのものの集計に近づいていく。

まず \(dx \wedge dy\) だけを曲面に食わせてみる：

\[(dx \wedge dy)(\Delta\mathbf{x}_S,\Delta\mathbf{y}_S)=\Delta x\,\Delta y\]

したがって上半球面での \(dx \wedge dy\) の積分は：

\[\iint_{S_{\text{upper}}} dx \wedge dy\]

の計算は §3.1.2 の \(y\) 積分と同じ方法でできて（\(a\) を \(R\)
と思えばまったく同じ）、結果は \(\pi R^2\) である。これは上半球を \(xy\)
平面に落とした影------半径 \(R\) の円板------の面積にほかならない。

下半球（\(z = -\sqrt{R^2 - x^2 - y^2}\)）でも同様に計算できる。ただし球面全体を外向きに測ると、下半球では面の向きが裏返るために符号が反転し、積分は
\(-\pi R^2\) になる。

したがって球面全体での \(dx \wedge dy\) の積分は
\(\pi R^2 + (-\pi R^2) = 0\)。上半球の影と下半球の影が打ち消し合ったのだ。これではっきりした：\textbf{$dx \wedge dy$ が測るのは「$xy$ 平面への影の符号付き面積」であり、曲面のスカラー面積ではない。}同様に、\(dy \wedge dz\)
は \(yz\) 平面への影、\(dz \wedge dx\) は \(zx\) 平面への影を測る。

\hypertarget{ux4e09ux3064ux306eux5f71ux304bux3089ux30b9ux30abux30e9ux30fcux9762ux7a4dux3078}{%
\subsubsection{3.2.3
三つの影からスカラー面積へ}\label{ux4e09ux3064ux306eux5f71ux304bux3089ux30b9ux30abux30e9ux30fcux9762ux7a4dux3078}}

では曲面のスカラー面積（小学校の面積）はどう測ればよいのか。各小区画で面素を構成する二本の変位ベクトルを、三つの面積測定器すべてに食わせ、その結果を合成する。

上半球面で残り二つも計算しよう。先ほどの二本の変位ベクトル

\[\Delta\mathbf{x}_S=\begin{pmatrix}\Delta x\\0\\-\dfrac{x}{z}\Delta x\end{pmatrix},\qquad
\Delta\mathbf{y}_S=\begin{pmatrix}0\\\Delta y\\-\dfrac{y}{z}\Delta y\end{pmatrix}\]

を三つの面積測定器に食わせると：

\[\begin{aligned}
(dy \wedge dz)(\Delta\mathbf{x}_S,\Delta\mathbf{y}_S) &= \frac{x}{z}\,\Delta x\,\Delta y
= \frac{x}{\sqrt{R^2-x^2-y^2}}\,\Delta x\,\Delta y \\[4pt]
(dz \wedge dx)(\Delta\mathbf{x}_S,\Delta\mathbf{y}_S) &= \frac{y}{z}\,\Delta x\,\Delta y
= \frac{y}{\sqrt{R^2-x^2-y^2}}\,\Delta x\,\Delta y \\[4pt]
(dx \wedge dy)(\Delta\mathbf{x}_S,\Delta\mathbf{y}_S) &= \Delta x\,\Delta y
\end{aligned}\]

三つの測定値 \(A_{yz},A_{zx},A_{xy}\)
は、面素の三つの座標平面への符号付き射影面積である。これら三成分をそろえると、向き付き微小面積の情報が記録される。これは第2章
§2.4.6
で見た、平行四辺形の向き付き面積を三つの射影面積として読む話の曲面版である。

ここから、向きを捨てた正の面積を取り出したい。そこで、この章では

\[
dS(\Delta\mathbf{a},\Delta\mathbf{b})
= \sqrt{A_{yz}^2+A_{zx}^2+A_{xy}^2}
\]

を、その小区画の\textbf{スカラー面素}と呼ぶ。

これは \(2\)-form ではない。\(dx\wedge dy\)
が面素を構成する二本の変位ベクトルを食べて一つの射影面積を返すのに対し、\(dS\)
は三つの射影面積を、直交デカルト座標でおなじみの二乗和平方根によって合成し、向きを捨てた正の面積を返す記法である。

ただし、すぐに \(dS\)
を足し上げる前に、三つの面積測定器をそれぞれ上半球面全体で積分すると何が起きるかを見ておこう。上半球が
\(xy\) 平面に落とす円板を \(D=\{(x,y)\mid x^2+y^2\le R^2\}\)
と書くと、三つの符号付き影面積は次の三本の積分になる：

\[\begin{aligned}
\iint_{S_{\text{upper}}} dy\wedge dz
&= \int_{-R}^{R}\int_{-\sqrt{R^2-x^2}}^{\sqrt{R^2-x^2}}
\frac{x}{\sqrt{R^2-x^2-y^2}}\,dy\,dx, \\[4pt]
\iint_{S_{\text{upper}}} dz\wedge dx
&= \int_{-R}^{R}\int_{-\sqrt{R^2-x^2}}^{\sqrt{R^2-x^2}}
\frac{y}{\sqrt{R^2-x^2-y^2}}\,dy\,dx, \\[4pt]
\iint_{S_{\text{upper}}} dx\wedge dy
&= \int_{-R}^{R}\int_{-\sqrt{R^2-x^2}}^{\sqrt{R^2-x^2}}
1\,dy\,dx.
\end{aligned}\]

この三本は、それぞれ \(yz\) 平面、\(zx\) 平面、\(xy\)
平面への符号付き影面積を測っている。実際、対称性から前二者は
\(0\)、最後は \(\pi R^2\) になる：

\[
\iint_{S_{\text{upper}}} dy\wedge dz=0,\qquad
\iint_{S_{\text{upper}}} dz\wedge dx=0,\qquad
\iint_{S_{\text{upper}}} dx\wedge dy=\pi R^2.
\]

ここで注意する。これら三本の積分値を後から二乗和平方根にしても、上半球面の面積にはならない。三つの影を\textbf{各小区画ごとに}合成してから足し上げなければ、曲面の傾きの情報が途中で消えてしまうからである。

したがって、表面積を出すには、各小区画で三つの射影面積を求め、そこからスカラー面素
\(dS\) を作り、それを足し上げればよい。先ほどの三つの射影面積を

\[\begin{aligned}
A_{yz} &= (dy \wedge dz)(\Delta\mathbf{x}_S,\Delta\mathbf{y}_S) \\
A_{zx} &= (dz \wedge dx)(\Delta\mathbf{x}_S,\Delta\mathbf{y}_S) \\
A_{xy} &= (dx \wedge dy)(\Delta\mathbf{x}_S,\Delta\mathbf{y}_S)
\end{aligned}\]

と書けば、スカラー面素は

\[
dS(\Delta\mathbf{x}_S,\Delta\mathbf{y}_S)
= \sqrt{A_{yz}^2+A_{zx}^2+A_{xy}^2}
\]

である。先ほど計算した三つの出力を代入すると、

\[\begin{aligned}
dS(\Delta\mathbf{x}_S,\Delta\mathbf{y}_S)
&= \sqrt{
\Bigl(\frac{x}{\sqrt{R^2-x^2-y^2}}\Delta x\,\Delta y\Bigr)^2
{}+ \Bigl(\frac{y}{\sqrt{R^2-x^2-y^2}}\Delta x\,\Delta y\Bigr)^2
{}+ (\Delta x\,\Delta y)^2} \\[4pt]
&= \sqrt{\frac{x^2+y^2}{R^2-x^2-y^2}+1}\;\Delta x\,\Delta y \\[4pt]
&= \frac{R}{\sqrt{R^2-x^2-y^2}}\,\Delta x\,\Delta y
\end{aligned}\]

つまり、上半球が \(xy\) 平面に落とす円板 \(x^2+y^2 \le R^2\)
を小区画に刻むと、面積のリーマン和は

\[\sum_i\sum_j
\frac{R}{\sqrt{R^2-x_i^2-y_j^2}}\,
\Delta x_i\,\Delta y_j\]

になる。つまり

\[\sum_i\sum_j dS(\Delta\mathbf{x}_{S,ij},\Delta\mathbf{y}_{S,ij})\]

という形で、各小区画のスカラー面素を愚直に足し上げているだけである。刻みを細かくする極限では：

\[
\operatorname{Area}(S_{\text{upper}})
= R\int_{-R}^{R}\biggl(\int_{-\sqrt{R^2-x^2}}^{\sqrt{R^2-x^2}}
\frac{dy}{\sqrt{R^2 - x^2 - y^2}}\biggr)\;dx
\]

内側の \(y\) 積分。\(a = \sqrt{R^2 - x^2}\) とおく：

\[\int_{-a}^{a} \frac{dy}{\sqrt{a^2 - y^2}}\]

置換積分 \(y = a\sin t\) を使う（§3.1.2
と同じ手順）。\(dy = a\cos t\,dt\)、\(\sqrt{a^2 - y^2} = a\cos t\)、\(y\)
が \(-a \to a\) のとき \(t\) は \(-\pi/2 \to \pi/2\)：

\[\int_{-a}^{a} \frac{dy}{\sqrt{a^2 - y^2}} = \int_{-\pi/2}^{\pi/2} \frac{a\cos t}{a\cos t}\,dt = \int_{-\pi/2}^{\pi/2} dt = \pi\]

積分結果は \(a\) に依らず \(\pi\) になる。したがって：

\[R\int_{-R}^{R} \pi\,dx = R \cdot \pi \cdot 2R = 2\pi R^2\]

これが上半球面の面積。下半球面でも（原点から見て「外向き」を正にとれば）符号は一部反転するが、二乗和は同じなので、同じく
\(2\pi R^2\)。よって球面全体の面積は：

\[2\pi R^2 + 2\pi R^2 = 4\pi R^2\]

見事に \(4\pi R^2\) が出た。第2章 §2.4.7
で平行四辺形の面積を三つの影から復元したのと、\textbf{まったく同じパターン}である。

\hypertarget{ux3053ux3053ux307eux3067ux3067ux308fux304bux3063ux305fux3053ux3068}{%
\subsubsection{3.2.4
ここまででわかったこと}\label{ux3053ux3053ux307eux3067ux3067ux308fux304bux3063ux305fux3053ux3068}}

係数1の基底 \(2\)-form
は、それ単独では曲面のスカラー面積を直接返さない。\(dx \wedge dy\)
は、曲面小片の \(xy\)
平面への符号付き射影面積を測る。曲面のスカラー面積を得るには、三つの射影成分を各小片ごとに合成してスカラー面素
\(dS\)
を作り、それを足し上げなければならない。これは第2章の平行四辺形の面積のときとまったく同じ状況だ。

しかし、物理ではしばしば「各点でどの方向の影をどれだけ重視するか」を制御したくなる。たとえば、曲面を貫く流体の流れを測るとき、流れが
\(x\) 方向に強ければ \(dy \wedge dz\)（\(yz\)
平面への影）に大きな重みをかけたい。そのような\textbf{場所ごとに変わる重み}------係数------の必要性が、ここで自然に姿を現す。これが
§3.4 への伏線である。

\begin{quote}
\textbf{【ここまでのチェックポイント】}

\begin{itemize}
\item
  曲面積分は、各小区画で面素を構成する二本の変位ベクトルを面積測定器に食わせ、得られたスカラーを集計する操作である。
\item
  \(dx \wedge dy\) は単独では \(xy\)
  平面への影の面積を測る。球面全体では正味ゼロ（上半球と下半球の影が打ち消し合う）。
\item
  三つの基底 \(2\)-form
  の測定値を各小片ごとに二乗和平方根で合成すると、スカラー面素 \(dS\)
  が得られる。これを足し上げれば曲面のスカラー面積になる。球面では
  \(4\pi R^2\)。
\item
  次節ではさらに次元を下げ、曲線の長さを測る。そこで係数1の限界が最もはっきりと姿を現す。
\end{itemize}
\end{quote}

\begin{center}\rule{0.5\linewidth}{0.5pt}\end{center}

\hypertarget{ux66f2ux7dda1ux6b21ux5143ux4fc2ux65701ux306eux9650ux754cux304cux3042ux3089ux308fux306bux306aux308b}{%
\subsection{§3.3
曲線------1次元、係数1の限界があらわになる}\label{ux66f2ux7dda1ux6b21ux5143ux4fc2ux65701ux306eux9650ux754cux304cux3042ux3089ux308fux306bux306aux308b}}

\hypertarget{ux76f4ux7ddaux304bux3089ux66f2ux7ddaux3078ux30eaux30fcux30deux30f3ux548cux3067ux5b9aux7fa9ux3059ux308b}{%
\subsubsection{3.3.1
直線から曲線へ------リーマン和で定義する}\label{ux76f4ux7ddaux304bux3089ux66f2ux7ddaux3078ux30eaux30fcux30deux30f3ux548cux3067ux5b9aux7fa9ux3059ux308b}}

第1章で、我々は \(dx, dy, dz\) を「ベクトルを食べて各成分を返す
\(1\)-form（横ベクトル）」として定義した。\(dx = \begin{pmatrix}1&0&0\end{pmatrix}\)、\(\mathbf{v} = \begin{pmatrix}\Delta x\\\Delta y\\\Delta z\end{pmatrix}\)
のとき \(dx(\mathbf{v}) = \Delta x\) である。\(dy, dz\) も同様。

では曲線に沿ってこれらを適用し、集計すると何が出るか。

まず、曲線を細かい小区間に刻む。各小区間には、始点から終点への変位ベクトル

\[
\Delta\mathbf{r}_i=(\Delta x_i,\Delta y_i,\Delta z_i)
\]

が対応する。この段階では、まだ曲線を式で表示する必要はない。各小区間で
\(dx\) をこの変位ベクトルに食わせれば、

\[dx(\Delta\mathbf{r}_i)=\Delta x_i\]

である。同様に、\(dy(\Delta\mathbf{r}_i)=\Delta y_i\)、\(dz(\Delta\mathbf{r}_i)=\Delta z_i\)
である。

全小区間で足し上げれば、

\[
\sum_i dx(\Delta\mathbf{r}_i)=\sum_i \Delta x_i
\]

となる。これは、曲線上の各小片が持つ \(x\)
方向の変位を、愚直に足し上げているだけである。刻みを無限に細かくする極限を

\[\int_\gamma dx\]

と\textbf{書くことにする}。これが \(1\)-form
の積分------線積分------の定義である。

さて、定義はできた。ではこの和を実際に計算するにはどうするか。

ここで \(t\) を導入するのは、これまでの方針を捨てるためではない。§3.1
では球の内部の小直方体を \(z,y,x\) の順に整理し、§3.2 では球面上の小片を
\(xy\)
平面への射影で整理した。同じように、曲線では小片を始点から終点へ順に並べるためのラベルが必要になる。そのラベルを
\(t\) と呼ぶだけである。

そこで、曲線上の小片に連続的な番号 \(t\) を振り、

\[\gamma(t) = \bigl(x(t),\, y(t),\, z(t)\bigr),\qquad t \in [t_0, t_1]\]

と書く。ここでの \(t\)
は、新しい空間座標ではない。曲線上の小片を始点から終点へ順番に並べるためのラベルである。このラベルを使うと、\(i\)
番目の小区間の変位は

\[
\Delta\mathbf{r}_i = \gamma(t_i+\Delta t)-\gamma(t_i)
\]

と書ける。

各小区間で\textbf{行列 $dx$ を変位ベクトル $\Delta\mathbf{r}_i$ に作用させる}：

\[dx(\Delta\mathbf{r}_i) = \begin{pmatrix}1&0&0\end{pmatrix} \begin{pmatrix}\Delta x_i \\ \Delta y_i \\ \Delta z_i\end{pmatrix} = \Delta x_i = \frac{\Delta x_i}{\Delta t}\,\Delta t\]

ここで等式なのは、固定した \(1\)-form \(dx\) を変位ベクトル
\(\Delta\mathbf{r}_i\) に作用させれば、成分 \(\Delta x_i\)
が正確に返る、という意味である。一方で、曲線全体の積分は、小区間ごとの代表点で測定器を評価して足し上げるリーマン和の極限として定義される。

\[\sum_i dx(\Delta\mathbf{r}_i) = \sum_i \frac{\Delta x_i}{\Delta t}\,\Delta t\]

第1章で扱った置換積分は、まさにこのような和の整理だった。ここでも同じことをする。曲線上の小片を
\(t\) で番号づけし、各成分の変化を「変化率 \(\times\)
小区間幅」として整理するのである。刻みを無限に細かくすると、換算率は導関数に収束する：\(\Delta x_i/\Delta t \to dx/dt\)。

この計算は第1章 §1.2.5 の置換積分とまったく同じ形であり：

\[\int_\gamma dx = \int_{t_0}^{t_1} \frac{dx}{dt}\,dt = x(t_1) - x(t_0)\]

同様に
\(\int_\gamma dy = y(t_1) - y(t_0)\)、\(\int_\gamma dz = z(t_1) - z(t_0)\)。すなわち\textbf{$dx, dy, dz$ という $1$-form の線積分は、曲線の正味の変位——始点と終点の座標の差——を返す。}

\hypertarget{ux5186ux5468ux3067ux8a66ux3059ux5f27ux9577ux306fux51faux306aux3044}{%
\subsubsection{3.3.2
円周で試す------弧長は出ない}\label{ux5186ux5468ux3067ux8a66ux3059ux5f27ux9577ux306fux51faux306aux3044}}

確認しよう。円周の場合も、まず考えているのは円を細かい弧に刻み、各小片の変位を足し上げる操作である。ただし、実際に一周分の和を計算するには、小片を順番に並べる番号が必要になる。ここではその番号として、角度
\(t\) を使う。半径 \(R\) の円を

\[\gamma(t) = (R\cos t,\; R\sin t,\; 0),\qquad t \in [0, 2\pi]\]

と表す。\(\gamma'(t) = (-R\sin t,\; R\cos t,\; 0)\) だから：

\[\begin{aligned}
\int_\gamma dx &= \int_0^{2\pi} dx(\gamma'(t))\,dt = \int_0^{2\pi} (-R\sin t)\,dt = R\bigl[\cos t\bigr]_0^{2\pi} = 0 \\[4pt]
\int_\gamma dy &= \int_0^{2\pi} dy(\gamma'(t))\,dt = \int_0^{2\pi} (R\cos t)\,dt = R\bigl[\sin t\bigr]_0^{2\pi} = 0
\end{aligned}\]

いずれもゼロである。閉曲線だから始点と終点が同じで、行きと帰りで打ち消し合う------当然の結果だ。

しかし、この円の弧長が \(2\pi R\) であることも我々は知っている。係数1の
\(dx, dy\)
だけではゼロになってしまうが、では\textbf{どうすれば弧長が出るのか}。

各番号 \(t\) で \(dx, dy\) が測った値はすでに手元にある。

\[
dx(\gamma'(t)) = -R\sin t,\qquad dy(\gamma'(t)) = R\cos t
\]

これらはスカラーだから、二乗して足すことができる：

\[dx(\gamma'(t))^2 + dy(\gamma'(t))^2 = R^2\sin^2 t + R^2\cos^2 t = R^2\]

その平方根をとれば、その小片あたりの長さの換算率 \(|\gamma'(t)| = R\)
が得られる。全区間で積分すれば：

\[\int_0^{2\pi} \!\sqrt{dx(\gamma'(t))^2 + dy(\gamma'(t))^2}\;dt = \int_0^{2\pi} \!R\,dt = 2\pi R\]

弧長 \(2\pi R\) が出た。ここで起きていることは重要だ------\(dx\) も
\(dy\)
も単独では閉曲線でゼロになったのに、それらの二乗和の平方根をとって積分すると、正しい弧長が姿を現した。係数1の
\(1\)-form
にはなかった「長さを測る力」が、二乗和という代数的な組み換えによって回復されたのである。

ここで起きていることを整理しよう。各番号 \(t\) で \(dx(\gamma'(t))\) と
\(dy(\gamma'(t))\) という二つのスカラーを得て、それらの二乗和の平方根
\(\sqrt{dx(\gamma'(t))^2 + dy(\gamma'(t))^2}\) をとり、\(t\)
で積分した。これはすなわち、直交座標での長さ測定を曲線上で積分したことにほかならない。三次元では、この線素を

\[ds(\mathbf{v}) = \sqrt{dx(\mathbf{v})^2 + dy(\mathbf{v})^2 + dz(\mathbf{v})^2}\]

と書く。いまの円周は \(z=0\) の平面内にあるので、\(dz(\gamma'(t))=0\)
となり、上の計算では \(dx,dy\) の二成分だけが残っている。

\begin{quote}
\textbf{注} （記法について）第1章 §1.1.6 の契約により、単独の \(dx\)
は横ベクトル（演算子）であり、それ自体を二乗することはできない。\(dx(\mathbf{v})^2\)
は、\(1\)-form \(dx\) をベクトル \(\mathbf{v}\)
に\textbf{作用させて得たスカラー}を二乗したものである。この区別を崩さないことが、本書の記法の一貫性を保つ。
\end{quote}

この \(ds\) は弧長を測るための便利な記法だが、本書でいう \(1\)-form
そのものではない。\(1\)-form は変位 \(\mathbf{v}\)
に対して線形に作用する測定器であるのに対し、\(ds(\mathbf{v})\)
は二乗和の平方根を含むため、一般に
\(ds(\mathbf{v}+\mathbf{w}) = ds(\mathbf{v})+ds(\mathbf{w})\)
を満たさない。したがって \(ds\) を \(dx,dy,dz\)
の線形結合------\(P\,dx+Q\,dy+R\,dz\) の形------で書くことはできない。

\hypertarget{ux3067ux306f-1-form-ux3067ux4f55ux304cux6e2cux308cux308bux306eux304b}{%
\subsubsection{\texorpdfstring{3.3.3 では \(1\)-form
で何が測れるのか}{3.3.3 では 1-form で何が測れるのか}}\label{ux3067ux306f-1-form-ux3067ux4f55ux304cux6e2cux308cux308bux306eux304b}}

弧長は
\(ds(\mathbf{v}) = \sqrt{dx(\mathbf{v})^2 + dy(\mathbf{v})^2 + dz(\mathbf{v})^2}\)
という線素を積分して測れた。ただし、この \(ds\) は \(dx,dy,dz\)
の線形結合では書けず、二乗和の平方根を必要とする。ここではまず、このような二乗和平方根を使わない
\(1\)-form------すなわち \(dx, dy, dz\) の線形結合で書ける
\(1\)-form------が、いったい何を測るのかを見ていこう。

空間の各点に力の場 \(\mathbf{F}(x,y,z) = (F_x, F_y, F_z)\)
が働いている。曲線に沿って質点を動かすとき、各小区間で力がする微小仕事は、力の変位方向成分と変位の大きさの積------すなわち
\(F_x\,\Delta x + F_y\,\Delta y + F_z\,\Delta z\)------である。これを全区間で足し上げたものが仕事
\(W = \int_\gamma \mathbf{F}\cdot d\mathbf{r}\) だ。

ここで重要なのは、力の成分 \(F_x, F_y, F_z\)
は\textbf{場所ごとに変わる}ことである。体積（§3.1）や表面積（§3.2）では係数1で十分な部分が多かったが、仕事を測るには場所ごとに変わる係数が不可欠だ。これが係数の必然性をもっとも鮮明に示す。

これが、次節で係数を導入する動機である。

\begin{quote}
\textbf{【ここまでのチェックポイント】}

\begin{itemize}
\item
  \(dx, dy, dz\) の線積分は「正味の変位」を返す。閉曲線ではゼロ。
\item
  弧長は
  \(ds(\mathbf{v}) = \sqrt{dx(\mathbf{v})^2 + dy(\mathbf{v})^2 + dz(\mathbf{v})^2}\)
  という線素を積分して計算できる。\(ds\) は \(1\)-form
  ではなく、\(dx, dy, dz\)
  の線形結合で書くことはできない（二乗和の平方根を必要とする）。
\item
  \(dx, dy, dz\) の線形結合で書ける \(1\)-form
  が測るのは「仕事」のような量である。→ §3.4 へ。
\end{itemize}
\end{quote}

\begin{center}\rule{0.5\linewidth}{0.5pt}\end{center}

\hypertarget{ux4fc2ux6570ux3092ux3064ux3051ux308bux5bc6ux5ea6ux3068ux5e7eux4f55ux306eux7a4d}{%
\subsection{§3.4
係数をつける------密度と幾何の積}\label{ux4fc2ux6570ux3092ux3064ux3051ux308bux5bc6ux5ea6ux3068ux5e7eux4f55ux306eux7a4d}}

\hypertarget{ux306aux305cux4fc2ux6570ux306aux306eux304bux7269ux7406ux304cux8981ux6c42ux3059ux308b}{%
\subsubsection{3.4.0
なぜ係数なのか------物理が要求する}\label{ux306aux305cux4fc2ux6570ux306aux306eux304bux7269ux7406ux304cux8981ux6c42ux3059ux308b}}

§3.1〜§3.3
では、係数1の形式で曲がった図形を測ってきた。しかし現実の物理では、測定器に\textbf{場所ごとに変わる重み}を掛けたくなる場面が無数にある：

\begin{itemize}
\item
  \textbf{力の場}（場所ごとに変わる）\(\times\) \textbf{変位} \(=\)
  \textbf{仕事}
\item
  \textbf{流速の場}（場所ごとに変わる）\(\times\) \textbf{断面} \(=\)
  \textbf{流量}
\item
  \textbf{密度の場}（場所ごとに変わる）\(\times\) \textbf{体積} \(=\)
  \textbf{質量}
\end{itemize}

ここに共通する構造は明白だ。\textbf{「場所ごとに変わる重み（$0$-form ＝ スカラー場）」}と\textbf{「幾何的な測り方（$k$-form）」}の積が、一般の
\(k\)-form を構成する。第2章 §2.4.8 で「\(0\)-form
はベクトルを0本食べる測定器、すなわちスカラー場」と述べたのは、まさにこの布石である。

以下、1次元・2次元・3次元の順に、係数付き形式の積分を定義していく。

\hypertarget{ux4e00ux822cux306e-1-form-ux306eux7ddaux7a4dux5206ux4ed5ux4e8b}{%
\subsubsection{\texorpdfstring{3.4.1 一般の \(1\)-form
の線積分（仕事）}{3.4.1 一般の 1-form の線積分（仕事）}}\label{ux4e00ux822cux306e-1-form-ux306eux7ddaux7a4dux5206ux4ed5ux4e8b}}

実空間で、係数が点 \((x,y,z)\) に依存する三つのスカラー場 \(P,Q,R\)
をとり、

\[\omega = P(x,y,z)\,dx + Q(x,y,z)\,dy + R(x,y,z)\,dz\]

を\textbf{$1$-form の一般形}と呼ぶ。各 \(dx, dy, dz\)
は第1章どおり、縦ベクトルから成分を読み取る横ベクトルである。係数
\(P,Q,R\) は場所ごとに変わる重み------\(0\)-form------である。もし
\(P=\partial f/\partial x,\;Q=\partial f/\partial y,\;R=\partial f/\partial z\)
と一つの関数 \(f\) から来るなら \(\omega=df\)
という特別な場合になるが、一般にはそうとは限らない。

曲線 \(\gamma(t)\) に沿った積分は、§3.3.1
のリーマン和に係数を乗せるだけだ。各小区間の代表点を \(\gamma(t_i)\)
とし、その点で係数を評価した \(\omega_{\gamma(t_i)}\) を変位
\(\Delta\mathbf{r}_i\) に食わせる。固定した点 \(\gamma(t_i)\) での
\(\omega_{\gamma(t_i)}\) は横ベクトル、\(\Delta\mathbf{r}_i\)
は縦ベクトルであり、その縮約は次の等式で書ける：

\[\omega_{\gamma(t_i)}(\Delta\mathbf{r}_i) = P(\gamma(t_i))\,\Delta x_i + Q(\gamma(t_i))\,\Delta y_i + R(\gamma(t_i))\,\Delta z_i\]

ここで等式なのは、一点で固定した横ベクトル \(\omega_{\gamma(t_i)}\)
と変位ベクトル \(\Delta\mathbf{r}_i\)
の縮約である、という意味である。一方で、小区間全体の寄与としては、代表点で係数を凍結したリーマン和の一項であり、刻みを細かくする極限で線積分になる。

各成分の変化を「有限の比 \(\times \Delta t\)」の形に整理する：

\[\omega(\Delta\mathbf{r}_i) = \bigl(P\frac{\Delta x_i}{\Delta t} + Q\frac{\Delta y_i}{\Delta t} + R\frac{\Delta z_i}{\Delta t}\bigr)\,\Delta t\]

\(P,Q,R\) は点 \(\gamma(t_i)\)
での値。全区間で足し上げて極限をとる。\(\Delta x_i/\Delta t \to dx/dt,\; \Delta y_i/\Delta t \to dy/dt,\; \Delta z_i/\Delta t \to dz/dt\)
だから：

\[\sum_i \omega(\Delta\mathbf{r}_i) \;\xrightarrow{\Delta t \to 0}\; \int_{t_0}^{t_1} \bigl(P\frac{dx}{dt} + Q\frac{dy}{dt} + R\frac{dz}{dt}\bigr)\,dt\]

この極限を

\[\int_\gamma \omega\]

と書く。§3.3.1
の係数1の場合とまったく同じ型で、各瞬間に\textbf{横ベクトル（係数付き $1$-form）}が\textbf{縦ベクトル（$\gamma'(t)$）}を食べてスカラーを返し、それを集計している。これは第1章
§1.2.5 の \(F(x)\,dx\) を3次元に自然拡張しただけである。

\begin{quote}
\textbf{注} （ベクトル解析との対応）ベクトル解析では、この量を
\(\int_\gamma \mathbf{F}\cdot d\mathbf{r}\) と書く。本書の
\(\int_\gamma \omega\) はそれと同じでありながら、\textbf{場 $\omega$}
と\textbf{経路 $\gamma$} が記号上も明確に分離されている。
\end{quote}

具体例を見よう。力場 \(\mathbf{F} = (y,\; x,\; 0)\) が、単位円
\(\gamma(t) = (\cos t,\; \sin t,\; 0),\; t \in [0, 2\pi]\)
に沿ってする仕事を計算する。

\(\omega = y\,dx + x\,dy\) だから、係数は
\(P=y,\;Q=x,\;R=0\)。曲線上では \(P=\sin t,\;Q=\cos t\)。また
\(dx/dt = -\sin t,\; dy/dt = \cos t\) より：

\[\omega(\gamma'(t)) = (\sin t)(-\sin t) + (\cos t)(\cos t) = -\sin^2\!t + \cos^2\!t = \cos 2t\]

したがって：

\[\int_\gamma \omega = \int_0^{2\pi} \cos 2t\,dt = \Bigl[\frac{1}{2}\sin 2t\Bigr]_0^{2\pi} = 0\]

仕事がゼロになった。この力場は単位円に沿って正味の仕事をしないのである。

\hypertarget{ux4e00ux822cux306e-2-form-ux306eux9762ux7a4dux5206ux6d41ux91cf}{%
\subsubsection{\texorpdfstring{3.4.2 一般の \(2\)-form
の面積分（流量）}{3.4.2 一般の 2-form の面積分（流量）}}\label{ux4e00ux822cux306e-2-form-ux306eux9762ux7a4dux5206ux6d41ux91cf}}

同様に、三つのスカラー場 \(P, Q, R\) を係数とする一般の \(2\)-form：

\[\eta = P\,dy \wedge dz + Q\,dz \wedge dx + R\,dx \wedge dy\]

曲面に沿った積分も、§3.2.1
のリーマン和に係数を乗せるだけだ。各小区画の代表点で係数を評価し、その小片を張る二本の変位ベクトル
\(\Delta\mathbf{a},\Delta\mathbf{b}\) に \(\eta\) を食わせる：

\[\eta(\Delta\mathbf{a},\Delta\mathbf{b})
= P\,(dy \wedge dz)(\Delta\mathbf{a},\Delta\mathbf{b})
{}+ Q\,(dz \wedge dx)(\Delta\mathbf{a},\Delta\mathbf{b})
{}+ R\,(dx \wedge dy)(\Delta\mathbf{a},\Delta\mathbf{b})\]

全区画で足し上げて極限をとる。この極限を

\[\iint_S \eta\]

と記す。各点で面素を構成する二本の変位ベクトルを面積測定器 \(\eta\)
が食べてスカラーを返し、それを曲面全体で集計する点は §3.2 と同じである。

\(dy \wedge dz,\; dz \wedge dx,\; dx \wedge dy\)
がそれぞれ各平面への「影の面積」を測ることは §3.2 で見た。係数
\(P, Q, R\)
は、それらの影に\textbf{場所ごとに変わる重み}をつける役割を果たす。物理的には、\((P, Q, R)\)
が流速場の成分であり、\(\iint_S \eta\)
が曲面を貫く流量（フラックス）に対応するが、これをベクトル解析の言葉に翻訳して厳密に整理するには追加の対応関係が必要である。詳しくは後の章で行う。

\begin{quote}
\textbf{注}
（ここではベクトル解析との対応は補助線）流量との対応は、既にベクトル解析を知っている読者のための橋渡しにすぎない。本書の本線はあくまで「\(2\)-form
を食わせて集計する」という操作であり、ベクトル解析の知識がなくともこの先の議論に支障はない。
\end{quote}

\hypertarget{ux4e00ux822cux306e-3-form-ux306eux4f53ux7a4dux5206ux7dcfux8ceaux91cf}{%
\subsubsection{\texorpdfstring{3.4.3 一般の \(3\)-form
の体積分（総質量）}{3.4.3 一般の 3-form の体積分（総質量）}}\label{ux4e00ux822cux306e-3-form-ux306eux4f53ux7a4dux5206ux7dcfux8ceaux91cf}}

\(3\)-form の一般形は、スカラー場 \(\rho(x,y,z)\) を係数として：

\[\Omega = \rho(x,y,z)\,dx \wedge dy \wedge dz\]

§3.1.1 のリーマン和に密度 \(\rho\) を乗せればよい。各小直方体の代表点で
\(\rho\) を評価し、体積素に \(\Omega\) を食わせる：

\[\Omega(\Delta x_i\,\hat{e}_x,\; \Delta y_i\,\hat{e}_y,\; \Delta z_i\,\hat{e}_z) = \rho(x_i,y_i,z_i)\,\Delta x_i\,\Delta y_i\,\Delta z_i\]

全区間で足し上げて極限をとる。この極限を

\[\iiint_V \Omega\]

と書く。\(\rho=1\) なら §3.1 の係数1の場合に戻る。\(\rho\)
が質量密度なら総質量、電荷密度なら総電荷を与える。\(3\)-form
の独立成分は \(dx \wedge dy \wedge dz\)
の1つだけなので、\textbf{「密度（$\rho$）」と「体積の測り方（$dx \wedge dy \wedge dz$）」がきれいに分離}して見通しがよい。

\hypertarget{ux7ddaux9762ux4f53ux306eux7d71ux4e00ux50cf}{%
\subsubsection{3.4.4
線・面・体の統一像}\label{ux7ddaux9762ux4f53ux306eux7d71ux4e00ux50cf}}

ここまでを俯瞰しよう。

\begin{longtable}[]{@{}
  >{\centering\arraybackslash}p{(\columnwidth - 6\tabcolsep) * \real{0.2941}}
  >{\raggedright\arraybackslash}p{(\columnwidth - 6\tabcolsep) * \real{0.2353}}
  >{\raggedright\arraybackslash}p{(\columnwidth - 6\tabcolsep) * \real{0.2353}}
  >{\raggedright\arraybackslash}p{(\columnwidth - 6\tabcolsep) * \real{0.2353}}@{}}
\toprule\noalign{}
\begin{minipage}[b]{\linewidth}\centering
対象
\end{minipage} & \begin{minipage}[b]{\linewidth}\raggedright
形式
\end{minipage} & \begin{minipage}[b]{\linewidth}\raggedright
積分の記法
\end{minipage} & \begin{minipage}[b]{\linewidth}\raggedright
物理的意味の例
\end{minipage} \\
\midrule\noalign{}
\endhead
\bottomrule\noalign{}
\endlastfoot
曲線 \(\gamma\) & \(1\)-form \(\omega\) &
\(\displaystyle\int_\gamma \omega\) & 仕事 \\
曲面 \(S\) & \(2\)-form \(\eta\) & \(\displaystyle\iint_S \eta\) &
流量 \\
領域 \(V\) & \(3\)-form \(\Omega\) & \(\displaystyle\iiint_V \Omega\) &
総質量 \\
\end{longtable}

リーマン和として書けば、それぞれ

\[
\int_\gamma \omega
= \lim_{\text{刻み}\to 0}
\sum_i \omega(\Delta\mathbf{r}_i),
\]

\[
\iint_S \eta
= \lim_{\text{刻み}\to 0}
\sum_i\sum_j\eta(\Delta\mathbf{a}_{ij},\Delta\mathbf{b}_{ij}),
\]

\[
\iiint_V \Omega
= \lim_{\text{刻み}\to 0}
\sum_i\sum_j\sum_k
\Omega(\Delta\mathbf{a}_{ijk},\Delta\mathbf{b}_{ijk},\Delta\mathbf{c}_{ijk})
\]

である。ここで「刻み
\(\to 0\)」とは、曲線・曲面・領域を構成する最も粗い小片の大きさを \(0\)
に近づける、という意味である。

どれも、\textbf{$k$-form（測定器）が $k$ 個の変位ベクトル、すなわち $k$ 次元の小片を食べてスカラーを返し、それを領域全体で集計する}という同じ型をしている。次数が上がるごとに食わせるベクトルの本数が増え、交代性が面積・体積の向きを司る。次数が違うだけで、原理はどこまでも同じだ。

\begin{quote}
\textbf{【ここまでのチェックポイント】}

\begin{itemize}
\item
  一般の \(k\)-form は「場所ごとに変わる重み（\(0\)-form）」による「基底
  \(k\)-form」の線形結合。
\item
  線積分 \(\int_\gamma \omega\) は仕事、面積分 \(\iint_S \eta\)
  は流量、体積分 \(\iiint_V \Omega\) は総質量に対応する。
\item
  いずれも「測定器と図形の縮約 →
  集計」という同一原理。次章では、これらの積分を別の座標で計算し直す方法------引き戻し------を導入する。
\end{itemize}
\end{quote}

\hypertarget{ux7a4dux5206ux306eux4e00ux822cux5f62}{%
\subsubsection{3.4.5
積分の一般形}\label{ux7a4dux5206ux306eux4e00ux822cux5f62}}

§3.1〜§3.3
で別々に定義した線積分、面積分、体積分は、実は一つのパターンにまとまる。\(k\)-form
\(\omega\) を \(k\) 次元領域 \(M\)
上で積分するという操作を、次元によらず

\[\int_M \omega\]

と書く。\(M\) が曲線（\(k=1\)）なら \(\int_\gamma\)、曲面（\(k=2\)）なら
\(\iint_S\)、立体（\(k=3\)）なら \(\iiint_V\) ------積分記号の重なりは
\(M\) の次元で決まるにすぎない。中身は常に「小区分に刻む → \(k\)-form に
\(k\) 個の変位ベクトルを食わせる → 和 \(\sum\) → 極限」である。

本書の舞台は 3 次元空間だから \(k=1,2,3\) で十分だが、この式が \(k\)
によらず成立するという事実は頭の片隅に置いておいて損はない。第5章で外微分
\(d\) を導入したとき、この統一的な見方がストークスの定理の一般形
\(\int_{\partial M} \omega = \int_M d\omega\) へと結実する。

\begin{center}\rule{0.5\linewidth}{0.5pt}\end{center}

\hypertarget{ux672cux7ae0ux306eux307eux3068ux3081ux3068ux6b21ux7ae0ux3078ux306eux5c55ux671b-1}{%
\subsection{§3.5
本章のまとめと次章への展望}\label{ux672cux7ae0ux306eux307eux3068ux3081ux3068ux6b21ux7ae0ux3078ux306eux5c55ux671b-1}}

ここまでの三章で、我々は次の三つをそろえた。

\begin{enumerate}
\def\labelenumi{\arabic{enumi}.}
\tightlist
\item
  \textbf{第1章}：\(dx\)
  を行列（\(1\)-form）と見なし、積分を行列作用の極限と読む
\item
  \textbf{第2章}：ウェッジ積 \(\wedge\) で \(2\)-form・\(3\)-form
  を構成し、面積・体積を代数で測る
\item
  \textbf{第3章（本章）}：曲線・曲面・領域に形式を適用し、その出力を集計する操作を確立した
\end{enumerate}

本章で確立したのは、次の統一原理である。

\textbf{$k$-form（測定器）が $k$ 個の変位ベクトル、すなわち $k$ 次元の小片を食べてスカラーを返し、それを領域全体で集計する。}次数が
\(0,1,2,3\)
と変わっても、この構図は不変だった。曲線上の仕事も、曲面上の流量も、領域内の総質量も------すべて同じ言語で書ける。

\begin{quote}
\textbf{【ここまでのチェックポイント——第3章】}

\begin{itemize}
\item
  \(k\)-form \(\omega\) の積分 \(\int_M \omega\) は「小区分に刻む →
  \(k\)-form に \(k\) 個の変位ベクトルを食わせる → 和 →
  極限」で定義される。
\item
  \(1\)-form の線積分 \(\int_\gamma \omega\)
  は、曲線小片の変位ベクトルに \(\omega\)
  を食わせて集計する操作である。仕事は \(\int_\gamma \omega\)
  で、\(\mathbf{F}\cdot d\mathbf{r}\) と同じ量。
\item
  \(2\)-form \(\eta\) の面積分 \(\iint_S \eta\)
  は、小片を構成する二本の変位ベクトルに \(\eta\)
  を食わせて集計する操作。係数1では各平面への「影」を測り、三つの影からスカラー面素
  \(dS\) を作って足し上げればスカラー面積が得られる。
\item
  \(3\)-form \(\Omega\) の体積分 \(\iiint_V \Omega\)
  は、小直方体を構成する三本の変位ベクトルに \(\Omega\)
  を食わせて集計する操作である。係数1でそのまま体積、係数 \(\rho\)
  で総質量・総電荷。
\item
  一般の \(k\)-form は、場所ごとに変わる係数（\(0\)-form）と基底
  \(k\)-form の組み合わせとして書ける。
\end{itemize}
\end{quote}

しかし、まだ重要な問題が残っている。

ここまでは、できるだけ標準デカルト座標 \(x,y,z\)
の測定器を保ったまま、小片を数えてきた。だが実際の計算では、別の変数で数えた方が圧倒的に楽なことがある。円周を角度でたどる。球を半径と角度で見る。物理の問題では、そのような変数変換を避けて通ることはできない。

そのとき、何を変え、何を保てばよいのか。

次章（第4章）では、変数変換を「点を動かす操作」ではなく、\textbf{測定器を作り替える操作}として読み直す。これが\textbf{引き戻し}である。

\newpage

\hypertarget{ux7b2c4ux7ae0ux5909ux6570ux5909ux63dbux3068ux306fux4f55ux304b-ux5f15ux304dux623bux3057-phiux6e2cux5b9aux5668ux306eux3064ux3058ux3064ux307eux5408ux308fux305b}{%
\chapter{\texorpdfstring{第4章：変数変換とは何か ------ 引き戻し
\(\Phi^*\)：測定器のつじつま合わせ}{第4章：変数変換とは何か ------ 引き戻し \textbackslash Phi\^{}*：測定器のつじつま合わせ}}\label{ux7b2c4ux7ae0ux5909ux6570ux5909ux63dbux3068ux306fux4f55ux304b-ux5f15ux304dux623bux3057-phiux6e2cux5b9aux5668ux306eux3064ux3058ux3064ux307eux5408ux308fux305b}}

\hypertarget{ux7269ux7406ux306fux66f2ux304cux308bux8a08ux7b97ux306fux56dbux89d2}{%
\subsection{§4.0
物理は曲がる、計算は四角}\label{ux7269ux7406ux306fux66f2ux304cux308bux8a08ux7b97ux306fux56dbux89d2}}

物理で本当に測りたい量は、座標の名前に依らない。

質点にした仕事は、軌道を \(x\) で書いても、時刻 \(t\)
で書いても同じでなければならない。惑星が掃く面積は、直交座標で測っても、極座標で測っても同じでなければならない。物体の質量は、直交座標で積分しても、円柱座標で積分しても同じでなければならない。

ところで、物理はたいてい曲がっている。
軌道は曲がり、面は曲がり、領域の境界も曲がる。

一方、計算は四角い。
区間に刻み、マス目に分け、箱に切る。積分とは、曲がった物理量を四角い計算の上で集計する技術である。

ここで問題が起こる。
計算に使う変数を変えると、測定器をそのままでは使えなくなる。

\(dx\) は \(x\) 方向の変位を測る測定器であり、\(dt\)
に差し替えただけでは仕事は正しく測れない。 \(dx \wedge dy\) は \(xy\)
平面の面積測定器であり、\(dr \wedge d\theta\)
に差し替えただけでは面積は正しく測れない。 \(dx \wedge dy \wedge dz\)
もまた、\(dr \wedge d\theta \wedge dz\)
に差し替えただけでは体積は正しく測れない。

物理量は変わってはいけない。 しかし、計算に使う四角は変わる。

そのとき、測定器をどう作り替えればよいか。

本章では、この問いを三つの例で追う。

まず有限の区間で、時間側に何を掛ければ仕事の小片が一致するかを見る。次に有限のマス目で、面積を合わせる係数を探す。最後に有限の箱で、体積を合わせる係数を探す。

分割を細かくすると、有限の比は微分係数になり、有限の行列式は偏微分でできた行列式になる。その極限で得られる測定器の作り替えを、引き戻しと呼び、\(\Phi^*\)
と書く。

つまり、本章の順番はこうである。

\begin{quote}
有限の区間で見つける。 有限のマス目で同じことをする。
有限の箱で同じことをする。 最後に、\(h\to0\) の極限として \(\Phi^*\)
が見える。
\end{quote}

本章では、公式を覚えるのではなく、つじつまを合わせる。
仕事、面積、体積という三つの物理量を通じて、\(dx\)、\(dx \wedge dy\)、\(dx \wedge dy \wedge dz\)
が、それぞれどのように作り替えられるかを見る。

\hypertarget{form-ux306eux5f15ux304dux623bux3057ux4ed5ux4e8bux3068ux904bux52d5ux30a8ux30cdux30ebux30aeux30fcux5b9aux7406}{%
\subsection{§4.1 1-form
の引き戻し------仕事と運動エネルギー定理}\label{form-ux306eux5f15ux304dux623bux3057ux4ed5ux4e8bux3068ux904bux52d5ux30a8ux30cdux30ebux30aeux30fcux5b9aux7406}}

\begin{quote}
\textbf{注（物理学未習の読者へ）} ここで使うのは、位置を \(x=\gamma(t)\)
と書けることだけである。\(F(x)\) は位置の関数、\(v(t)\)
は時刻の関数と読めればよい。
\end{quote}

\textbf{① 物理量から始める。}
いきなり一般論には行かない。まずは、よく知られた力学の事実から始めよう。

質点に力 \(F\) が働き、質点が \(x\) 軸上を運動しているとする。仕事は

\[W = \int F\,dx\]

で与えられる。一方、運動方程式は

\[F = m\frac{dv}{dt}\]

である。この二つから、おなじみの仕事-エネルギー定理

\[W = \frac{1}{2}mv^2\,\Big|_{t_0}^{t_1}\]

が出る。本節ではこの導出を、「測定器をどう作り替えれば同じ仕事が出るか」という目で読み直す。有限区間ごとに、時間側の測定器に必要な係数を求めるところから始める。

\textbf{② 素朴な置き換え——そのままでは合わない。} 質点の位置を
\(x = \gamma(t)\) と書く。速度は \(v(t)=\frac{d\gamma}{dt}(t)\)
である。仕事の測定器は \(F\,dx\)
------これは実空間の変位を食べて仕事の小片を返す \(1\)-form だ。

しかし運動は時刻 \(t\)
によって与えられている。そこで、仕事を時間で計算したい。

まず実空間側で、等加速度運動 \(x(t) = \frac{1}{2}at^2\)、一定力
\(F = ma\) の場合を計算する。終点 \(X = \frac{1}{2}aT^2\) だから

\[W_{\text{実}} = \int_0^X ma\,dx = ma\,X = \frac{1}{2}ma^2T^2\]

である。一方、時間側で素朴に------つまり \(dx\) をそのまま \(dt\)
に置き換えて------積分してみると

\[W_{\text{素朴}} = \int_0^T ma\,dt = ma\,T\]

となる。\(\frac{1}{2}ma^2T^2 \neq ma\,T\)
だから、明らかに一致しない。\(dt\)
のままでは仕事の値が変わってしまうのだ。

\textbf{③ 有限区間で原因を見る。} 時刻の小区間 \([t_i, t_{i+1}]\)
を考える。その幅を \(\Delta t_i = t_{i+1} - t_i\)、対応する位置の変化を
\(\Delta x_i = \gamma(t_{i+1}) - \gamma(t_i)\)
とする。実空間側の仕事の一項は
\(F_i\,\Delta x_i\)。時間側で素朴に作った一項は
\(F_i\,\Delta t_i\)。この二つが一致しないのは、\(\Delta t_i\) と
\(\Delta x_i\) のつじつま係数を無視したからだ。

そこで、時間側の一項に係数 \(c_i\) を掛けて \(F_i\,c_i\,\Delta t_i\)
としてみる。これが \(F_i\,\Delta x_i\)
と一致するには、\(c_i\,\Delta t_i = \Delta x_i\)
でなければならない。したがって

\[c_i = \frac{\Delta x_i}{\Delta t_i}\]

である。

ここまでは、ただの割り算である。けれども、この割り算には意味がある。時間側の区間
\(\Delta t_i\) に、どんな測定器を食わせれば、実空間側の変位
\(\Delta x_i\) と同じ値が出るかを探していた。その答えが

\[\frac{\Delta x_i}{\Delta t_i}\,dt\]

である。実際、この測定器に \(\Delta t_i\) を食わせると

\[\left(\frac{\Delta x_i}{\Delta t_i}\,dt\right)(\Delta t_i)=\Delta x_i\]

となる。これで、時間側の有限区間でも、実空間側と同じ変位を読めるようになった。

この有限区間上の仮の測定器を、記号で

\[\gamma_h^\square(dx):=\frac{\Delta x_i}{\Delta t_i}\,dt\]

と書くことにする。ここで \(h\)
は分割の最大幅を表す。この測定器を有限区間に食わせると、

\[\gamma_h^\square(dx)\bigl(\Delta t_i\bigr)=\Delta x_i\]

となる。

分割を細かくすると、有限比は速度

\[\frac{\Delta x_i}{\Delta t_i}\to \frac{d\gamma}{dt}(t)\]

へ近づく。したがって極限では

\[
\gamma_h^\square(dx)
\xrightarrow{h\to0}
\gamma^*(dx)
=
\frac{d\gamma}{dt}(t)\,dt
\]

が得られる。この極限後の測定器を
\(\gamma^*(dx)=\frac{d\gamma}{dt}(t)\,dt\)
と書く。有限区間で見つけた係数 \(\Delta x_i/\Delta t_i\) が、極限で速度
\(\frac{d\gamma}{dt}(t)\) に変わったのである。

以後、有限のマス目や箱に対する同じ作り替えを、まとめて
\(\Phi_h^\square\) と書く。

\begin{quote}
\textbf{注}（\(\Phi_h^\square\) と \(\Phi^*\)）

\(\gamma_h^\square\) や \(\Phi_h^\square\)
は、有限区間・有限セル上で使うための、本書内の仮の測定器である。\(h\to0\)
後の \(\Phi^*\) とは区別する。

有限セル版では、有限差分ベクトルが張る図形の測定値を、計算側のセル幅で割ったものを係数にする。したがって、この測定器を有限セルに食わせると、もとの測定値が戻る。

この係数はセルごとに定まる。記号を軽くするため、セル番号への依存は
\(\Phi_h^\square\) には書き込まない。

ここで測っているのは、写された有限セルそのものではなく、選んだ差分ベクトルが張る図形である。1次元なら変位、2次元なら平行四辺形、3次元なら平行六面体である。たとえば極座標の有限マス目の像は曲がった扇形状になるが、ここで測っているのはその扇形そのものではなく、選んだ二本の差分ベクトルが張る平行四辺形である。

分割幅 \(h\to0\) の極限で、\(\gamma_h^\square\) や \(\Phi_h^\square\)
の係数は、極限後の引き戻し \(\Phi^*\) の係数に収束する。
\end{quote}

\textbf{④ 具体例の引き戻しを書く。}
極限を取ったあと、時間軸上での測定器（引き戻された結果）は

\[\gamma^*(F\,dx) = F(\gamma(t))\,v(t)\,dt\]

である。これにより、実空間での仕事の測定器 \(F\,dx\)
を、時間軸用の測定器へと焼き直すことができた。

\textbf{⑤ 係数付き $1$-form の引き戻し。} 任意の係数 \(F(x)\)
に対して、同じ構造が成り立つ。位置が \(x = \gamma(t)\) で与えられるとき

\[
\gamma^*(F(x)\,dx)
=
F(\gamma(t))\,\frac{d\gamma}{dt}(t)\,dt
\]

である。

\textbf{⑥ この場合の変換で定義する。}
ここまでで、有限区間版から標準版へ移る道筋が見えた。そこで、滑らかな曲線
\(\gamma: t \mapsto x\) に対しては、\(1\)-form \(\omega = F(x)\,dx\)
の作り替えを

\[
\gamma^*(F(x)\,dx)
=
F(\gamma(t))\,\frac{d\gamma}{dt}(t)\,dt
\]

で定義する。この式が、時間側で使う \(1\)-form の引き戻しである。速度
\(\frac{d\gamma}{dt}(t)\) は、有限区間で見つけた
\(\Delta x_i/\Delta t_i\) の極限であり、「実空間の測定器 \(dx\)
を時間軸の測定器 \(dt\) に焼き直す」ためのつじつま係数になっている。

次は同じ手順を、2次元の面積測定器 \(dx\wedge dy\)
で繰り返す。1次元では有限区間を写して、その像の変位を測った。2次元では有限四角を写して、その像の二本の差分ベクトルが張る面積を測る。

\textbf{⑦ 物理量へ戻る。} ここで運動方程式 \(F = m\,dv/dt\) を使うと、

\[\gamma^*(F\,dx) = m\frac{dv}{dt}\,v\,dt = m v\,dv\]

と読める（\(m\frac{dv}{dt}\,dt\) は速度軸の測定器 \(dv\)
を時間軸へ引き戻した \(\gamma^*(dv)\) に対応している）。よって

\[W = \int F\,dx = \int_{\gamma} F\,dx = \int_{t_0}^{t_1} \gamma^*(F\,dx) = \int_{v(t_0)}^{v(t_1)} m v\,dv = \frac{1}{2}mv^2\,\Big|_{t_0}^{t_1}\]

となる------仕事が運動エネルギーの変化に等しい、という式である。

\begin{quote}
\textbf{注}
（なぜ「引き戻し」という名前なのか）実空間にある測定器をパラメータ空間に「送る」のなら、それは「押し出し」ではないのか------筆者はそう感じる。筆者にとって、実空間が「現実」で、パラメータ空間は「計算用の道具」だからだ。しかし数学の写像
\(\gamma\) は、つねに「パラメータ空間 \(\to\)
実空間」という向きで定義される。数学の世界では、パラメータ空間のほうが出発地なのだ。だから、実空間からパラメータ空間へ逆行する測定器の移動は、写像の向きに逆らって「引き戻し（Pull-back）」と呼ばれる。そういう命名規則なのだと割り切ることにしている。
\end{quote}

\begin{quote}
\textbf{注} （\(dx\) は横ベクトル、\(dt\) も横ベクトル）§4.1
の計算では、\(dx\) も \(dt\) も \(dv\)
も、すべてが「変位を食べてスカラーを返す」という同じ型を保っている。第1章
§1.1.3
以来の「横ベクトル×縦ベクトル→スカラー」の原則は、引き戻しのただ中でも崩れない。\(\gamma^*(F\,dx)\)
の結果も再び時間軸上の
\(1\)-form（横ベクトル）であり、それを時間側の変位 \(\Delta t\)
に食わせればスカラー（仕事の小片）が得られる------この一貫性こそが、引き戻しの理解を支える。
\end{quote}

\hypertarget{form-ux306eux5f15ux304dux623bux3057ux89d2ux904bux52d5ux91cfux4fddux5b58ux3068ux9762ux7a4dux901fux5ea6}{%
\subsection{§4.2 2-form
の引き戻し------角運動量保存と面積速度}\label{form-ux306eux5f15ux304dux623bux3057ux89d2ux904bux52d5ux91cfux4fddux5b58ux3068ux9762ux7a4dux901fux5ea6}}

\textbf{① 物理量から始める。} §4.1 では 1
次元の有限区間を写して測った。今度は同じ発想を、2
次元の有限マス目に対して行う。

\begin{quote}
\textbf{注}
（物理学未習の読者へ）以下では角運動量保存やケプラーの第二法則という言葉が出るが、本節で本当に使うのは「平面上の面積を、別の変数で測り直す」という事実である。物理の意味は後から読めばよい。
\end{quote}

中心力場の中を運動する質点を考える。力が常に原点に向かうとき、角運動量

\[
L
=
m\left(
x\frac{dy}{dt}
-
y\frac{dx}{dt}
\right)
\]

は時間に依らず一定である（角運動量保存則）。この保存則には美しい幾何学的意味がある------質点の位置ベクトルが掃く面積の速度（\textbf{面積速度}）が一定になるのだ（ケプラーの第二法則）。

質点の運動は \(z=0\) 平面上で起こるとしてよい。時刻 \(t_0\) から \(t_1\)
までに位置ベクトルが掃く扇形領域 \(D\) の面積は

\[A = \iint_D dx \wedge dy\]

と書ける。\(dx \wedge dy\) は \(xy\)
平面上の面積測定器（\(2\)-form）である。

\begin{quote}
\textbf{注}
（物理学者の二次元）ここで確認しておこう。本書は一貫して三次元空間の実在を前提としている。以下の角運動量の議論で
\(z=0\)
と置くのは、計算を見やすくするための一時的な簡略化であり、我々は依然として三次元空間の中の一つの平面を眺めているにすぎない。
\end{quote}

\textbf{② 素朴な置き換え——そのままでは合わない。} 右辺の扇形は \((x,y)\)
の式では書きづらい。そこで \((r,\theta)\)
で書きたくなる。しかし、\(dx \wedge dy\) をそのまま
\(dr \wedge d\theta\) に置き換えるとどうなるか。半径 \(C\)
の完全な円の面積を計算してみよう。正しい値は \(\pi C^2\)
である。ところが、素朴に置き換えると：

\[
A_{\text{素朴}}
=
\iint dr \wedge d\theta
=
\int_0^{2\pi}
\biggl(
  \int_0^C dr
\biggr)
d\theta
=
2\pi C
\]

これは \(\pi C^2\) とはまったく違う。\(dr \wedge d\theta\)
のままでは、正しい面積が出ない。

\textbf{③ 有限マス目で原因を見る。} 今度は、\((r,\theta)\)
側の有限マス目に何を食わせれば、\(xy\) 側の面積と同じ値になるかを見る。

\((r,\theta)\) 空間に、有限の大きさ \(\Delta r \times \Delta\theta\)
のグリッドを切る。座標 \((r, \theta)\) にある一マスが、\((x,y)\)
空間でどのような二つの変位ベクトルを作るかを求める。

\(\Delta r\) 方向の辺：\(r\) だけ変えて \(\theta\)
を固定するから、差は正確に求まる：

\[\Delta x_r = \Delta r \cos\theta,\qquad \Delta y_r = \Delta r \sin\theta\]

\(\Delta\theta\) 方向の辺：三角関数の差の公式を用いる：

\[\begin{aligned}
\Delta x_\theta &= r\bigl(\cos(\theta+\Delta\theta) - \cos\theta\bigr)
= -2r\sin\biggl(\theta+\frac{\Delta\theta}{2}\biggr)\sin\biggl(\frac{\Delta\theta}{2}\biggr) \\[4pt]
\Delta y_\theta &= r\bigl(\sin(\theta+\Delta\theta) - \sin\theta\bigr)
= 2r\cos\biggl(\theta+\frac{\Delta\theta}{2}\biggr)\sin\biggl(\frac{\Delta\theta}{2}\biggr)
\end{aligned}\]

二辺 \(\mathbf{v}_r = (\Delta x_r, \Delta y_r)\) と
\(\mathbf{v}_\theta = (\Delta x_\theta, \Delta y_\theta)\)
が張る平行四辺形の面積を行列式で求める：

\[\begin{aligned}
\det(\mathbf{v}_r,\mathbf{v}_\theta) &= \Delta x_r\,\Delta y_\theta - \Delta x_\theta\,\Delta y_r \\
&= 2r\Delta r\,\sin\biggl(\frac{\Delta\theta}{2}\biggr)\biggl[ \cos\theta\cos\biggl(\theta+\frac{\Delta\theta}{2}\biggr) + \sin\theta\sin\biggl(\theta+\frac{\Delta\theta}{2}\biggr) \biggr] \\
&= r\Delta r\,\sin\Delta\theta
\end{aligned}\]

結果は

\[\det = r\Delta r\,\sin\Delta\theta\]

である。

この測定値を、\((r,\theta)\) 側のセル幅 \(\Delta r_i\,\Delta\theta_j\)
で割ると、有限セル上で使う仮の測定器の係数が得られる。つまり

\[
\Phi_h^\square(dx\wedge dy)
:=
r_i\frac{\sin\Delta\theta_j}{\Delta\theta_j}\,dr\wedge d\theta
\]

と置く。この測定器を有限マス目に食わせると、

\[
\begin{aligned}
\Phi_h^\square(dx\wedge dy)(\Delta r_i,\Delta\theta_j)
&=
\left(
r_i\frac{\sin\Delta\theta_j}{\Delta\theta_j}\,dr\wedge d\theta
\right)(\Delta r_i,\Delta\theta_j) \\
&=
r_i\,\Delta r_i\,\sin\Delta\theta_j
\end{aligned}
\]

となる。右辺は、二本の有限差分ベクトルが張る平行四辺形の符号付き面積である。

有限セル版の面積総和を \(A_h^\square\) と置く。したがって

\[
A_h^\square
:=
\sum_i\sum_j
\Phi_h^\square(dx\wedge dy)(\Delta r_i,\Delta\theta_j)
=
\sum_i\sum_j r_i\,\Delta r_i\,\sin\Delta\theta_j
\]

である。

\(\sin\Delta\theta_j\) は有限の \(\Delta\theta_j\)
に対して厳密な値である。ここで測っているのは、極座標の有限マス目そのものの曲がった扇形面積ではなく、二つの有限差分ベクトルが張る平行四辺形の符号付き面積である。実空間側で質点の位置ベクトルが掃く扇形領域を
\(D\) と書く。一方、極座標側でその領域に対応する \((r,\theta)\) の範囲を
\(D'\) と書く。

総和の極限をとる最終段階で

\[
\frac{\sin\Delta\theta_j}{\Delta\theta_j}\to 1
\]

となる。したがって、真の面積 \(A\) は

\[
A
=
\lim_{h\to0} A_h^\square
=
\iint_{D'} r\,dr\wedge d\theta
=
\int_{\theta_0}^{\theta_1}
\biggl(
  \int_0^{R(\theta)} r\,dr
\biggr)
d\theta
=
\int_{\theta_0}^{\theta_1}
\frac{1}{2}R(\theta)^2\,d\theta
\]

として得られる。

有限セル上の測定器の係数は
\(r_i\frac{\sin\Delta\theta_j}{\Delta\theta_j}\)
である。分割を細かくすると、この係数は \(r\) へ近づく。こうして、極限で

\[
\Phi_h^\square(dx\wedge dy)
\xrightarrow{h\to0}
\Phi^*(dx\wedge dy)
=
r\,dr\wedge d\theta
\]

が現れる。

\begin{quote}
\textbf{注} （\(\Delta r,\Delta\theta\)
は微小である必要はない）ここまでの差分ベクトルの計算に、\(\Delta r\) や
\(\Delta\theta\)
が「十分小さい」という仮定は一切使っていない。上で求めた
\(\Delta x_r, \Delta y_r, \Delta x_\theta, \Delta y_\theta\)
は、\(\Delta r,\Delta\theta\)
がどれほど大きくても厳密に成り立つ式である。ここで「厳密」と言っているのは、有限差分ベクトルが張る平行四辺形を測るという意味である。微小量への移行は、\(\Phi^*\)
と積分へ移る最終段階------総和の極限をとるとき------で初めて行う。
\end{quote}

\textbf{④ 極限で得られた標準形を書く。} 極座標への変換
\(\Phi(r,\theta) = (r\cos\theta,\; r\sin\theta)\) に対して：

\[\Phi^*(dx \wedge dy) = r\,dr \wedge d\theta\]

これが、有限セル版 \(\Phi_h^\square\) から \(h\to0\) で得られる
\(2\)-form の作り替えである。この極限で得られた作り替えを \(\Phi^*\)
と書く。この \(r\) が、§4.1 の速度 \(d\gamma/dt\)
に対応するつじつま係数である。

\textbf{⑤ 係数付き $2$-form の引き戻し。} 任意の係数 \(G(x,y)\) を持つ
\(2\)-form \(G(x,y)\,dx \wedge dy\) に対して、同じ構造が成り立つ：

\[\Phi^*\bigl(G(x,y)\,dx \wedge dy\bigr) = G(r\cos\theta,\, r\sin\theta)\; r\,dr \wedge d\theta\]

係数 \(G\) は座標を極座標に焼き直すだけ（\(0\)-form
の引き戻し）で、\(dx \wedge dy\) の部分だけが \(r\,dr \wedge d\theta\)
に変わる。引き戻し \(\Phi^*\)
は「係数の焼き直し」と「測定器の焼き直し」に自然に分解されるのである。

\textbf{⑥ 同じ発見を一般の変換へ移す。}
極座標だけに特別なことはない。平面間の変換 \(\Phi(u,v)=(x(u,v),y(u,v))\)
を考え、\((u,v)\) 側の有限四角を写す。\(u\) 方向と \(v\)
方向の二本の差分ベクトルを \(dx\wedge dy\)
で測ると、有限四角ごとの面積値が出る。

その面積値を \(\Delta u\,\Delta v\)
で割り、分割を細かくしていく。すると、差分比は偏微分係数へ移り、\(2\times2\)
の行列式が現れる。したがって、\(h\to0\) 後の作り替えは

\[\Phi^*\bigl(G(x,y)\,dx \wedge dy\bigr) = G(x(u,v), y(u,v))\;
\det\!\begin{pmatrix}
\frac{\partial x}{\partial u} & \frac{\partial x}{\partial v} \\[6pt]
\frac{\partial y}{\partial u} & \frac{\partial y}{\partial v}
\end{pmatrix}
\,du \wedge dv\]

となる。つまり

\[\Phi^*\bigl(G(x,y)\,dx \wedge dy\bigr)
=G(\Phi(u,v))\left(
\frac{\partial x}{\partial u}\frac{\partial y}{\partial v}
-\frac{\partial x}{\partial v}\frac{\partial y}{\partial u}
\right)\,du\wedge dv\]

である。ここで現れる \(2\times2\) 行列式が、実空間の面積測定器
\(dx\wedge dy\) を、\((u,v)\) 側で使える測定器 \(du\wedge dv\)
に焼き直すためのつじつま係数である。

有限マス目で見た

\[\Delta x_u\,\Delta y_v-\Delta x_v\,\Delta y_u\]

は、この偏微分行列式を発見するための有限セル版である。有限差分ベクトルが張る平行四辺形については等式だが、\(\Phi^*\)
と書く極限後の作り替えでは偏微分係数を使う。有限セル上の測定値を
\(\Delta u\,\Delta v\) で割った係数は、分割幅 \(h\) を \(0\)
に近づけると、上の \(2\times2\) 行列式へ収束する。

\textbf{⑦ 物理量へ戻る。} 扇形領域 \(D\)
の面積を、引き戻した測定器で対応する極座標側の領域 \(D'\)
上で積分すれば：

\[
A
=
\iint_D dx\wedge dy
=
\iint_{D'} \Phi^*(dx\wedge dy)
=
\iint_{D'} r\,dr\wedge d\theta
\]

\[
=
\int_{\theta_0}^{\theta_1}
\biggl(
  \int_0^{R(\theta)} r\,dr
\biggr)
d\theta
=
\int_{\theta_0}^{\theta_1}
\frac{1}{2}R(\theta)^2\,d\theta
\]

\(r\) で積分した結果、\(\frac{1}{2}r^2\) が現れ、面積が \(\theta\)
軸上の \(1\)-form の線積分に還元された。さらに時間パラメータ \(t\)
へ引き戻す。質点の軌道を \(r = R(t),\, \theta = \Theta(t)\) とすれば：

\[
\frac{dA}{dt}
=
\frac{1}{2}r(t)^2\,\frac{d\theta}{dt}(t)
\]

角運動量の定義

\[
L
=
m r^2\frac{d\theta}{dt}
\]

より、面積速度は

\[\frac{dA}{dt} = \frac{L}{2m}\]

中心力場では \(L\)
が保存されるから、面積速度も一定------これがケプラーの第二法則であり、角運動量保存則の幾何学的表現である。§4.1
で \(1\)-form のつじつま係数 \(d\gamma/dt\)
がエネルギー保存を記述したのとまったく同じ構造が、\(2\)-form
のつじつま係数 \(r\) を通じて角運動量保存を記述している。

\hypertarget{form-ux306eux5f15ux304dux623bux3057ux8ceaux91cfux4fddux5b58ux3068ux4f53ux7a4dux5206}{%
\subsection{§4.3 3-form
の引き戻し------質量保存と体積分}\label{form-ux306eux5f15ux304dux623bux3057ux8ceaux91cfux4fddux5b58ux3068ux4f53ux7a4dux5206}}

\textbf{① 物理量から始める。} 最後は体積分である。今度は 3
次元の有限箱を写して測る。

密度 \(\rho\) が一様な円柱（半径 \(C\)、高さ
\(H\)）の質量を考えよう。\((x,y,z)\) 空間で直接計算すれば

\[M = \rho \cdot (\text{体積}) = \rho \cdot \pi C^2 H\]

である。この \(M\) は、どんな座標で計算しても変化してはならない。

\textbf{② 素朴な置き換え——そのままでは合わない。} 積分を円柱座標
\((r,\theta,z)\) で書きたい。円柱の範囲は
\(r \in [0, C],\; \theta \in [0, 2\pi],\; z \in [0, H]\)
だから、\(dx \wedge dy \wedge dz\) をそのまま
\(dr \wedge d\theta \wedge dz\) に置き換えて：

\[
M_{\text{素朴}}
=
\rho \iiint dr\wedge d\theta\wedge dz
=
\rho
\int_0^H
\biggl(
  \int_0^{2\pi}
  \biggl(
    \int_0^C dr
  \biggr)
  d\theta
\biggr)
dz
=
2\pi\rho C H
\]

これを正しい値 \(\rho \pi C^2 H\)
と比べると、\(2\pi\rho C H \neq \rho \pi C^2 H\)。明らかに一致しない。

\textbf{③ 有限マス目で原因を見る。} 3次元では、\((r,\theta,z)\)
側の有限の箱に何を食わせれば、\((x,y,z)\)
側の体積と同じ値になるかを探す。

\((r,\theta,z)\) 空間に、有限の大きさ
\(\Delta r \times \Delta\theta \times \Delta z\)
のグリッドを切る。変換は
\(x = r\cos\theta,\; y = r\sin\theta,\; z = z\)。\(z\)
方向は変換を受けないから、体積は「\(xy\)
平面内のマス目面積」\(\times \Delta z\) である。\(xy\)
平面内の面積は、§4.2 の手順そのままだ：

\[\det = r\Delta r\,\sin\Delta\theta\]

したがって、一マスの体積は

\[\text{体積} = \bigl(r\Delta r\,\sin\Delta\theta\bigr) \times \Delta z\]

円柱座標の有限セルそのものは \(\theta\)
方向に曲がっているが、ここで測っているのは、三本の差分ベクトルが張る平行六面体の符号付き体積である。

この測定値を、\((r,\theta,z)\) 側のセル幅
\(\Delta r_i\,\Delta\theta_j\,\Delta z_k\)
で割ると、有限セル上で使う仮の測定器の係数が得られる。つまり

\[
\Phi_h^\square(dx\wedge dy\wedge dz)
:=
r_i\frac{\sin\Delta\theta_j}{\Delta\theta_j}\,dr\wedge d\theta\wedge dz
\]

と置く。この測定器を有限箱に食わせると、

\[
\begin{aligned}
&\Phi_h^\square(dx\wedge dy\wedge dz)
(\Delta r_i,\Delta\theta_j,\Delta z_k)\\
&\qquad =
r_i\frac{\sin\Delta\theta_j}{\Delta\theta_j}
\Delta r_i\Delta\theta_j\Delta z_k \\
&\qquad =
r_i\Delta r_i\sin\Delta\theta_j\Delta z_k
\end{aligned}
\]

となる。右辺は、三本の有限差分ベクトルが張る平行六面体の符号付き体積である。

有限セル版の質量総和を \(M_h^\square\) と置くと、

\[
M_h^\square
:=
\sum_i\sum_j\sum_k
\rho\,
\Phi_h^\square(dx\wedge dy\wedge dz)(\Delta r_i,\Delta\theta_j,\Delta z_k)
=
\sum_i\sum_j\sum_k
\rho\,r_i\,\Delta r_i\,\sin\Delta\theta_j\,\Delta z_k
\]

である。総和の極限をとる最終段階で

\[
\frac{\sin\Delta\theta_j}{\Delta\theta_j}\to 1
\]

となる。したがって、真の質量 \(M\) は

\[
M
=
\lim_{h\to0} M_h^\square
=
\rho\iiint_{V'} r\,dr\wedge d\theta\wedge dz
=
\rho
\int_0^H
\biggl(
  \int_0^{2\pi}
  \biggl(
    \int_0^C r\,dr
  \biggr)
  d\theta
\biggr)
dz
=
\rho\cdot\pi C^2H
\]

として得られる。

有限セル上の測定器の係数は
\(r_i\frac{\sin\Delta\theta_j}{\Delta\theta_j}\) である。\(h\to0\)
で、この係数はまた \(r\) へ近づく。

\textbf{④ 極限で得られた標準形を書く。} 円柱座標への変換
\(\Phi(r,\theta,z) = (r\cos\theta,\; r\sin\theta,\; z)\) に対して：

\[
\Phi_h^\square(dx\wedge dy\wedge dz)
\xrightarrow{h\to0}
\Phi^*(dx \wedge dy \wedge dz)
=
r\,dr \wedge d\theta \wedge dz
\]

これが、有限セル版 \(\Phi_h^\square\) から \(h\to0\) で得られる
\(3\)-form の作り替えである。この極限で得られた作り替えを \(\Phi^*\)
と書く。この \(r\) が、§4.2 の \(r\)
をそのまま引き継いだつじつま係数である。

\textbf{⑤ 係数付き $3$-form の引き戻し。} 任意の密度 \(\rho(x,y,z)\)
を持つ \(3\)-form \(\rho\,dx \wedge dy \wedge dz\)
に対して、同じ構造が成り立つ：

\[\Phi^*\bigl(\rho(x,y,z)\,dx \wedge dy \wedge dz\bigr) = \rho(r\cos\theta,\, r\sin\theta,\, z)\; r\,dr \wedge d\theta \wedge dz\]

係数 \(\rho\) は座標を焼き直すだけ（\(0\)-form
の引き戻し）で、\(dx \wedge dy \wedge dz\) の部分だけが
\(r\,dr \wedge d\theta \wedge dz\) に変わる。

\textbf{⑥ 同じ発見を一般の変換へ移す。} より一般に、空間の変換
\(\Phi(u,v,w) = (x(u,v,w), y(u,v,w), z(u,v,w))\) を考える。\((u,v,w)\)
側の有限箱を写すと、三本の差分ベクトルができる。それを
\(dx\wedge dy\wedge dz\)
で測れば、写された平行六面体の向き付き体積が出る。

その体積値を \(\Delta u\,\Delta v\,\Delta w\)
で割り、分割を細かくしていく。すると差分比は偏微分係数へ移り、\(3\times3\)
の行列式が現れる。したがって、\(h\to0\) 後の引き戻しは：

\[
\begin{aligned}
&\Phi^*\bigl(\rho(x,y,z)\,dx \wedge dy \wedge dz\bigr) \\
&= \rho(x(u,v,w), y(u,v,w), z(u,v,w))\,
\det\!\begin{pmatrix}
\frac{\partial x}{\partial u} & \frac{\partial x}{\partial v} & \frac{\partial x}{\partial w} \\[6pt]
\frac{\partial y}{\partial u} & \frac{\partial y}{\partial v} & \frac{\partial y}{\partial w} \\[6pt]
\frac{\partial z}{\partial u} & \frac{\partial z}{\partial v} & \frac{\partial z}{\partial w}
\end{pmatrix} \\
&\qquad\qquad {}\wedge du \wedge dv \wedge dw
\end{aligned}
\]

この \(3 \times 3\) 行列式を \(J(u,v,w)\) と書けば

\[\Phi^*(\rho\,dx \wedge dy \wedge dz) = \rho(\Phi(u,v,w))\,J\,du \wedge dv \wedge dw\]

となる。この \(3\times3\) 行列式 \(J(u,v,w)\) が、\(3\)-form
のつじつま係数である。

有限マス目で三個の差分ベクトルを並べれば、

\[\det\!\begin{pmatrix}
\Delta x_u & \Delta x_v & \Delta x_w \\[2pt]
\Delta y_u & \Delta y_v & \Delta y_w \\[2pt]
\Delta z_u & \Delta z_v & \Delta z_w
\end{pmatrix}\]

が出る。これは有限セル上の測定器を有限箱に食わせた値であり、差分ベクトルが張る平行六面体については等式である。分割幅
\(h\) を \(0\) に近づけると、この有限行列式を
\(\Delta u\,\Delta v\,\Delta w\)
で割った係数が、偏微分からなるヤコビアン \(J\) へ収束する。

\(J\)
が負なら向きが反転していることを意味する。形式の引き戻しでは、この符号を持った
\(J\)
をそのまま使う。一方、「正の体積」や密度を積分するときには、向きを忘れた大きさが必要になるため
\(|J|\) が現れる。ここで \(J\) と \(|J|\) を混同しないことが重要である。

\(1\)-form なら \(1 \times 1\) 行列式（速度 \(dx/dt\)）、\(2\)-form なら
\(2 \times 2\) 行列式、\(3\)-form なら \(3 \times 3\)
行列式（\(J\)）------次数が上がるごとに、つじつま係数を司る行列式のサイズが一つずつ大きくなる。上の偏微分行列式の式こそが
\textbf{$3$-form の引き戻しの一般的な定義} である。

\begin{quote}
\textbf{注} （ベクトル解析を学んだ読者へ）変数変換で「極座標にすると
\(r\,dr\,d\theta\) が出る」と習ったとき、多くの教科書は「\(\theta\)
方向の弧の長さ \(r\Delta\theta\)
を考えて、\(\Delta V \approx \Delta r \cdot r\Delta\theta \cdot \Delta z\)」という幾何学的な説明をする。その説明に間違いはない。しかし本書の計算では弧の長さを一切考えていない。変換式と三角関数の公式、そして
\(\wedge\) に相当する行列式の計算だけで、つじつま係数 \(r\) が出てくる。
\end{quote}

\textbf{⑦ 物理量へ戻る。} 引き戻した測定器で積分すれば

\[
M
=
\rho\iiint_V dx\wedge dy\wedge dz
=
\rho\iiint_{V'} r\,dr\wedge d\theta\wedge dz
\]

\[
=
\rho
\int_0^H
\biggl(
  \int_0^{2\pi}
  \biggl(
    \int_0^C r\,dr
  \biggr)
  d\theta
\biggr)
dz
=
\rho\cdot\pi C^2H
\]

正しい値に一致した。積分値を保つために、測定器
\(dx \wedge dy \wedge dz\) は \((r,\theta,z)\) 空間では
\(r\,dr \wedge d\theta \wedge dz\)
に作り替えられなければならなかったのである。§4.1 の速度 \(d\gamma/dt\)
がエネルギー保存を、§4.2 の \(r\)
が角運動量保存を記述したのとまったく同じ構造で、§4.3 の \(r\)（一般には
\(J\)）が質量保存を記述している。

\hypertarget{ux5f15ux304dux623bux3057ux306eux6027ux8cea-ux3053ux3053ux307eux3067ux306bux78baux7acbux3057ux305fux3053ux3068}{%
\subsection{§4.4 引き戻しの性質 ------
ここまでに確立したこと}\label{ux5f15ux304dux623bux3057ux306eux6027ux8cea-ux3053ux3053ux307eux3067ux306bux78baux7acbux3057ux305fux3053ux3068}}

§4.1〜§4.3
で三つの具体例を通じて引き戻しを構成してきた。ここで、これらに共通する代数的な性質を整理しておく。あわせて、§4.3
で現れた \(3 \times 3\) 行列式 \(J\) を
\textbf{ヤコビアン（ヤコビ行列式）} と定義する。もととなる行列

\[\begin{pmatrix}
\frac{\partial x}{\partial u} & \frac{\partial x}{\partial v} & \frac{\partial x}{\partial w} \\[4pt]
\frac{\partial y}{\partial u} & \frac{\partial y}{\partial v} & \frac{\partial y}{\partial w} \\[4pt]
\frac{\partial z}{\partial u} & \frac{\partial z}{\partial v} & \frac{\partial z}{\partial w}
\end{pmatrix}\]

を \textbf{ヤコビ行列} と呼び、その行列式 \(\det\) がヤコビアン \(J\)
である。

\textbf{① 0-form（スカラー場）の引き戻し。} スカラー場 \(f(x,y,z)\)
を変換 \(\Phi\)
で引き戻すとは、単に座標をパラメータ表示に焼き直すことである：

\[\Phi^*(f) = f(\Phi(u,v,w))\]

たとえば §4.1 で \(F(x)\,dx\) の \(F(x)\) を \(F(\gamma(t))\)
に焼き直した操作は、\(F\) という \(0\)-form の引き戻しにほかならない。

\textbf{② 線形性。} 二つの \(k\)-form \(\omega, \eta\) とスカラー
\(a,b\) に対し

\[\Phi^*(a\omega + b\eta) = a\,\Phi^*(\omega) + b\,\Phi^*(\eta)\]

が成り立つ。これは測定器の和や定数倍が、引き戻しの前後で整合することを保証する。

\textbf{③ ウェッジ積との整合性。} \(k\)-form \(\omega\) と \(\ell\)-form
\(\eta\) に対し

\[\Phi^*(\omega \wedge \eta) = \Phi^*(\omega) \wedge \Phi^*(\eta)\]

が成り立つ。すなわち「組み立ててから引き戻す」のと「引き戻してから組み立てる」のは同じ結果を与える。§4.2
で \(dx\) と \(dy\) を個別に展開してから \(\wedge\)
で組み立てた操作が、まさにこの性質を実行していた。

\textbf{④ 次数ごとのつじつま係数。} 有限セルから \(h\to0\)
の極限を取った後の引き戻し \(\Phi^*\) では、変換 \(\Phi\)
の偏微分を並べたヤコビ行列から、必要な \(k\) 個の方向を取り出し、その
\(k \times k\)
行列式がつじつま係数となる。本書の範囲では、次の形で現れる。

\begin{itemize}
\tightlist
\item
  \(k=1\)（\(1\)-form）：\(1 \times 1\) 行列式
\end{itemize}

\[
\det
\begin{pmatrix}
\frac{dx}{dt}
\end{pmatrix}
=
\frac{dx}{dt}
\]

\begin{itemize}
\tightlist
\item
  \(k=2\)（\(2\)-form）：\(2 \times 2\) 行列式
\end{itemize}

\[
\det
\begin{pmatrix}
\frac{\partial x}{\partial u} & \frac{\partial x}{\partial v}\\[4pt]
\frac{\partial y}{\partial u} & \frac{\partial y}{\partial v}
\end{pmatrix}
=
\frac{\partial x}{\partial u}\frac{\partial y}{\partial v}
-
\frac{\partial x}{\partial v}\frac{\partial y}{\partial u}
\]

\begin{itemize}
\tightlist
\item
  \(k=3\)（\(3\)-form）：\(3 \times 3\) 行列式
\end{itemize}

\[
J
=
\det
\begin{pmatrix}
\frac{\partial x}{\partial u} & \frac{\partial x}{\partial v} & \frac{\partial x}{\partial w}\\[4pt]
\frac{\partial y}{\partial u} & \frac{\partial y}{\partial v} & \frac{\partial y}{\partial w}\\[4pt]
\frac{\partial z}{\partial u} & \frac{\partial z}{\partial v} & \frac{\partial z}{\partial w}
\end{pmatrix}
\]

こうして、\(1\)-form では \(1\times1\) 行列式、\(2\)-form では
\(2\times2\) 行列式、\(3\)-form では \(3\times3\)
行列式が現れる。次数が上がるごとに、測定器のつじつま係数を決める行列式のサイズが一つずつ大きくなる。この構造は、より高い次元の
\(k\)-form でも同じである。変換式を微分して得られる係数を並べ、必要な
\(k\) 本の方向を取り出して行列式を作る。その行列式が、\(k\)
次元の測定器を計算側の変数で使える形へ焼き直すためのつじつま係数になる。

\begin{quote}
\textbf{注}
（ヤコビアンという呼び名）一般には、変数変換の微分係数を並べた行列をヤコビ行列、その行列式をヤコビアンと呼ぶ。\(1\)
次元なら \(\frac{dx}{dt}\)、\(2\) 次元なら上の \(2\times2\)
行列式、\(3\) 次元なら \(J\) がそれにあたる。

本書では、特に体積換算率として現れる \(3\times3\) 行列式を \(J\)
と書き、ヤコビアンと呼ぶことが多い。
\end{quote}

\textbf{⑤ 有限セル版と極限後の引き戻しの関係。}
本章の導出では、いきなり偏微分を書かなかった。まず §4.1
で、有限区間上の測定器の作り替えを \(\gamma_h^\square\)
と書いた。そこでは有限比 \(\Delta x/\Delta t\)
が区間ごとの係数になり、\(h\to0\) で \(\gamma^*\) へ移った。

§4.2 と §4.3 では、この考えを 2 次元・3 次元へ拡張した。有限セル版
\(\Phi_h^\square\)
は、有限セル上で使う仮の測定器である。係数は、「有限差分ベクトルが張る図形の測定値
\(\div\)
計算側のセル幅」で定める。この測定器を有限セルに食わせると、もとの測定値が戻る。

たとえば \(2\)-form では、二本の差分ベクトルが張る平行四辺形の面積を
\(\Delta u\,\Delta v\) で割ったものが係数になる。\(3\)-form
では、三本の差分ベクトルが張る平行六面体の体積を
\(\Delta u\,\Delta v\,\Delta w\) で割ったものが係数になる。

分割幅 \(h\to0\)
の極限で、これらの係数は偏微分からなる行列式へ収束する。これが極限後の引き戻し
\(\Phi^*\)
の係数である。したがって本章の流れは、「有限セルで仮の測定器を作る →
\(h\to0\) で \(\Phi^*\) を得る」という二段階になっている。

\begin{quote}
\textbf{注} （外微分 \(d\)
との可換性------第5章への伏線）これらに加えて、引き戻しと外微分 \(d\)
のあいだには重要な関係 \[\Phi^*(d\omega) = d(\Phi^*\omega)\]
が成り立つ。すなわち「微分してから引き戻す」のと「引き戻してから微分する」のは同じ結果を与える。この性質は第5章で
\(d\) を正式に導入したのち、§5.9 で改めて取り上げる。
\end{quote}

これ以上深く立ち入ることはしない。後の章では、曲がった空間そのものを扱う言葉へ進む。そのときにも、本章で見た有限セルと行列式の経験は足場になる。

\begin{center}\rule{0.5\linewidth}{0.5pt}\end{center}

\hypertarget{ux672cux7ae0ux306eux307eux3068ux3081ux3068ux7b2c5ux7ae0ux5916ux5faeux5206ux3078ux306eux5c55ux671b}{%
\subsection{§4.5
本章のまとめと第5章（外微分）への展望}\label{ux672cux7ae0ux306eux307eux3068ux3081ux3068ux7b2c5ux7ae0ux5916ux5faeux5206ux3078ux306eux5c55ux671b}}

ここまでの四章で、我々は次の四つをそろえた：

\begin{enumerate}
\def\labelenumi{\arabic{enumi}.}
\tightlist
\item
  \textbf{第1章}：\(dx\)
  を行列（\(1\)-form）と見なし、積分を行列作用の極限と読む
\item
  \textbf{第2章}：ウェッジ積 \(\wedge\) で \(2\)-form・\(3\)-form
  を構成し、面積・体積を代数で測る
\item
  \textbf{第3章}：曲線・曲面・領域に形式を適用し、その出力を集計する操作を確立した
\item
  \textbf{第4章（本章）}：引き戻しにより「積分値を保つために測定器を作り替える」という統一原理を、仕事-エネルギー定理・角運動量保存・質量保存の三つの具体例を通じて確立した
\end{enumerate}

本章で確立したのは、次の統一原理である：

\textbf{別の変数で数えるときは、積分値が変わらないように測定器を作り替える。}これが引き戻しである。曲線上の仕事も、曲面上の面積や流量も、領域内の総質量も、この同じ考え方で扱える。

引き戻しは「積分値を保つために測定器を作り替える」技術であり、その核心は、§4.1〜§4.3
で見た発見的導出にある。すなわち、「素朴にやって合わない →
有限セルで仮の測定器を作る → \(h\to0\)
で一般形を得る」という流れである。有限セル版 \(\Phi_h^\square\)
では、有限差分ベクトルの行列式をセル幅で割ったものが係数になり、\(\Phi^*\)
ではその極限として偏微分係数やヤコビアンが現れる。\(1\)-form では速度
\(d\gamma/dt\)、\(2\)-form では \(r\)、\(3\)-form では \(3\times3\)
行列式 \(J\) が、有限セルを測る計算の中から現れた。

\begin{quote}
\textbf{【ここまでのチェックポイント——第4章】}

\begin{itemize}
\item
  引き戻し \(\Phi^*\)
  は、積分値が変わらないように、実空間側の測定器を計算側の変数で使える形へ焼き直す操作である。
\item
  \(1\)-form では、実空間の測定器 \(F(x)\,dx\) が、時間側では
  \(\gamma^*(F(x)\,dx)=F(\gamma(t))\,\frac{d\gamma}{dt}(t)\,dt\)
  に作り替えられる。有限区間では \(\Delta x/\Delta t\)、極限では速度
  \(d\gamma/dt\) がつじつま係数になる。
\item
  \(2\)-form では、\(dx\wedge dy\) が極座標側で
  \(\Phi^*(dx\wedge dy)=r\,dr\wedge d\theta\)
  に作り替えられる。有限マス目では \(r\Delta r\sin\Delta\theta\)
  が現れ、\(h\to0\) の極限で係数 \(r\) が残る。
\item
  \(3\)-form では、\(dx\wedge dy\wedge dz\) が円柱座標側で
  \(\Phi^*(dx\wedge dy\wedge dz)=r\,dr\wedge d\theta\wedge dz\)
  に作り替えられる。一般の変換では
  \(\Phi^*(dx\wedge dy\wedge dz)=J\,du\wedge dv\wedge dw\) となる。
\item
  有限セル版 \(\Phi_h^\square\)
  は、有限セル上で使う仮の測定器である。その係数は「有限差分ベクトルが張る図形の測定値
  \(\div\)
  計算側セル幅」で定める。この測定器を有限セルに食わせると、もとの測定値が戻る。\(h\to0\)
  で、この係数が \(\Phi^*\) の係数になる。
\item
  \(1\)-form では \(1\times1\) 行列式、\(2\)-form では \(2\times2\)
  行列式、\(3\)-form では \(3\times3\)
  行列式が、測定器を作り替えるためのつじつま係数になる。
\item
  第1章 §1.4 の \(dx = \cos\theta\,dr - r\sin\theta\,d\theta\)
  は、すでに \(1\)-form の引き戻しの原型だった。
\end{itemize}
\end{quote}

\begin{center}\rule{0.5\linewidth}{0.5pt}\end{center}

次章では、いよいよ\textbf{外微分 $d$}
を導入する。測定器を作り、積分し、変数に合わせて作り替えるところまで来た。次に見るのは、その測定器自身の変化である。場所ごとの変化を、一つ上の次数の形式として記録する操作------それが外微分
\(d\) である。

\newpage

\hypertarget{ux7b2c5ux7ae0ux5faeux5206ux3059ux308bux3068ux306fux4f55ux304b-ux5916ux5faeux5206-dux7a4dux5206ux91cfux3068ux5c40ux6240ux6cd5ux5247}{%
\chapter{\texorpdfstring{第5章：微分するとは何か ------ 外微分
\(d\)：積分量と局所法則}{第5章：微分するとは何か ------ 外微分 d：積分量と局所法則}}\label{ux7b2c5ux7ae0ux5faeux5206ux3059ux308bux3068ux306fux4f55ux304b-ux5916ux5faeux5206-dux7a4dux5206ux91cfux3068ux5c40ux6240ux6cd5ux5247}}

\hypertarget{ux5faeux5206ux5074ux3078ux306eux6a4bux89b3ux6e2cux306fux7a4dux5206ux6cd5ux5247ux306fux5faeux5206}{%
\subsection{§5.0
微分側への橋------観測は積分、法則は微分}\label{ux5faeux5206ux5074ux3078ux306eux6a4bux89b3ux6e2cux306fux7a4dux5206ux6cd5ux5247ux306fux5faeux5206}}

第3章で我々は、曲線・曲面・領域上の積分を \(k\)-form
の言葉で定義した。曲線 \(\gamma\) に沿った仕事は
\(\int_\gamma \omega\)、曲面 \(S\) を貫く流量は \(\iint_S \eta\)、領域
\(V\) の総質量は \(\iiint_V \Omega\)
である。いずれも「\(k\)-form（測定器）に \(k\)
個の変位ベクトルを食わせてスカラーを得、領域全体で集計する」という同一の原理に従っていた。

第4章では、同じ積分を別の変数で計算するために、測定器を作り替える操作------引き戻し
\(\Phi^*\)------を導入した。これで、積分値を保ったまま計算側の変数を変える道具も手に入った。

ここで我々が手にした積分は、すべて\textbf{領域全体にわたる}量である（以下こうした量を\textbf{大域的}と呼ぶ）。仕事は曲線全体にわたる総和であり、流量は曲面全体を貫く総量であり、質量は領域全体での総計だ。それらは測定器の当て方と領域の選び方に依存する。

しかし、物理学者が最終的に欲しがるのは、これだけではない。積分量として観測される法則の背後に、各点で成り立つ局所的な関係を見たい。マクスウェル方程式にせよナビエ--ストークス方程式にせよ、自然界の基本法則は\textbf{空間の各点・各瞬間において何が成り立つか}を記述する\textbf{局所的}な関係------微分方程式------として書かれる。大域的な積分量は領域に依存するが、局所的な法則は領域に依存しない普遍性を持つからだ。

\begin{quote}
\textbf{注}
（大域的／局所的という用語）他書では似た対比を\textbf{マクロ／ミクロ}と呼ぶこともある。ただし本章では、スケールの大小というより、領域全体に積分された量と、各点で成り立つ関係の対比を指して、大域的／局所的という語を使う。
\end{quote}

ここに根本的な問いが生まれる。

\begin{quote}
\textbf{積分された量から、どのようにして局所法則を取り出すのか。}
\end{quote}

この問いに答えるのが本章の主題------\textbf{外微分 $d$}------である。\(d\)
は \(k\)-form に作用して \((k+1)\)-form
を返す演算子であり、第3章で定義した積分と組み合わせることで「境界で測った大域的な積分」を「内部の局所的な変化の集積」へと翻訳する。いわば、\textbf{$d$ は積分法則を微分法則に変換する装置}なのである。

本章の構成はこうだ。まず第1章で導入した \(df\) を思い出し、一般の
\(1\)-form を閉ループで測ったときに残るズレから \(d\omega\)
を作る。次に同じ考えを \(2\)-form
まで拡張し、ストークスの定理とガウスの定理を一つの形へまとめる。最後に
\(d^2=0\)
の構造と、外微分が物理法則をどのように局所化するかを見て、次章のホッジ・スターへつなぐ。

\begin{center}\rule{0.5\linewidth}{0.5pt}\end{center}

\hypertarget{df-ux518dux8a2aux5faeux5206ux3068ux7a4dux5206ux306fux9006ux6f14ux7b97ux304b}{%
\subsection{\texorpdfstring{§5.1 \(df\)
再訪------微分と積分は逆演算か}{§5.1 df 再訪------微分と積分は逆演算か}}\label{df-ux518dux8a2aux5faeux5206ux3068ux7a4dux5206ux306fux9006ux6f14ux7b97ux304b}}

\hypertarget{df-ux306fux6b21ux5143ux4e0aux3052ux3092ux3057ux3066ux3044ux308b}{%
\subsubsection{\texorpdfstring{5.1.1 \(df\)
は「次元上げ」をしている}{5.1.1 df は「次元上げ」をしている}}\label{df-ux306fux6b21ux5143ux4e0aux3052ux3092ux3057ux3066ux3044ux308b}}

第1章 §1.2 で、我々は関数 \(f(x,y,z)\)
の全微分を横ベクトルとして定義した：

\[df = \frac{\partial f}{\partial x}\,dx + \frac{\partial f}{\partial y}\,dy + \frac{\partial f}{\partial z}\,dz = \begin{pmatrix} \frac{\partial f}{\partial x} & \frac{\partial f}{\partial y} & \frac{\partial f}{\partial z} \end{pmatrix}\]

これは
\textbf{$0$-form（スカラー場 $f$）を入力として $1$-form（横ベクトル $df$）を出力する}操作である。次数が
0 から 1 へ上がった。

ここで思い出してほしいのは、\(df(\mathbf{v})\) が変位 \(\mathbf{v}\)
に対する \(f\)
の一次変化量を返す、ということだ。\(\mathbf{v} = \begin{pmatrix}\Delta x\\\Delta y\\\Delta z\end{pmatrix}\)
が十分小さいとき：

\[df(\mathbf{v}) = \frac{\partial f}{\partial x}\,\Delta x + \frac{\partial f}{\partial y}\,\Delta y + \frac{\partial f}{\partial z}\,\Delta z \approx f(\mathbf{r}+\mathbf{v}) - f(\mathbf{r})\]

\(df\) 単独では「変化量そのもの」ではない------第1章 §1.1.6
の契約どおり、\(df\)
は演算子であり、変位ベクトルに作用して初めて数値が得られる。

\hypertarget{ux7ddaux7a4dux5206ux3059ux308bux3068ux7aefux70b9ux306eux5deeux3060ux3051ux304cux6b8bux308b}{%
\subsubsection{5.1.2
線積分すると端点の差だけが残る}\label{ux7ddaux7a4dux5206ux3059ux308bux3068ux7aefux70b9ux306eux5deeux3060ux3051ux304cux6b8bux308b}}

この \(df\) を曲線 \(\gamma\) に沿って積分する。第3章 §3.3.1
の定義どおり、曲線を小区間に刻み、各ステップの変位
\(\Delta\mathbf{r}_i\) に \(df\) を食わせて足し上げる。

各区間で
\(df(\Delta\mathbf{r}_i) \approx f(\mathbf{r}_i) - f(\mathbf{r}_{i-1})\)
だから、全区間の和は：

\[\sum_{i=1}^n df(\Delta\mathbf{r}_i) \approx \sum_{i=1}^n \bigl(f(\mathbf{r}_i) - f(\mathbf{r}_{i-1})\bigr)\]

この和を展開すると：

\[\bigl(f(\mathbf{r}_1)-f(\mathbf{r}_0)\bigr) + \bigl(f(\mathbf{r}_2)-f(\mathbf{r}_1)\bigr) + \cdots + \bigl(f(\mathbf{r}_n)-f(\mathbf{r}_{n-1})\bigr)\]

隣り合う項
\(f(\mathbf{r}_1), f(\mathbf{r}_2), \dots, f(\mathbf{r}_{n-1})\)
がプラスとマイナスで打ち消し合う------\textbf{望遠和（telescoping sum）}だ。生き残るのは最初と最後だけ。刻みを無限に細かくする極限で：

\[\int_\gamma df = f(\mathbf{r}_n) - f(\mathbf{r}_0) = f(B) - f(A)\]

ここで \(A\) は曲線の始点、\(B\) は終点である。

\begin{quote}
\textbf{【ここまでのチェックポイント】}

\begin{itemize}
\item
  \(\int_\gamma df = f(B) - f(A)\)。積分の結果は経路の形状に一切依存せず、端点の値だけで決まる。
\item
  これは \(df\) という特別な \(1\)-form
  だからこそ成立する。力学でいえば「保存力の仕事は経路によらない」という事実の数学的根拠だ。
\end{itemize}
\end{quote}

\hypertarget{ux7b2c3ux7ae0ux306eux4f8bux3067ux78baux304bux3081ux308b}{%
\subsubsection{5.1.3
第3章の例で確かめる}\label{ux7b2c3ux7ae0ux306eux4f8bux3067ux78baux304bux3081ux308b}}

第3章 §3.4.1 で我々は、\(\omega = y\,dx + x\,dy\) を単位円
\(\gamma(t) = (\cos t,\; \sin t,\; 0),\; t \in [0,2\pi]\)
に沿って積分し、結果が \(0\)
になることを計算した。あのときは係数に沿って愚直に積分したが、いまの言葉で見直すと話はもっと単純だ。

実は \(y\,dx + x\,dy\) は \(d(xy)\) にほかならない。なぜなら：

\[d(xy) = \frac{\partial (xy)}{\partial x}\,dx + \frac{\partial (xy)}{\partial y}\,dy + \frac{\partial (xy)}{\partial z}\,dz = y\,dx + x\,dy\]

したがって \(\omega = df\)（\(f = xy\)）であり、§5.1.2
の望遠和の理屈がそのまま使える。単位円は閉曲線だから \(B = A\)、よって：

\[\oint_\gamma \omega = \oint_\gamma d(xy) = xy\Big|_A^A = \cos 0\cdot\sin 0 - \cos 0\cdot\sin 0 = 0\]

あのときの \(\int_0^{2\pi} \cos 2t\,dt = 0\) という計算は、実は
\(\omega\) が \(d(xy)\)
という\textbf{完全形式}だったからこそ、機械的にゼロになったのである。

\begin{quote}
\textbf{注} （完全形式）\(\omega = df\) と書ける \(1\)-form
を\textbf{完全形式}（exact
form）と呼ぶ。完全形式を閉曲線に沿って積分すると、\(f\)
が一価関数である限り必ずゼロになる。一方「ある閉曲線で積分がゼロになったからといって、その形式が完全形式だとは限らない」という点には注意が必要である（すべての閉曲線でゼロになるかどうかは領域の位相とも関係する）。この点は
§5.2 で詳しく見る。
\end{quote}

\begin{quote}
\textbf{注} （\(d\) と \(\int\)
の関係------逆写像ではなく境界公式）\(\int_\gamma df = f(B)-f(A)\)
は、積分が外微分 \(d\)
の大域的な逆写像である、という意味ではない。積分は曲線 \(\gamma\)
の選び方に依存する汎関数であり、\(df\) から元の関数 \(f\)
全体を復元する操作ではない。この式は「完全形式を曲線に沿って積分すると、端点という境界だけが残る」という微積分学の基本定理である。後で見る
Stokes 型公式の \(0\)-form 版だ、と考えるのが安全である。
\end{quote}

\begin{center}\rule{0.5\linewidth}{0.5pt}\end{center}

\hypertarget{ux9589ux30ebux30fcux30d7ux3067ux59ffux3092ux73feux3059ux30baux30ec}{%
\subsection{§5.2
閉ループで姿を現す「ズレ」}\label{ux9589ux30ebux30fcux30d7ux3067ux59ffux3092ux73feux3059ux30baux30ec}}

\hypertarget{ux4e00ux822cux306e-1-form-ux3067ux306fux3069ux3046ux304b}{%
\subsubsection{\texorpdfstring{5.2.1 一般の \(1\)-form
ではどうか}{5.2.1 一般の 1-form ではどうか}}\label{ux4e00ux822cux306e-1-form-ux3067ux306fux3069ux3046ux304b}}

では、一般の \(1\)-form

\[\omega = P(x,y,z)\,dx + Q(x,y,z)\,dy + R(x,y,z)\,dz\]

ではどうか。ここで \(P, Q, R\)
は場所ごとに変わる係数（スカラー場）であり、第3章 §3.4.1
で導入した「一般の \(1\)-form」である。

\begin{quote}
\textbf{注} （「一般の」とはどういうことか）§5.1 で扱った \(df\)
は、係数が
\(\frac{\partial f}{\partial x}, \frac{\partial f}{\partial y}, \frac{\partial f}{\partial z}\)
という特別な組み合わせだった------\(f\)
という一つの関数から偏微分で導かれた三つ組である。一方、ここでの「一般の」\(\omega\)
は \(P, Q, R\)
が\textbf{互いに無関係な、任意の三つのスカラー場}である。言い換えれば、\(\omega = df\)
と書けるとは限らない \(1\)-form のことを「一般の
\(1\)-form」と呼んでいる。§5.1.3 の \(y\,dx + x\,dy\) はたまたま
\(d(xy)\)
と書ける完全形式だったが、そうでないもののほうがむしろ普通なのである。
\end{quote}

\(\omega\) は力を表すとは限らず、任意の「線に沿って測る測定器」を表す。

この \(\omega\) に対して、§5.1
と同じ望遠和が成り立つ保証はない。それを判定する最も簡潔な方法は\textbf{閉ループ}------始点と終点が同じ曲線------で積分してみることだ。もし
\(\int_\gamma \omega\) が端点の値の差だけで書けるなら、閉ループでは
\(f(A)-f(A)=0\) とならなければならない。

閉ループに沿った線積分を \(\oint_\gamma \omega\) と書く。一般には：

\[\oint_\gamma \omega \neq 0\]

である。摩擦力に逆らって物体を一周させれば仕事はゼロにならず、渦巻く水流に沿って一周すれば正味の仕事を受ける。この\textbf{「閉じても帳尻が合わない」という事実こそが、循環や渦のような局所構造の証拠}だ。

\begin{quote}
\textbf{注}
（物理を知らなくても大丈夫）摩擦力や水流をまだ物理学で扱っていない読者も、ここで立ち止まる必要はない。いま重要なのは具体的な物理現象ではなく、「閉ループの積分がゼロにならないなら、その場には循環のような局所構造がある」という感覚だけである。物理での詳しい対応は後の章で触れる。
\end{quote}

\hypertarget{ux5c40ux6240ux60c5ux5831ux306fux3069ux3053ux306bux5bbfux308bux304b}{%
\subsubsection{5.2.2
局所情報はどこに宿るか}\label{ux5c40ux6240ux60c5ux5831ux306fux3069ux3053ux306bux5bbfux308bux304b}}

では、この「一周して残るズレ」はどこから来るのか。閉ループをいきなり大域的な大きさで考えていては、中のどの点がどれだけズレに寄与したのかわからない。第1章で
\(f'(x)\) を求めたときのことを思い出そう------我々は
\(\Delta f / \Delta x\) の \(\Delta x \to 0\)
の極限をとることで、\textbf{一点での}変化率を取り出した。

同じ発想をここでも使う。ループをどんどん小さくしていったとき、一周の積分
\(\oint \omega\) の値はどうなるか。もし \(\oint \omega\)
の主項がループの\textbf{囲む面積}に比例して小さくなるなら、面積あたりのズレ------単位面積あたりどれだけ帳尻が合わないか------という量が、その点に固有の値として決まるはずだ。

逆に、もし比例しなければ（たとえばループの周の長さに比例するなら）、面積あたりの量としては決まらず、「その点での性質」とは呼べない。つまり問題はこうだ：

\begin{quote}
\textbf{微小ループを一周したとき、$\oint \omega$ の主項は囲んだ面積に比例するのか。比例するなら、その比例係数は何か。}
\end{quote}

次節で実際に、\(xy\) 平面に置いた微小長方形でこの問いに答える。

\begin{center}\rule{0.5\linewidth}{0.5pt}\end{center}

\hypertarget{ux5faeux5c0fux30ebux30fcux30d7ux306eux89e3ux4f53ux30baux30ecux306fux9762ux7a4dux306bux6bd4ux4f8bux3059ux308b}{%
\subsection{§5.3
微小ループの解体------ズレは面積に比例する}\label{ux5faeux5c0fux30ebux30fcux30d7ux306eux89e3ux4f53ux30baux30ecux306fux9762ux7a4dux306bux6bd4ux4f8bux3059ux308b}}

\hypertarget{xy-ux5e73ux9762ux306eux5faeux5c0fux9577ux65b9ux5f62}{%
\subsubsection{\texorpdfstring{5.3.1 \(xy\)
平面の微小長方形}{5.3.1 xy 平面の微小長方形}}\label{xy-ux5e73ux9762ux306eux5faeux5c0fux9577ux65b9ux5f62}}

空間内の点 \((x, y, z)\) を左下の頂点とし、\(xy\)
平面に平行な微小長方形を考える。\(x\) 方向の幅を \(\Delta x\)、\(y\)
方向の幅を \(\Delta y\)
とする。この長方形の四辺を\textbf{反時計回り（右手系）}に一周し、\(\omega = P\,dx + Q\,dy + R\,dz\)
を積分する。

各辺での積分を、テイラー展開の1次近似（\(\Delta x, \Delta y\)
が十分小さいとして）で評価しよう。

\begin{quote}
\textbf{注}
（テイラー展開について）「テイラー展開」という言葉に身構える必要はない。ここで使うのは多変数関数の1次近似------第1章
§1.2.1 で \(\Delta f \approx f'(x)\,\Delta x\)
とやったのと同じことだ。たとえば \(Q(x+\Delta x, y)\) を \(x\) 方向に
\(\Delta x\) だけずらした値は、\(Q\) の \(x\) に関する偏微分
\(\frac{\partial Q}{\partial x}\) を使って
\(Q + \frac{\partial Q}{\partial x}\,\Delta x\)
と近似できる。\(\Delta x\) が十分小さければ、\(\Delta x\)
の2次以上の項（\((\Delta x)^2\) やそれ以上）は \(\Delta x\)
に比べて圧倒的に小さいため無視してよい------これが1次近似の意味である。
\end{quote}

\(z\) は一定なので \(dz=0\)、\(R\) の項は寄与しない。

\textbf{辺1（下辺、右向き）}：\((x, y)\) から \((x+\Delta x, y)\)
へ。\(y\) 一定なので \(dy=0\)。\(P\) はおよそ \(P(x, y, z)\)。寄与は
\(P(x, y, z)\,\Delta x\)。

\textbf{辺2（右辺、上向き）}：\((x+\Delta x, y)\) から
\((x+\Delta x, y+\Delta y)\) へ。\(dx=0\)。\(Q\) を \(x\) 方向に
\(\Delta x\) だけずれた位置で評価する：
\[Q(x+\Delta x, y, z) \approx Q(x, y, z) + \frac{\partial Q}{\partial x}\Delta x\]
寄与は
\(\bigl(Q + \frac{\partial Q}{\partial x}\Delta x\bigr)\,\Delta y\)。

\textbf{辺3（上辺、左向き）}：\((x+\Delta x, y+\Delta y)\) から
\((x, y+\Delta y)\) へ。負の方向なので符号が反転。\(P\) を \(y\) 方向に
\(\Delta y\) だけずれた位置で：
\[P(x, y+\Delta y, z) \approx P(x, y, z) + \frac{\partial P}{\partial y}\Delta y\]
寄与は
\(-\bigl(P + \frac{\partial P}{\partial y}\Delta y\bigr)\,\Delta x\)。

\textbf{辺4（左辺、下向き）}：\((x, y+\Delta y)\) から \((x, y)\)
へ戻る。寄与は \(-Q(x, y, z)\,\Delta y\)。

四辺の寄与を合計する：

\[\begin{aligned}
\oint \omega &\approx P\Delta x + (Q + \frac{\partial Q}{\partial x}\Delta x)\Delta y - (P + \frac{\partial P}{\partial y}\Delta y)\Delta x - Q\Delta y \\
&= P\Delta x + Q\Delta y + \frac{\partial Q}{\partial x}\Delta x\Delta y - P\Delta x - \frac{\partial P}{\partial y}\Delta x\Delta y - Q\Delta y \\
&= (\frac{\partial Q}{\partial x} - \frac{\partial P}{\partial y})\,\Delta x\,\Delta y
\end{aligned}\]

\(P\Delta x\) と \(Q\Delta y\) の項が美しく相殺し、残ったのは
\((\frac{\partial Q}{\partial x} - \frac{\partial P}{\partial y})\,\Delta x\,\Delta y\)
だけだ。

\hypertarget{ux6c7aux5b9aux7684ux306aux4e8bux5b9f}{%
\subsubsection{5.3.2
決定的な事実}\label{ux6c7aux5b9aux7684ux306aux4e8bux5b9f}}

この結果が告げていることはただ一つ：

\begin{quote}
\textbf{一周して残るズレの主項は、囲んだ面積 $\Delta x\,\Delta y$ に比例する。}
\end{quote}

そして比例係数
\(\frac{\partial Q}{\partial x} - \frac{\partial P}{\partial y}\)
こそが、点 \((x,y,z)\)
における\textbf{単位面積あたりのズレ}------その点での「渦の強さ」を表す局所量------である。より正確には、微小長方形
\(R_{\Delta x,\Delta y}\) について

\[\lim_{\Delta x,\Delta y\to 0}\frac{1}{\Delta x\,\Delta y}\oint_{\partial R_{\Delta x,\Delta y}}\omega = \frac{\partial Q}{\partial x} - \frac{\partial P}{\partial y}\]

という極限で取り出される量である。

これは第1章で \(f'(x) = \lim_{\Delta x\to 0} \frac{\Delta f}{\Delta x}\)
とやったのと同じ発想の「微分」だ。分母が「動いた距離」から「囲んだ面積」に変わっただけで、\textbf{変化率を局所的に取り出す}という精神は変わっていない。

\begin{quote}
\textbf{注} （なぜ
\(\frac{\partial Q}{\partial x}-\frac{\partial P}{\partial y}\)
だけが残るか）\(\frac{\partial P}{\partial x}\) や
\(\frac{\partial Q}{\partial y}\) は出てこない。\(P\) の \(x\)
方向の変化（\(\frac{\partial P}{\partial x}\)）は下辺と上辺で同じ向きに効くため相殺し、\(Q\)
の \(y\)
方向の変化（\(\frac{\partial Q}{\partial y}\)）も右辺と左辺で相殺する。生き残るのは「\(P\)
の \(y\) 方向の変化」と「\(Q\) の \(x\)
方向の変化」------互いに\textbf{直交する方向への偏微分の差}だけだ。この交叉性が、次節以降のすべての鍵になる。
\end{quote}

\begin{center}\rule{0.5\linewidth}{0.5pt}\end{center}

\hypertarget{d-ux306eux8a95ux751fux9762ux7a4dux3042ux305fux308aux306eux30baux30ecux3092ux6e2cux308bux65b0ux3057ux3044ux6e2cux5b9aux5668}{%
\subsection{\texorpdfstring{§5.4 \(d\)
の誕生------面積あたりのズレを測る新しい測定器}{§5.4 d の誕生------面積あたりのズレを測る新しい測定器}}\label{d-ux306eux8a95ux751fux9762ux7a4dux3042ux305fux308aux306eux30baux30ecux3092ux6e2cux308bux65b0ux3057ux3044ux6e2cux5b9aux5668}}

\hypertarget{ux30a6ux30a7ux30c3ux30b8ux7a4dux306eux518dux78baux8a8d}{%
\subsubsection{5.4.1
ウェッジ積の再確認}\label{ux30a6ux30a7ux30c3ux30b8ux7a4dux306eux518dux78baux8a8d}}

§5.3 で得た
\((\frac{\partial Q}{\partial x} - \frac{\partial P}{\partial y})\,\Delta x\,\Delta y\)
を、第2章の言葉で言い直そう。第2章 §2.4.4 で我々は、\(dx \wedge dy\)
という面積測定器を定義した。これは2本のベクトル
\(\mathbf{v}_1, \mathbf{v}_2\) を食わせると、それらが \(xy\)
平面に落とす影の符号付き面積を返す装置だった：

\[(dx \wedge dy)(\mathbf{v}_1, \mathbf{v}_2) = \det\begin{pmatrix} dx(\mathbf{v}_1) & dx(\mathbf{v}_2) \\ dy(\mathbf{v}_1) & dy(\mathbf{v}_2) \end{pmatrix}\]

特に、\(x\) 方向の変位 \(\Delta x\,\hat{e}_x\) と \(y\) 方向の変位
\(\Delta y\,\hat{e}_y\) を食わせれば
\((dx \wedge dy)(\Delta x\,\hat{e}_x,\; \Delta y\,\hat{e}_y) = \Delta x\,\Delta y\)
である。

\hypertarget{domega-ux306eux5b9aux7fa9}{%
\subsubsection{\texorpdfstring{5.4.2 \(d\omega\)
の定義}{5.4.2 d\textbackslash omega の定義}}\label{domega-ux306eux5b9aux7fa9}}

§5.3 の結果をこの言葉で書けば、微小長方形の二辺
\(\Delta x\,\hat{e}_x,\; \Delta y\,\hat{e}_y\)
に対して、ある\textbf{新しい $2$-form} が
\((\frac{\partial Q}{\partial x} - \frac{\partial P}{\partial y})\,\Delta x\,\Delta y\)
という主項を返したことになる。この \(2\)-form を \(\omega\)
の\textbf{外微分}と呼び、\(d\omega\) と書く：

\[d\omega := (\frac{\partial Q}{\partial x} - \frac{\partial P}{\partial y})\,dx \wedge dy\]

\(d\omega\) は \(2\)-form
である------2本のベクトルを食べてスカラーを返す。その値は「食わせた2本が張る微小平行四辺形に対する、単位面積あたりの閉ループのズレ」を表す。

ここで起きていることを整理しよう：

\begin{itemize}
\item
  \textbf{入力}：\(\omega = P\,dx + Q\,dy + R\,dz\)（\(1\)-form、ベクトル1本を食べる線の測定器）
\item
  \textbf{出力}：\(d\omega = (\frac{\partial Q}{\partial x} - \frac{\partial P}{\partial y})\,dx \wedge dy\)（\(2\)-form、ベクトル2本を食べる面の測定器）
\item
  \textbf{操作}：\(d\) は次数を 1 から 2 へ\textbf{一つ上げる}
\end{itemize}

これは §5.1 で見た \(df\) とまったく同じ構造だ。\(d\) は \(0\)-form
\(f\) に作用して \(1\)-form \(df\) を返し、\(1\)-form \(\omega\)
に作用して \(2\)-form \(d\omega\)
を返す。\textbf{$d$ の本質は「次数を1つ上げる」こと}にある。

\begin{quote}
\textbf{【ここまでのチェックポイント】}

\begin{itemize}
\item
  \(df\) は \(0\)-form → \(1\)-form
  の外微分だった。\(\int_\gamma df = f(B)-f(A)\)（端点のみに依存）。
\item
  一般の \(\omega\) では閉ループで
  \(\oint \omega \neq 0\)。このズレを調べるためにループを微小化する。
\item
  \(xy\) 平面の微小長方形では
  \(\oint \omega = (\frac{\partial Q}{\partial x} - \frac{\partial P}{\partial y})\,\Delta x\,\Delta y +\)
  高次の項。ズレの主項は面積に比例する。
\item
  この比例係数を測る \(2\)-form を
  \(d\omega := (\frac{\partial Q}{\partial x} - \frac{\partial P}{\partial y})\,dx \wedge dy\)
  と定義する。\(d\) は次数を上げる演算子。
\end{itemize}
\end{quote}

\begin{center}\rule{0.5\linewidth}{0.5pt}\end{center}

\hypertarget{ux4e00ux822cux306e-1-form-ux306eux5916ux5faeux52063ux6b21ux5143ux3078ux306eux62e1ux5f35}{%
\subsection{\texorpdfstring{§5.5 一般の \(1\)-form
の外微分------3次元への拡張}{§5.5 一般の 1-form の外微分------3次元への拡張}}\label{ux4e00ux822cux306e-1-form-ux306eux5916ux5faeux52063ux6b21ux5143ux3078ux306eux62e1ux5f35}}

\hypertarget{yz-ux5e73ux9762ux3068-zx-ux5e73ux9762}{%
\subsubsection{\texorpdfstring{5.5.1 \(yz\) 平面と \(zx\)
平面}{5.5.1 yz 平面と zx 平面}}\label{yz-ux5e73ux9762ux3068-zx-ux5e73ux9762}}

§5.3--§5.4 では \(xy\) 平面のループだけを考えた。しかし
3次元空間には面の向きが3種類ある。\(yz\)
平面の微小長方形では同様の計算で
\((\frac{\partial R}{\partial y} - \frac{\partial Q}{\partial z})\,\Delta y\,\Delta z\)
が、\(zx\) 平面では
\((\frac{\partial P}{\partial z} - \frac{\partial R}{\partial x})\,\Delta z\,\Delta x\)
が得られる。それぞれ \(dy \wedge dz\) と \(dz \wedge dx\) に対応する。

直感から導かれる完全な式はこうだ：

\[d\omega = (\frac{\partial R}{\partial y} - \frac{\partial Q}{\partial z})\,dy \wedge dz + (\frac{\partial P}{\partial z} - \frac{\partial R}{\partial x})\,dz \wedge dx + (\frac{\partial Q}{\partial x} - \frac{\partial P}{\partial y})\,dx \wedge dy\]

\begin{quote}
\textbf{注} （\(2\)-form は傾いた平行四辺形を測る）この \(d\omega\) は
\(2\)-form------\textbf{2本のベクトルを食べる測定器}------である。食わせる2本は、\(xy\)
平面に平行とは限らない。3次元空間の中に斜めに浮かぶ微小平行四辺形の二辺
\(\mathbf{v}_1, \mathbf{v}_2\)
を食わせれば、その面に対する単位面積あたりの閉ループのズレを返す。§5.3
でやったのとまったく同じこと------測定する面が座標平面に平行とは限らなくなっただけだ。2本のベクトルから平行四辺形の面積（と向き）を取り出すしくみは、第2章
§2.4 で \(dy \wedge dz\)
等を反対称行列として構成したときにすでに確かめてある。
\end{quote}

\hypertarget{ux4ee3ux6570ux7684ux5c0eux51faux30e9ux30a4ux30d7ux30cbux30c3ux30c4ux5247ux3068-ddx0}{%
\subsubsection{\texorpdfstring{5.5.2 代数的導出------ライプニッツ則と
\(d(dx)=0\)}{5.5.2 代数的導出------ライプニッツ則と d(dx)=0}}\label{ux4ee3ux6570ux7684ux5c0eux51faux30e9ux30a4ux30d7ux30cbux30c3ux30c4ux5247ux3068-ddx0}}

上記の式は、各平面で別々にループ計算をしても当然導ける。しかし、毎回ループを描いて四辺を足すのは面倒だ。ここでやることは、§5.3
で幾何的に導入した外微分を計算用の代数ルールに焼き直し、これが確かに外微分を再現することの確認である。すなわち、\(xy\)平面の例から得た直感を使って、機械的に計算できる形へ移そう。

\begin{quote}
\textbf{注} （公理的定義との関係）多くの数学書では、ライプニッツ則と
\(d(dx)=0\) を \(d\)
の\textbf{定義}として置くことも多い。その定義は簡潔だが、なぜそのルールになるのかは見えにくい。本書では、微小ループの解体から出発する。これから確かめるのは、その幾何的な
\(d\)
と代数的な計算ルールが整合することだ。ルールは計算の便宜、\textbf{意味は §5.3 にある}。もっとも、高次元まで見据えれば、この代数的定義が非常に簡潔であることも確かである。
\end{quote}

まず、\(\omega = P\,dx + Q\,dy + R\,dz\) に対し、\(d\omega\) を
\(d(P\,dx) + d(Q\,dy) + d(R\,dz)\)
であることは問題ないだろう。これは傾いた平行四辺形のループを考えるということだ。問題は
\(d(P\,dx)\) ------係数 \(P\) のついた \(1\)-form に \(d\)
をどう作用させるか------である。

ここで、§5.3--§5.4 で得た幾何的な \(d\)
がすでに教えてくれていることがある。\(\omega = P\,dx\)（\(Q=R=0\)）の場合、§5.3
の微小ループ計算を \(xy\) 平面だけでなく \(zx\)
平面でも行えば、\(d(P\,dx)\) は \(yz\) 面に成分を持たず、

\[d(P\,dx) = \frac{\partial P}{\partial z}\,(dz \wedge dx) - \frac{\partial P}{\partial y}\,(dx \wedge dy)\]

となるはずだ（\(xy\) 面のループでは
\((0 - \frac{\partial P}{\partial y}) = -\frac{\partial P}{\partial y}\)、\(zx\)
面では
\((\frac{\partial P}{\partial z} - 0) = \frac{\partial P}{\partial z}\)）。

さて、ここで
\(dP = \frac{\partial P}{\partial x}\,dx + \frac{\partial P}{\partial y}\,dy + \frac{\partial P}{\partial z}\,dz\)
を思い出そう（§5.1）。これと \(dx\) のウェッジ積をとると：

\[dP \wedge dx = (\frac{\partial P}{\partial x}\,dx + \frac{\partial P}{\partial y}\,dy + \frac{\partial P}{\partial z}\,dz) \wedge dx = \frac{\partial P}{\partial z}\,(dz \wedge dx) - \frac{\partial P}{\partial y}\,(dx \wedge dy)\]

（\(dx \wedge dx = 0\) により \(\frac{\partial P}{\partial x}\)
の項は消えた）

幾何的計算と\textbf{完全に一致する}。すなわち、少なくとも \(P\,dx\)
という形の項について：

\[d(P\,dx) = dP \wedge dx\]

が成り立つ。つまり、係数 \(P\) の場所による変化だけが、\(dx\)
方向の測定器に新しい面積測定器として貼りつく。

幾何的結果とライプニッツ則 \(d(P\,dx) = dP \wedge dx + P\,d(dx)\)
を見比べれば、余分な項 \(P\,d(dx)\) はゼロでなければならない。\(P\)
は任意の関数だから、これが常に成り立つためには \(d(dx)=0\)
でなければならない------すなわち \(d(dx) = 0\) が要請される。

この観察を一般化すると、次の\textbf{次数付きライプニッツ則}に到達する。\(P\)
は \(0\)-form、\(dx\) は \(1\)-form であり、\(0\)-form と \(1\)-form
の「積」に対する \(d\) の振る舞いは：

\[d(P\,dx) = (dP) \wedge dx + P\,d(dx)\]

\(dP\) は既知------§5.1 より
\(dP = \frac{\partial P}{\partial x}\,dx + \frac{\partial P}{\partial y}\,dy + \frac{\partial P}{\partial z}\,dz\)。問題は
\(d(dx)\) の値だ。

ここで、座標関数 \(x\) の外微分は \(dx\)
である（\(d(x) = dx\)）。デカルト座標の \(dx\)
は、場所によって係数が変わらない基準の測定器である。つまり \(dx\)
自身を微小ループで一周させても、面積あたりのズレは出てこない。同じことが
\(dy,dz\)
にも成り立つ。したがって座標基底については次を計算ルールとして置く：

\[d(dx) = d(dy) = d(dz) = 0\]

すると：

\[d(P\,dx) = (\frac{\partial P}{\partial x}\,dx + \frac{\partial P}{\partial y}\,dy + \frac{\partial P}{\partial z}\,dz) \wedge dx + P\,0\]

\(\wedge\)
の反対称性（\(dx \wedge dx = 0\)、\(dy \wedge dx = -dx \wedge dy\)）を使って：

\[d(P\,dx) = \frac{\partial P}{\partial x}\,(dx \wedge dx) + \frac{\partial P}{\partial y}\,(dy \wedge dx) + \frac{\partial P}{\partial z}\,(dz \wedge dx) = \frac{\partial P}{\partial z}\,(dz \wedge dx) - \frac{\partial P}{\partial y}\,(dx \wedge dy)\]

同様に：

\[\begin{aligned}
d(Q\,dy) &= (\frac{\partial Q}{\partial x}\,dx + \frac{\partial Q}{\partial y}\,dy + \frac{\partial Q}{\partial z}\,dz) \wedge dy = \frac{\partial Q}{\partial x}\,(dx \wedge dy) - \frac{\partial Q}{\partial z}\,(dy \wedge dz) \\
d(R\,dz) &= (\frac{\partial R}{\partial x}\,dx + \frac{\partial R}{\partial y}\,dy + \frac{\partial R}{\partial z}\,dz) \wedge dz = -\frac{\partial R}{\partial x}\,(dz \wedge dx) + \frac{\partial R}{\partial y}\,(dy \wedge dz)
\end{aligned}\]

三つを足し合わせ、基底 \(dy \wedge dz,\; dz \wedge dx,\; dx \wedge dy\)
の順に整理する（この巡回順は次章の右手系との整合のための布石である）：

\[\begin{aligned}
d\omega &= (-\frac{\partial Q}{\partial z} + \frac{\partial R}{\partial y})\,dy \wedge dz + (\frac{\partial P}{\partial z} - \frac{\partial R}{\partial x})\,dz \wedge dx + (-\frac{\partial P}{\partial y} + \frac{\partial Q}{\partial x})\,dx \wedge dy \\
&= (\frac{\partial R}{\partial y} - \frac{\partial Q}{\partial z})\,dy \wedge dz + (\frac{\partial P}{\partial z} - \frac{\partial R}{\partial x})\,dz \wedge dx + (\frac{\partial Q}{\partial x} - \frac{\partial P}{\partial y})\,dx \wedge dy
\end{aligned}\]

これは §5.3 で幾何的に導いた \(xy\) 成分
\((\frac{\partial Q}{\partial x}-\frac{\partial P}{\partial y})\)
を含み、\(yz\) と \(zx\) の成分も同じ交叉パターンで揃っている。

\hypertarget{ux4fc2ux6570ux306eux30d1ux30bfux30fcux30f3}{%
\subsubsection{5.5.3
係数のパターン}\label{ux4fc2ux6570ux306eux30d1ux30bfux30fcux30f3}}

係数の並び方に注目してほしい。\((\frac{\partial R}{\partial y} - \frac{\partial Q}{\partial z},\; \frac{\partial P}{\partial z} - \frac{\partial R}{\partial x},\; \frac{\partial Q}{\partial x} - \frac{\partial P}{\partial y})\)
は、\((P, Q, R)\)
の偏微分を\textbf{自分以外の軸方向に関して交叉させ、差をとった}ものである。巡回的な対称性がある：

\[\begin{aligned}
dy \wedge dz &\longleftrightarrow \frac{\partial R}{\partial y} - \frac{\partial Q}{\partial z} \\
dz \wedge dx &\longleftrightarrow \frac{\partial P}{\partial z} - \frac{\partial R}{\partial x} \\
dx \wedge dy &\longleftrightarrow \frac{\partial Q}{\partial x} - \frac{\partial P}{\partial y}
\end{aligned}\]

\begin{quote}
\textbf{注}
（既習者への補助線）ベクトル解析を既習の読者には、この3つの係数の並びに見覚えがあるだろう。本書ではその名前に依存せず、まず
\(d\) という操作そのものを組み立てる。対応関係は後の章で整理する。
\end{quote}

\begin{quote}
\textbf{【ここまでのチェックポイント】}

\begin{itemize}
\item
  一般の \(1\)-form \(\omega\) の外微分 \(d\omega\) は \(2\)-form
  であり、係数は偏微分の交叉差
  \((\frac{\partial R}{\partial y}-\frac{\partial Q}{\partial z},\; \frac{\partial P}{\partial z}-\frac{\partial R}{\partial x},\; \frac{\partial Q}{\partial x}-\frac{\partial P}{\partial y})\)。
\item
  計算ルールは (1) 線形性、(2) ライプニッツ則
  \(d(P\,dx) = dP \wedge dx + P\,d(dx)\)、(3) \(d(dx)=d(dy)=d(dz)=0\)
  の三つ。
\item
  結果は §5.3 の微小ループ計算と完全に整合する。
\end{itemize}
\end{quote}

\begin{center}\rule{0.5\linewidth}{0.5pt}\end{center}

\hypertarget{ux96c6ux7a4dux3059ux308cux3070ux5883ux754cux3060ux3051ux304cux6b8bux308bux30b9ux30c8ux30fcux30afux30b9ux306eux5b9aux7406}{%
\subsection{§5.6
集積すれば境界だけが残る------ストークスの定理}\label{ux96c6ux7a4dux3059ux308cux3070ux5883ux754cux3060ux3051ux304cux6b8bux308bux30b9ux30c8ux30fcux30afux30b9ux306eux5b9aux7406}}

\hypertarget{ux9762ux3092ux30bfux30a4ux30ebux8cbcux308aux3059ux308b}{%
\subsubsection{5.6.1
面をタイル貼りする}\label{ux9762ux3092ux30bfux30a4ux30ebux8cbcux308aux3059ux308b}}

§5.3 では\textbf{ひとつの}微小長方形の周りで \(\oint \omega\)
を計算し、その主項が \(d\omega\)
と面積の積になることを見た。では、この微小長方形を\textbf{面全体に敷き詰めて}、全部足し合わせたらどうなるか。

曲面 \(S\)
を小さな座標パッチに分け、それぞれのパラメータ平面上で微小長方形に刻む（第3章
§3.2
の面素の考え方そのものだ）。各微小長方形について、面の向きに合う向きで境界を一周した
\(\omega\) の積分は、面素に \(d\omega\) を食わせた値を主項として持つ。

ここでは、曲面はなめらかにパラメータ表示でき、境界にも自然な向きが誘導される場合を考える。面の向きを逆にすれば、境界をたどる向きも同時に逆になる。この「面の向きと境界の向きの組」を固定して初めて、左右の符号が一致する。

\hypertarget{ux5185ux90e8ux306eux8fbaux306fux76f8ux6bbaux3059ux308b}{%
\subsubsection{5.6.2
内部の辺は相殺する}\label{ux5185ux90e8ux306eux8fbaux306fux76f8ux6bbaux3059ux308b}}

ここで決定的な幾何学的観察がある。\textbf{隣り合う二つの長方形が共有する辺}に注目しよう。左の長方形にとってその辺は「上向き」、右の長方形にとっては「下向き」にたどられる。経路が逆向きだから、\(\omega\)
の線積分はこの辺の上で\textbf{完全に打ち消し合う}。

この相殺は面 \(S\)
の内部の\textbf{すべての共有辺}で起こる。いくら細かく分割して足し合わせても、内部の辺の寄与はプラスマイナスで消えるのだ。

\hypertarget{ux751fux304dux6b8bux308bux306eux306fux5883ux754cux3060ux3051}{%
\subsubsection{5.6.3
生き残るのは境界だけ}\label{ux751fux304dux6b8bux308bux306eux306fux5883ux754cux3060ux3051}}

すると、最終的に相殺されずに生き残る辺はどこか------\textbf{隣の長方形が存在しない辺}、すなわち面
\(S\) の\textbf{境界 $\partial S$} に沿った辺だけである。

つまり：

\[\oint_{\partial S} \omega = \iint_S d\omega\]

左辺は「面の境界という1次元の閉曲線に沿った \(\omega\)
の線積分」、右辺は「面の内部に敷き詰めた無数の微小ループのズレ密度の総和」である。分割を細かくする極限で、内部辺の相殺だけが残り、両者が一致する。

これが\textbf{ケルビン–ストークスの定理}（単にストークスの定理とも呼ばれる）だ。

\(\omega = P\,dx + Q\,dy + R\,dz\) の成分をあらわに書けば：

\[\oint_{\partial S} \bigl(P\,dx + Q\,dy + R\,dz\bigr) = \iint_S \Bigl( (\frac{\partial R}{\partial y} - \frac{\partial Q}{\partial z})\,dy \wedge dz + (\frac{\partial P}{\partial z} - \frac{\partial R}{\partial x})\,dz \wedge dx + (\frac{\partial Q}{\partial x} - \frac{\partial P}{\partial y})\,dx \wedge dy \Bigr)\]

第2章以来の行列の言葉で書くなら、右辺の \(2\)-form は反対称行列
\(\mathbf{J}^T - \mathbf{J}\) で表される。\(\mathbf{v}_1, \mathbf{v}_2\)
を面素の二辺とすれば
\((d\omega)(\mathbf{v}_1, \mathbf{v}_2) = \mathbf{v}_1^T(\mathbf{J}^T - \mathbf{J})\mathbf{v}_2\)。付録B.3に全成分を書き下してあるので、必要に応じて参照されたい。

この定理の構造は驚くほど単純だ。\textbf{大域的な境界で測った積分} \(=\)
\textbf{内部の局所的なズレ（$d\omega$）の集積}。\(d\) とは、境界
\(\partial S\) における \(\omega\) の情報を、内部 \(S\) における
\(d\omega\) の情報へと\textbf{翻訳}する装置なのである。

\begin{quote}
\textbf{注}
（局所化に使っている仮定）ここでのタイル貼りの説明は、十分なめらかな曲面と場を、十分細かい座標パッチに分けられる状況を前提にしている。境界の尖り、自己交差、特異点などがある場合には、パッチ分割や向きの扱いを別途整える必要がある。本書では物理で通常使う滑らかな場合を扱う。
\end{quote}

\begin{quote}
\textbf{注} （第3章との対応）第3章 §3.2.1
で曲面積分を定義したとき、我々は面素
\((\mathbf{r}_u\Delta u,\;\mathbf{r}_v\Delta v)\) に \(dx \wedge dy\)
を食わせていた。ここでも同じく、曲面を小さな面素に分け、各面素で測定器を作用させ、内部で相殺する寄与を整理している。
\end{quote}

\begin{quote}
\textbf{【ここまでのチェックポイント】}

\begin{itemize}
\item
  面 \(S\) を微小ループで敷き詰めると、内部の辺は相殺し、境界
  \(\partial S\) の寄与だけが残る。
\item
  \(\int_{\partial S} \omega = \int_S d\omega\)。大域的な境界積分が局所的な外微分の積分に等しい。
\item
  第3章の曲面積分と同じく、面を小さなパッチに分け、各パッチで測定器を作用させて集計している。
\end{itemize}
\end{quote}

\begin{center}\rule{0.5\linewidth}{0.5pt}\end{center}

\hypertarget{ux540cux3058ux3053ux3068ux3092ux3082ux3046ux4e00ux6bb52-form-ux306eux5916ux5faeux5206ux3068ux767aux6563}{%
\subsection{\texorpdfstring{§5.7 同じことをもう一段------\(2\)-form
の外微分と発散}{§5.7 同じことをもう一段------2-form の外微分と発散}}\label{ux540cux3058ux3053ux3068ux3092ux3082ux3046ux4e00ux6bb52-form-ux306eux5916ux5faeux5206ux3068ux767aux6563}}

\hypertarget{form-ux306fux9762ux306eux6e2cux5b9aux5668}{%
\subsubsection{\texorpdfstring{5.7.1 \(2\)-form
は面の測定器}{5.7.1 2-form は面の測定器}}\label{form-ux306fux9762ux306eux6e2cux5b9aux5668}}

第3章 §3.4.2 で、一般の \(2\)-form を次の形で書いた：

\[\eta = A\,dy \wedge dz + B\,dz \wedge dx + C\,dx \wedge dy\]

これは「面に沿って流量を測る測定器」である。第2章 §2.4.5
で見たように、3つの基底 \(2\)-form がそれぞれ \(yz\) 平面、\(zx\)
平面、\(xy\) 平面への影の面積を測る。

\hypertarget{ux5faeux5c0fux76f4ux65b9ux4f53ux306eux8868ux9762ux3067ux6e2cux308b}{%
\subsubsection{5.7.2
微小直方体の表面で測る}\label{ux5faeux5c0fux76f4ux65b9ux4f53ux306eux8868ux9762ux3067ux6e2cux308b}}

§5.3 と同じ精神で、今度は空間内の微小直方体
\([x, x+\Delta x] \times [y, y+\Delta y] \times [z, z+\Delta z]\)
の\textbf{6つの面すべて}で \(\eta\) を積分する。向きは外向きに合わせる。

たとえば \(C\,dx \wedge dy\) の項について、\(z\)
に垂直な底面と上面だけを考える：

\begin{itemize}
\item
  \textbf{底面}（\(z\)、外向きは \(-z\) 方向）：\(C(x,y,z)\) に
  \(-dx \wedge dy\) の向きで評価 → およそ
  \(-C(x,y,z)\,\Delta x\,\Delta y\)
\item
  \textbf{上面}（\(z+\Delta z\)、外向きは \(+z\)
  方向）：\(C(x,y,z+\Delta z) \approx C + \frac{\partial C}{\partial z}\Delta z\)
  → およそ
  \((C + \frac{\partial C}{\partial z}\Delta z)\,\Delta x\,\Delta y\)
\end{itemize}

和をとると \(C\Delta x\Delta y\)
が相殺し、\(\frac{\partial C}{\partial z}\,\Delta x\,\Delta y\,\Delta z\)
が残る。\(A\,dy \wedge dz\) の項は \(x\) に垂直な面から
\(\frac{\partial A}{\partial x}\,\Delta x\,\Delta y\,\Delta z\)
を、\(B\,dz \wedge dx\) の項は \(y\) に垂直な面から
\(\frac{\partial B}{\partial y}\,\Delta x\,\Delta y\,\Delta z\)
を与える。

全6面の合計は：

\[\iint_{\partial (\text{直方体})} \eta \approx (\frac{\partial A}{\partial x} + \frac{\partial B}{\partial y} + \frac{\partial C}{\partial z})\,\Delta x\,\Delta y\,\Delta z\]

ここでも同じ構造だ------\textbf{表面で測ったズレの主項は、囲んだ体積に比例する}。比例係数
\(\frac{\partial A}{\partial x} + \frac{\partial B}{\partial y} + \frac{\partial C}{\partial z}\)
が「単位体積あたりの湧き出し」を表す。

\hypertarget{deta-ux306eux5b9aux7fa9ux3068ux4ee3ux6570ux7684ux5c0eux51fa}{%
\subsubsection{\texorpdfstring{5.7.3 \(d\eta\)
の定義と代数的導出}{5.7.3 d\textbackslash eta の定義と代数的導出}}\label{deta-ux306eux5b9aux7fa9ux3068ux4ee3ux6570ux7684ux5c0eux51fa}}

直方体の三辺
\(\Delta x\,\hat{e}_x,\; \Delta y\,\hat{e}_y,\; \Delta z\,\hat{e}_z\)
を食わせる \(3\)-form として：

\[d\eta := (\frac{\partial A}{\partial x} + \frac{\partial B}{\partial y} + \frac{\partial C}{\partial z})\,dx \wedge dy \wedge dz\]

この式は、§5.5.2 の代数的ルールからも導ける。その前に、§5.5.2
でやったのと同じく、幾何的な結果がライプニッツ則と整合することを
\(A\,dy \wedge dz\) の項で確かめておこう。

§5.7.2 の幾何計算によれば、\(A\,dy \wedge dz\) の寄与は
\(\frac{\partial A}{\partial x}\,\Delta x\,\Delta y\,\Delta z\)
だった。これを \(dx \wedge dy \wedge dz\) の言葉で書けば
\(\frac{\partial A}{\partial x}\,dx \wedge dy \wedge dz\) である。

一方、\(d(A\,dy \wedge dz)\) にライプニッツ則を適用すると：

\[d(A\,dy \wedge dz) = dA \wedge dy \wedge dz + A\,d(dy \wedge dz)\]

ここで \(d(dy \wedge dz)\) は、§5.5.2 と同じく \(d(dy)=d(dz)=0\)
を使って：

\[d(dy \wedge dz) = d(dy) \wedge dz - dy \wedge d(dz) = 0 \wedge dz - dy \wedge 0 = 0\]

したがって第2項は消え、第1項だけが残る：

\[dA \wedge dy \wedge dz = (\frac{\partial A}{\partial x}\,dx + \frac{\partial A}{\partial y}\,dy + \frac{\partial A}{\partial z}\,dz) \wedge dy \wedge dz\]

\(\wedge\)
の反対称性（\(dy \wedge dy = 0\)、\(dz \wedge dz = 0\)）により
\(\frac{\partial A}{\partial y}\) と \(\frac{\partial A}{\partial z}\)
の項は消え、\(\frac{\partial A}{\partial x}\,dx \wedge dy \wedge dz\)
だけが残る。これは §5.7.2 の幾何計算と\textbf{完全に一致する}。

\(B\,dz \wedge dx\)、\(C\,dx \wedge dy\)
の項も同様で、三つを足し合わせれば：

\[d\eta = (\frac{\partial A}{\partial x} + \frac{\partial B}{\partial y} + \frac{\partial C}{\partial z})\,dx \wedge dy \wedge dz\]

\(d\) は \(2\)-form → \(3\)-form の昇格を果たした。

\hypertarget{ux30acux30a6ux30b9ux306eux5b9aux7406}{%
\subsubsection{5.7.4
ガウスの定理}\label{ux30acux30a6ux30b9ux306eux5b9aux7406}}

§5.6 と同じタイル貼りの論法------今度は立体 \(V\)
を微小直方体で埋め尽くす------により、内部の面は相殺し、外面
\(\partial V\) だけが残る：

\[\iint_{\partial V} \eta = \iiint_V d\eta\]

これが\textbf{ガウスの定理}である。ストークスの定理と、次元が一つずつずれただけで\textbf{まったく同じ形}をしている。並べてみれば一目瞭然だ：\(\omega\)（\(1\)-form）\(\to\)
\(d\omega\)（\(2\)-form）\(\to\) 面 \(S\)、\(\eta\)（\(2\)-form）\(\to\)
\(d\eta\)（\(3\)-form）\(\to\) 立体 \(V\)。

\(\eta = A\,dy \wedge dz + B\,dz \wedge dx + C\,dx \wedge dy\)
の成分をあらわに書けば：

\[\iint_{\partial V} \bigl(A\,dy \wedge dz + B\,dz \wedge dx + C\,dx \wedge dy\bigr) = \iiint_V (\frac{\partial A}{\partial x} + \frac{\partial B}{\partial y} + \frac{\partial C}{\partial z})\,dx \wedge dy \wedge dz\]

\begin{quote}
\textbf{注}
（既習者への補助線）ベクトル解析を既習の読者には、\(\frac{\partial A}{\partial x}+\frac{\partial B}{\partial y}+\frac{\partial C}{\partial z}\)
にも見覚えがあるだろう。本書ではまず \(2\)-form
の外微分としてこの量を得た、という順序を大切にする。名前との対応は後の章で整理する。
\end{quote}

\begin{quote}
\textbf{【ここまでのチェックポイント】}

\begin{itemize}
\item
  \(2\)-form \(\eta\) の外微分 \(d\eta\) は \(3\)-form で、係数は
  \(\frac{\partial A}{\partial x} + \frac{\partial B}{\partial y} + \frac{\partial C}{\partial z}\)。
\item
  導出は §5.5 と同じ代数的ルール（ライプニッツ則 \(+\)
  \(d(dx)=0\)）で機械的にできる。
\item
  \(\iint_{\partial V} \eta = \iiint_V d\eta\)。ストークス（\(1\)-form→面）とガウス（\(2\)-form→立体）は次数が違うだけで同じ構造。
\end{itemize}
\end{quote}

\begin{center}\rule{0.5\linewidth}{0.5pt}\end{center}

\hypertarget{d2-0ux30baux30ecux306eux30baux30ecux306fux6b8bux3089ux306aux3044}{%
\subsection{\texorpdfstring{§5.8
\(d^2 = 0\)------ズレのズレは残らない}{§5.8 d\^{}2 = 0------ズレのズレは残らない}}\label{d2-0ux30baux30ecux306eux30baux30ecux306fux6b8bux3089ux306aux3044}}

次節でストークスの定理とガウスの定理を一本化する前に、外微分が持つもう一つの基本構造を確認しておこう。\(d\)
は次数を一つ上げるが、それを二度続けると新しいズレは残らない。この事実は、後で見る「境界の境界は向き付き和として消える」という幾何的構造ともつながっている。

\hypertarget{ddf-0}{%
\subsubsection{\texorpdfstring{5.8.1
\(d(df) = 0\)}{5.8.1 d(df) = 0}}\label{ddf-0}}

§5.1 で
\(df = \frac{\partial f}{\partial x}\,dx + \frac{\partial f}{\partial y}\,dy + \frac{\partial f}{\partial z}\,dz\)
を定義した。これをさらにもう一度外微分してみる。§5.5 の公式で
\(P = \frac{\partial f}{\partial x},\; Q = \frac{\partial f}{\partial y},\; R = \frac{\partial f}{\partial z}\)
とすれば：

\[d(df) = \left(\frac{\partial^2 f}{\partial y \partial z} - \frac{\partial^2 f}{\partial z \partial y}\right) dy \wedge dz + \left(\frac{\partial^2 f}{\partial z \partial x} - \frac{\partial^2 f}{\partial x \partial z}\right) dz \wedge dx + \left(\frac{\partial^2 f}{\partial x \partial y} - \frac{\partial^2 f}{\partial y \partial x}\right) dx \wedge dy\]

混合偏微分の対称性（\(f\) が \(C^2\) 級（2階連続微分可能）なら
\(\frac{\partial^2 f}{\partial x \partial y} = \frac{\partial^2 f}{\partial y \partial x}\)
など。第1章 §1.2 参照）により、すべての係数がゼロになる：

\[d(df) = 0\]

\begin{quote}
\textbf{注} （公理的な立場での \(d^2=0\) と \(C^2\)
条件）座標表示で具体的に計算すると、\(d^2=0\)
は混合偏微分の対称性とウェッジ積の反対称性によって必然的に現れる性質であることがわかる。微分形式の公理的な立場では、逆に
\(d^2=0\)
を要請（定義の一部）とすることで、暗黙のうちに混合偏微分の対称性が成り立つ（\(C^2\)
級の）関数クラスを扱っている、と読むこともできる。本書ではまず具体計算を通じて、物理法則の階層に矛盾が生じないための仕組みとして理解する。
\end{quote}

\hypertarget{ddomega-0}{%
\subsubsection{\texorpdfstring{5.8.2
\(d(d\omega) = 0\)}{5.8.2 d(d\textbackslash omega) = 0}}\label{ddomega-0}}

§5.5 の \(d\omega\) に §5.7
の公式を適用する。\(A = \frac{\partial R}{\partial y} - \frac{\partial Q}{\partial z},\; B = \frac{\partial P}{\partial z} - \frac{\partial R}{\partial x},\; C = \frac{\partial Q}{\partial x} - \frac{\partial P}{\partial y}\)
として：

\[d(d\omega) = \left(\frac{\partial^2 R}{\partial x \partial y} - \frac{\partial^2 Q}{\partial x \partial z} + \frac{\partial^2 P}{\partial y \partial z} - \frac{\partial^2 R}{\partial y \partial x} + \frac{\partial^2 Q}{\partial z \partial x} - \frac{\partial^2 P}{\partial z \partial y}\right) dx \wedge dy \wedge dz\]

括弧内を展開して並べ替えると：

\[\left(\frac{\partial^2 R}{\partial x \partial y} - \frac{\partial^2 Q}{\partial x \partial z}\right) + \left(\frac{\partial^2 P}{\partial y \partial z} - \frac{\partial^2 R}{\partial y \partial x}\right) + \left(\frac{\partial^2 Q}{\partial z \partial x} - \frac{\partial^2 P}{\partial z \partial y}\right)\]

混合偏微分の対称性（\(\frac{\partial^2 R}{\partial y \partial x}=\frac{\partial^2 R}{\partial x \partial y}\)
など）により、6項が3組の相殺を起こし、合計はゼロ：

\[d(d\omega) = 0\]

\begin{quote}
\textbf{【ここまでのチェックポイント】}

\begin{itemize}
\item
  \(d^2 = 0\)：外微分を2回続けると必ずゼロ。混合偏微分の対称性とウェッジの反対称性が相殺を起こす。
\item
  既習者が知るいくつかの「二度微分すると消える」恒等式は、この \(d^2=0\)
  の別表現として後の章で整理される。
\end{itemize}
\end{quote}

\begin{center}\rule{0.5\linewidth}{0.5pt}\end{center}

\hypertarget{ux5916ux5faeux5206ux306eux7d71ux5408ux4e00ux3064ux306eux5f0fux4e00ux3064ux306eux30ebux30fcux30eb}{%
\subsection{§5.9
外微分の統合------一つの式、一つのルール}\label{ux5916ux5faeux5206ux306eux7d71ux5408ux4e00ux3064ux306eux5f0fux4e00ux3064ux306eux30ebux30fcux30eb}}

\hypertarget{ux30b9ux30c8ux30fcux30afux30b9ux3068ux30acux30a6ux30b9ux306fux540cux3058ux5f0fux3060ux3063ux305f}{%
\subsubsection{5.9.1
ストークスとガウスは同じ式だった}\label{ux30b9ux30c8ux30fcux30afux30b9ux3068ux30acux30a6ux30b9ux306fux540cux3058ux5f0fux3060ux3063ux305f}}

§5.6 で得たストークスの定理
\(\oint_{\partial S}\omega = \iint_S d\omega\) と、§5.7
で得たガウスの定理
\(\iint_{\partial V}\eta = \iiint_V d\eta\)。並べてみると、\(\oint\)、\(\iint\)、\(\iiint\)
と積分記号の数が違うだけで、構造は同一だ。\(M\) が曲線なら
\(\int\)（1次元）、曲面なら \(\iint\)（2次元）、立体なら
\(\iiint\)（3次元）------つまり積分記号の重なりは \(M\)
の次元で決まるにすぎない。

そこで、次元によらず統一的に：

\[\int_{\partial M} \omega = \int_M d\omega\]

と書く。\(\omega\) は \(k\)-form、\(M\) は \((k+1)\)
次元領域、\(\partial M\) はその \(k\) 次元の境界である。\(k=0\) なら
\(\int_{\partial M} f = f(B)-f(A) = \int_M df\)
という微積分学の基本定理に、\(k=1\) ならストークスの定理に、\(k=2\)
ならガウスの定理になる。すべてがこの一本の式に吸収される。

この統一形の威力は §5.8 の \(d^2=0\)
との組み合わせでさらに際立つ。\(\int_{\partial M} \omega = \int_M d\omega\)
の \(\omega\) を \(d\omega\) に置き換えれば：

\[\int_{\partial(\partial M)} \omega = \int_{\partial M} d\omega = \int_M d(d\omega) = 0\]

つまり
\textbf{「境界の境界は向き付き和としてゼロになる」}（\(\partial^2 M = 0\)）という幾何学的事実と、\textbf{$d^2=0$}
は互いに呼応している。集合として二度目の境界が常に空になる、という意味ではない。たとえば多角形の境界をさらに境界に分ければ頂点は現れるが、向き付きで数えると隣り合う辺から来た寄与が相殺する。\(d^2=0\)
はこの位相的な相殺の、微分形式による代数的な写し絵なのである。

\begin{quote}
\textbf{注} （3次元までで十分）本書の舞台は 3次元空間だから、\(M\)
の次元は1・2・3のいずれかであり、\(\omega\) は
\(0\)-form・\(1\)-form・\(2\)-form のいずれかである。\(k=3\)（\(M\)
が4次元）には立ち入らない。しかしこの式が \(k\)
によらず成立するという事実は、頭の片隅に置いておいて損はない。
\end{quote}

\hypertarget{ux4e00ux822cux306e-k-form-ux306eux5916ux5faeux5206}{%
\subsubsection{\texorpdfstring{5.9.2 一般の \(k\)-form
の外微分}{5.9.2 一般の k-form の外微分}}\label{ux4e00ux822cux306e-k-form-ux306eux5916ux5faeux5206}}

同様に、\(d\) の作用も次数によらず一つのパターンにまとまる。任意の
\(k\)-form \(\omega\) が係数 \(a_{i_1\cdots i_k}\) と基底
\(dx_{i_1}\wedge\cdots\wedge dx_{i_k}\) の和で書かれているとき：

\[d\omega = \sum \Bigl( \sum_j \frac{\partial a_{i_1\cdots i_k}}{\partial x_j}\,dx_j \Bigr) \wedge dx_{i_1}\wedge\cdots\wedge dx_{i_k}\]

本書で必要なのは \(k=0,1,2\) の三つの場合だけであり、それぞれ
§5.1（\(df\)）、§5.5（\(d\omega\)）、§5.7（\(d\eta\)）で具体的に書き下した。上の一般形は「どの次数でも、係数を全微分してウェッジでつなぐ」という単一の原理に従っていることの確認にすぎない。

\hypertarget{ux8a08ux7b97ux30ebux30fcux30ebux307eux3068ux3081}{%
\subsubsection{5.9.3
計算ルールまとめ}\label{ux8a08ux7b97ux30ebux30fcux30ebux307eux3068ux3081}}

本章で構築した外微分 \(d\) の全計算は、以下の 4 つのルールに集約される：

\begin{enumerate}
\def\labelenumi{\arabic{enumi}.}
\tightlist
\item
  \textbf{$0$-form への作用}：\(df = \frac{\partial f}{\partial x}\,dx + \frac{\partial f}{\partial y}\,dy + \frac{\partial f}{\partial z}\,dz\)（第1章
  §1.2 以来の定義）
\item
  \textbf{線形性}：\(d(\omega_1 + \omega_2) = d\omega_1 + d\omega_2\)
\item
  \textbf{次数付きライプニッツ則}：\(d(f\,\omega) = df \wedge \omega + f\,d\omega\)
\item
  \textbf{座標基底の外微分はゼロ}：\(d(dx) = d(dy) = d(dz) = 0\)
\end{enumerate}

ルール4 は、座標基底 \(dx,dy,dz\)
自身には位置による係数の変化がなく、微小ループで測っても局所的なズレを生まない、という幾何的事実を反映している。§5.8
の \(d^2=0\) とも整合する。

これら4つのルールさえあれば、\(0\)-form から \(3\)-form
までの任意の形式の外微分が機械的に計算できる。実際に、本章で扱ったすべての計算はこのルールの組み合わせにすぎない。

\begin{quote}
\textbf{注}
（引き戻しとの可換性------第4章の伏線回収）第4章で予告したように、引き戻し
\(\Phi^*\) と外微分 \(d\) のあいだには
\[\Phi^*(d\omega) = d(\Phi^*\omega)\]
という可換関係が成り立つ。「微分してから引き戻す」のと「引き戻してから微分する」のは同じ結果を与える。この性質は、これまでの4ルールからは直接導けないが、\(d\)
が「局所的なズレを測る」操作であり、\(\Phi^*\)
が「座標を焼き直す」操作であることを思い出せば自然に納得できる------局所的なズレの構造は、座標のとり方に依存しないのだ。この事実は、第11章で一般の多様体上の微分形式を考える際の重要な足場となる。
\end{quote}

\begin{quote}
\textbf{【ここまでのチェックポイント】}

\begin{itemize}
\item
  \(\int_{\partial M} \omega = \int_M d\omega\)
  一本で、微積分学の基本定理・ストークスの定理・ガウスの定理を統合できる。
\item
  \(d\) の計算は 4
  ルール（\(df\)、線形性、ライプニッツ則、\(d(dx)=0\)）に集約される。
\item
  \(d^2=0\) と \(\partial^2 M=0\) は代数的真理と幾何学的真理の呼応。
\end{itemize}
\end{quote}

\begin{center}\rule{0.5\linewidth}{0.5pt}\end{center}

\hypertarget{ux7a4dux5206ux304bux3089ux5faeux5206ux65b9ux7a0bux5f0fux3078ux7269ux7406ux6cd5ux5247ux306eux5c40ux6240ux5316}{%
\subsection{§5.10
積分から微分方程式へ------物理法則の局所化}\label{ux7a4dux5206ux304bux3089ux5faeux5206ux65b9ux7a0bux5f0fux3078ux7269ux7406ux6cd5ux5247ux306eux5c40ux6240ux5316}}

\hypertarget{ux7a4dux5206ux6cd5ux5247ux3092ux5faeux5206ux6cd5ux5247ux306bux5909ux3048ux308b}{%
\subsubsection{5.10.1
積分法則を微分法則に変える}\label{ux7a4dux5206ux6cd5ux5247ux3092ux5faeux5206ux6cd5ux5247ux306bux5909ux3048ux308b}}

本章の冒頭に掲げた問いに戻ろう。

\begin{quote}
\textbf{積分された量から、どのようにして局所法則を取り出すのか。}
\end{quote}

答えは、これまでに得た二つの道具------\textbf{ストークスの定理}と\textbf{外微分 $d$}
------にある。

物理学の多くの基本法則は、まず\textbf{積分の形}で発見される。ある領域の境界で測った量が、内部にある源の総和に等しい、という形だ：

\[\int_{\partial M} \omega = \int_M \eta\]

左辺は境界で測った積分（大域的な観測量）、右辺は内部の源の積分（大域的な総量）。\(\omega\)
と \(\eta\) は、それぞれ適切な次数の微分形式である。

ここにストークスの定理を左辺に適用する。\(\int_{\partial M} \omega = \int_M d\omega\)
だから：

\[\int_M d\omega = \int_M \eta\]

移項して：

\[\int_M (d\omega - \eta) = 0\]

ここが決定的な瞬間だ。この等式が、十分小さい領域を含む\textbf{任意の滑らかな領域 $M$}
について成り立つならば、被積分形式は各点で恒等的にゼロでなければならない。なぜなら、もし
\(d\omega - \eta \neq 0\)
となる点があれば、その点の周りの微小領域で積分がゼロでなくなり矛盾するからだ。ここでも、場が十分滑らかで特異な集中源を含まない場合を考えている。

したがって：

\[d\omega = \eta\]

これが\textbf{積分法則から抽出された局所微分法則}である。\(d\)
はまさに「積分方程式を微分方程式に変換する演算子」として機能したのだ。

つまり、外微分 \(d\)
は単に形式の次数を上げる記号ではない。境界で観測された法則を、内部の各点で成り立つ法則へ翻訳するための演算子である。

\hypertarget{ux7269ux7406ux306eux5b9fux4f8b}{%
\subsubsection{5.10.2 物理の実例}\label{ux7269ux7406ux306eux5b9fux4f8b}}

この構図の典型例は電磁気学の基本法則である。

\begin{quote}
\textbf{注}
（電磁気学を未習の読者へ）この節では、電磁気学の詳しい内容ではなく、積分形の法則が局所形へ変わるしくみを見る。\(E,D,B,H,J,\rho\)
は、ここでは物理量の名前というより、「線で測る量」「面で測る量」「体積で測る量」を表す記号だと思って読めばよい。
\end{quote}

\begin{quote}
\textbf{注}
（ベクトル解析で電磁気学を学んだ読者へ）ここでは、通常の三成分ベクトルとしての場ではなく、どの次元の対象に積分するかに合わせて形式の次数を選んでいる。\(E,H\)
は線に沿って積分されるので \(1\)-form、\(D,B,J\)
は面を貫いて積分されるので \(2\)-form、\(\rho\)
は体積にわたって積分されるので \(3\)-form
として扱う。通常のベクトル解析記法との対応は、ホッジ・スターを導入した後で整理する。
\end{quote}

電磁気では、量の種類によって積分する対象が異なる。線に沿って測る量、面を貫いて測る量、体積にわたって測る量がある。そこで模式的に、次のように読むことができる。

\begin{longtable}[]{@{}lll@{}}
\toprule\noalign{}
量 & 測る対象 & 形式 \\
\midrule\noalign{}
\endhead
\bottomrule\noalign{}
\endlastfoot
\(E,H\) & 線に沿って測る & \(1\)-form \\
\(D,B,J\) & 面を貫いて測る & \(2\)-form \\
\(\rho\) & 体積にわたって測る & \(3\)-form \\
\end{longtable}

たとえば、電荷に関するガウスの法則は

\[
\int_{\partial V}D=\int_V\rho
\]

という形をしている。左辺は閉曲面 \(\partial V\)
を通って出ていく電束の総量であり、右辺は内部 \(V\)
に含まれる電荷の総量である。ストークスの定理を使えば、

\[
\int_{\partial V}D=\int_V dD
\]

だから、

\[
\int_V dD=\int_V\rho
\]

となる。これが任意の領域 \(V\) で成り立つなら、

\[
dD=\rho
\]

である。

同じように、閉曲面を貫く正味の磁束がゼロである、という法則は

\[
\int_{\partial V}B=0
\]

と書ける。ストークスの定理より、

\[
\int_V dB=0
\]

であり、任意の \(V\) について成り立つなら、

\[
dB=0
\]

となる。

時間変化を含む法則も同じ型を持つ。曲面 \(S\) に対して、

\[
\int_{\partial S}E
=
-\frac{\partial}{\partial t}\int_S B
\]

なら、

\[
dE=-\frac{\partial B}{\partial t}
\]

であり、

\[
\int_{\partial S}H
=
\int_SJ+\frac{\partial}{\partial t}\int_S D
\]

なら、

\[
dH=J+\frac{\partial D}{\partial t}
\]

である。ここで \(d\) は空間方向の外微分であり、時間変化は
\(\partial/\partial t\) で表している。

したがって、電磁気の基礎方程式は、この章の言葉では

\[
dD=\rho,\qquad
dB=0,\qquad
dE=-\frac{\partial B}{\partial t},\qquad
dH=J+\frac{\partial D}{\partial t}
\]

という同じ型の式として現れる。

ここで重要なのは、これらを電磁気学の新しい公式として覚えることではない。境界で測る積分法則が、外微分
\(d\) によって内部の局所法則へ移る、という一点である。

一方で、\(E\) と \(D\)、あるいは \(B\) と \(H\)
を関係づけるには、空間や媒質の測り方が必要になる。また、線で測る量と面で測る量を対応させるにも、長さ・面積・体積をどう測るかという追加の規則が必要になる。次章で導入するホッジ・スターは、この対応を与えるための道具である。

成分表示を含む詳しい展開は、付録C「電磁気の積分形と微分形式」に回す。

\begin{quote}
\textbf{【ここまでのチェックポイント】}

\begin{itemize}
\item
  物理法則はしばしば \(\int_{\partial M} \omega = \int_M \eta\)
  の形で与えられる。
\item
  ストークスの定理で \(\int_{\partial M} \omega = \int_M d\omega\)
  と書き換え、任意の \(M\) で成立することから
  \(d\omega = \eta\)（局所法則）を得る。
\item
  ただしこの局所化は、場が十分滑らかで、特異な集中源を含まない領域での議論である。
\item
  \(d\) は積分法則を微分法則に変換する演算子である。
\end{itemize}
\end{quote}

\begin{center}\rule{0.5\linewidth}{0.5pt}\end{center}

\hypertarget{ux7b2ciiux90e8ux3078ux306eux5c55ux671bux30dbux30c3ux30b8ux30b9ux30bfux30fcux3078ux306eux4f0fux7dda}{%
\subsection{§5.11
第II部への展望------ホッジ・スターへの伏線}\label{ux7b2ciiux90e8ux3078ux306eux5c55ux671bux30dbux30c3ux30b8ux30b9ux30bfux30fcux3078ux306eux4f0fux7dda}}

本章で我々は、外微分 \(d\) という強力な演算子を手に入れた。\(d\) は：

\begin{itemize}
\item
  \(0\)-form \(\to\)
  \(1\)-form：\(df = \frac{\partial f}{\partial x}\,dx + \frac{\partial f}{\partial y}\,dy + \frac{\partial f}{\partial z}\,dz\)
\item
  \(1\)-form \(\to\)
  \(2\)-form：\(d\omega = (\frac{\partial R}{\partial y}-\frac{\partial Q}{\partial z})\,dy\wedge dz + (\frac{\partial P}{\partial z}-\frac{\partial R}{\partial x})\,dz\wedge dx + (\frac{\partial Q}{\partial x}-\frac{\partial P}{\partial y})\,dx\wedge dy\)
\item
  \(2\)-form \(\to\)
  \(3\)-form：\(d\eta = (\frac{\partial A}{\partial x}+\frac{\partial B}{\partial y}+\frac{\partial C}{\partial z})\,dx\wedge dy\wedge dz\)
\end{itemize}

次数を一段ずつ上げるたびに、係数は異なるパターンで組み替わる。そして
\(d^2=0\) が、これらの階層の間に矛盾が入り込めないことを保証する。

しかし、ここで一つの非対称性が目につく。\(0\)-form と \(3\)-form
は独立成分が 1 つ、\(1\)-form と \(2\)-form は独立成分が 3
つ------\(1, 3, 3, 1\) という対称性がある（第2章 §2.5.9--§2.5.10
で見た）。この対称性は偶然ではない。デカルト座標の 3
次元空間では、長さ・面積・体積を測る規則を加えることで、同じ情報量の形式どうしを結びつける「辞書」が現れる。

その辞書が、次章で導入するホッジ・スター \(\ast\)
である。ただし、\(\ast\)
は単なる記号の置き換えではない。どの方向をどの面に対応させるか、面積や体積をどう測るか、向きをどう決めるかに依存する。そのため本章では、具体的な対応式はまだ使わない。

ここまでで構築したのは、あくまで \(d\) である。\(d\)
は次数を一つ上げ、局所的なズレを取り出し、積分法則を局所法則へ変換する。しかし
\(d\)
だけでは、これまで見慣れていたベクトル解析の各演算をまだ再構成できない。\(1\)-form
と \(2\)-form
を対応させるには、空間の長さ・面積・体積をどう測るかという追加規則------計量------が必要になる。

次章では、その「測り方」の規則をホッジ・スター \(\ast\)
として定義する。第5章で作った \(d\) に、第6章で \(\ast\)
を加え、その二つを両輪として後の章で見慣れたベクトル解析の演算を組み立て直す。ここから第II部の核心------ナブラの解体------へと進む。

\begin{center}\rule{0.5\linewidth}{0.5pt}\end{center}

\begin{quote}
\textbf{【ここまでのチェックポイント — 第5章全体】}

\begin{itemize}
\item
  外微分 \(d\) は \(k\)-form を \((k+1)\)-form に写す演算子。\(0\)-form
  \(\to\) \(1\)-form（\(df\)）、\(1\)-form \(\to\)
  \(2\)-form、\(2\)-form \(\to\) \(3\)-form。
\item
  計算ルールは 4 つ：\(df\) の定義、線形性、ライプニッツ則
  \(d(f\omega) = df\wedge\omega + f\,d\omega\)、\(d(dx)=d(dy)=d(dz)=0\)。
\item
  幾何学的起源は §5.3
  の微小ループ解体：閉ループのズレが面積に比例し、その比例係数が
  \(d\omega\)。
\item
  ストークスの定理 \(\int_{\partial M} \omega = \int_M d\omega\)
  は、境界と内部をつなぐ普遍的な橋。
\item
  \(d^2 = 0\)
  は混合偏微分の対称性とウェッジ積の反対称性に由来し、幾何学的には「境界の境界は向き付き和としてゼロになる」ことに対応する。
\item
  \(d\)
  は積分法則を局所微分法則へ変換する演算子------物理学の定式化に不可欠な役割を果たす。
\item
  次章では、長さ・面積・体積を測る規則を加え、ホッジ・スター \(\ast\)
  を通じて \(d\) とベクトル解析の演算との対応を整理する。
\end{itemize}
\end{quote}

\begin{center}\rule{0.5\linewidth}{0.5pt}\end{center}

\hypertarget{ux4ed8ux9332bux5916ux5faeux5206ux306eux884cux5217ux8868ux793a}{%
\section{付録B：外微分の行列表示}\label{ux4ed8ux9332bux5916ux5faeux5206ux306eux884cux5217ux8868ux793a}}

本章本文は基底 \(dx, dy, dz\)
とウェッジ積の代数で進めた。ここでは第2章以来の本書の流儀------\textbf{成分をひとつ残らず行列に並べる}------によって、外微分を行列表示の言葉で書き直す。

\hypertarget{b.1-0-formdf-ux306e-1-times-3-ux884cux30d9ux30afux30c8ux30eb}{%
\subsection{\texorpdfstring{B.1 \(0\)-form：\(df\) の \(1 \times 3\)
行ベクトル}{B.1 0-form：df の 1 \textbackslash times 3 行ベクトル}}\label{b.1-0-formdf-ux306e-1-times-3-ux884cux30d9ux30afux30c8ux30eb}}

\(f = f(x,y,z)\) に対し、第1章 §1.2 の定義どおり：

\[df = \frac{\partial f}{\partial x}\,dx + \frac{\partial f}{\partial y}\,dy + \frac{\partial f}{\partial z}\,dz\]

行ベクトル（\(1 \times 3\) 行列）として：

\[df = \begin{pmatrix} \frac{\partial f}{\partial x} & \frac{\partial f}{\partial y} & \frac{\partial f}{\partial z} \end{pmatrix}\]

である。

\hypertarget{b.2-1-formux4fc2ux6570ux306eux30e4ux30b3ux30d3ux884cux5217-mathbfj}{%
\subsection{\texorpdfstring{B.2 \(1\)-form：係数のヤコビ行列
\(\mathbf{J}\)}{B.2 1-form：係数のヤコビ行列 \textbackslash mathbf\{J\}}}\label{b.2-1-formux4fc2ux6570ux306eux30e4ux30b3ux30d3ux884cux5217-mathbfj}}

\(\omega = P\,dx + Q\,dy + R\,dz\) に対し、係数 \((P,Q,R)\)
の偏微分を\textbf{すべて}並べた \(3 \times 3\) 行列を：

\[\mathbf{J} := \begin{pmatrix}
\frac{\partial P}{\partial x} & \frac{\partial P}{\partial y} & \frac{\partial P}{\partial z} \\
\frac{\partial Q}{\partial x} & \frac{\partial Q}{\partial y} & \frac{\partial Q}{\partial z} \\
\frac{\partial R}{\partial x} & \frac{\partial R}{\partial y} & \frac{\partial R}{\partial z}
\end{pmatrix}\]

とする。これは \((P,Q,R)\)
を「縦に積んだ」ベクトル場のヤコビ行列にあたる。§5.1
で「単なる偏微分の行列では足りない」と述べたのは、まさにこの
\(\mathbf{J}\) が反対称でないからだ。

\hypertarget{b.3-domega-mathbfjt---mathbfj}{%
\subsection{\texorpdfstring{B.3
\(d\omega = \mathbf{J}^T - \mathbf{J}\)}{B.3 d\textbackslash omega = \textbackslash mathbf\{J\}\^{}T - \textbackslash mathbf\{J\}}}\label{b.3-domega-mathbfjt---mathbfj}}

§5.5 の結果は：

\[d\omega = (\frac{\partial R}{\partial y} - \frac{\partial Q}{\partial z})\,dy \wedge dz + (\frac{\partial P}{\partial z} - \frac{\partial R}{\partial x})\,dz \wedge dx + (\frac{\partial Q}{\partial x} - \frac{\partial P}{\partial y})\,dx \wedge dy\]

第2章 §2.4.4 の流儀に従い、任意の縦ベクトル
\(\mathbf{v}_1, \mathbf{v}_2\) に対して
\((d\omega)(\mathbf{v}_1, \mathbf{v}_2) = \mathbf{v}_1^T \mathbf{M} \mathbf{v}_2\)
となる反対称行列 \(\mathbf{M}\) を求めよう。

\(\mathbf{J}^T - \mathbf{J}\) の成分を書き下す：

\[\mathbf{J}^T - \mathbf{J} = \begin{pmatrix}
0 & \frac{\partial Q}{\partial x} - \frac{\partial P}{\partial y} & \frac{\partial R}{\partial x} - \frac{\partial P}{\partial z} \\
\frac{\partial P}{\partial y} - \frac{\partial Q}{\partial x} & 0 & \frac{\partial R}{\partial y} - \frac{\partial Q}{\partial z} \\
\frac{\partial P}{\partial z} - \frac{\partial R}{\partial x} & \frac{\partial Q}{\partial z} - \frac{\partial R}{\partial y} & 0
\end{pmatrix}\]

すなわち、2章で求めた\(2\)-formの行列表示と照らし合わせて：

\[d\omega = \mathbf{J}^T - \mathbf{J}\]

となり、\textbf{外微分 $d\omega$ の行列表示は、転置から係数のヤコビ行列を引いた反対称行列である。}

\hypertarget{b.4-2-formdeta-ux3068ux30e4ux30b3ux30d3ux884cux5217ux306eux30c8ux30ecux30fcux30b9}{%
\subsection{\texorpdfstring{B.4 \(2\)-form：\(d\eta\)
とヤコビ行列のトレース}{B.4 2-form：d\textbackslash eta とヤコビ行列のトレース}}\label{b.4-2-formdeta-ux3068ux30e4ux30b3ux30d3ux884cux5217ux306eux30c8ux30ecux30fcux30b9}}

\(\eta = A\,dy \wedge dz + B\,dz \wedge dx + C\,dx \wedge dy\)
に対し、係数 \((A,B,C)\) のヤコビ行列を：

\[\mathbf{J}_\eta := \begin{pmatrix}
\frac{\partial A}{\partial x} & \frac{\partial A}{\partial y} & \frac{\partial A}{\partial z} \\
\frac{\partial B}{\partial x} & \frac{\partial B}{\partial y} & \frac{\partial B}{\partial z} \\
\frac{\partial C}{\partial x} & \frac{\partial C}{\partial y} & \frac{\partial C}{\partial z}
\end{pmatrix}\]

とすると、§5.7 の結果
\(d\eta = (\frac{\partial A}{\partial x} + \frac{\partial B}{\partial y} + \frac{\partial C}{\partial z})\,dx \wedge dy \wedge dz\)
の係数は \(\mathbf{J}_\eta\)
の\textbf{トレース}（対角成分の和）に一致する：

\[\frac{\partial A}{\partial x} + \frac{\partial B}{\partial y} + \frac{\partial C}{\partial z} = \operatorname{tr}(\mathbf{J}_\eta)\]

\hypertarget{b.4.1-deta-ux306e27ux6210ux52063ux679aux306eux884cux5217ux306bux5c55ux958bux3059ux308b}{%
\subsubsection{\texorpdfstring{B.4.1 \(d\eta\)
の27成分------3枚の行列に展開する}{B.4.1 d\textbackslash eta の27成分------3枚の行列に展開する}}\label{b.4.1-deta-ux306e27ux6210ux52063ux679aux306eux884cux5217ux306bux5c55ux958bux3059ux308b}}

\(\eta\) の反対称成分は
\(\eta_{yz}=A,\; \eta_{zx}=B,\; \eta_{xy}=C\)（他は符号反転または
\(0\)）である。\(d\eta\) の成分を
\((d\eta)_{abc} = \partial_a \eta_{bc}\) と書く。ここで
\(\partial_x=\frac{\partial}{\partial x}\)、\(\partial_y=\frac{\partial}{\partial y}\)、\(\partial_z=\frac{\partial}{\partial z}\)
であり、\(a,b,c \in \{x,y,z\}\) だから成分は \(3^3=27\) 個ある。\(a\)
を固定した \(3 \times 3\) 行列3枚に並べると：

\[
\begin{aligned}
d\eta_{x,\cdot,\cdot}
&= \begin{pmatrix}
0 & \frac{\partial C}{\partial x} & -\frac{\partial B}{\partial x} \\
-\frac{\partial C}{\partial x} & 0 & \frac{\partial A}{\partial x} \\
\frac{\partial B}{\partial x} & -\frac{\partial A}{\partial x} & 0
\end{pmatrix}, \\
d\eta_{y,\cdot,\cdot}
&= \begin{pmatrix}
0 & \frac{\partial C}{\partial y} & -\frac{\partial B}{\partial y} \\
-\frac{\partial C}{\partial y} & 0 & \frac{\partial A}{\partial y} \\
\frac{\partial B}{\partial y} & -\frac{\partial A}{\partial y} & 0
\end{pmatrix}, \\
d\eta_{z,\cdot,\cdot}
&= \begin{pmatrix}
0 & \frac{\partial C}{\partial z} & -\frac{\partial B}{\partial z} \\
-\frac{\partial C}{\partial z} & 0 & \frac{\partial A}{\partial z} \\
\frac{\partial B}{\partial z} & -\frac{\partial A}{\partial z} & 0
\end{pmatrix}
\end{aligned}
\]

この \(27\) 成分を、対応する基底 \(1\)-form のウェッジ積（たとえば
\(a=x,\;b=y,\;c=z\) なら
\(dx \wedge dy \wedge dz\)）と縮約するとき、添字の重複がある項（対角成分や
\(b=c\) 等）は \(dx \wedge dx = 0\) で消え、生き残るのは \(a,b,c\)
がすべて異なる \(6\) 項（\(3!\)
の順列）だけ。それぞれを符号つきで和をとる。ここで各成分を反対称化して行列に並べて二重にカウントしているため、第2章
§2.4.4 や付録Dと同様に因子 \(\frac{1}{2}\) の補正が入り：

\[d\eta = \left( \frac{\partial A}{\partial x} + \frac{\partial B}{\partial y} + \frac{\partial C}{\partial z} \right) dx \wedge dy \wedge dz\]

すなわち \(d\eta\) は \(27\) 成分のうち \(6\)
項が生き残り、ペアで足し合わさって \(\mathbf{J}_\eta\)
のトレースへ収束する。

\(3\)-form は基底が \(dx \wedge dy \wedge dz\) の 1
つだけなので、行列としてはスカラー係数に縮退している。

\hypertarget{b.5-d2-f-0-ux3068ux30d8ux30c3ux30bbux884cux5217}{%
\subsection{\texorpdfstring{B.5 \(d^2 f = 0\)
とヘッセ行列}{B.5 d\^{}2 f = 0 とヘッセ行列}}\label{b.5-d2-f-0-ux3068ux30d8ux30c3ux30bbux884cux5217}}

\(0\)-form \(f\) のヘッセ行列：

\[\mathbf{H}_f := \begin{pmatrix}
\frac{\partial^2 f}{\partial x^2} & \frac{\partial^2 f}{\partial x \partial y} & \frac{\partial^2 f}{\partial x \partial z} \\
\frac{\partial^2 f}{\partial y \partial x} & \frac{\partial^2 f}{\partial y^2} & \frac{\partial^2 f}{\partial y \partial z} \\
\frac{\partial^2 f}{\partial z \partial x} & \frac{\partial^2 f}{\partial z \partial y} & \frac{\partial^2 f}{\partial z^2}
\end{pmatrix}\]

は \(f\) の 2 階偏導関数を並べたものだ。\(f\) が \(C^2\) 級なら
\(\mathbf{H}_f\)
は対称行列（\(\frac{\partial^2 f}{\partial x \partial y} = \frac{\partial^2 f}{\partial y \partial x}\)
など）である。

\(d(df) = 0\) は §5.5 の \(d\omega\) 公式で
\((P,Q,R) = (\frac{\partial f}{\partial x}, \frac{\partial f}{\partial y}, \frac{\partial f}{\partial z})\)
とおいたもの。\(\mathbf{J} = \mathbf{H}_f\) であり、\(\mathbf{H}_f\)
が対称だから \(\mathbf{J}^T - \mathbf{J} = 0\)、したがって
\(d(df) = 0\)。行列の言葉では「ヘッセ行列の反対称成分がゼロ」と言い換えられる。

\hypertarget{ux4ed8ux9332cux96fbux78c1ux6c17ux306eux7a4dux5206ux5f62ux3068ux5faeux5206ux5f62ux5f0f}{%
\section{付録C：電磁気の積分形と微分形式}\label{ux4ed8ux9332cux96fbux78c1ux6c17ux306eux7a4dux5206ux5f62ux3068ux5faeux5206ux5f62ux5f0f}}

この付録では、電磁気の四つの基礎方程式を、積分形から局所形へ移す過程を成分表示で確認する。ここで使うのは、外微分
\(d\)
と微分形式の次数である。空間の計量は、これらの局所化そのものには現れない。

\hypertarget{c.1-ux7269ux7406ux91cfux306eux6b21ux6570}{%
\subsection{C.1
物理量の次数}\label{c.1-ux7269ux7406ux91cfux306eux6b21ux6570}}

線に沿って測る量を \(1\)-form、面を貫いて測る量を
\(2\)-form、体積にわたって測る量を \(3\)-form として置く。成分で書けば：

\[
E = E_x\,dx+E_y\,dy+E_z\,dz,
\qquad
H = H_x\,dx+H_y\,dy+H_z\,dz
\]

\[D = D_x\,dy\wedge dz + D_y\,dz\wedge dx + D_z\,dx\wedge dy\]

\[B = B_x\,dy\wedge dz + B_y\,dz\wedge dx + B_z\,dx\wedge dy\]

\[J = J_x\,dy\wedge dz + J_y\,dz\wedge dx + J_z\,dx\wedge dy\]

\[
\rho = \rho\,dx\wedge dy\wedge dz
\]

ここでは記号を簡単にするため、電荷密度を表す係数と、それを体積形式に乗せた
\(3\)-form の両方を同じ \(\rho\) で書いている。

\(E,H\) は線に沿って測る \(1\)-form、\(D,B,J\) は面を貫いて測る
\(2\)-form、\(\rho\) は体積にわたって測る \(3\)-form である。

\hypertarget{c.2-ux96fbux8377ux306bux95a2ux3059ux308bux30acux30a6ux30b9ux306eux6cd5ux5247}{%
\subsection{C.2
電荷に関するガウスの法則}\label{c.2-ux96fbux8377ux306bux95a2ux3059ux308bux30acux30a6ux30b9ux306eux6cd5ux5247}}

積分形は：

\[
\int_{\partial V}D=\int_V\rho
\]

左辺に一般ストークスの定理を使うと：

\[
\int_{\partial V}D=\int_V dD
\]

したがって任意の領域 \(V\) で

\[
\int_V dD=\int_V\rho
\]

が成り立つなら、局所形は：

\[
dD=\rho
\]

である。成分展開すると：

\[
dD
=
\left(\frac{\partial D_x}{\partial x}
{}+\frac{\partial D_y}{\partial y}
{}+\frac{\partial D_z}{\partial z}\right)
dx\wedge dy\wedge dz
\]

ゆえに：

\[
\frac{\partial D_x}{\partial x}
{}+\frac{\partial D_y}{\partial y}
{}+\frac{\partial D_z}{\partial z}
=\rho
\]

\hypertarget{c.3-ux78c1ux675fux306eux6cd5ux5247}{%
\subsection{C.3 磁束の法則}\label{c.3-ux78c1ux675fux306eux6cd5ux5247}}

閉曲面を貫く正味の磁束がゼロである、という積分形は：

\[
\int_{\partial V}B=0
\]

一般ストークスの定理より：

\[
\int_V dB=0
\]

任意の領域 \(V\) について成り立つなら、局所形は：

\[
dB=0
\]

である。成分展開すると：

\[
dB
=
\left(\frac{\partial B_x}{\partial x}
{}+\frac{\partial B_y}{\partial y}
{}+\frac{\partial B_z}{\partial z}\right)
dx\wedge dy\wedge dz
\]

したがって：

\[
\frac{\partial B_x}{\partial x}
{}+\frac{\partial B_y}{\partial y}
{}+\frac{\partial B_z}{\partial z}
=0
\]

\hypertarget{c.4-ux30d5ux30a1ux30e9ux30c7ux30fcux306eux6cd5ux5247}{%
\subsection{C.4
ファラデーの法則}\label{c.4-ux30d5ux30a1ux30e9ux30c7ux30fcux306eux6cd5ux5247}}

積分形は：

\[
\int_{\partial S}E
=
-\frac{\partial}{\partial t}\int_S B
\]

左辺に一般ストークスの定理を使うと：

\[
\int_S dE
=
-\frac{\partial}{\partial t}\int_S B
\]

任意の曲面 \(S\) について成り立つなら、局所形は：

\[
dE=-\frac{\partial B}{\partial t}
\]

である。成分展開すると：

\[dE = \left(\frac{\partial E_z}{\partial y}-\frac{\partial E_y}{\partial z}\right)dy\wedge dz + \left(\frac{\partial E_x}{\partial z}-\frac{\partial E_z}{\partial x}\right)dz\wedge dx + \left(\frac{\partial E_y}{\partial x}-\frac{\partial E_x}{\partial y}\right)dx\wedge dy\]

一方、

\[
-\frac{\partial B}{\partial t}
=
-\frac{\partial B_x}{\partial t}\,dy\wedge dz
-\frac{\partial B_y}{\partial t}\,dz\wedge dx
-\frac{\partial B_z}{\partial t}\,dx\wedge dy
\]

だから、各成分は：

\[
\frac{\partial E_z}{\partial y}
-
\frac{\partial E_y}{\partial z}
=
-\frac{\partial B_x}{\partial t}
\]

\[
\frac{\partial E_x}{\partial z}
-
\frac{\partial E_z}{\partial x}
=
-\frac{\partial B_y}{\partial t}
\]

\[
\frac{\partial E_y}{\partial x}
-
\frac{\partial E_x}{\partial y}
=
-\frac{\partial B_z}{\partial t}
\]

となる。

\hypertarget{c.5-ux30a2ux30f3ux30daux30fcux30ebux30deux30afux30b9ux30a6ux30a7ux30ebux306eux6cd5ux5247}{%
\subsection{C.5
アンペール・マクスウェルの法則}\label{c.5-ux30a2ux30f3ux30daux30fcux30ebux30deux30afux30b9ux30a6ux30a7ux30ebux306eux6cd5ux5247}}

積分形は：

\[
\int_{\partial S}H
=
\int_SJ+\frac{\partial}{\partial t}\int_S D
\]

左辺に一般ストークスの定理を使うと：

\[
\int_S dH
=
\int_SJ+\frac{\partial}{\partial t}\int_S D
\]

任意の曲面 \(S\) について成り立つなら、局所形は：

\[
dH=J+\frac{\partial D}{\partial t}
\]

である。成分展開すると：

\[dH = \left(\frac{\partial H_z}{\partial y}-\frac{\partial H_y}{\partial z}\right)dy\wedge dz + \left(\frac{\partial H_x}{\partial z}-\frac{\partial H_z}{\partial x}\right)dz\wedge dx + \left(\frac{\partial H_y}{\partial x}-\frac{\partial H_x}{\partial y}\right)dx\wedge dy\]

また、

\[J+\frac{\partial D}{\partial t} = \left(J_x+\frac{\partial D_x}{\partial t}\right)dy\wedge dz + \left(J_y+\frac{\partial D_y}{\partial t}\right)dz\wedge dx + \left(J_z+\frac{\partial D_z}{\partial t}\right)dx\wedge dy\]

だから、各成分は：

\[
\frac{\partial H_z}{\partial y}
-
\frac{\partial H_y}{\partial z}
=
J_x+\frac{\partial D_x}{\partial t}
\]

\[
\frac{\partial H_x}{\partial z}
-
\frac{\partial H_z}{\partial x}
=
J_y+\frac{\partial D_y}{\partial t}
\]

\[
\frac{\partial H_y}{\partial x}
-
\frac{\partial H_x}{\partial y}
=
J_z+\frac{\partial D_z}{\partial t}
\]

となる。

\hypertarget{c.6-ux4f55ux306bux8a08ux91cfux304cux5fc5ux8981ux304b}{%
\subsection{C.6
何に計量が必要か}\label{c.6-ux4f55ux306bux8a08ux91cfux304cux5fc5ux8981ux304b}}

以上の4つの局所形は、外微分 \(d\)
と微分形式の次数だけで書ける。この段階では、長さや角度を測る計量は使っていない。

計量が必要になるのは、\(E\) と \(D\)、\(B\) と \(H\)
を関係づけるときである。また、\(1\)-form と \(2\)-form
を対応させて、通常の三成分の場として同じ種類の対象のように扱うときにも、空間の測り方が必要になる。

次章で導入するホッジ・スターは、その対応を与える道具である。

\newpage

\hypertarget{ux7b2c6ux7ae0ux8a08ux91cf-g-ux3068ux30dbux30c3ux30b8ux30b9ux30bfux30fc-ast-ux5185ux7a4dux306eux53ecux559aux3068ux6b21ux6570ux306eux53cdux8ee2}{%
\chapter{\texorpdfstring{第6章：計量 \(g\) とホッジ・スター \(\ast\) ---
内積の召喚と次数の反転}{第6章：計量 g とホッジ・スター \textbackslash ast --- 内積の召喚と次数の反転}}\label{ux7b2c6ux7ae0ux8a08ux91cf-g-ux3068ux30dbux30c3ux30b8ux30b9ux30bfux30fc-ast-ux5185ux7a4dux306eux53ecux559aux3068ux6b21ux6570ux306eux53cdux8ee2}}

\hypertarget{ux8a00ux3044ux8a33ux306eux7d42ux7109ux5185ux7a4dux3092ux89e3ux653eux3059ux308b}{%
\subsection{§6.0
言い訳の終焉------内積を解放する}\label{ux8a00ux3044ux8a33ux306eux7d42ux7109ux5185ux7a4dux3092ux89e3ux653eux3059ux308b}}

　第2章で平行四辺形のスカラー面積を三つの射影面積の二乗和平方根として復元したとき、第3章で球の表面積
\(4\pi R^2\) や円周の弧長 \(2\pi R\)
を同じ発想で測ったとき------本書ではこれまで、長さや面積を計算するたびに「計量（内積）をまだ正式には入れていない」と何度か断ってきた。

読者もそろそろくどいと感じている頃だろう。正直に言えば、筆者もそろそろ飽きてきた。

そもそも我々は、実空間 \((x,y,z)\)
という直交デカルト座標の上で計算している。この空間ではピタゴラスの定理が成り立つ。そろそろ気兼ねなく、\textbf{内積}を使うことにしよう。

内積とは何か。すでに我々はこれと何度も遭遇している。第2章で向き付き面積を得たとき、3つの影の成分
\(A_{yz}, A_{zx}, A_{xy}\)
の二乗和の平方根をとれば小学校の面積に一致すると言った------あの「二乗和」こそ、自分自身との内積（の平方根）である。より一般に、縦ベクトル同士、あるいは横ベクトルとして表された
\(1\)-form
同士について、対応する成分同士を掛けて足し合わせたスカラー量を、ここでは\textbf{内積}と呼ぶ。

\ldots\ldots{}\textbf{以上だ}。おそらく多くの読者は高校数学で内積を習っているだろうから、「何を当たり前のことを」と思っていることだろう。しかし、本書では居心地が悪い思いをしながらも、今まで内積を正面から定義してこなかった。それには理由がある。

その理由は、内積を早く出しすぎると、これまで大事にしてきた区別が崩れてしまうからだ。

より進んだ立場では、「ベクトルの足し算・引き算」と「内積」は不可分ではなく、分けて考えることができる、と見る。この立場では「内積が入った空間」「入らない空間」を区別し、内積を定義するための行列を「計量テンソル」と呼ぶ。筆者もこの区別には敬意を払っている。だからこそ、これまで計量や内積をまだ正式には導入していない、と何度か断ってきたわけだ。

つまり、ここで本当に区別したいのは、「横ベクトル×縦ベクトルでスカラーを得る操作」と「同じ種類のベクトル同士で内積をとる操作」である。

\begin{quote}
\textbf{注}
（なぜ横と縦をここまで分けるのか）計量や内積を前面に出さなくても、かなり強力に物理を記述できる、という立場がある。たとえば、初等的な
\(n\)-form
の積分や、多くの学部レベルの物理計算では、その見方がよく働く。筆者はその立場も尊重している。一方で、長さ・角度・内積・ホッジ・スター
\(\ast\)
を使う計量つきの見方も、もちろん強力である。筆者個人としては、問題に応じてこの二つの見方を行き来できることが理想だと考えている。そして行き来するためには、まず「横ベクトル×縦ベクトルで測る操作」と「同種ベクトル同士で内積をとる操作」を混同しないことが必要になる。本書が横ベクトルと縦ベクトルをしつこく分けてきたのは、そのためである。
\end{quote}

第1章以来、我々は横ベクトル×縦ベクトルでスカラーを得る計算を山ほどやってきた。では、それと内積は何が違うのか。

結論から言えば、\textbf{まったく別物}だ。横ベクトル掛ける縦ベクトルは「横と縦という異種のベクトル間」の計算である。対して内積は「横と横」あるいは「縦と縦」という\textbf{同種のベクトル間}の計算だ。

そう、事は読者が思ったよりも深遠だ。実空間という特別な舞台では、この二つがたまたま同じ数値を返すことが多いために混同されてきた------その混同の上に、ベクトル解析という巨大な体系が築かれている。本章の目標は、この区別を明確にした上で、内積の自然な一般化として計量
\(g\) を導き、そこからホッジ・スター \(\ast\) の正体を暴くことである。

\begin{center}\rule{0.5\linewidth}{0.5pt}\end{center}

\hypertarget{ux30d1ux30e9ux30e1ux30fcux30bfux7a7aux9593ux306eux5185ux7a4dux8a08ux91cf-g-ux306eux6b63ux4f53}{%
\subsection{\texorpdfstring{§6.1 パラメータ空間の内積------計量 \(g\)
の正体}{§6.1 パラメータ空間の内積------計量 g の正体}}\label{ux30d1ux30e9ux30e1ux30fcux30bfux7a7aux9593ux306eux5185ux7a4dux8a08ux91cf-g-ux306eux6b63ux4f53}}

\hypertarget{ux5b9fux7a7aux9593ux306eux5185ux7a4d}{%
\subsubsection{6.1.1
実空間の内積}\label{ux5b9fux7a7aux9593ux306eux5185ux7a4d}}

実空間 \((x,y,z)\) における二つの縦ベクトルを成分で書こう。

\[\mathbf{v}_1 = \begin{pmatrix} x_1 \\ y_1 \\ z_1 \end{pmatrix},\qquad \mathbf{v}_2 = \begin{pmatrix} x_2 \\ y_2 \\ z_2 \end{pmatrix}\]

同じ添字の成分同士を掛けて総和をとる操作を、本書では\textbf{内積}（ドット積）と定義し、記号
\(\cdot\) で表す。

\[\mathbf{v}_1 \cdot \mathbf{v}_2 = x_1 x_2 + y_1 y_2 + z_1 z_2\]

\(xyz\)
座標での内積は、対応する成分同士を掛けて足し合わせるだけで計算できる。

\hypertarget{ux5185ux7a4dux3092ux6e2cux5b9aux5668ux3068ux3057ux3066ux8aadux307fux76f4ux3059}{%
\subsubsection{6.1.2
内積を測定器として読み直す}\label{ux5185ux7a4dux3092ux6e2cux5b9aux5668ux3068ux3057ux3066ux8aadux307fux76f4ux3059}}

ここで、同じ計算を「測定器」の言葉で読み直しておこう。内積は一見すると、これまで扱ってきた測定器とは別系統の演算に見える。しかし実は、内積もまた「あるものを測定器に変え、別のものに作用させてスカラーを得る」操作として読むことができる。

その入り口が転置である。\(\mathbf{v}_1\)
を転置すると、縦ベクトルは横ベクトルになる。

\[\mathbf{v}_1^T = \begin{pmatrix} x_1 & y_1 & z_1 \end{pmatrix}\]

これは、ペアになる縦ベクトル \(\mathbf{v}_2\)
を食べて、「\(\mathbf{v}_1\) との内積はいくらか」を測る測定器である。

\[\mathbf{v}_1^T\mathbf{v}_2
= \begin{pmatrix} x_1 & y_1 & z_1 \end{pmatrix}
\begin{pmatrix} x_2 \\ y_2 \\ z_2 \end{pmatrix}
= x_1x_2 + y_1y_2 + z_1z_2\]

つまり、実空間の直交デカルト座標では、縦ベクトル \(\mathbf{v}_1\)
を転置するだけで、そのベクトルと内積をとるための横ベクトル------測定器------が得られる。

\begin{quote}
\textbf{注}
（添字の上げ下げの正体）テンソル解析では、ベクトル成分と測定器側の成分を、上付き添字・下付き添字で区別する流儀がある。そして計量
\(g\)
を使って添字を上げ下げする、と説明される。本書の言葉で言えば、その正体は「縦ベクトルを、内積を測る横ベクトルへ変換する」操作である。直交デカルト座標では
\(g=I\) なので転置だけで済むが、次に見るようにパラメータ空間では途中に
\(\mathbf{g}=J^T J\) を挟む必要が出てくる。
\end{quote}

\hypertarget{ux30d1ux30e9ux30e1ux30fcux30bfux7a7aux9593ux3067ux5185ux7a4dux3092ux3068ux308bux306bux306f}{%
\subsubsection{6.1.3
パラメータ空間で内積をとるには}\label{ux30d1ux30e9ux30e1ux30fcux30bfux7a7aux9593ux3067ux5185ux7a4dux3092ux3068ux308bux306bux306f}}

では、同じ読み替えをパラメータ空間 \((r,\theta,z)\) で行うとどうなるか。

実空間の直交デカルト座標では、\(\mathbf{v}_1^T\)
がそのまま「\(\mathbf{v}_1\)
との内積を測る測定器」だった。だからパラメータ空間でも、\(\mathbf{v}_1^T\)
をそのまま測定器にしたくなる。しかしそれでは、\(\Delta r_1 \Delta r_2 + \Delta\theta_1 \Delta\theta_2 + \Delta z_1 \Delta z_2\)
を計算してしまうだけで、実空間での正しい内積にはならない。実空間とパラメータ空間のメモリの対応が場所によって異なるからである。

では、パラメータ空間の縦ベクトルを、実空間の内積を正しく測る横ベクトルに変えるには、何を挟めばよいか。

第4章で、座標変換のヤコビ行列が成分をどのように変換するかを見た。ここではその同じ行列
\(J\) を使って、パラメータ空間の変位を実空間の変位として読む。

\[J = \begin{pmatrix}
\frac{\partial x}{\partial r} & \frac{\partial x}{\partial \theta} & \frac{\partial x}{\partial z} \\
\frac{\partial y}{\partial r} & \frac{\partial y}{\partial \theta} & \frac{\partial y}{\partial z} \\
\frac{\partial z}{\partial r} & \frac{\partial z}{\partial \theta} & \frac{\partial z}{\partial z}
\end{pmatrix}\]

円柱座標の換算式 \(x = r\cos\theta,\; y = r\sin\theta,\; z = z\)
から偏微分を計算すれば：

\[J = \begin{pmatrix}
\cos\theta & -r\sin\theta & 0 \\
\sin\theta & r\cos\theta & 0 \\
0 & 0 & 1
\end{pmatrix}\]

パラメータ空間のベクトル \(\mathbf{v}_1, \mathbf{v}_2\) を成分で書く。

\[\mathbf{v}_1 = \begin{pmatrix} \Delta r_1 \\ \Delta\theta_1 \\ \Delta z_1 \end{pmatrix},\qquad \mathbf{v}_2 = \begin{pmatrix} \Delta r_2 \\ \Delta\theta_2 \\ \Delta z_2 \end{pmatrix}\]

これらを実空間の変位として読むと、それぞれ
\(J\mathbf{v}_1, J\mathbf{v}_2\) になる。実空間で「\(J\mathbf{v}_1\)
との内積」を測る測定器は、前節の読み方では \((J\mathbf{v}_1)^T\)
である。したがって、求めたい内積は次の形で書ける。

\[\text{内積} = (J\mathbf{v}_1)^T (J\mathbf{v}_2)\]

ここで転置の規則を使うと、\((J\mathbf{v}_1)^T = \mathbf{v}_1^T J^T\)
だから、

\[= \mathbf{v}_1^T (J^T J) \mathbf{v}_2\]

ここで \(J^T J\) を実際に計算してみよう。\(J^T\) は \(J\)
の行と列を入れ替えたものだ。

\[J^T = \begin{pmatrix}
\cos\theta & \sin\theta & 0 \\
-r\sin\theta & r\cos\theta & 0 \\
0 & 0 & 1
\end{pmatrix}\]

したがって、

\[J^T J = \begin{pmatrix}
\cos\theta & \sin\theta & 0 \\
-r\sin\theta & r\cos\theta & 0 \\
0 & 0 & 1
\end{pmatrix}
\begin{pmatrix}
\cos\theta & -r\sin\theta & 0 \\
\sin\theta & r\cos\theta & 0 \\
0 & 0 & 1
\end{pmatrix}
= \begin{pmatrix}
1 & 0 & 0 \\
0 & r^2 & 0 \\
0 & 0 & 1
\end{pmatrix}\]

\(\cos^2\theta + \sin^2\theta = 1\)
という代数的恒等式がすべての幾何を代行し、対角成分に \(1, r^2, 1\)
が残る。\(\theta\)
方向の弧の長さなど一度も考えていない------行列の積と転置の性質だけを使った、純粋に代数的な変形だ。

\textbf{これで、新しい測定器が出てきた。} パラメータ空間の縦ベクトル
\(\mathbf{v}_1\)
から、実空間での内積を正しく測る横ベクトルを作るには、単に転置するのではなく、右側に
\(J^T J\) を掛ければよい。

\[\omega_{\mathbf{v}_1} = \mathbf{v}_1^T (J^T J)\]

そしてこの測定器を \(\mathbf{v}_2\)
に作用させれば、実空間での内積と同じ値が得られる。

この \(J^T J\) を、\textbf{計量行列 $\mathbf{g}$}と呼ぶ。

\[\mathbf{g} = J^T J\]

\(xyz\)
座標では、そもそも「引き戻す」という操作を通らず、成分の積和だけで内積を計算していた。パラメータ空間に来ると、まず
\(J\)
で実空間の変位として読み、その後で成分の積和をとる。この二段階を行列の形に畳み込むと、中央に
\(\mathbf{g}=J^T J\) が現れる------そしてその中身は、ヤコビ行列 \(J\)
さえ知っていれば、偏微分と行列の積だけで機械的に求まる。新しい幾何学の公理など、どこにも入り込む余地はない。

念のため、具体的な数値で確認しておこう。\(r=2,\; \theta=\pi/3\)
の地点を考える。このとき \(\cos\theta = 1/2,\; \sin\theta = \sqrt{3}/2\)
だから、

\[J = \begin{pmatrix}
1/2 & -\sqrt{3} & 0 \\
\sqrt{3}/2 & 1 & 0 \\
0 & 0 & 1
\end{pmatrix}\]

パラメータ空間で
\(\mathbf{v} = \begin{pmatrix}0 \\ 1 \\ 0\end{pmatrix}\) という
\(\theta\)
方向の単位変位を考える。\(\mathbf{g} = \begin{pmatrix} 1 & 0 & 0 \\ 0 & r^2 & 0 \\ 0 & 0 & 1 \end{pmatrix}\)
であり、いま \(r=2\) だから
\(\mathbf{g} = \begin{pmatrix} 1 & 0 & 0 \\ 0 & 4 & 0 \\ 0 & 0 & 1 \end{pmatrix}\)。パラメータ空間で内積（長さの二乗）を計算すると：

\[\mathbf{v}^T \mathbf{g} \,\mathbf{v} = \begin{pmatrix} 0 & 1 & 0 \end{pmatrix}
\begin{pmatrix} 1 & 0 & 0 \\ 0 & 4 & 0 \\ 0 & 0 & 1 \end{pmatrix}
\begin{pmatrix} 0 \\ 1 \\ 0 \end{pmatrix} = 4\]

一方、同じ \(\mathbf{v}\) を実空間の変位として読むと：

\[J\mathbf{v} = \begin{pmatrix}
1/2 & -\sqrt{3} & 0 \\
\sqrt{3}/2 & 1 & 0 \\
0 & 0 & 1
\end{pmatrix}
\begin{pmatrix} 0 \\ 1 \\ 0 \end{pmatrix} = \begin{pmatrix} -\sqrt{3} \\ 1 \\ 0 \end{pmatrix}\]

実空間での内積（長さの二乗）は
\((-\sqrt{3})^2 + 1^2 + 0^2 = 4\)。確かに一致している。\(\mathbf{v} = \begin{pmatrix}1 \\ 0 \\ 0\end{pmatrix}\)（\(r\)
方向の単位変位）でも同様に、\(\mathbf{v}^T \mathbf{g} \,\mathbf{v} = 1\)、\((J\mathbf{v})^T (J\mathbf{v}) = \cos^2\theta + \sin^2\theta = 1\)
で一致する。\(\mathbf{g} = J^T J\)
はこうして、パラメータ空間と実空間の内積を矛盾なく繫いでいる。

なお、ここでは円柱座標 \((r,\theta,z)\) で計算したが、球座標
\((r,\theta,\phi)\) の場合も同様に \(J\) から \(\mathbf{g}\) が求まる。

\[\mathbf{g}_{\text{球座標}} = \begin{pmatrix} 1 & 0 & 0 \\ 0 & r^2 & 0 \\ 0 & 0 & r^2\sin^2\theta \end{pmatrix}\]

対角成分に \(r^2\) と \(r^2\sin^2\theta\)
が現れるが、これらもすべて座標変換の一次係数を並べた \(J\)
から機械的に導かれる換算係数である。

この \(g\)
を仲立ちにすることで、どの座標系でも一貫して内積（長さ・角度）を計算できる。また、後半で扱うホッジ・スター
\(\ast\) の働きも、この内積のルールから自然に読み解ける。

\begin{center}\rule{0.5\linewidth}{0.5pt}\end{center}

\hypertarget{g-ux306bux3088ux308bux7e26ux30d9ux30afux30c8ux30ebux3068ux6a2aux30d9ux30afux30c8ux30ebux306eux5909ux63db}{%
\subsection{\texorpdfstring{§6.2 \(g\)
による縦ベクトルと横ベクトルの変換}{§6.2 g による縦ベクトルと横ベクトルの変換}}\label{g-ux306bux3088ux308bux7e26ux30d9ux30afux30c8ux30ebux3068ux6a2aux30d9ux30afux30c8ux30ebux306eux5909ux63db}}

\hypertarget{mathbfvt-mathbfg-ux306fux6a2aux30d9ux30afux30c8ux30ebux3067ux3042ux308b}{%
\subsubsection{\texorpdfstring{6.2.1 \(\mathbf{v}^T \mathbf{g}\)
は横ベクトルである}{6.2.1 \textbackslash mathbf\{v\}\^{}T \textbackslash mathbf\{g\} は横ベクトルである}}\label{mathbfvt-mathbfg-ux306fux6a2aux30d9ux30afux30c8ux30ebux3067ux3042ux308b}}

§6.1 で導いた内積の式 \(\mathbf{v}_1^T \mathbf{g} \,\mathbf{v}_2\)
を、もう一度見てみよう。行列の積は左から順に実行できるから、この式は「まず
\(\mathbf{v}_1^T \mathbf{g}\) を計算し、その結果に \(\mathbf{v}_2\)
を掛ける」と読むことができる。

\(\mathbf{v}_1\) は縦ベクトルだが、転置された \(\mathbf{v}_1^T\)
は横ベクトルだ。横ベクトルに正方行列を掛けると、結果は再び横ベクトルになる。つまり、\(\mathbf{v}_1^T \mathbf{g}\)
という部分だけを取り出せば、

\[\omega = \mathbf{v}^T \mathbf{g}\]

これは \(1 \times 3\) の\textbf{横ベクトル}である。計量 \(g\)
とは、内積の行列形式から自然に現れた、\textbf{縦ベクトルを横ベクトルに変換する装置}なのだ。

逆向きも書いておこう。横ベクトルとしての \(1\)-form

\[
\omega =
\begin{pmatrix}
\omega_1 & \omega_2 & \omega_3
\end{pmatrix}
\]

が与えられたとする。これに対応する縦ベクトル \(\mathbf{v}_\omega\) は、

\[
\mathbf{v}_\omega = \mathbf{g}^{-1}\omega^T
\]

である。

\begin{quote}
\textbf{注} （逆行列）正方行列 \(A\) に対して、 \[
A^{-1}A = AA^{-1} = I
\] を満たす行列 \(A^{-1}\) を、\(A\) の逆行列と呼ぶ。直感的には、\(A\)
による変換を元に戻す行列である。ここでは、\(\mathbf{g}\)
を使って作った横ベクトルをもとの縦ベクトルへ戻すために、\(\mathbf{g}^{-1}\)
を使っている。
\end{quote}

実際、\(\omega = \mathbf{v}^T\mathbf{g}\) なら、転置して

\[
\omega^T = \mathbf{g}^T\mathbf{v}
\]

となる。\(\mathbf{g}\) は対称行列なので
\(\mathbf{g}^T = \mathbf{g}\)。したがって、

\[
\mathbf{v} = \mathbf{g}^{-1}\omega^T
\]

で元に戻る。

実空間の直交デカルト座標では \(\mathbf{g}=I\)
なので、この逆向きの変換も単なる転置に見える。しかし一般の座標では、横ベクトルを縦ベクトルへ戻すために
\(\mathbf{g}^{-1}\) が必要になる。

\begin{quote}
\textbf{注} （第2章で密輸していたもの）第2章で \(2\)-form
の行列表現を探したとき、我々はすでに \(\mathbf{v}_1^T M \mathbf{v}_2\)
という形を使っていた。あのときは「縦ベクトルを横に寝かせる操作は数学的には行儀が悪い」とだけ断って、深入りせずに押し通した。いま見ている
\(\mathbf{v}^T\mathbf{g}\)
は、そのとき密輸していた操作の正体である。実空間では \(\mathbf{g}=I\)
なので、転置だけで縦ベクトルが横ベクトルに見えてしまう。しかしパラメータ空間では、実空間の内積を再現するように縦ベクトルを測定器としての横ベクトルへ変えるには、計量
\(\mathbf{g}\) を明示的に通す必要がある。
\end{quote}

この変換の意味を考えよう。第1章以来、我々は「測られる図形（縦ベクトル
\(\mathbf{v}\)）」と「測る装置（横ベクトル
\(\omega\)）」を区別してきた。両者は異なる世界の住人だった。しかし \(g\)
を経由することで、図形の側にいた縦ベクトルが、はかりの側------横ベクトル------へと姿を変える。同じ矢印でも、場所によって「はかりとしての感度」が変わるのは、\(\mathbf{g}\)
がその場所ごとの目盛りの違いをすべて吸収してくれるからだ。

実空間（\(\mathbf{g}=I\)）では \(\omega = \mathbf{v}^T\)
であり、縦ベクトルの成分がそのまま横ベクトルの成分になる。パラメータ空間では
\(\mathbf{g} = \begin{pmatrix} 1 & 0 & 0 \\ 0 & r^2 & 0 \\ 0 & 0 & 1 \end{pmatrix}\)
なので、\(\omega = \begin{pmatrix} \Delta r & r^2\Delta\theta & \Delta z \end{pmatrix}\)。\(\theta\)
成分だけ \(r^2\) 倍されるのは、\(\theta\) 方向の目盛りが \(r\)
倍に伸びていることの反映だ。

\(g\)
がなければ、「測られる図形」と「測る装置」は互いに変換不能である。\(g\)
は両者を繫ぐ唯一の仲介者なのだ。

\hypertarget{ux8a08ux91cfux304cux4e0eux3048ux308bux3082ux306eux5e73ux884c-k-ux9762ux4f53ux306eux5e7eux4f55ux5b66ux7684ux306aux5927ux304dux3055}{%
\subsubsection{\texorpdfstring{6.2.2 計量が与えるもの------平行 \(k\)
面体の幾何学的な大きさ}{6.2.2 計量が与えるもの------平行 k 面体の幾何学的な大きさ}}\label{ux8a08ux91cfux304cux4e0eux3048ux308bux3082ux306eux5e73ux884c-k-ux9762ux4f53ux306eux5e7eux4f55ux5b66ux7684ux306aux5927ux304dux3055}}

計量 \(g\) の最も基本的な役割は、\(k\)
本のベクトルが貼る\textbf{平行 $k$ 面体の実際の幾何学的な大きさ}（長さ・面積・体積）を定義することである。

本書のこれまでの章では、\(k\)-form に \(k\)
本のベクトルを食わせることで「向き付きの数え上げ」を行ってきた。\(dx(\mathbf{v})\)
は \(\mathbf{v}\) の \(x\)
成分という数、\(dx \wedge dy(\mathbf{v}_1, \mathbf{v}_2)\) は
\(\mathbf{v}_1, \mathbf{v}_2\)
が張る平行四辺形の向き付き面積という数である。しかしこれらはあくまで影の測量であり、ベクトルたちが実際に何メートルの線分を張り、何平方メートルの平行四辺形を張っているか------\textbf{幾何学的な本当の大きさ}------は、計量
\(g\) が与えられて初めて決まる。

\(k=1\)（1本のベクトル \(\mathbf{v}\)
が張る線分）の場合、その長さの二乗は自分自身との内積で定義される。これは
§6.1 ですでに導いた通りだ。

\[|\mathbf{v}|^2 = \mathbf{v}^T \mathbf{g} \,\mathbf{v}\]

実空間 \((\mathbf{g}=I)\) では
\(|\mathbf{v}|^2 = \Delta x^2 + \Delta y^2 + \Delta z^2\)。円柱座標のパラメータ空間では
\(|\mathbf{v}|^2 = \Delta r^2 + r^2\Delta\theta^2 + \Delta z^2\)。

\(k=0\)（0本のベクトル、すなわち点）の場合、\(0\)-form であるスカラー場
\(f\) の「大きさ」は単にその絶対値 \(|f|\) である。

\(k=2\)（2本のベクトルが張る平行四辺形の面積）および
\(k=3\)（3本のベクトルが張る平行六面体の体積）についても、同じことだ。第2章
§2.4.7 で見たように、\(dx \wedge dy, dy \wedge dz, dz \wedge dx\)
という3つの基底 \(2\)-form の係数 \(A_{xy}, A_{yz}, A_{zx}\)
の二乗和の平方根 \(\sqrt{A_{yz}^2 + A_{zx}^2 + A_{xy}^2}\)
が、平行四辺形の符号なし面積を与える。\(xyz\)
座標でこの計算に必要なのは基底 \(2\)-form
の係数だけであり、それ以外の情報はいらない。

\begin{quote}
\textbf{注} （なぜ \(k=2\) ではベクトル表示しないのか）\(k=1\) では
\(\mathbf{v}\) という縦ベクトルの内積
\(\mathbf{v}^T\mathbf{g}\,\mathbf{v}\) として長さを書いた。\(k=2\)
での同様の内積------行列同士の内積（付録Dのフロベニウス積）------はまだ定義していない。さらに、\(2\)-form
の係数を \(3\)
成分のベクトルとして表示するには、次数を変換するホッジ・スター \(\ast\)
も必要になる。ここではまだ、面を面のまま、すなわち \(2\)-form
の係数として扱っている。
\end{quote}

計量 \(g\) の役割は、この「\(xyz\)
での計算」をパラメータ空間のような目盛りの歪んだ座標系へ一般化することである。いずれの場合も、\textbf{$g$ が決まれば、あらゆる次数の平行 $k$ 面体に一貫した幾何学的な大きさが定義できる}------これが計量の最も根本的な役割である。

\hypertarget{ux5185ux7a4dux306eux6027ux8cea}{%
\subsubsection{6.2.3 内積の性質}\label{ux5185ux7a4dux306eux6027ux8cea}}

第5章で外微分 \(d\)
の性質（線形性・次数を上げる・\(d^2=0\)・ライプニッツ則）をまとめたのと同じように、内積についても、ここまでの計算から浮かび上がる性質を整理しておこう。これらは特定の座標系に依存しない。

\begin{enumerate}
\def\labelenumi{\arabic{enumi}.}
\item
  \textbf{対称性}：\(\mathbf{v}_1 \cdot \mathbf{v}_2 = \mathbf{v}_2 \cdot \mathbf{v}_1\)。これを
  §6.1 で導いた行列の形 \(\mathbf{v}_1^T \mathbf{g} \,\mathbf{v}_2\)
  に当てはめてみよう。\(\mathbf{v}_2 \cdot \mathbf{v}_1 = \mathbf{v}_2^T \mathbf{g} \,\mathbf{v}_1\)
  だが、スカラーは転置しても変わらないから
  \(\mathbf{v}_2^T \mathbf{g} \,\mathbf{v}_1 = (\mathbf{v}_2^T \mathbf{g} \,\mathbf{v}_1)^T = \mathbf{v}_1^T \mathbf{g}^T \mathbf{v}_2\)。これが
  \(\mathbf{v}_1^T \mathbf{g} \,\mathbf{v}_2\) と等しくなるには
  \(\mathbf{g} = \mathbf{g}^T\)
  でなければならない。つまり、内積の対称性から、\(\mathbf{g}\)
  は必ず\textbf{対称行列}になることが導かれる。
\item
  \textbf{線形性}（双線形性）：\((a\mathbf{v}_1 + b\mathbf{v}_2) \cdot \mathbf{v}_3 = a(\mathbf{v}_1 \cdot \mathbf{v}_3) + b(\mathbf{v}_2 \cdot \mathbf{v}_3)\)。各引数について線形。これは
  \((a\mathbf{v}_1 + b\mathbf{v}_2)^T \mathbf{g} \,\mathbf{v}_3 = a\mathbf{v}_1^T \mathbf{g} \,\mathbf{v}_3 + b\mathbf{v}_2^T \mathbf{g} \,\mathbf{v}_3\)
  という、行列の積の分配法則そのものである。
\item
  \textbf{正定値性}：\(\mathbf{v} \cdot \mathbf{v} > 0\)（\(\mathbf{v} \neq \mathbf{0}\)）。「長さの二乗」が正であることを保証する。特殊相対論のように符号が混ざる場合はこの限りではないが、本書の舞台である3次元実空間ではつねに成り立つ。
\end{enumerate}

\begin{quote}
\textbf{注} （対称行列と反対称行列の対比）
第2章で面積の計量器が反対称行列（\(M = -M^T\)）を要請したことと、
ここで内積が対称行列（\(\mathbf{g} = \mathbf{g}^T\)）を要請することは、
ちょうど対をなしている。前者は「重なればゼロ」、後者は「順番によらず同じ値」という、どちらもごく自然な要請から導かれる。行列の世界では、反対称が向き付きの数え上げ（符号付きの面積・体積）を、対称が向きなしの大きさ（符号なし面積・体積）と角度を司る。
\end{quote}

\hypertarget{ux8a08ux91cfux884cux5217ux306eux5e7eux4f55ux3092ux6210ux5206ux306bux5bbfux3089ux305bux308b}{%
\subsubsection{6.2.4
計量行列の幾何を成分に宿らせる}\label{ux8a08ux91cfux884cux5217ux306eux5e7eux4f55ux3092ux6210ux5206ux306bux5bbfux3089ux305bux308b}}

\(\mathbf{g}\)
の対角成分は「その軸方向の一歩が長さにどれだけ効くか」、非対角成分は「異なる軸方向の一歩が互いにどれだけ干渉するか」を表す。円柱座標の
\(\mathbf{g}\) が対角行列だったのは、\(r,\theta,z\)
の各軸が互いに直交しているからだ。いずれにせよ、\(\mathbf{g}\)
という一個の行列が、その場所での「ものさしの性質」のすべてを成分として担っている。

\begin{quote}
\textbf{注}
（公理的な内積------公理と行列の順序逆転）標準的な線形代数の教科書では、まず上記の3性質を公理として定義し、その後に「これらを満たす演算は、ある基底を選べば必ず対称行列
\(\mathbf{g}\) を用いた \(\mathbf{v}^T \mathbf{g} \mathbf{w}\)
の形で書ける」という定理を導く。論理の順番が我々とは逆だ。公理的な立場では行列表示は一表現に過ぎないが、測定器として読む本書の立場では\textbf{この行列 $\mathbf{g}$ こそが空間の目盛りそのもの}であり、すべての測定の出発点である。どちらが正しいかではない。抽象的な性質から入るか、具体的なシンボルの操作による測り方から入るか。向いている方向が違うだけだ。
\end{quote}

\begin{quote}
\textbf{【ここまでのチェックポイント — §6.0–§6.2】}

\begin{itemize}
\item
  内積とは、同種のベクトル間（横×横、または縦×縦）で成分の積和をとる演算。横×縦とは別物。
\item
  パラメータ空間での内積は、実空間へ引き戻してから計算する。その過程で
  \(J^T J\) が自然に現れる。これが計量 \(g\) の正体。
\item
  \(g\) は縦ベクトルを横ベクトルに変換する。具体的には、縦ベクトル
  \(\mathbf{v}\) から対応する \(1\)-form を作るには
  \(\omega_{\mathbf{v}} = \mathbf{v}^T\mathbf{g}\) とする。
\item
  逆に、横ベクトルとしての \(1\)-form \(\omega\)
  から対応する縦ベクトルを作るには、\(\mathbf{v}_\omega = \mathbf{g}^{-1}\omega^T\)
  とする。デカルト座標では \(\mathbf{g}=I\)
  なので、どちらもほとんど転置だけに見える。
\item
  \(g\) を仲立ちに、あらゆる次数の平行 \(k\)
  面体の幾何学的な大きさが定義できる。
\item
  内積の性質（対称性・線形性・正定値性）から \(g\)
  の対称性が導かれる。これは面積測定器 \(\wedge\) の反対称性と対をなす。
\end{itemize}
\end{quote}

\begin{center}\rule{0.5\linewidth}{0.5pt}\end{center}

\hypertarget{ux30dbux30c3ux30b8ux30b9ux30bfux30fc-ast-ux4e8cux3064ux306eux65b9ux6cd5ux3092ux7e6bux3050ux5bfeux5fdc}{%
\subsection{\texorpdfstring{§6.3 ホッジ・スター \(\ast\)
------二つの方法を繫ぐ対応}{§6.3 ホッジ・スター \textbackslash ast ------二つの方法を繫ぐ対応}}\label{ux30dbux30c3ux30b8ux30b9ux30bfux30fc-ast-ux4e8cux3064ux306eux65b9ux6cd5ux3092ux7e6bux3050ux5bfeux5fdc}}

\hypertarget{ux30b9ux30abux30e9ux30fcux3092ux5f97ux308bux4e8cux3064ux306eux65b9ux6cd5}{%
\subsubsection{6.3.1
スカラーを得る二つの方法}\label{ux30b9ux30abux30e9ux30fcux3092ux5f97ux308bux4e8cux3064ux306eux65b9ux6cd5}}

ここまでで \(g = J^T J\)
の正体が明らかになり、内積をどの座標系でも計算できるようになった。ここで、より大きな構図を見ておこう。

物理の式を眺めると、スカラー量を作るやり方が一つではないことに気づく。ある場を線や面に沿って積分するときは、「測るもの」が「測られるもの」に作用して数を返す、という書き方が自然に現れる。一方で、仕事率やエネルギー密度のような量は、しばしば
\(\mathbf{E}\cdot\mathbf{J}\) や
\(\mathbf{A}\cdot(\nabla\times\mathbf{A})\)
のように、同じ種類の矢印同士の内積として書かれる。

つまり世の中の物理記法には、少なくとも見かけ上、スカラーを得る二つの流儀が併存している。本書の言葉で言えば、それは次の二つである。

\begin{itemize}
\item
  \textbf{方法①（微分形式の方法 / dの方法）}：測る装置（横ベクトル：\(1\)-form
  など）と、測られる図形（縦ベクトル：ベクトル場など）を用意し、横×縦の計算でスカラーを得る。この計算で使うのは、形式がベクトルを食べてスカラーを返すという、作用と評価の論理である。
\item
  \textbf{方法②（ベクトル解析の方法 / ∇の方法）}：すべての量を「同じ種類の矢印」に統一し、同種ベクトル間の内積（縦×縦、または横×横）でスカラーを得る。この計算では、同種ベクトル同士を比べるために計量
  \(g\)
  を挟む。多くの読者は学校教育を通じてこちらの方法に慣れ親しんでいるだろう。
\end{itemize}

\begin{quote}
\textbf{注}
（二つの方法の呼び方------筆者の苦悩）筆者はこの二つの流儀を教科書でどう呼ぶか長らく悩んだ。方法①を「微分形式の方法」と呼ぶのが素直だろうが、本書では「dの方法」とも併記することにした。なぜなら、微分形式を知らない読者にとっては、dという記号の方が親しみやすいからだ。筆者は個人的に「トポロジーの方法」とも呼んでいる。トポロジーや幾何学の文脈で育った読者には、そちらの呼び方の方がしっくりくるだろう。同様に、方法②を「ベクトル解析の方法」と呼ぶのが素直だろうが、本書では「∇の方法」とも併記することにした。∇（ナブラ）という記号は、物理学を学んだ読者にとって最も馴染み深い記号だからだ。筆者は個人的に「内積の方法」とも呼んでいる。内積の構造が核心に据わっていることを強調したいときに使う。どれも同じ二つの流儀を指している。
\end{quote}

\begin{quote}
\textbf{注}
（二つの記法の系譜）方法①はグラスマンの外積代数やエリ・カルタンの流れに近く、方法②はギブズやヘヴィサイドが整備したベクトル解析の流れに近い。どちらかが正しく、どちらかが誤りという話ではない。
\end{quote}

実は、どちらの方法を採用しても、最終的に得られる物理的な結果は同じになる。これは単なる表現の違いだ。

しかしここに実務上の問題が生じる。方法②の枠組みでは、本来「面」で測るべき量（\(2\)-form）も「線」で測る量（\(1\)-form）も、すべて同じ「\(3\)
成分の矢印」として扱われる。方法①における面積や体積の概念を、方法②の「矢印同士の内積」の枠組みで扱うには、次数の異なる形式どうしを結びつける対応が必要になる。

\(k\)-form の独立成分数には、特徴的な対称性があった。第2章 §2.5.9
で見たように、\(3\) 次元空間では \(0\)-form が \(1\) 成分、\(1\)-form が
\(3\) 成分、\(2\)-form も \(3\) 成分、\(3\)-form が \(1\)
成分------真ん中を折り目にして、\(1, 3, 3, 1\)
と左右対称になっている。独立成分数が同じ \(1\)-form と
\(2\)-form、\(0\)-form と \(3\)-form
のあいだには、何らかの対応関係が存在するはずだ。その対応を与える装置こそが、\textbf{ホッジ・スター $\ast$}である。

\(\ast\) は、内積の情報（すなわち
\(g\)）を使って、ある形式に対し、それと組み合わせれば体積になる相手を対応させる。まず本文では実空間の基底でこの対応を作り、成分表示としての確認は付録Dに回そう。

\begin{quote}
\textbf{注}
（二つの方法のあいだで）方法①と方法②のどちらか一方だけを使う必要はない。筆者の理想は、両方を自由に行き来できることだ。本書が方法①（微分形式）から出発したのは、空間の歪みに依存しない強固な土台を築くためである。その上で、必要に応じて方法②の言葉でも同じ量を読めるようにする------その橋渡しをするのが
\(\ast\) だ。
\end{quote}

\begin{quote}
\textbf{注}
（ベクトル解析を知っている読者へ）通常のベクトル解析で面の向きや外積として扱われる操作の多くは、この
\(\ast\)
が与える対応として読み直せる。ただし本章では、それらを既知の公式として使うのではなく、\(g\)
と第2章の基底 \(2\)-form から作り直す。
\end{quote}

\hypertarget{ux5b9fux7a7aux9593ux3067ux306e-ast-ux8f9eux66f8ux4f53ux7a4dux306bux306aux308bux76f8ux624bux3092ux9078ux3076}{%
\subsubsection{\texorpdfstring{6.3.2 実空間での \(\ast\)
辞書------体積になる相手を選ぶ}{6.3.2 実空間での \textbackslash ast 辞書------体積になる相手を選ぶ}}\label{ux5b9fux7a7aux9593ux3067ux306e-ast-ux8f9eux66f8ux4f53ux7a4dux306bux306aux308bux76f8ux624bux3092ux9078ux3076}}

では、\(\ast\) はどのように決まるのか。

ここではまず、実空間 \((x,y,z)\)
の直交デカルト座標で考える。この空間では、\(dx,dy,dz\)
は互いに直交し、長さ \(1\) の基準測定器である。また、空間の向きを

\[
dx \wedge dy \wedge dz
\]

が正の体積形式になるように決めておく。

\(\ast\)
は、ある形式に対して、それと組み合わせれば体積形式になる相手を返す操作だと考えればよい。

たとえば \(dx\) と組み合わせて正の体積形式を作る \(2\)-form
は何か。それは \(dy \wedge dz\) である。実際、

\[
dx \wedge dy \wedge dz
\]

が得られる。したがって、

\[
\ast dx = dy \wedge dz
\]

と決めるのが自然である。

同じように、

\[
dy \wedge dz \wedge dx = dx \wedge dy \wedge dz
\]

なので、

\[
\ast dy = dz \wedge dx
\]

であり、

\[
dz \wedge dx \wedge dy = dx \wedge dy \wedge dz
\]

なので、

\[
\ast dz = dx \wedge dy
\]

である。

つまり、\(\ast\)
はそれぞれの線の測定器に対して、それと組み合わせると正の体積形式になる面の測定器を対応させている。同じ考えで、スカラー
\(1\) は体積形式 \(dx \wedge dy \wedge dz\) と対応し、体積形式はスカラー
\(1\) と対応する。

したがって、実空間での \(\ast\) の辞書は次のようになる。

\[\begin{aligned}
\ast(1) &= dx \wedge dy \wedge dz, \qquad &\ast(dx \wedge dy \wedge dz) &= 1 \\
\ast(dx) &= dy \wedge dz, \qquad &\ast(dy \wedge dz) &= dx \\
\ast(dy) &= dz \wedge dx, \qquad &\ast(dz \wedge dx) &= dy \\
\ast(dz) &= dx \wedge dy, \qquad &\ast(dx \wedge dy) &= dz
\end{aligned}\]

この時点では、\(\ast\)
を「対応表」として読めば十分である。ただし、この対応は単なる暗記表ではない。基底を選べば、\(\ast\)
は係数を並べ替える線形変換として配列表示できる。全次数を含む配列表現は付録Dで確認する。

この辞書の意味するところは単純だ。\(x\)
軸方向の線分（\(dx\)）に対応するのは、\(x\)
軸に垂直な面（\(dy \wedge dz\)）である。同様に
\(y \leftrightarrow dz \wedge dx\)、\(z \leftrightarrow dx \wedge dy\)。そしてスカラー
\(1\)（\(0\)-form）と体積
\(dx \wedge dy \wedge dz\)（\(3\)-form）が互いに対応する。これがまさに、第2章で見た独立成分数
\(1, 3, 3, 1\) の対称性の現れである。

\begin{quote}
\textbf{注} （なぜこの辞書でよいのか）\(\ast dx = dy \wedge dz\)
となるのは、\(dx,dy,dz\) が正規直交基底（互いに直交し、長さ
\(1\)）であり、かつ空間の向きを \(dx \wedge dy \wedge dz\)
が正となる右手系にとっているからだ。この条件のもとでは、上の辞書が一意的に定まる。パラメータ空間まで考慮すると、\(\ast\)
の辞書は計量 \(g\)
と空間の向きの選び方に依存する------言い換えると、\(g\) が \(I\)
でなければ、辞書の係数はもっと複雑になる。
\end{quote}

\hypertarget{astast-mathrmid}{%
\subsubsection{\texorpdfstring{6.3.3
\(\ast\ast = \mathrm{id}\)}{6.3.3 \textbackslash ast\textbackslash ast = \textbackslash mathrm\{id\}}}\label{astast-mathrmid}}

上の辞書から直ちにわかるように、\(\ast\)
を二度続けて作用させると元に戻る：

\[\ast(\ast(dx)) = \ast(dy \wedge dz) = dx\]

一般に、3次元デカルト実空間では
\(\ast\ast = \mathrm{id}\)（恒等演算子）が成り立つ。\(\ast\) を \(1\)
回適用すれば形式の次数が反転し、\(2\) 回適用すれば元に戻る。

\begin{quote}
\textbf{【ここまでのチェックポイント — §6.3】}

\begin{itemize}
\item
  スカラーを得る方法には二つある。方法①は横×縦の作用でスカラーを得る。方法②は同種ベクトルの内積として読み、そこに計量
  \(g\) が現れる。
\item
  第2章で見た \(1,3,3,1\) の対称性は、独立成分数が同じ \(1\)-form と
  \(2\)-form のあいだに対応をつけられることを示唆している。
\item
  ホッジ・スター \(\ast\)
  は、ある形式に対して、それと組み合わせれば正の体積形式になる相手を返す。実空間での辞書は
  \(\ast dx = dy \wedge dz\) などである。
\item
  この辞書は単なる暗記表ではなく、基底を選べば線形変換として配列表示できる。配列表現の詳細は付録Dで確認する。
\item
  \(\ast\ast = \mathrm{id}\) は、この辞書から直ちに従う。
\end{itemize}
\end{quote}

\hypertarget{ast-ux306eux6027ux8cea}{%
\subsubsection{\texorpdfstring{6.3.4 \(\ast\)
の性質}{6.3.4 \textbackslash ast の性質}}\label{ast-ux306eux6027ux8cea}}

§6.2.3 で内積の性質を、第5章で外微分 \(d\)
の性質をまとめたのに倣い、\(\ast\)
についてもここまでの構成から浮かび上がる性質を整理しておく。

\begin{enumerate}
\def\labelenumi{\arabic{enumi}.}
\item
  \textbf{次数の反転}：3次元空間において、\(\ast\) は \(0\)-form
  \(\leftrightarrow\) \(3\)-form を、\(1\)-form \(\leftrightarrow\)
  \(2\)-form を対応させる。一般に \(k\)-form を \((3-k)\)-form へ送る。
\item
  \textbf{線形性}：\(\ast(\omega_1 + \omega_2) = \ast\omega_1 + \ast\omega_2\)。\(\ast\)
  は基底に対する辞書を係数つきで線形に延長する操作であり、線形性は自明である。配列表現の詳細は付録Dで確認する。
\item
  \textbf{内積との対応}：同じ次数の形式 \(\alpha,\beta\)
  に対して、\(\alpha\) と \(\ast\beta\)
  をウェッジ積で組み合わせると、\(\alpha\) と \(\beta\)
  の内積を係数にもつ \(3\)-form が得られる。デカルト座標では、
\end{enumerate}

\[
\alpha \wedge \ast\beta
=
(\alpha\cdot\beta)\,dx\wedge dy\wedge dz
\]

である。ここで \(\alpha\cdot\beta\)
は、同じ次数の形式どうしの内積を表す。\(1\)-form
どうしなら成分の積和、\(2\)-form
どうしなら付録Dで見るフロベニウス積である。この式は、\(\ast\)
が「内積で得られるスカラー」と「ウェッジ積で得られる
\(3\)-form」を繫いでいることを表している。

\begin{quote}
\textbf{注}
（より進んだ定義との関係）より進んだ教科書では、\(\alpha\wedge\ast\beta=(\alpha\cdot\beta)\,dx\wedge dy\wedge dz\)
を \(\ast\)
の定義として採用することが多い。本書では先に基底ごとの辞書と配列表示を作ったが、それはこの定義をデカルト座標で具体化したものになっている。
\end{quote}

\begin{enumerate}
\def\labelenumi{\arabic{enumi}.}
\setcounter{enumi}{3}
\item
  \textbf{$\ast\ast = \mathrm{id}$}：二度作用させると元に戻る（§6.3.3）。これは正規直交基底かつ右手系という条件のもとで成り立ち、付録Dでは反対称行列とフロベニウス積を通じて代数的にも確認できる。
\item
  \textbf{計量 $g$ への依存}：\(\ast\) の辞書は空間の計量 \(g\)
  によって決まる。デカルト座標（\(g=I\)）では §6.3.2
  の簡潔な辞書で済むが、一般の座標系では \(\ast\)
  の係数は場所の関数になる。この一般の場合については第9章で詳しく扱う。
\end{enumerate}

これらの性質は、\(\ast\)
が単なる「辞書」ではなく、計量と向きを備えた空間における\textbf{次数の折り返し構造}そのものであることを示している。

\begin{center}\rule{0.5\linewidth}{0.5pt}\end{center}

\hypertarget{ux5bfeux5fdcux306eux5b9fux4f8bux5faeux5206ux5f62ux5f0fux3068ux30d9ux30afux30c8ux30ebux89e3ux6790}{%
\subsection{§6.4
対応の実例------微分形式とベクトル解析}\label{ux5bfeux5fdcux306eux5b9fux4f8bux5faeux5206ux5f62ux5f0fux3068ux30d9ux30afux30c8ux30ebux89e3ux6790}}

\hypertarget{ux30b8ux30e5ux30fcux30ebux71b1-pvi-ux4e8cux3064ux306eux7d4cux8def}{%
\subsubsection{\texorpdfstring{6.4.1 ジュール熱 \(P=VI\)
------二つの経路}{6.4.1 ジュール熱 P=VI ------二つの経路}}\label{ux30b8ux30e5ux30fcux30ebux71b1-pvi-ux4e8cux3064ux306eux7d4cux8def}}

\(\ast\)
が具体的にどのように微分形式と同種ベクトルの内積表示を繫いでいるか、実際の計算で確かめよう。この節の目的は、電磁気学やベクトル解析の公式を学ぶことではない。\(\ast\)
を使うと、微分形式で自然に得られる量が、スカラー積やベクトル解析の記法に読み替えられることを、二つの例で眺めることである。

\begin{quote}
\textbf{注}
（電磁気学やベクトル解析を知らなくてよい）ここから先、説明のために電磁気学やベクトル解析の記号を少し先取りして使う。\(E\)
や \(J\) が何を表すか、あるいは
\(\mathbf{A}\cdot(\nabla\times\mathbf{A})\)
が何を意味するかを、今すぐ理解する必要はない。ここで見たいのは、微分形式で自然に出てくる
\(3\)-form が、ホッジ・スター \(\ast\)
を通じて、スカラーや内積の形に読み替えられるという一点である。ベクトル解析の記法との関係は、このあと少しだけ見えるが、本格的な整理は後の章に回す。
\end{quote}

高校物理で学ぶ「電力 = 電圧 ×
電流（\(P=VI\)）」を、空間の各点で起こる現象として見直す。

\begin{itemize}
\item
  電圧の起源である\textbf{電場 $E$} は、線分に沿って数え上げる
  \textbf{$1$-form} である。
\item
  電流の起源である\textbf{電流密度 $J$} は、面を貫く電荷を数え上げる
  \textbf{$2$-form} である。
\end{itemize}

まず、\textbf{方法①（微分形式）}で計算してみよう。\(1\)-form \(E\) と
\(2\)-form \(J\) のウェッジ積をとる。

\[E \wedge J\]

\(1\)-form と \(2\)-form のウェッジ積は
\(3\)-form（体積密度）になる。これは「単位体積あたりの電力（W/m³）」を意味する。これを空間領域で積分（\(\int_V E \wedge J\)）すれば、空間全体の消費電力（ワット）が得られる。この経路で使っているのは、\(1\)-form
と \(2\)-form をウェッジ積で \(3\)-form にする操作だけである。

次に、\textbf{方法②（同種ベクトルの内積）}で同じ計算をどう行うか見てみよう。\(E\)
と \(J\) をどちらも \(1\)-form のように扱って内積をとるには、面を測る
\(J\) を \(\ast\) で線を測る形式へ変換する：

\[\text{ベクトル } J \longleftrightarrow \ast J\]

これで \(E\) と \(\ast J\) という二つの \(1\)-form
が揃った。これらの内積（ドット積）をとる。微分形式で \(1\)-form
の内積を書くには、一方に \(\ast\) をかけて \(2\)-form にし、ウェッジ積で
\(3\)-form にしてから、全体に \(\ast\) で \(0\)-form に落とす：

\[E \cdot J \longleftrightarrow \ast(E \wedge \ast(\ast J))\]

ここで \(\ast\ast = \mathrm{id}\) により、内側の \(\ast(\ast J) = J\)
となる：

\[\ast(E \wedge J)\]

カッコの中身 \(E \wedge J\) は、方法①で得た \(3\)-form
と全く同じものだ。外側の \(\ast\) は、\(3\)-form から体積
\(dx \wedge dy \wedge dz\)
を剥ぎ取り、スカラー値を取り出している。スカラー \(E \cdot J\)
に微小体積 \(dV\) を掛け直して積分する書き方は、この \(3\)-form
が持っていた体積を後から補っている、と読める。

\textbf{つまり、$\ast$ は方法①の $3$-form と方法②のスカラー表示を矛盾なく繫いでいる。}
\(\ast\ast = \mathrm{id}\)
があるからこそ、どちらの経路を通っても同じ結果に至るのだ。

\hypertarget{ux30d8ux30eaux30b7ux30c6ux30a3-a-cdot-nabla-times-a-ast-ux306eux5f80ux5fa9}{%
\subsubsection{\texorpdfstring{6.4.2 ヘリシティ
\(A \cdot (\nabla \times A)\) ------\(\ast\)
の往復}{6.4.2 ヘリシティ A \textbackslash cdot (\textbackslash nabla \textbackslash times A) ------\textbackslash ast の往復}}\label{ux30d8ux30eaux30b7ux30c6ux30a3-a-cdot-nabla-times-a-ast-ux306eux5f80ux5fa9}}

もう一つの例で、\(\ast\) の働きをさらに見てみよう。\(1\)-form \(\alpha\)
と、その外微分 \(d\alpha\) を組み合わせた次の \(3\)-form を考える：

\[\alpha \wedge d\alpha\]

この形は、方法②の言葉で書けば、ベクトル場 \(\mathbf{A}\)
とその回転の内積として現れる。

\[\mathbf{A} \cdot (\nabla \times \mathbf{A})\]

ここで \(\mathbf{A}\) に対応する \(1\)-form を \(\alpha\)
と書く。\(d\alpha\) はまず \(2\)-form として得られ、方法②の言葉では
\(\ast d\alpha\) という \(1\)-form と対応する。

まず\textbf{方法①}で計算する。\(1\)-form \(\alpha\) と \(2\)-form
\(d\alpha\) のウェッジ積：

\[\alpha \wedge d\alpha\]

\(1\)-form と \(2\)-form のウェッジ積は
\(3\)-form。この経路では、\(\alpha\) と \(d\alpha\)
をそのままウェッジ積で組み合わせている。

次に\textbf{方法②}で同じ計算の内部構造を見る。同種ベクトルの内積として読むには、\(2\)-form
\(d\alpha\) をそのまま置けないので、\(\ast d\alpha\) という \(1\)-form
を対応させる：

\[\nabla \times \mathbf{A} \longleftrightarrow \ast d\alpha\]

これで \(\alpha\) と \(\ast d\alpha\) という二つの \(1\)-form
が揃った。内積をとる：

\[\mathbf{A} \cdot (\nabla \times \mathbf{A}) \longleftrightarrow \ast(\alpha \wedge \ast(\ast d\alpha))\]

\(\ast\ast = \mathrm{id}\) により内側の
\(\ast(\ast d\alpha) = d\alpha\)。したがって、

\[\ast(\alpha \wedge d\alpha)\]

カッコの中身は、方法①の \(3\)-form \(\alpha \wedge d\alpha\)
と同じものだ。外側の \(\ast\)
体積を剥ぎ取り、スカラー値を取り出している。

\textbf{ここで注目すべきは、$\ast$ が $2$ 回現れて、$\ast\ast = \mathrm{id}$ により打ち消し合っていることだ。}
方法②の内部では、\(d\alpha\)（\(2\)-form）に対応する \(1\)-form として
\(\ast d\alpha\) を使い、内積をとるために再び \(\ast\)
を呼び出している。この往復が、\(\ast\ast = \mathrm{id}\)
によって帳消しになるのである。

\begin{quote}
\textbf{注} （\(\ast\)
の往復が起きる理由）ベクトル解析は「形式の次数」という概念を持たず、すべてを「\(3\)
成分の矢印」として統一的に扱う。\(2\)-form を矢印として読むには \(\ast\)
で対応する \(1\)-form を作る必要があり、内積計算のためにはまた \(\ast\)
で元の次数へ戻る必要がある。\(\ast\)
とは、この二つの表し方を結ぶ一つの対応なのである。
\end{quote}

\begin{quote}
\textbf{【ここまでのチェックポイント — §6.4】}

\begin{itemize}
\item
  ジュール熱の計算では、方法①（\(E \wedge J\)）と方法②（\(E \cdot J = \ast(E \wedge J)\)）が
  \(\ast\) を通じて一致する。
\item
  ヘリシティの計算では、\(\ast\) が \(2\) 回使われ
  \(\ast\ast = \mathrm{id}\)
  で打ち消し合う。方法①（\(\alpha \wedge d\alpha\)）と方法②（\(\mathbf{A} \cdot (\nabla \times \mathbf{A})\)）は同じ結果を与える。
\item
  \(\ast\)
  は微分形式とベクトル解析を繫ぐ対応であり、\(\ast\ast = \mathrm{id}\)
  が往復の整合性を保証している。
\end{itemize}
\end{quote}

\begin{center}\rule{0.5\linewidth}{0.5pt}\end{center}

\hypertarget{ux4e09ux3064ux306eux6f14ux7b97ux306eux578bgrad-rot-div}{%
\subsection{§6.5 三つの演算の型------grad, rot,
div}\label{ux4e09ux3064ux306eux6f14ux7b97ux306eux578bgrad-rot-div}}

ここまでで \(d\)（第5章）と
\(\ast\)（本章）という二つの演算子が揃った。この二つを組み合わせると、形式の次数をさまざまに移動させる操作を作ることができる。

その中でも、三次元のベクトル解析で特に名前を持つものがある。

\[
0 \to 1,\qquad
1 \to 2 \to 1,\qquad
1 \to 2 \to 3 \to 0
\]

という三つの型である。

本書ではこれらを、それぞれ

\[
\mathrm{grad},\qquad
\mathrm{rot},\qquad
\mathrm{div}
\]

と呼ぶことにする。

本章では、これら三つの型に名前をつけるところまでに留める。通常の
\(\nabla f\)、\(\nabla\times\mathbf F\)、\(\nabla\cdot\mathbf F\)
との対応は、後の章で改めて扱う。

\begin{quote}
\textbf{注} （ベクトル解析を知っている読者へ）ここまで来ると、通常の
\(\nabla f\)、\(\nabla\times\mathbf F\)、\(\nabla\cdot\mathbf F\)
が背後に見えているかもしれない。ただし本章では、通常のベクトル解析記法を正式には展開しない。ここでは、「\(d\)
と \(\ast\) を組み合わせた中から、grad, rot, div
に相当する三つの型を取り出す」という道筋を確認する。通常のナブラ記法、ストークスの定理との対応、曲線座標で公式が長くなる理由は、後の章で改めて扱う。
\end{quote}

以降、本節ではまずデカルト座標で考える。\(1\)-form

\[
\omega = P\,dx + Q\,dy + R\,dz
\]

を用いる。ベクトル解析を知っている読者は、これを成分 \((P,Q,R)\)
を持つ場に対応するものとして読んでよい。ただし、この対応の本格的な整理は後の章に回す。

\hypertarget{ux52feux914d-mathrmgradf-df}{%
\subsubsection{\texorpdfstring{6.5.1 勾配
\(\mathrm{grad}\,f = df\)}{6.5.1 勾配 \textbackslash mathrm\{grad\}\textbackslash,f = df}}\label{ux52feux914d-mathrmgradf-df}}

\(0\)-form（スカラー場）\(f\) の外微分 \(df\) は \(1\)-form である。

\[
df = \frac{\partial f}{\partial x}\,dx
{}+ \frac{\partial f}{\partial y}\,dy
{}+ \frac{\partial f}{\partial z}\,dz
\]

本書ではこれを、勾配の \(1\)-form 表示と呼び、

\[
\mathrm{grad}\,f = df
\]

と書く。

\(d\) は次数を \(0 \to 1\)
に上げる。これが、この章で必要な勾配の姿である。通常のベクトル解析で
\(\nabla f\)
と書かれる縦ベクトル表示との関係は、後の章で改めて整理する。

\hypertarget{ux56deux8ee2-mathrmrotomega-astdomega}{%
\subsubsection{\texorpdfstring{6.5.2 回転
\(\mathrm{rot}\,\omega = \ast\,d\,\omega\)}{6.5.2 回転 \textbackslash mathrm\{rot\}\textbackslash,\textbackslash omega = \textbackslash ast\textbackslash,d\textbackslash,\textbackslash omega}}\label{ux56deux8ee2-mathrmrotomega-astdomega}}

\(1\)-form

\[
\omega = P\,dx + Q\,dy + R\,dz
\]

に \(d\) を作用させると、\(2\)-form \(d\omega\) が得られる。

\[d\omega = (\frac{\partial R}{\partial y} - \frac{\partial Q}{\partial z})\,dy \wedge dz + (\frac{\partial P}{\partial z} - \frac{\partial R}{\partial x})\,dz \wedge dx + (\frac{\partial Q}{\partial x} - \frac{\partial P}{\partial y})\,dx \wedge dy\]

ここに \(\ast\) を作用させると、\(2\)-form が再び \(1\)-form に戻る。

\[\ast(d\omega) = (\frac{\partial R}{\partial y} - \frac{\partial Q}{\partial z})\,dx + (\frac{\partial P}{\partial z} - \frac{\partial R}{\partial x})\,dy + (\frac{\partial Q}{\partial x} - \frac{\partial P}{\partial y})\,dz\]

本書ではこの \(1\)-form を回転の \(1\)-form 表示と呼び、

\[\mathrm{rot}\,\omega = \ast\,d\,\omega\]

と書く。

\(\ast\,d\) は

\[
1 \to 2 \to 1
\]

という次数の遷移を持つ。この「一度面に上がってから線に戻る」構造が、回転に対応する。通常のベクトル解析で
\(\nabla\times\mathbf F\) と書かれるものとの関係は、後の章で改めて扱う。

\hypertarget{ux767aux6563-mathrmdivomega-astdastomega}{%
\subsubsection{\texorpdfstring{6.5.3 発散
\(\mathrm{div}\,\omega = \ast\,d\,\ast\,\omega\)}{6.5.3 発散 \textbackslash mathrm\{div\}\textbackslash,\textbackslash omega = \textbackslash ast\textbackslash,d\textbackslash,\textbackslash ast\textbackslash,\textbackslash omega}}\label{ux767aux6563-mathrmdivomega-astdastomega}}

\(1\)-form \(\omega\) にまず \(\ast\) を作用させると、\(2\)-form
になる。そこに \(d\) を作用させると、\(3\)-form になる。最後にもう一度
\(\ast\) を作用させると、\(0\)-form、つまりスカラー場に戻る。

\[
\ast d\ast\omega
\]

デカルト座標で計算すれば、

\[\ast(d(\ast\omega)) = \frac{\partial P}{\partial x} + \frac{\partial Q}{\partial y} + \frac{\partial R}{\partial z}\]

となる。

本書ではこの \(0\)-form を発散と呼び、

\[\mathrm{div}\,\omega = \ast\,d\,\ast\,\omega\]

と書く。

\(\ast\,d\,\ast\) は

\[
1 \to 2 \to 3 \to 0
\]

という次数の遷移を持つ。通常のベクトル解析で \(\nabla\cdot\mathbf F\)
と書かれるものとの関係は、後の章で改めて扱う。

\hypertarget{d2-0-ux306eux5f71}{%
\subsubsection{\texorpdfstring{6.5.4 \(d^2 = 0\)
の影}{6.5.4 d\^{}2 = 0 の影}}\label{d2-0-ux306eux5f71}}

第5章で学んだ \(d^2 = 0\)
を思い出そう。この一つの事実から、三つの演算のあいだに関係が生まれる。

まず、

\[
\mathrm{rot}(\mathrm{grad}\,f)
=
\ast d(df)
=
\ast(d^2 f)
=
0
\]

である。

また、

\[
\mathrm{div}(\mathrm{rot}\,\omega)
=
\ast d\ast(\ast d\omega)
=
\ast d(d\omega)
=
\ast(d^2\omega)
=
0
\]

である。

ここでは、これらを形式側の恒等式として眺めるに留める。ベクトル解析を知っている読者には、通常の恒等式との対応が見えているだろう。しかし本書では、通常のナブラ記法としての意味づけは後の章で改めて行う。

\begin{center}\rule{0.5\linewidth}{0.5pt}\end{center}

\hypertarget{ux30d9ux30afux30c8ux30ebux89e3ux6790ux3078ux5411ux3051ux3066}{%
\subsection{§6.6
ベクトル解析へ向けて}\label{ux30d9ux30afux30c8ux30ebux89e3ux6790ux3078ux5411ux3051ux3066}}

第5章で外微分 \(d\) を、本章で計量 \(g\) とホッジ・スター \(\ast\)
を手に入れた。これで、形式の世界からベクトル解析へ向かうための主要な道具が揃った。

本章で確認したのは、\(d\) と \(\ast\)
の組み合わせの中から、三次元ベクトル解析で grad, rot, div
と呼ばれる三つの型を取り出せる、ということである。

ここで得た対応を表にしておこう。

\begin{longtable}[]{@{}llll@{}}
\toprule\noalign{}
演算 & 本書の分解 & 次数の遷移 & \(\ast\) の使用回数 \\
\midrule\noalign{}
\endhead
\bottomrule\noalign{}
\endlastfoot
\(\mathrm{grad}\) & \(d\) & \(0 \to 1\) & 0回 \\
\(\mathrm{rot}\) & \(\ast\,d\) & \(1 \to 2 \to 1\) & 1回 \\
\(\mathrm{div}\) & \(\ast\,d\,\ast\) & \(1 \to 2 \to 3 \to 0\) & 2回 \\
\end{longtable}

ここまでで、ナブラ記法の背後にある構造のかなりの部分が見えてきた。ただし、本章ではまだ通常のベクトル解析を正式には展開していない。

ベクトル解析を知っている読者には、ここから通常の \(\nabla\)
記法が立ち上がってくることが見えているかもしれない。しかし本書では、ここで急がない。通常のナブラ記法としての整理は、次章以降で改めて行う。

次章以降では、ここで得た \(d\)、\(g\)、\(\ast\)
を使って、通常のベクトル解析の記法へ戻っていく。ストークスの定理、ガウスの発散定理、曲線座標で公式が長くなる理由も、この道具立てから改めて読み直す。

\begin{quote}
\textbf{注}
（計量の符号を一つ変えるだけで）本章では正定値計量（\(\mathbf{v}^T \mathbf{g} \,\mathbf{v} > 0\)
for
\(\mathbf{v} \neq 0\)）を前提とした。しかし特殊相対性理論の舞台である
\(4\) 次元時空では、時間成分だけ符号が反転した
\(\mathbf{g} = \begin{pmatrix} -1 & 0 & 0 & 0 \\ 0 & 1 & 0 & 0 \\ 0 & 0 & 1 & 0 \\ 0 & 0 & 0 & 1 \end{pmatrix}\)
のような不定計量が登場する。本書は \(3\)
次元デカルト実空間に固定しているためこれ以上踏み込まないが、本章で学んだ「行列
\(\mathbf{g}\) を挟む」という発想は、そのまま \(4\)
次元時空へと拡張可能である。必要になったとき、読者はこの章に立ち返ってほしい。
\end{quote}

\begin{center}\rule{0.5\linewidth}{0.5pt}\end{center}

\begin{quote}
\textbf{【ここまでのチェックポイント — 第6章全体】}

\begin{itemize}
\item
  内積とは、同種のベクトル間（横×横、または縦×縦）で成分の積和をとる演算。横×縦の計算とは別物である。
\item
  実空間では内積の行列は単位行列 \(I\)。パラメータ空間では
  \(\mathbf{g}=J^T J\)
  により、座標変換の一次係数から求まる。すなわち計量 \(g\)
  とは、実空間へ引き戻して内積をとったとき、中央に現れる行列に他ならない。
\item
  \(\mathbf{v}^T \mathbf{g}\)
  は縦ベクトルを横ベクトルに変換し、\(\mathbf{g}^{-1}\omega^T\)
  は横ベクトルを縦ベクトル表示へ戻す。\(g\)
  がなければ物体と計器は互いに変換不能である。
\item
  スカラーを得る方法には二つある：方法①（横×縦）と方法②（同種ベクトルの内積）。ホッジ・スター
  \(\ast\) は両者を繫ぐ対応である。
\item
  \(\ast\)
  は、ある形式に対して、それと組み合わせれば正の体積形式になる相手を返す。実空間での辞書は
  \(\ast dx = dy \wedge dz\) などであり、\(\ast\ast = \mathrm{id}\)
  を満たす。
\item
  ジュール熱とヘリシティの計算例は、方法①と方法②が \(\ast\) と
  \(\ast\ast = \mathrm{id}\) を通じて同じ結果を与えることを示している。
\item
  第6章では、\(\mathrm{grad}=d\)、\(\mathrm{rot}=\ast\,d\)、\(\mathrm{div}=\ast\,d\,\ast\)
  という三つの型が現れることを確認した。通常のナブラ記法としての整理は次章以降で行う。
\item
  \(d^2=0\) から、形式側では \(\mathrm{rot}(\mathrm{grad}\,f)=0\) と
  \(\mathrm{div}(\mathrm{rot}\,\omega)=0\)
  が現れる。通常のベクトル解析での意味づけは後の章で改めて扱う。
\end{itemize}
\end{quote}

\begin{center}\rule{0.5\linewidth}{0.5pt}\end{center}

\hypertarget{ux4ed8ux9332dux30dbux30c3ux30b8ux30b9ux30bfux30fcux306eux914dux5217ux8868ux793a}{%
\section{付録D：ホッジ・スターの配列表示}\label{ux4ed8ux9332dux30dbux30c3ux30b8ux30b9ux30bfux30fcux306eux914dux5217ux8868ux793a}}

§6.3.2 では実空間の \(\ast\)
辞書を与えた。この付録では、その辞書が配列操作としてどう見えるかを確認する。ホッジ・スターは抽象的な記号ではなく、基底を選べば具体的な配列として表示できる線形変換である。まず
\(1\)-form \(\leftrightarrow\) \(2\)-form
の対応を反対称行列で見て、最後に \(0\)-form \(\leftrightarrow\)
\(3\)-form の対応を三階配列で見る。

\hypertarget{d.1-ast_1to2-ux4e09ux3064ux306eux4fc2ux6570ux3092ux53cdux5bfeux79f0ux884cux5217ux3078ux7f6eux304f}{%
\subsection{\texorpdfstring{D.1 \(\ast_{1\to2}\)
------三つの係数を反対称行列へ置く}{D.1 \textbackslash ast\_\{1\textbackslash to2\} ------三つの係数を反対称行列へ置く}}\label{d.1-ast_1to2-ux4e09ux3064ux306eux4fc2ux6570ux3092ux53cdux5bfeux79f0ux884cux5217ux3078ux7f6eux304f}}

\(1\)-form

\[
\omega =
\begin{pmatrix}P & Q & R\end{pmatrix}
=
P\,dx + Q\,dy + R\,dz
\]

に本文の辞書

\[
\ast dx = dy \wedge dz,\qquad
\ast dy = dz \wedge dx,\qquad
\ast dz = dx \wedge dy
\]

を適用すると、

\[
\ast\omega
=
P\,dy\wedge dz
+
Q\,dz\wedge dx
+
R\,dx\wedge dy
\]

である。第2章の反対称行列表現では、\(dy\wedge dz\) を表す行列を

\[
E_1 =
\begin{pmatrix}
0 & 0 & 0 \\
0 & 0 & 1 \\
0 & -1 & 0
\end{pmatrix}
\]

とおく。同様に、\(dz\wedge dx\) を表す行列を

\[
E_2 =
\begin{pmatrix}
0 & 0 & -1 \\
0 & 0 & 0 \\
1 & 0 & 0
\end{pmatrix}
\]

とおき、\(dx\wedge dy\) を表す行列を

\[
E_3 =
\begin{pmatrix}
0 & 1 & 0 \\
-1 & 0 & 0 \\
0 & 0 & 0
\end{pmatrix}
\]

とおく。

したがって、

\[
\ast_{1\to2}
=
\begin{pmatrix}
E_1\\
E_2\\
E_3
\end{pmatrix}
\]

という「行列の縦ベクトル」として見られる。\(1\)-form
\(\omega=\begin{pmatrix}P&Q&R\end{pmatrix}\) を左から掛けると、

\[
\omega\,\ast_{1\to2}
=
\begin{pmatrix}P&Q&R\end{pmatrix}
\begin{pmatrix}
E_1\\
E_2\\
E_3
\end{pmatrix}
=
P E_1+Q E_2+R E_3
\]

である。よって

\[
\ast_{1\to2}(\omega)
=
\begin{pmatrix}
0 & R & -Q \\
-R & 0 & P \\
Q & -P & 0
\end{pmatrix}
\]

が得られる。\(1\)-form の三つの係数 \(P,Q,R\) が、\(2\)-form
の反対称行列の三つの独立成分へ配置された。これが \(\ast_{1\to2}\)
の最も見える姿である。

\begin{quote}
\textbf{注} （第2章の \(\widehat{\epsilon}\) との関係）この
\(E_1,E_2,E_3\) は、付録Aで
\(\varepsilon_{1,\cdot,\cdot},\varepsilon_{2,\cdot,\cdot},\varepsilon_{3,\cdot,\cdot}\)
として書き下した行列と同一である。\(\ast_{1\to2}\)
の本質は、エディントンのイプシロン \(\varepsilon_{ijk}\) の第一添字を
\(1\)-form の成分方向に、残りの二添字を \(3\times3\)
行列の方向に配置したものだ。第2章で体積測定器 \(\widehat{\epsilon}\)
として導入したあの三枚組が、ここでホッジ・スターとして再登場している。
\end{quote}

\hypertarget{d.2-ast_2to1-ux30d5ux30edux30d9ux30cbux30a6ux30b9ux7a4dux306bux3088ux308bux4fc2ux6570ux62bdux51fa}{%
\subsection{\texorpdfstring{D.2 \(\ast_{2\to1}\)
------フロベニウス積による係数抽出}{D.2 \textbackslash ast\_\{2\textbackslash to1\} ------フロベニウス積による係数抽出}}\label{d.2-ast_2to1-ux30d5ux30edux30d9ux30cbux30a6ux30b9ux7a4dux306bux3088ux308bux4fc2ux6570ux62bdux51fa}}

逆向きの \(\ast_{2\to1}\)
は、反対称行列から三つの独立成分を取り出す操作である。行列成分だけを見れば、\(A=M_{23}\)、\(B=M_{31}\)、\(C=M_{12}\)
を読むだけでよい。つまり、\(2\)-form

\[
\eta
=
A\,dy\wedge dz
+
B\,dz\wedge dx
+
C\,dx\wedge dy
\]

を表す反対称行列

\[
M=
\begin{pmatrix}
0 & C & -B \\
-C & 0 & A \\
B & -A & 0
\end{pmatrix}
\]

から \(A\,dx+B\,dy+C\,dz\)
へ戻すには、この三成分を読めばよい。しかし、これを単なる「読み取り」として終わらせると、\(\ast_{1\to2}\)
との転置関係が見えにくい。そこで、係数抽出そのものを行列同士の内積として書き直す。

\(3\times3\) 行列 \(A,B\) の内積を

\[
A\cdot B
=
\frac{1}{2}\operatorname{tr}(A^T B)
=
\frac{1}{2}
\sum_{i=1}^{3}
\sum_{j=1}^{3}
A_{ij}B_{ij}
\]

で定義する。

\begin{quote}
\textbf{注} （トレース）\(\operatorname{tr}\)
は行列の\textbf{トレース}（対角成分の和）である。付録Bで外微分の行列表現を扱った際に登場した。\(\operatorname{tr}(A^T B)\)
は「\(A\) を転置して \(B\)
を掛け、対角成分を足し合わせる」三手順の操作だ。\(\frac{1}{2}\sum A_{ij}B_{ij}\)
と等価だから、成分で考えたいときは後者の表式を使えばよい。
\end{quote}

§6.2.3 で縦ベクトル同士の内積を \(\mathbf{v}_1 \cdot \mathbf{v}_2\)
と書いたのと同じ記法である。行列もまた「同じ種類のもの」どうしであり、対応する成分同士を掛けて総和をとる。\(i\)
と \(j\)
の二つの添字にわたって次々と縮約していくことから、これを\textbf{フロベニウス積}または\textbf{連続縮約}と呼ぶ。

\begin{quote}
\textbf{注} （係数 \(1/2\) について）フロベニウス積には
\(\operatorname{tr}(A^T B)\)（\(1/2\)
なし）を採用する流儀もある。しかし反対称行列では \(M_{23}\) と
\(M_{32}=-M_{23}\) のようなペア成分の両方を拾ってしまい、内積が \(2\)
倍される。本書では \(1/2\) を定義に織り込む。これにより、\(2\)-form
\(M\) の自分自身との内積 \(M\cdot M=M_{23}^2+M_{31}^2+M_{12}^2\)
がそのまま §6.2.2 の「平行四辺形の符号なし面積の二乗」に一致する。
\end{quote}

\begin{quote}
\textbf{注} （抽象化の引力）§6.2.3
で公理的な内積に触れたが、こうして実際に手を動かしてみると、その簡潔さも分からなくはない。縦ベクトルの内積には
\(\mathbf{v}_1^T \mathbf{v}_2\) を、行列の内積には
\(\frac{1}{2}\operatorname{tr}(A^T B)\)
を------表示が変わるたびに内積を一から定義し直すのは、考えてみれば相当に大変だ。これらをまとめて「内積とは、ある公理を満たす演算である」と一言で済ませたい欲求に飲まれそうな自分に気付く。それでも本書は最後まで表示を明示する方針を崩さない。ここまで来ると意地である。
\end{quote}

D.1 の \(E_1,E_2,E_3\) は、この内積について正規直交になる。実際、各
\(E_k\) の非ゼロ成分は二つだけで、同じ \(E_k\) どうしでは
\(\frac{1}{2}(1^2+(-1)^2)=1\)、異なる \(E_i,E_j\)
どうしでは非ゼロ成分の位置が重ならない。したがって、

\[
E_i\cdot E_j=\delta_{ij}
\]

である。

この正規直交性を使えば、逆向きの変換 \(\ast_{2\to1}\)
は、\(E_1,E_2,E_3\) との内積として書ける。\(M\) を任意の \(3\times3\)
行列とする。実際に \(2\)-form から来る場合、\(M\) は反対称行列である。

\[
\ast_{2\to1}(M)
=
\begin{pmatrix}
E_1\cdot M & E_2\cdot M & E_3\cdot M
\end{pmatrix}
\]

\(E_k\cdot M=\frac{1}{2}\operatorname{tr}(E_k^T M)\)
は、上で定義したフロベニウス積である。成分で見ると、

\[
\begin{aligned}
E_1 \cdot M &= \frac{1}{2}(M_{23}-M_{32}) \\
E_2 \cdot M &= \frac{1}{2}(M_{31}-M_{13}) \\
E_3 \cdot M &= \frac{1}{2}(M_{12}-M_{21})
\end{aligned}
\]

である。\(M\) が反対称行列なら
\(M_{32}=-M_{23}\)、\(M_{13}=-M_{31}\)、\(M_{21}=-M_{12}\) なので、

\[
\ast_{2\to1}(M)
=
\begin{pmatrix}
M_{23} & M_{31} & M_{12}
\end{pmatrix}
\]

となる。したがって、\(\ast_{2\to1}\)
は反対称行列から独立成分を取り出す操作であると同時に、\(E_k\)
とのフロベニウス積による係数抽出でもある。

成分をすべて書き下せば、\(\ast_{2\to1}\)
は次のような「行列の横ベクトル」である。

\[
\ast_{2\to1} = \begin{pmatrix}
\begin{pmatrix} 0 & 0 & 0 \\ 0 & 0 & 1 \\ 0 & -1 & 0 \end{pmatrix} &
\begin{pmatrix} 0 & 0 & -1 \\ 0 & 0 & 0 \\ 1 & 0 & 0 \end{pmatrix} &
\begin{pmatrix} 0 & 1 & 0 \\ -1 & 0 & 0 \\ 0 & 0 & 0 \end{pmatrix}
\end{pmatrix}
\]

\(\ast_{1\to2}\) では係数倍して足し合わせ、\(\ast_{2\to1}\)
では内積で係数を抽出する。縦と横、加重和と内積抽出が対応している。

\hypertarget{d.3-1-form-ux3068-2-form-ux306eux5834ux5408ux306e-astastmathrmid}{%
\subsection{\texorpdfstring{D.3 \(1\)-form と \(2\)-form の場合の
\(\ast\ast=\mathrm{id}\)}{D.3 1-form と 2-form の場合の \textbackslash ast\textbackslash ast=\textbackslash mathrm\{id\}}}\label{d.3-1-form-ux3068-2-form-ux306eux5834ux5408ux306e-astastmathrmid}}

D.1 と D.2 を合わせると、\(1\)-form と \(2\)-form の場合の
\(\ast\ast=\mathrm{id}\) は、配列の正規直交性として見える。

\[
\omega=P\,dx+Q\,dy+R\,dz
\]

に対して、D.1 より

\[
\ast_{1\to2}(\omega)=P E_1+Q E_2+R E_3
\]

である。これに \(\ast_{2\to1}\) を作用させると、

\[
\begin{aligned}
\ast_{2\to1}(\ast_{1\to2}(\omega))
&=
\begin{pmatrix}
E_1\cdot (P E_1+Q E_2+R E_3) \\
E_2\cdot (P E_1+Q E_2+R E_3) \\
E_3\cdot (P E_1+Q E_2+R E_3)
\end{pmatrix}^{T} \\
&=
\begin{pmatrix}
P & Q & R
\end{pmatrix}
=
\omega
\end{aligned}
\]

となる。最後の等号では \(E_i\cdot E_j=\delta_{ij}\)
を使った。つまり、\(\ast_{1\to2}\) で \(E_1,E_2,E_3\)
へ配置した三つの係数を、\(\ast_{2\to1}\) が同じ \(E_1,E_2,E_3\)
との内積でそのまま取り戻している。

逆向きも同じである。反対称行列

\[
M=P E_1+Q E_2+R E_3
\]

に対して、D.2 より

\[
\ast_{2\to1}(M)
=
\begin{pmatrix}
E_1\cdot M & E_2\cdot M & E_3\cdot M
\end{pmatrix}
=
\begin{pmatrix}
P & Q & R
\end{pmatrix}
\]

であり、さらに D.1 の \(\ast_{1\to2}\) を作用させれば、

\[
\ast_{1\to2}(\ast_{2\to1}(M))
=
P E_1+Q E_2+R E_3
=
M
\]

となる。

したがって、\(1\)-form と \(2\)-form の場合には

\[
\ast_{2\to1}(\ast_{1\to2}(\omega))=\omega,
\qquad
\ast_{1\to2}(\ast_{2\to1}(M))=M
\]

が成り立つ。\(\ast_{1\to2}\) は三つの係数を \(E_1,E_2,E_3\)
へ配置し、\(\ast_{2\to1}\) は \(E_1,E_2,E_3\)
との内積で係数を取り出す。正規化されたフロベニウス積を使えば、両者は転置的な関係にあることが、この配置と係数抽出の対応として見えている。

\hypertarget{d.4-ast_0to3-ux3068-ast_3to0-ux306eux914dux5217ux8868ux793a}{%
\subsection{\texorpdfstring{D.4 \(\ast_{0\to3}\) と \(\ast_{3\to0}\)
の配列表示}{D.4 \textbackslash ast\_\{0\textbackslash to3\} と \textbackslash ast\_\{3\textbackslash to0\} の配列表示}}\label{d.4-ast_0to3-ux3068-ast_3to0-ux306eux914dux5217ux8868ux793a}}

ここまでで、\(\ast_{1\to2}\) と \(\ast_{2\to1}\)
を、係数配列を移す変換として表示した。残るのは、\(\ast_{0\to3}\) と
\(\ast_{3\to0}\) である。

\(0\)-form は一つの係数を持つ。一方、\(3\)-form
は、完全に配列として見れば、三つの添字を持つ完全反対称な \(3\)
階配列である。したがって、\(\ast_{0\to3}\) は \(1\) 成分を
\(3\times3\times3\) 成分へ広げる変換であり、\(\ast_{3\to0}\) は
\(3\times3\times3\) 成分から \(1\) 成分を取り出す変換である。

ここでは第2章でも現れた \(\varepsilon_{ijk}\)
を前面に押し出してその姿を見てみよう。

\begin{quote}
\textbf{注} （\(\varepsilon_{ijk}\) の正体）\(\varepsilon_{ijk}\)
は、第2章で現れた完全反対称な記号である。添字 \((i,j,k)\) が \((1,2,3)\)
の偶置換なら \(+1\)、奇置換なら \(-1\)、同じ添字を含めば \(0\)
になる。直交デカルト基底では、完全反対称テンソルの成分として読める。本書では、その成分をホッジ・スターの表示として使う。
\end{quote}

まず、\(0\)-form \(f\) に \(\ast_{0\to3}\) を作用させる。出力は
\(3\)-form
なので、三つの添字を持つ成分で書ける。このとき、\(\ast_{0\to3}\)
そのものの成分を

\[
(\ast_{0\to3})_{ijk}
=
\varepsilon_{ijk}
\]

とおく。

これを、\(i\) を固定した三枚の \(3\times3\)
行列として表示しよう。すなわち、\(i=1,2,3\)
のそれぞれについて、\((j,k)\) 成分を行列として並べる。

\[
(\ast_{0\to3})_{1jk}
=
\begin{pmatrix}
0&0&0\\
0&0&1\\
0&-1&0
\end{pmatrix},
\]

\[
(\ast_{0\to3})_{2jk}
=
\begin{pmatrix}
0&0&-1\\
0&0&0\\
1&0&0
\end{pmatrix},
\]

\[
(\ast_{0\to3})_{3jk}
=
\begin{pmatrix}
0&1&0\\
-1&0&0\\
0&0&0
\end{pmatrix}.
\]

見ての通り、これは D.1 で使った反対称行列 \(E_1,E_2,E_3\)
そのものである。つまり、

\[
(\ast_{0\to3})_{1jk}= (E_1)_{jk},
\qquad
(\ast_{0\to3})_{2jk}= (E_2)_{jk},
\qquad
(\ast_{0\to3})_{3jk}= (E_3)_{jk}
\]

である。

ここで新しいことはほとんど起きていない。D.1 で見た反対称行列
\(E_1,E_2,E_3\) が、今度は \(i=1,2,3\)
の三枚のスライスとして並んでいるだけである。\(1\)-form と \(2\)-form
の場合は、三つの係数を三つの反対称行列の係数として読んだ。\(0\)-form と
\(3\)-form
の場合は、一つの係数から、その三枚すべてをまとめた完全反対称な \(3\)
階配列を作っている。

実際、\(0\)-form \(f\) に作用させると、

\[
(\ast_{0\to3}f)_{ijk}
=
(\ast_{0\to3})_{ijk}f
=
\varepsilon_{ijk}f
\]

である。これは、係数 \(f\) を完全反対称な \(3\)
階配列へ広げる操作である。

次に、逆向きの \(\ast_{3\to0}\) を見る。これは \(3\)
階配列を受け取り、一つの係数を返す変換である。その成分は

\[
(\ast_{3\to0})_{ijk}
=
\frac{1}{3!}\varepsilon_{ijk}
\]

である。

したがって、同じく \(i\) を固定した三枚の \(3\times3\)
行列として表示すれば、

\[
(\ast_{3\to0})_{1jk}
=
\frac{1}{3!}
\begin{pmatrix}
0&0&0\\
0&0&1\\
0&-1&0
\end{pmatrix},
\]

\[
(\ast_{3\to0})_{2jk}
=
\frac{1}{3!}
\begin{pmatrix}
0&0&-1\\
0&0&0\\
1&0&0
\end{pmatrix},
\]

\[
(\ast_{3\to0})_{3jk}
=
\frac{1}{3!}
\begin{pmatrix}
0&1&0\\
-1&0&0\\
0&0&0
\end{pmatrix}.
\]

つまり、\(\ast_{3\to0}\)
でも同じ三枚の反対称行列が使われる。違いは、全体に \(1/3!\)
が掛かっていることだけである。

添字 \((i,j,k)\) を一列に並べて一つの番号だと思えば、\(\ast_{0\to3}\) は
\(1\) 成分を \(27\) 成分へ広げる列ベクトルであり、\(\ast_{3\to0}\) は
\(27\) 成分から \(1\)
成分を取り出す行ベクトルである。したがって、正規化係数 \(1/3!\)
を除けば、両者は転置の関係にある。

\(3\)-form

\[
\eta
=
\eta_{ijk}
\]

に \(\ast_{3\to0}\) を作用させると、三つの添字をすべて足し合わせる。

\[
\ast_{3\to0}\eta
=
\sum_{i=1}^{3}
\sum_{j=1}^{3}
\sum_{k=1}^{3}
(\ast_{3\to0})_{ijk}\eta_{ijk}
\]

したがって、

\[
\ast_{3\to0}\eta
=
\frac{1}{3!}
\sum_{i=1}^{3}
\sum_{j=1}^{3}
\sum_{k=1}^{3}
\varepsilon_{ijk}\eta_{ijk}
\]

である。これは、完全反対称な \(3\)
階配列から一つの係数を取り出す三重縮約である。

係数 \(1/3!\) は、D.2 のフロベニウス積に現れた \(1/2\)
と同じ役割を持つ。反対称行列では独立成分が二回ずつ現れるので、\(1/2\)
を掛けて重複を補正した。ここでは、完全反対称な \(3\) 階配列の独立成分が
\(3!=6\) 回現れるので、\(1/3!\) を掛けて重複を補正している。

実際に、\(\ast\)
を二度続けて作用させると元に戻ることを確認しておこう。まず、\(0\)-form
\(f\) から始める。

\[
\ast_{3\to0}(\ast_{0\to3}f)
=
\frac{1}{3!}
\sum_{i=1}^{3}
\sum_{j=1}^{3}
\sum_{k=1}^{3}
\varepsilon_{ijk}(\ast_{0\to3}f)_{ijk}
\]

ここで \((\ast_{0\to3}f)_{ijk}=\varepsilon_{ijk}f\) だから、

\[
\ast_{3\to0}(\ast_{0\to3}f)
=
\frac{1}{3!}
\sum_{i=1}^{3}
\sum_{j=1}^{3}
\sum_{k=1}^{3}
\varepsilon_{ijk}\varepsilon_{ijk}f
\]

非ゼロになるのは、\((i,j,k)\) が \((1,2,3)\) の順列になっている \(3!=6\)
通りだけである。そのとき \(\varepsilon_{ijk}\varepsilon_{ijk}=1\)
なので、

\[
\ast_{3\to0}(\ast_{0\to3}f)
=
\frac{1}{3!}(3!)f
=
f
\]

である。

逆向きも確認しておく。任意の \(3\)-form
は、三次元では独立成分を一つしか持たない。したがって、完全反対称な成分
\(\eta_{ijk}\) は、ある係数 \(h\) を使って

\[
\eta_{ijk}
=
h\,\varepsilon_{ijk}
\]

と書ける。

このとき、

\[
\ast_{3\to0}\eta
=
\frac{1}{3!}
\sum_{i=1}^{3}
\sum_{j=1}^{3}
\sum_{k=1}^{3}
\varepsilon_{ijk}(h\varepsilon_{ijk})
=
h
\]

である。したがって、

\[
(\ast_{0\to3}(\ast_{3\to0}\eta))_{ijk}
=
(\ast_{0\to3}h)_{ijk}
=
h\varepsilon_{ijk}
=
\eta_{ijk}
\]

となる。これで、\(0\)-form と \(3\)-form の場合にも
\(\ast\ast=\mathrm{id}\) が配列表現で確認できた。

結局、ここで見たことは D.1〜D.3 と同じである。\(\ast_{1\to2}\)
では三つの係数を三枚の反対称行列へ分配した。\(\ast_{0\to3}\)
では一つの係数を三枚の反対称行列すべてへ同時に分配した。\(\ast_{2\to1}\)
では反対称行列から係数を取り出し、\(\ast_{3\to0}\)
では三枚の反対称行列との三重縮約によって係数を取り出した。ホッジ・スターは、どの次数でも、基底を選べば具体的な配列として表示できる線形変換なのである。

\hypertarget{d.5-ux307eux3068ux3081}{%
\subsection{D.5 まとめ}\label{d.5-ux307eux3068ux3081}}

\begin{quote}
\textbf{【ここまでのチェックポイント — 付録D】}

\begin{itemize}
\item
  \(\ast_{1\to2}\) は、三つの係数を反対称行列 \(E_1,E_2,E_3\)
  へ配置する線形変換である。
\item
  \(\ast_{2\to1}\)
  は、反対称行列から係数を取り出す線形変換であり、\(E_k\)
  とのフロベニウス積によって書ける。
\item
  フロベニウス積 \(A\cdot B=\frac{1}{2}\operatorname{tr}(A^TB)\)
  を使うと、\(E_i\cdot E_j=\delta_{ij}\) となる。
\item
  したがって、\(1\)-form と \(2\)-form
  の場合には、\(\ast_{2\to1}(\ast_{1\to2}(\omega))=\omega\)、\(\ast_{1\to2}(\ast_{2\to1}(M))=M\)
  が成り立つ。
\item
  \(\ast_{0\to3}\) は、三枚の反対称行列 \(E_1,E_2,E_3\) を重ねた \(3\)
  階配列として表示できる。成分では
  \((\ast_{0\to3})_{ijk}=\varepsilon_{ijk}\) である。
\item
  \(\ast_{3\to0}\) は、同じ配列に正規化係数 \(1/3!\)
  を掛けたものとして表示できる。成分では
  \((\ast_{3\to0})_{ijk}=\frac{1}{3!}\varepsilon_{ijk}\) である。
\item
  \(1/2\) と \(1/3!\) は、反対称成分の重複を補正する係数である。
\item
  \(0\)-form/\(3\)-form
  の場合にも、\(\ast_{3\to0}(\ast_{0\to3}f)=f\)、\(\ast_{0\to3}(\ast_{3\to0}\eta)=\eta\)
  が成り立つ。
\item
  したがって、ホッジ・スター \(\ast\) は、\(0\leftrightarrow3\) と
  \(1\leftrightarrow2\)
  の両方について、基底を選べば具体的な配列として表示できる線形変換である。
\end{itemize}
\end{quote}

\newpage

\hypertarget{ux7b2c7ux7ae0ux30d9ux30afux30c8ux30ebux89e3ux6790-ux30caux30d6ux30e9ux306eux767bux5834}{%
\chapter{第7章：ベクトル解析 ------
ナブラの登場}\label{ux7b2c7ux7ae0ux30d9ux30afux30c8ux30ebux89e3ux6790-ux30caux30d6ux30e9ux306eux767bux5834}}

\hypertarget{ux30caux30d6ux30e9ux306eux767bux5834}{%
\subsection{§7.0
ナブラの登場}\label{ux30caux30d6ux30e9ux306eux767bux5834}}

　本書のタイトルは「ナブラ解体新書」である。しかしここに至るまで、ナブラ
\(\nabla\)
を正面から定義したことは一度もない、という建前だ。第6章では、\(d\) と
\(\ast\) を組み合わせることで grad, rot, div
の型だけを先に見た。本章では、同じ対象を標準的なベクトル解析の記法で書き直す。

代わりに何をしてきたか。\(dx\) は横ベクトル、\(2\)-form
は反対称行列、\(d\)
は次数を上げる演算------いわば「裏方」の道具をひたすら積み上げてきた。読者の中には「いつになったら
\(\nabla\) が出てくるんだ」と思っていた人もいるだろう。

筆者もずっとタイトル詐欺が気になっている。

本章では、ついに \(\nabla\)
を定義する。そして標準的なベクトル解析------勾配・発散・回転・ラプラシアン・諸恒等式------を、天下りのまま一気に展開する。本文の計算では、\(d\)
も \(\ast\) も \(\wedge\)
もいったん脇に置く。これは多くの読者が大学1-2年で（おそらく挫折しながら）学んだ、あのベクトル解析そのものである。

\begin{quote}
\textbf{注} （なぜ今なのか）第6章で我々は、\(d\) と \(\ast\) を使って
grad, rot, div
の型を先に見た。本章では、同じ対象を標準的なベクトル解析の記法で、あえて公式だらけに展開する。本章を読み終えたとき「公式が多すぎる」と感じたら、それでよい。その違和感が、次章で二つの記法を対応させるための推進力になる。
\end{quote}

本章で使う道具は、ここまでに揃えたものだけである------縦ベクトル、横ベクトル、行列、内積、ヤコビ行列、そして偏微分。

\begin{center}\rule{0.5\linewidth}{0.5pt}\end{center}

\hypertarget{ux30c9ux30c3ux30c8ux7a4dux3068ux30afux30edux30b9ux7a4d}{%
\subsection{§7.1
ドット積とクロス積}\label{ux30c9ux30c3ux30c8ux7a4dux3068ux30afux30edux30b9ux7a4d}}

\hypertarget{ux30c9ux30c3ux30c8ux7a4d}{%
\subsubsection{7.1.1 ドット積}\label{ux30c9ux30c3ux30c8ux7a4d}}

二つの縦ベクトル

\[
\mathbf{a}
=
\begin{pmatrix}
a_x\\
a_y\\
a_z
\end{pmatrix},
\qquad
\mathbf{b}
=
\begin{pmatrix}
b_x\\
b_y\\
b_z
\end{pmatrix}
\]

の\textbf{ドット積}（内積）を次で定義する。

\[\mathbf{a} \cdot \mathbf{b} = a_x b_x + a_y b_y + a_z b_z\]

§6.1 で見たように、これは実空間 \((x,y,z)\) で \(\mathbf{g} = I\)
とした場合の内積 \(\mathbf{a}^T \mathbf{g} \,\mathbf{b}\)
に等しい。本書ではすでに縦ベクトルと横ベクトルの積として何度も使ってきた操作だ。

\hypertarget{ux30afux30edux30b9ux7a4d}{%
\subsubsection{7.1.2 クロス積}\label{ux30afux30edux30b9ux7a4d}}

二つの縦ベクトル \(\mathbf{a}, \mathbf{b}\)
の\textbf{クロス積}（外積）を次で定義する。

\[\mathbf{a} \times \mathbf{b} = \begin{pmatrix}
a_y b_z - a_z b_y \\
a_z b_x - a_x b_z \\
a_x b_y - a_y b_x
\end{pmatrix}\]

この結果は再び縦ベクトルである。クロス積の成分の並びには規則性があり、次の反対称行列を使うと簡潔に書ける。

\[\mathbf{a} \times = \begin{pmatrix}
0 & -a_z & a_y \\
a_z & 0 & -a_x \\
-a_y & a_x & 0
\end{pmatrix}\]

この \(3\times 3\) 行列をベクトル \(\mathbf{b}\) に左から掛ければ、

\[\mathbf{a} \times \mathbf{b} = (\mathbf{a} \times)\,\mathbf{b} = \begin{pmatrix}
0 & -a_z & a_y \\
a_z & 0 & -a_x \\
-a_y & a_x & 0
\end{pmatrix}\begin{pmatrix}
b_x \\ b_y \\ b_z
\end{pmatrix} = \begin{pmatrix}
a_y b_z - a_z b_y \\
a_z b_x - a_x b_z \\
a_x b_y - a_y b_x
\end{pmatrix}\]

となる。この \((\mathbf{a} \times)\) を \(\mathbf{a}\)
の\textbf{外積行列}と呼ぶ。第2章で \(2\)-form
の行列表現として現れ、第6章の付録Dではホッジ・スターの配列表示として再登場した行列と、まったく同じ形である。本章ではこれを純粋に「ベクトルにベクトルを掛けてベクトルを得る演算」として使う。

クロス積には次の基本的な性質がある。

\begin{enumerate}
\def\labelenumi{\arabic{enumi}.}
\tightlist
\item
  \textbf{反可換性}：\(\mathbf{a} \times \mathbf{b} = -\,\mathbf{b} \times \mathbf{a}\)
\item
  \textbf{自分自身とのクロス積はゼロ}：\(\mathbf{a} \times \mathbf{a} = \mathbf{0}\)
\item
  \textbf{直交性}：\((\mathbf{a} \times \mathbf{b}) \cdot \mathbf{a} = 0\)、\((\mathbf{a} \times \mathbf{b}) \cdot \mathbf{b} = 0\)
\end{enumerate}

性質3は、クロス積が \(\mathbf{a}\) にも \(\mathbf{b}\)
にも直交するベクトルを生み出すことを意味する。これがいわゆる「右手の法則」の正体だ。

\begin{quote}
\textbf{【ここまでのチェックポイント — §7.1】}

\begin{itemize}
\item
  ドット積
  \(\mathbf{a}\cdot\mathbf{b} = a_x b_x + a_y b_y + a_z b_z\)。実空間では、縦ベクトル
  \(\mathbf{a}\) を転置して \(\mathbf{a}^T \mathbf{b}\) と書ける。
\item
  クロス積 \(\mathbf{a}\times\mathbf{b}\) は外積行列
  \((\mathbf{a}\times)\) を使って \((\mathbf{a}\times)\,\mathbf{b}\)
  と書ける。
\item
  クロス積は反可換、自分自身とはゼロ、結果は二つのベクトルに直交する。
\end{itemize}
\end{quote}

\begin{center}\rule{0.5\linewidth}{0.5pt}\end{center}

\hypertarget{nabla-ux3068ux52feux914d}{%
\subsection{\texorpdfstring{§7.2 \(\nabla\)
と勾配}{§7.2 \textbackslash nabla と勾配}}\label{nabla-ux3068ux52feux914d}}

\hypertarget{nabla-ux306eux5b9aux7fa9}{%
\subsubsection{\texorpdfstring{7.2.1 \(\nabla\)
の定義}{7.2.1 \textbackslash nabla の定義}}\label{nabla-ux306eux5b9aux7fa9}}

ナブラ \(\nabla\)
を、偏微分演算子を縦に並べた形式的なベクトルとして定義する。

\[\nabla = \begin{pmatrix}
\frac{\partial}{\partial x} \\[0.5em]
\frac{\partial}{\partial y} \\[0.5em]
\frac{\partial}{\partial z}
\end{pmatrix}\]

\(\nabla\)
は「ベクトル」だが、その成分は数ではなく\textbf{微分演算子}である。\(\nabla\)
は単独では意味を持たず、何かに作用して初めて値が定まる。

\hypertarget{ux52feux914d}{%
\subsubsection{7.2.2 勾配}\label{ux52feux914d}}

スカラー場 \(f(x,y,z)\) に \(\nabla\)
を作用させたものを\textbf{勾配}と呼び、\(\mathrm{grad}\,f\) または
\(\nabla f\) と書く。

\[\mathrm{grad}\,f = \nabla f = \begin{pmatrix}
\frac{\partial f}{\partial x} \\[0.5em]
\frac{\partial f}{\partial y} \\[0.5em]
\frac{\partial f}{\partial z}
\end{pmatrix}\]

勾配の各成分は、\(f\)
の各座標方向の変化率である。幾何学的には、\(\nabla f\) は \(f\)
が最も急に増加する方向を指すベクトル場になる。

\begin{quote}
\textbf{注} （勾配は横ベクトルか縦ベクトルか）§6.5 では
\(\mathrm{grad}\,f = df\) と書き、\(df\)
は係数の横ベクトルだった。本章では \(\nabla f\)
を縦ベクトルとして扱う。第6章で導入したユークリッド計量を用いて、\(1\)-form
\(df\) を対応する縦ベクトルとして読むと \(\nabla f\) が得られる。§6.5
では「微分そのもの（form）が空間の性質を測る」立場に立ち、本章では標準的なベクトル解析に合わせて「各点に矢印（ベクトル）が立っている」立場に立っている。出発点の絵は異なるが、線積分で集計すれば同じスカラー量を与える。この一致は偶然ではない。
\end{quote}

\textbf{例}：\(f(x,y,z) = x^2 + y^2 + z^2\) の勾配。

\[\nabla f = \begin{pmatrix} 2x \\ 2y \\ 2z \end{pmatrix}\]

このベクトル場は原点から放射状に外向きで、原点から離れるほど大きくなる。

\begin{center}\rule{0.5\linewidth}{0.5pt}\end{center}

\hypertarget{ux767aux6563}{%
\subsection{§7.3 発散}\label{ux767aux6563}}

ベクトル場

\[
\mathbf{F}(x,y,z)
=
\begin{pmatrix}
F_x\\
F_y\\
F_z
\end{pmatrix}
\]

に対して、\(\nabla\) とのドット積をとったものを\textbf{発散}と呼ぶ。

\[\mathrm{div}\,\mathbf{F} = \nabla \cdot \mathbf{F} = \frac{\partial F_x}{\partial x} + \frac{\partial F_y}{\partial y} + \frac{\partial F_z}{\partial z}\]

形式的には「\(\nabla\) と \(\mathbf{F}\) のドット積」だが、\(\nabla\)
の各成分が微分演算子であるため、結果はスカラー場になる。

発散は行列の言葉でも書ける。\(\mathbf{F}\) のヤコビ行列
\(\mathbf{J}_\mathbf{F}\) を考える。

\[\mathbf{J}_\mathbf{F} = \begin{pmatrix}
\frac{\partial F_x}{\partial x} & \frac{\partial F_x}{\partial y} & \frac{\partial F_x}{\partial z} \\[0.3em]
\frac{\partial F_y}{\partial x} & \frac{\partial F_y}{\partial y} & \frac{\partial F_y}{\partial z} \\[0.3em]
\frac{\partial F_z}{\partial x} & \frac{\partial F_z}{\partial y} & \frac{\partial F_z}{\partial z}
\end{pmatrix}\]

この行列の\textbf{トレース}（対角成分の和）が発散に一致する。

\[\mathrm{div}\,\mathbf{F} = \operatorname{tr}(\mathbf{J}_\mathbf{F}) = \frac{\partial F_x}{\partial x} + \frac{\partial F_y}{\partial y} + \frac{\partial F_z}{\partial z}\]

物理的には、\(\mathrm{div}\,\mathbf{F}\)
はその点での「湧き出し」の強さを表す。正なら湧き出し、負なら吸い込み、ゼロならその点では湧き出しも吸い込みもない。

\textbf{例}：

\[
\mathbf{F}
=
\begin{pmatrix}
x\\
y\\
z
\end{pmatrix}
\]

の発散。

\[\nabla \cdot \mathbf{F} = \frac{\partial x}{\partial x} + \frac{\partial y}{\partial y} + \frac{\partial z}{\partial z} = 1 + 1 + 1 = 3\]

\textbf{例}：

\[
\mathbf{F}
=
\begin{pmatrix}
-y\\
x\\
0
\end{pmatrix}
\]

の発散。

\[\nabla \cdot \mathbf{F} = \frac{\partial (-y)}{\partial x} + \frac{\partial x}{\partial y} + \frac{\partial 0}{\partial z} = 0 + 0 + 0 = 0\]

この場は回転しているが、湧き出してはいない。

\begin{center}\rule{0.5\linewidth}{0.5pt}\end{center}

\hypertarget{ux56deux8ee2}{%
\subsection{§7.4 回転}\label{ux56deux8ee2}}

ベクトル場 \(\mathbf{F}\) に対して、\(\nabla\)
とのクロス積をとったものを\textbf{回転}（ローテーション）と呼ぶ。

\[\mathrm{rot}\,\mathbf{F} = \nabla \times \mathbf{F} = \begin{pmatrix}
\frac{\partial F_z}{\partial y} - \frac{\partial F_y}{\partial z} \\[0.5em]
\frac{\partial F_x}{\partial z} - \frac{\partial F_z}{\partial x} \\[0.5em]
\frac{\partial F_y}{\partial x} - \frac{\partial F_x}{\partial y}
\end{pmatrix}\]

外積行列を使えば次のようにも書ける。

\[\mathrm{rot}\,\mathbf{F} = (\nabla \times)\,\mathbf{F} = \begin{pmatrix}
0 & -\frac{\partial}{\partial z} & \frac{\partial}{\partial y} \\[0.3em]
\frac{\partial}{\partial z} & 0 & -\frac{\partial}{\partial x} \\[0.3em]
-\frac{\partial}{\partial y} & \frac{\partial}{\partial x} & 0
\end{pmatrix}\begin{pmatrix}
F_x \\ F_y \\ F_z
\end{pmatrix}\]

物理的には、\(\mathrm{rot}\,\mathbf{F}\)
はその点での「渦」の強さと向きを表す。渦の軸方向を向き、大きさが渦の強さに対応する。

\textbf{例}：

\[
\mathbf{F}
=
\begin{pmatrix}
-y\\
x\\
0
\end{pmatrix}
\]

の回転。

\[\nabla \times \mathbf{F} = \begin{pmatrix}
\frac{\partial 0}{\partial y} - \frac{\partial x}{\partial z} \\[0.3em]
\frac{\partial (-y)}{\partial z} - \frac{\partial 0}{\partial x} \\[0.3em]
\frac{\partial x}{\partial x} - \frac{\partial (-y)}{\partial y}
\end{pmatrix} = \begin{pmatrix}
0 - 0 \\ 0 - 0 \\ 1 - (-1)
\end{pmatrix} = \begin{pmatrix} 0 \\ 0 \\ 2 \end{pmatrix}\]

この場は \(z\) 軸まわりに一様に回転しており、渦の強さは \(2\) である。

\textbf{例}：

\[
\mathbf{F}
=
\begin{pmatrix}
x\\
y\\
z
\end{pmatrix}
\]

の回転。

\[\nabla \times \mathbf{F} = \begin{pmatrix}
\frac{\partial z}{\partial y} - \frac{\partial y}{\partial z} \\[0.3em]
\frac{\partial x}{\partial z} - \frac{\partial z}{\partial x} \\[0.3em]
\frac{\partial y}{\partial x} - \frac{\partial x}{\partial y}
\end{pmatrix} = \begin{pmatrix} 0 \\ 0 \\ 0 \end{pmatrix}\]

放射状の場には渦がない------直感とも一致する。

\begin{quote}
\textbf{【ここまでのチェックポイント — §7.2–§7.4】}

\begin{itemize}
\item
  \(\nabla\) は \(\partial_x,\partial_y,\partial_z\)
  を縦に並べた形式的ベクトル。
\item
  \(\mathrm{grad}\,f = \nabla f\)（勾配）：スカラー場 → ベクトル場。
\item
  \(\mathrm{div}\,\mathbf{F} = \nabla \cdot \mathbf{F}\)（発散）：ベクトル場
  → スカラー場。ヤコビ行列のトレースに等しい。
\item
  \(\mathrm{rot}\,\mathbf{F} = \nabla \times \mathbf{F}\)（回転）：ベクトル場
  → ベクトル場。外積行列で書ける。
\end{itemize}
\end{quote}

\begin{center}\rule{0.5\linewidth}{0.5pt}\end{center}

\hypertarget{ux30e9ux30d7ux30e9ux30b7ux30a2ux30f3}{%
\subsection{§7.5
ラプラシアン}\label{ux30e9ux30d7ux30e9ux30b7ux30a2ux30f3}}

勾配（スカラー→ベクトル）と発散（ベクトル→スカラー）を連続して作用させると、スカラー場からスカラー場への演算が得られる。これを\textbf{ラプラシアン}と呼ぶ。

\[\nabla^2 f = \mathrm{div}(\mathrm{grad}\,f) = \nabla \cdot (\nabla f) = \frac{\partial^2 f}{\partial x^2} + \frac{\partial^2 f}{\partial y^2} + \frac{\partial^2 f}{\partial z^2}\]

記号 \(\nabla^2\) は \(\nabla \cdot \nabla\) の略記であり、\(\Delta f\)
と書く流儀もある。

\textbf{例}：\(f(x,y,z) = x^2 + y^2 + z^2\) のラプラシアン。

\[\nabla^2 f = \frac{\partial^2}{\partial x^2}(x^2) + \frac{\partial^2}{\partial y^2}(y^2) + \frac{\partial^2}{\partial z^2}(z^2) = 2 + 2 + 2 = 6\]

ベクトル場 \(\mathbf{F}\)
に対しても、成分ごとにラプラシアンを作用させたものを\textbf{ベクトルラプラシアン}と呼ぶ。

\[\nabla^2 \mathbf{F} = \begin{pmatrix} \nabla^2 F_x \\ \nabla^2 F_y \\ \nabla^2 F_z \end{pmatrix}\]

ラプラシアンは物理の至るところに現れる。熱伝導方程式、波動方程式、ポアソン方程式------いずれも
\(\nabla^2\) を中心に据えた方程式である。

\begin{center}\rule{0.5\linewidth}{0.5pt}\end{center}

\hypertarget{ux6052ux7b49ux5f0f}{%
\subsection{§7.6 恒等式}\label{ux6052ux7b49ux5f0f}}

ここまでに定義した三つの演算の間には、二つの重要な恒等式が成り立つ。どちらも「二度作用させるとゼロになる」という形をしている。

\hypertarget{mathrmrotmathrmgradf-mathbf0}{%
\subsubsection{\texorpdfstring{7.6.1
\(\mathrm{rot}(\mathrm{grad}\,f) = \mathbf{0}\)}{7.6.1 \textbackslash mathrm\{rot\}(\textbackslash mathrm\{grad\}\textbackslash,f) = \textbackslash mathbf\{0\}}}\label{mathrmrotmathrmgradf-mathbf0}}

勾配でベクトル場を得て、その回転をとると必ずゼロベクトルになる。

\[\nabla \times (\nabla f) = \begin{pmatrix}
\frac{\partial}{\partial y}\frac{\partial f}{\partial z} - \frac{\partial}{\partial z}\frac{\partial f}{\partial y} \\[0.5em]
\frac{\partial}{\partial z}\frac{\partial f}{\partial x} - \frac{\partial}{\partial x}\frac{\partial f}{\partial z} \\[0.5em]
\frac{\partial}{\partial x}\frac{\partial f}{\partial y} - \frac{\partial}{\partial y}\frac{\partial f}{\partial x}
\end{pmatrix} = \begin{pmatrix} 0 \\ 0 \\ 0 \end{pmatrix}\]

各成分で、偏微分の順序交換
\(\frac{\partial^2 f}{\partial y\partial z} = \frac{\partial^2 f}{\partial z\partial y}\)
により打ち消し合う。これは \(f\)
がなめらか（二階連続微分可能）であれば常に成り立つ。

物理的意味：「勾配場には渦がない」。保存的な重力場や静電場がこれに当たる。

\hypertarget{mathrmdivmathrmrotmathbff-0}{%
\subsubsection{\texorpdfstring{7.6.2
\(\mathrm{div}(\mathrm{rot}\,\mathbf{F}) = 0\)}{7.6.2 \textbackslash mathrm\{div\}(\textbackslash mathrm\{rot\}\textbackslash,\textbackslash mathbf\{F\}) = 0}}\label{mathrmdivmathrmrotmathbff-0}}

回転でベクトル場を得て、その発散をとると必ずゼロになる。

\[\nabla \cdot (\nabla \times \mathbf{F}) = \frac{\partial}{\partial x}\!\left(\frac{\partial F_z}{\partial y} - \frac{\partial F_y}{\partial z}\right) + \frac{\partial}{\partial y}\!\left(\frac{\partial F_x}{\partial z} - \frac{\partial F_z}{\partial x}\right) + \frac{\partial}{\partial z}\!\left(\frac{\partial F_y}{\partial x} - \frac{\partial F_x}{\partial y}\right)\]

展開すると、\(\frac{\partial^2 F_z}{\partial x\partial y}\) と
\(-\frac{\partial^2 F_z}{\partial y\partial x}\)
のように、やはり偏微分の順序交換により項がペアで打ち消し合い、総和は
\(0\) になる。

物理的意味：「渦場に湧き出しはない」。磁場がこれに当たる（モノポールが存在しないことの数学的表現）。

\begin{quote}
\textbf{【ここまでのチェックポイント — §7.5–§7.6】}

\begin{itemize}
\item
  \(\nabla^2 f = \mathrm{div}(\mathrm{grad}\,f)\)
  はスカラー場のラプラシアン。熱伝導や波動を記述する。
\item
  \(\mathrm{rot}(\mathrm{grad}\,f) = \mathbf{0}\)：勾配場に渦はない。
\item
  \(\mathrm{div}(\mathrm{rot}\,\mathbf{F}) = 0\)：渦場に湧き出しはない。
\item
  いずれも偏微分の順序交換
  \(\partial_i\partial_j = \partial_j\partial_i\) からの当然の帰結。
\end{itemize}
\end{quote}

\begin{center}\rule{0.5\linewidth}{0.5pt}\end{center}

\hypertarget{ux30caux30d6ux30e9ux306eux516cux5f0fux96c6}{%
\subsection{§7.7
ナブラの公式集}\label{ux30caux30d6ux30e9ux306eux516cux5f0fux96c6}}

実用的な計算のために、\(\nabla\)
を含む積の微分則と代表的なベクトル恒等式をまとめておく。以下、\(f, g\)
はスカラー場、\(\mathbf{F}, \mathbf{G}\) はベクトル場とする。

\hypertarget{ux7a4dux306eux5faeux5206ux5247}{%
\subsubsection{7.7.1 積の微分則}\label{ux7a4dux306eux5faeux5206ux5247}}

\textbf{勾配の積}：

\[\nabla(fg) = f\,\nabla g + g\,\nabla f\]

\textbf{発散の積}：

\[\nabla \cdot (f\mathbf{F}) = (\nabla f) \cdot \mathbf{F} + f\,(\nabla \cdot \mathbf{F})\]

\textbf{回転の積}：

\[\nabla \times (f\mathbf{F}) = (\nabla f) \times \mathbf{F} + f\,(\nabla \times \mathbf{F})\]

\textbf{クロス積の発散}：

\[\nabla \cdot (\mathbf{F} \times \mathbf{G}) = (\nabla \times \mathbf{F}) \cdot \mathbf{G} - \mathbf{F} \cdot (\nabla \times \mathbf{G})\]

\hypertarget{ux4e8cux56deux4f5cux7528ux306eux6052ux7b49ux5f0f}{%
\subsubsection{7.7.2
二回作用の恒等式}\label{ux4e8cux56deux4f5cux7528ux306eux6052ux7b49ux5f0f}}

スカラー場の勾配の回転と、ベクトル場の回転の発散は §7.6
で導いた。公式集として再掲する。

\textbf{勾配の回転はゼロ}：

\[\nabla \times (\nabla f) = \mathbf{0}\]

\textbf{回転の発散はゼロ}：

\[\nabla \cdot (\nabla \times \mathbf{F}) = 0\]

\textbf{回転の回転}：

\[\nabla \times (\nabla \times \mathbf{F}) = \nabla(\nabla \cdot \mathbf{F}) - \nabla^2 \mathbf{F}\]

これはラプラシアンを発散と回転に分解する式である。導出には BAC-CAB
則（§7.7.4）を使う。

\textbf{スカラー場の勾配の発散}（= ラプラシアン）：

\[\nabla \cdot (\nabla f) = \nabla^2 f\]

§7.5 の定義そのもの。

\hypertarget{ux30b9ux30abux30e9ux30fcux4e09ux91cdux7a4d}{%
\subsubsection{7.7.3
スカラー三重積}\label{ux30b9ux30abux30e9ux30fcux4e09ux91cdux7a4d}}

三つの縦ベクトル \(\mathbf{A}, \mathbf{B}, \mathbf{C}\)
のスカラー三重積は、\(\mathbf{A}\) と \(\mathbf{B}\times\mathbf{C}\)
のドット積で定義される。

\[\mathbf{A} \cdot (\mathbf{B} \times \mathbf{C})\]

この値は \(\mathbf{A}, \mathbf{B}, \mathbf{C}\)
の張る平行六面体の符号付き体積に等しい。巡回置換に対して不変であり、二つのベクトルを入れ替えると符号が反転する。

\[\mathbf{A} \cdot (\mathbf{B} \times \mathbf{C}) = \mathbf{B} \cdot (\mathbf{C} \times \mathbf{A}) = \mathbf{C} \cdot (\mathbf{A} \times \mathbf{B})\]

\hypertarget{ux30d9ux30afux30c8ux30ebux4e09ux91cdux7a4dbac-cab-ux5247}{%
\subsubsection{7.7.4 ベクトル三重積（BAC-CAB
則）}\label{ux30d9ux30afux30c8ux30ebux4e09ux91cdux7a4dbac-cab-ux5247}}

\[\mathbf{A} \times (\mathbf{B} \times \mathbf{C}) = \mathbf{B}(\mathbf{A} \cdot \mathbf{C}) - \mathbf{C}(\mathbf{A} \cdot \mathbf{B})\]

クロス積を二重に使うときに現れる。右辺の各項はスカラー倍されたベクトルであり、BAC-CAB
の語呂で記憶される。

\begin{quote}
\textbf{注}
（なぜ公式を列挙したのか）ここに並べた公式は、大学のベクトル解析の試験で暗記を強いられる代表格である。では、\(d\)
と \(\ast\)
で書き直せば公式の数が劇的に減るかというと、そうでもない。微分形式の側でも、\(d(f\omega) = df\wedge\omega + f d\omega\)
のような規則はそれなりにある。

ベクトル解析の真の問題は、公式の数ではない。\textbf{すべてが同じ「$3$ 成分の矢印」に見えてしまうこと}だ。\(\nabla f\)
も \(\nabla \times \mathbf{F}\) も \(\nabla^2\mathbf{F}\)
も、どれも画面上ではただの矢印の束である。しかし \(\nabla f\)
は勾配（ポテンシャルの傾き）、\(\nabla \times \mathbf{F}\)
は渦の軸------矢印の幾何的意味はまったく異なる。お絵かきと直感に頼っていると、問題が複雑になった瞬間に区別がつかなくなる。微分形式は、\(0\)-form、\(1\)-form、\(2\)-form、\(3\)-form
と測定器の次数を明示することで、この混同を原理的に防ぐ。\(d\)
は常に次数を \(+1\) し、\(\ast\)
は常に次数を反転させる------少数の規則に型の情報が載っているから、矢印を見なくても式の形だけで何をしているかがわかる。これが、第8章以降で
form の言葉に立ち戻る理由である。
\end{quote}

\begin{quote}
\textbf{注} （ここだけの話）\(\mathrm{grad}=d\),
\(\mathrm{rot}=\ast d\), \(\mathrm{div}=\ast d\ast\)
と対応を書き下すのは、実はベクトル解析を解体するための妥協案という側面がある。

一部の微分形式の啓蒙的な説明では、きれいに書ける範囲が強調されることが多い。しかし、ベクトル解析を本気で翻訳しようとすると、\(\ast\)
が何度も顔を出し、必ずしも見た目が簡単になるわけではない。たとえば
§7.7.4 のベクトル三重積 \(\mathbf{A}\times(\mathbf{B}\times\mathbf{C})\)
をこの枠組みで書き直すと、\(\ast(\alpha \wedge \ast(\beta \wedge \gamma))\)
のように \(\ast\)
の入れ子が発生し、むしろベクトル解析よりも不格好な姿を晒すことになる。

それでも一部の啓蒙書が美しさを強調するのは、計量という「強力な制限」を召喚する前の世界------すなわち
\(\ast\) を使わずに \(d\) と \(\wedge\)
だけで編まれた風景------があまりに純粋で美しいからだ。計量（内積）を定義する前からマクスウェル方程式のような物理法則の骨格が定まる、という事実は確かに衝撃的である。

その「純粋な微分形式の世界」へ向かうために、あえて不格好な \(\ast\)
まみれの翻訳作業を一度やり遂げる。それが本書のこの段階の狙いである。この「微分形式によるベクトル解析」の計算を汚い、不自然だ、と読者が思えたなら、それは正しい。その時、君はもうベクトル解析の初学者を卒業している。
\end{quote}

\begin{center}\rule{0.5\linewidth}{0.5pt}\end{center}

\hypertarget{ux7a4dux5206ux5b9aux7406ux30b9ux30c8ux30fcux30afux30b9ux30acux30a6ux30b9ux30b0ux30eaux30fcux30f3}{%
\subsection{§7.8
積分定理------ストークス・ガウス・グリーン}\label{ux7a4dux5206ux5b9aux7406ux30b9ux30c8ux30fcux30afux30b9ux30acux30a6ux30b9ux30b0ux30eaux30fcux30f3}}

ベクトル解析の最も重要な仕事は、微分演算（grad, div,
rot）と積分を結びつけることにある。以下に三つの積分定理を、\(\nabla\)
の言葉で示す。

\hypertarget{ux30b9ux30c8ux30fcux30afux30b9ux306eux5b9aux7406}{%
\subsubsection{7.8.1
ストークスの定理}\label{ux30b9ux30c8ux30fcux30afux30b9ux306eux5b9aux7406}}

曲面 \(S\) とその境界である閉曲線 \(C = \partial S\) に対し、

\[\oint_C \mathbf{F} \cdot d\mathbf{r} = \iint_S (\nabla \times \mathbf{F}) \cdot \mathbf{n}\,dS\]

左辺は曲線 \(C\) に沿った \(\mathbf{F}\) の線積分（循環）、右辺は曲面
\(S\)
上の回転の法線成分の面積分である。「境界での循環は、内部の渦の総和に等しい」という意味を持つ。\(\mathbf{n}\)
は曲面の単位法線ベクトルで、\(C\) の向きに対して右ねじの関係で選ぶ。

\begin{quote}
\textbf{注} （グリーンの定理）\(S\) が \(xy\) 平面内の領域 \(D\)
で、\(\mathbf{F}\) の成分が \(F_x=P,\;F_y=Q,\;F_z=0\)
の場合、ストークスの定理は次の形になる。これはグリーンの定理と呼ばれる。

\[\oint_{\partial D} (P\,dx + Q\,dy) = \iint_D \left(\frac{\partial Q}{\partial x} - \frac{\partial P}{\partial y}\right) dx\,dy\]
\end{quote}

\hypertarget{ux30acux30a6ux30b9ux306eux767aux6563ux5b9aux7406}{%
\subsubsection{7.8.2
ガウスの発散定理}\label{ux30acux30a6ux30b9ux306eux767aux6563ux5b9aux7406}}

立体 \(V\) とその境界である閉曲面 \(S = \partial V\) に対し、

\[\oiint_S \mathbf{F} \cdot \mathbf{n}\,dS = \iiint_V (\nabla \cdot \mathbf{F})\,dV\]

左辺は閉曲面 \(S\) を通る \(\mathbf{F}\) の外向きの流束、右辺は立体
\(V\)
内の発散の体積分である。「境界を通る流れの総和は、内部の湧き出しの総和に等しい」という意味を持つ。

\hypertarget{ux7a4dux5206ux5b9aux7406ux306eux5171ux901aux69cbux9020}{%
\subsubsection{7.8.3
積分定理の共通構造}\label{ux7a4dux5206ux5b9aux7406ux306eux5171ux901aux69cbux9020}}

ストークスの定理とガウスの発散定理を並べて見ると、共通のパターンが浮かぶ。

\begin{longtable}[]{@{}
  >{\raggedright\arraybackslash}p{(\columnwidth - 6\tabcolsep) * \real{0.2500}}
  >{\raggedright\arraybackslash}p{(\columnwidth - 6\tabcolsep) * \real{0.2500}}
  >{\raggedright\arraybackslash}p{(\columnwidth - 6\tabcolsep) * \real{0.2500}}
  >{\raggedright\arraybackslash}p{(\columnwidth - 6\tabcolsep) * \real{0.2500}}@{}}
\toprule\noalign{}
\begin{minipage}[b]{\linewidth}\raggedright
定理
\end{minipage} & \begin{minipage}[b]{\linewidth}\raggedright
境界での積分
\end{minipage} & \begin{minipage}[b]{\linewidth}\raggedright
内部での積分
\end{minipage} & \begin{minipage}[b]{\linewidth}\raggedright
微分演算
\end{minipage} \\
\midrule\noalign{}
\endhead
\bottomrule\noalign{}
\endlastfoot
ストークス & \(\displaystyle\oint_C \mathbf{F}\cdot d\mathbf{r}\) &
\(\displaystyle\iint_S (\nabla\times\mathbf{F})\cdot\mathbf{n}\,dS\) &
rot \\
ガウス & \(\displaystyle\oiint_S \mathbf{F}\cdot\mathbf{n}\,dS\) &
\(\displaystyle\iiint_V (\nabla\cdot\mathbf{F})\,dV\) & div \\
\end{longtable}

いずれも「境界で測った量 = 内部の微分量の総和」という形をしている。§7.6
の \(\mathrm{div}(\mathrm{rot})=0\)
は、この枠組みでは「渦場には湧き出しがないから、閉曲面で囲んでも中から出てくるものはない」と読める。

しかし、ここには一つの不満が残る。ストークスとガウスは、grad / rot / div
という異なる演算子を使うため、\textbf{別々の定理として暗記しなければならない}。第5章で
\(\int_{\partial M}\omega = \int_M d\omega\)
という一本の式がこれらすべてを統合していたことを思い出してほしい。第8章では、\(\nabla\)
の世界の積分定理を、この一本の形に読み替える対応を整理する。

\begin{quote}
\textbf{【ここまでのチェックポイント — §7.8】}

\begin{itemize}
\item
  ストークスの定理：\(\oint_C \mathbf{F}\cdot d\mathbf{r} = \iint_S (\nabla\times\mathbf{F})\cdot\mathbf{n}\,dS\)
\item
  ガウスの発散定理：\(\oiint_S \mathbf{F}\cdot\mathbf{n}\,dS = \iiint_V (\nabla\cdot\mathbf{F})\,dV\)
\item
  グリーンの定理はストークスの定理の平面版。
\item
  これらの積分定理は「境界の積分 =
  内部の微分の積分」という共通構造を持つが、\(\nabla\)
  の言葉では別々の定理として扱わざるを得ない。
\end{itemize}
\end{quote}

\begin{center}\rule{0.5\linewidth}{0.5pt}\end{center}

\hypertarget{ux7b2c8ux7ae0ux3078-ux77e2ux5370ux3092ux898bux308bux306aux6e2cux5b9aux5668ux3092ux898bux308d}{%
\subsection{§7.9 第8章へ ------
矢印を見るな、測定器を見ろ}\label{ux7b2c8ux7ae0ux3078-ux77e2ux5370ux3092ux898bux308bux306aux6e2cux5b9aux5668ux3092ux898bux308d}}

第6章で、我々は外微分 \(d\) とホッジ・スター \(\ast\)
を組み合わせることで、\(\mathrm{grad}=d\)、\(\mathrm{rot}=\ast d\)、\(\mathrm{div}=\ast d\ast\)
という三つの型が現れることを見た。

本章では、同じ対象を標準的なベクトル解析の言葉で書き直した。\(\nabla f\)、\(\nabla\times\mathbf F\)、\(\nabla\cdot\mathbf F\)
は便利だが、画面上ではすべて矢印やスカラーに見える。しかし、それぞれが担う幾何的意味はまったく異なる。そして問題が複雑になるほど、矢印だけを見て区別するのは難しくなる。

次章では、第6章で得た \(d,\ast\) の型と、本章で導入した \(\nabla\)
記法を対応させる。\(\nabla f\) を \(df\) に、\(\nabla\times\mathbf F\)
を \(\ast d\omega\) に、\(\nabla\cdot\mathbf F\) を \(\ast d\ast\omega\)
に読み替えることで、標準ベクトル解析の式が、測定器の次数を持った式として整理される。矢印の絵に頼らず、測定器の代数で考える------その視点を、次章で標準ベクトル解析の記法と接続する。

\begin{quote}
\textbf{【ここまでのチェックポイント — 第7章全体】}

\begin{itemize}
\item
  \(\nabla\)
  を偏微分演算子を縦に並べた形式的ベクトルとして定義し、勾配・発散・回転・ラプラシアンを導入した。
\item
  \(\mathrm{grad}\,f = \nabla f\),
  \(\mathrm{div}\,\mathbf{F} = \nabla \cdot \mathbf{F} = \operatorname{tr}(\mathbf{J}_\mathbf{F})\),
  \(\mathrm{rot}\,\mathbf{F} = \nabla \times \mathbf{F}\)。
\item
  \(\nabla^2 = \mathrm{div}\,\mathrm{grad}\),
  \(\mathrm{rot}(\mathrm{grad}) = 0\),
  \(\mathrm{div}(\mathrm{rot}) = 0\)。
\item
  クロス積の外積行列表示と BAC-CAB 則を確認した。
\item
  ベクトル解析の難しさは公式の数もさることながら、すべてが同じ矢印に見えることにある。微分形式は測定器（\(n\)-form）の次数で区別する。第8章でこの二つの言語を繫ぐ。
\end{itemize}
\end{quote}

\newpage

\hypertarget{ux7b2c8ux7ae0ux4e8cux3064ux306eux8a00ux8a9e-ux6e2cux5b9aux5668ux306eux5faeux5206ux3068ux5834ux306eux5faeux5206}{%
\chapter{第8章：二つの言語 ------
測定器の微分と、場の微分}\label{ux7b2c8ux7ae0ux4e8cux3064ux306eux8a00ux8a9e-ux6e2cux5b9aux5668ux306eux5faeux5206ux3068ux5834ux306eux5faeux5206}}

\hypertarget{ux672cux66f8ux306eux30cfux30a4ux30e9ux30a4ux30c8}{%
\subsection{§8.0
本書のハイライト}\label{ux672cux66f8ux306eux30cfux30a4ux30e9ux30a4ux30c8}}

　ここまで長かった。第1章で \(dx\)
を行ベクトルと定義し、第2章で面積測定器と体積測定器を組み立て、第3章で積分を再構築し、第5章で外微分
\(d\) を得て、第6章で計量 \(g\) とホッジ・スター \(\ast\)
を手に入れ、そして第7章ではついに \(\nabla\)
を正面から定義してベクトル解析を一気に展開した。

ここから先は、もう遠慮しない。\(d\) も \(\ast\) も \(\nabla\)
も、すべての道具が揃った。ベクトル解析の用語を引き合いに出すのに、もはや「ベクトル解析を知っている読者へ」と断る必要もない。第7章で
\(\nabla\) は正式に本書の住人になったのだから。

本章は、この本の\textbf{ハイライト}である。中心となる翻訳そのものは、驚くほど短い。ただし最後に、次章の実践編へ進むために、曲線座標でこの辞書がどう動くかを一度だけ見ておく。そのために第1章から第7章までの入念な準備があった------行列、ウェッジ積、外微分、計量、ホッジ・スター、そしてベクトル解析。これらの道具がすべて揃っているからこそ、本章ではただ「翻訳する」だけで、二つの世界が一つの絵に収まる。我々は二つの言語を手に入れた------測定器の微分（\(d\)）と、場の微分（\(\nabla\)）。この二つは、やっていることが根本的に違う。しかし積分を通じて繋がるとき、まったく同じ結果を与える。その仕組みを、これから徹底的に掘り下げる。

\begin{center}\rule{0.5\linewidth}{0.5pt}\end{center}

\hypertarget{ux4e8cux3064ux306eux5faeux5206ux4e8cux3064ux306eux4e16ux754c}{%
\subsection{§8.1
二つの微分、二つの世界}\label{ux4e8cux3064ux306eux5faeux5206ux4e8cux3064ux306eux4e16ux754c}}

\hypertarget{ux6e2cux5b9aux5668ux306bux5faeux5206ux3092ux4e57ux305bux308bd-ux306eux7acbux5834}{%
\subsubsection{\texorpdfstring{8.1.1 測定器に微分を乗せる------\(d\)
の立場}{8.1.1 測定器に微分を乗せる------d の立場}}\label{ux6e2cux5b9aux5668ux306bux5faeux5206ux3092ux4e57ux305bux308bd-ux306eux7acbux5834}}

第1章以来、我々は \(dx = \begin{pmatrix}1&0&0\end{pmatrix}\),
\(dy = \begin{pmatrix}0&1&0\end{pmatrix}\),
\(dz = \begin{pmatrix}0&0&1\end{pmatrix}\)
を横ベクトルの測定器として扱ってきた。これらは縦ベクトル（変位）を食わせると、その方向の変位量をスカラーで返す。

外微分 \(d\) は、この測定器に\textbf{微分を乗せる}操作である。

\(f(x,y,z)\) がスカラー場のとき、\(df\) は \(f\)
の変化を測る新しい測定器だ。各偏微分係数を成分として持つ横ベクトルになる。

\[df = \frac{\partial f}{\partial x}\,dx + \frac{\partial f}{\partial y}\,dy + \frac{\partial f}{\partial z}\,dz = \begin{pmatrix} \frac{\partial f}{\partial x} & \frac{\partial f}{\partial y} & \frac{\partial f}{\partial z} \end{pmatrix}\]

\(\omega = P\,dx + Q\,dy + R\,dz\) が \(1\)-form のとき、\(d\omega\) は
\(\omega\) の局所的なズレを測る \(2\)-form（反対称行列）になる。

\[d\omega = \left(\frac{\partial R}{\partial y}-\frac{\partial Q}{\partial z}\right)dy\wedge dz + \left(\frac{\partial P}{\partial z}-\frac{\partial R}{\partial x}\right)dz\wedge dx + \left(\frac{\partial Q}{\partial x}-\frac{\partial P}{\partial y}\right)dx\wedge dy\]

\(\eta = A\,dy\wedge dz + B\,dz\wedge dx + C\,dx\wedge dy\) が
\(2\)-form のとき、\(d\eta\) は \(3\)-form（三階反対称テンソル）になる。

\[d\eta = \left(\frac{\partial A}{\partial x} + \frac{\partial B}{\partial y} + \frac{\partial C}{\partial z}\right) dx\wedge dy\wedge dz\]

このように、\(d\) が作用するたびに測定器の次数が \(0\to1\to2\to3\)
と一つずつ上がっていく。核心は\textbf{場はそのまま、測定器の側が変化する}ことだ。\(f\)
はスカラー場のままだが、\(d\)
が作用するたびに新しい測定器------\(df\)（\(1\)-form）、\(d\omega\)（\(2\)-form）、\(d\eta\)（\(3\)-form）------が生み出される。

\hypertarget{ux5f62ux5f0fux7684ux306aux5faeux5206ux6f14ux7b97ux5b50ux306fux540cux3058ux307eux307eux65b0ux305fux306aux5834ux3092ux8003ux3048ux308bnabla-ux306eux7acbux5834}{%
\subsubsection{\texorpdfstring{8.1.2
形式的な微分演算子は同じまま、新たな場を考える------\(\nabla\)
の立場}{8.1.2 形式的な微分演算子は同じまま、新たな場を考える------\textbackslash nabla の立場}}\label{ux5f62ux5f0fux7684ux306aux5faeux5206ux6f14ux7b97ux5b50ux306fux540cux3058ux307eux307eux65b0ux305fux306aux5834ux3092ux8003ux3048ux308bnabla-ux306eux7acbux5834}}

一方、第7章で導入した

\[
\nabla
=
\begin{pmatrix}
\frac{\partial}{\partial x}\\[0.3em]
\frac{\partial}{\partial y}\\[0.3em]
\frac{\partial}{\partial z}
\end{pmatrix}
\]

の世界では、\textbf{形式的な微分演算子は同じまま、作用する相手に応じて新たな場が生まれる}。\(\nabla\)
はつねに同じ演算子だが、スカラー場 \(f\) に作用すればベクトル場
\(\nabla f\) を、ベクトル場 \(\mathbf{F}\) に作用すればスカラー場
\(\nabla\cdot\mathbf{F}\) やベクトル場 \(\nabla\times\mathbf{F}\)
を返す。

スカラー場 \(f(x,y,z)\) に \(\nabla\)
を作用させると、各方向の変化率を成分とする縦ベクトル場が得られる。

\[\nabla f = \begin{pmatrix} \frac{\partial f}{\partial x} \\[0.3em] \frac{\partial f}{\partial y} \\[0.3em] \frac{\partial f}{\partial z} \end{pmatrix}\]

これは空間の各点に立つ矢印であり、もはや測定器ではない。\(f\)
が最も急に増加する方向とその強さを表す。

成分が \(F_x,F_y,F_z\) であるベクトル場 \(\mathbf{F}\) に
\(\nabla\cdot\) を作用させると、スカラー場が得られる。

\[\nabla \cdot \mathbf{F} = \frac{\partial F_x}{\partial x} + \frac{\partial F_y}{\partial y} + \frac{\partial F_z}{\partial z}\]

各方向の変化率の和------湧き出しの強さ------を返す。

同じ \(\mathbf{F}\) に \(\nabla\times\)
を作用させると、再び縦ベクトル場が得られる。

\[\nabla \times \mathbf{F} = \begin{pmatrix} \frac{\partial F_z}{\partial y} - \frac{\partial F_y}{\partial z} \\[0.3em] \frac{\partial F_x}{\partial z} - \frac{\partial F_z}{\partial x} \\[0.3em] \frac{\partial F_y}{\partial x} - \frac{\partial F_x}{\partial y} \end{pmatrix}\]

渦の軸方向と強さを表す。

\(\nabla\)
の世界では、結果はつねに実空間の中の矢印かスカラーとして返される。\(\nabla f\)
も \(\nabla\times\mathbf{F}\)
も紙面上では同じ「矢印」であり、型による区別はつかない。

\begin{quote}
\textbf{【ここまでのチェックポイント — §8.1】}

\begin{itemize}
\item
  \(d\)
  の立場：場はそのまま、測定器の側が変化する。新しい測定器が生まれる。
\item
  \(\nabla\) の立場：形式的な微分演算子は同じまま、新たな場が生まれる。
\end{itemize}
\end{quote}

\begin{center}\rule{0.5\linewidth}{0.5pt}\end{center}

\hypertarget{ux7ffbux8a33ux8f9eux66f8ux306eux5b8cux6210}{%
\subsection{§8.2
翻訳辞書の完成}\label{ux7ffbux8a33ux8f9eux66f8ux306eux5b8cux6210}}

第6章 §6.5 で \(\mathrm{grad}=d\), \(\mathrm{rot}=\ast d\),
\(\mathrm{div}=\ast d\ast\) という対応を導入した。ここで、この辞書を
\(\nabla\)
の言葉と突き合わせて、完全な形で整理する。ここで完成するのは、三次元ユークリッド空間・本書の行列表現の範囲で使う翻訳辞書である。

以下の辞書では、ベクトル場

\[
\mathbf{F}
=
\begin{pmatrix}
F_x\\
F_y\\
F_z
\end{pmatrix}
\]

に対応する \(1\)-form を、計量 \(g\) を使って

\[
\omega=\mathbf{F}^T g
\]

と読む。デカルト座標では \(g=I\) なので、

\[
\omega
=
\begin{pmatrix}
F_x & F_y & F_z
\end{pmatrix}
=
F_x\,dx+F_y\,dy+F_z\,dz
\]

である。直交デカルト座標では成分は同じに見えるが、これは「同じ対象だ」という意味ではない。ベクトル場は縦ベクトル、\(1\)-form
は横ベクトルであり、両者の対応には計量 \(g\) が入っている。

\hypertarget{ux52feux914d-1}{%
\subsubsection{8.2.1 勾配}\label{ux52feux914d-1}}

\[
\mathrm{grad}\,f = \nabla f
\quad\longleftrightarrow\quad
df = \frac{\partial f}{\partial x}dx + \frac{\partial f}{\partial y}dy + \frac{\partial f}{\partial z}dz
\]

\(\nabla f\) は\textbf{縦ベクトル}、\(df\)
は\textbf{横ベクトル}である。デカルト座標では成分が同じに見えるが、型は異なる。第6章の記法では、\(df\)
を対応する縦ベクトルとして読むには \(g^{-1}(df)^T\)
を使う。デカルト座標では \(g=I\) なので、その成分が通常の \(\nabla f\)
と一致する。

\hypertarget{ux56deux8ee2-1}{%
\subsubsection{8.2.2 回転}\label{ux56deux8ee2-1}}

\[
\mathrm{rot}\,\mathbf{F} = \nabla \times \mathbf{F}
\quad\longleftrightarrow\quad
\ast d\omega
\]

\(\nabla\times\mathbf{F}\) は\textbf{縦ベクトル}、\(d\omega\)
は\textbf{$2$-form（反対称行列）}、そして \(\ast d\omega\)
は\textbf{$1$-form（横ベクトル）}である。回転は「外微分で次数を上げてから、ホッジ・スターで次数を戻す」という二段階の操作だ。

\(\nabla\times\mathbf{F}\) の \(x\) 成分
\(\frac{\partial F_z}{\partial y} - \frac{\partial F_y}{\partial z}\)
は、\(\ast d\omega\) の \(dx\)
の係数と一致する。これは第6章で確認した対応そのものである。

\hypertarget{ux767aux6563-1}{%
\subsubsection{8.2.3 発散}\label{ux767aux6563-1}}

\[
\mathrm{div}\,\mathbf{F} = \nabla \cdot \mathbf{F}
\quad\longleftrightarrow\quad
\ast d \ast \omega
\]

\(\nabla\cdot\mathbf{F}\) は\textbf{スカラー場}、\(\ast d\ast\omega\)
も\textbf{$0$-form（スカラー場）}である。発散は「ホッジ・スターで
\(2\)-form に変え、外微分で \(3\)-form
に上げ、再びホッジ・スターでスカラーに落とす」という三段階の操作だ。

\hypertarget{ux8f9eux66f8ux306eux8868}{%
\subsubsection{8.2.4 辞書の表}\label{ux8f9eux66f8ux306eux8868}}

\begin{longtable}[]{@{}
  >{\centering\arraybackslash}p{(\columnwidth - 6\tabcolsep) * \real{0.2500}}
  >{\centering\arraybackslash}p{(\columnwidth - 6\tabcolsep) * \real{0.2500}}
  >{\centering\arraybackslash}p{(\columnwidth - 6\tabcolsep) * \real{0.2500}}
  >{\centering\arraybackslash}p{(\columnwidth - 6\tabcolsep) * \real{0.2500}}@{}}
\toprule\noalign{}
\begin{minipage}[b]{\linewidth}\centering
演算
\end{minipage} & \begin{minipage}[b]{\linewidth}\centering
ベクトル解析側
\end{minipage} & \begin{minipage}[b]{\linewidth}\centering
微分形式（測定器）側
\end{minipage} & \begin{minipage}[b]{\linewidth}\centering
ベクトル場として戻すなら
\end{minipage} \\
\midrule\noalign{}
\endhead
\bottomrule\noalign{}
\endlastfoot
勾配 & \(\mathrm{grad}\,f\) & \(df\) & \(g^{-1}(df)^T\) \\
回転 & \(\mathrm{rot}\,\mathbf{F}\) & \(\ast d\omega\) &
\(g^{-1}(\ast d\omega)^T\) \\
発散 & \(\mathrm{div}\,\mathbf{F}\) & \(\ast d\ast\omega\) & そのまま
\(0\)-form \\
\end{longtable}

\(\ast\) の使用回数が増えるほど、\(\nabla\) 表記と \(d,\ast\)
表記の間の「翻訳コスト」が上がる。 ここで \(\omega=\mathbf{F}^Tg\)
である。勾配は \(\ast\) がゼロ回だが、\(df\)
を通常の勾配ベクトルとして読むには計量による変換 \(g^{-1}(df)^T\)
が入る。回転は \(\ast\) が \(1\) 回、発散は \(2\) 回である。この計量
\(g\) と \(\ast\) による翻訳辞書が、第7章 §7.7
の注で「微分形式のほうが汚い」と感じた原因だった。

\begin{quote}
\textbf{【ここまでのチェックポイント — §8.2】}

\begin{itemize}
\item
  \(\nabla f \leftrightarrow df\),
  \(\nabla\times\mathbf{F} \leftrightarrow \ast d\omega\),
  \(\nabla\cdot\mathbf{F} \leftrightarrow \ast d\ast\omega\)。
\item
  \(\ast\) の使用回数が勾配 \(0\)、回転 \(1\)、発散
  \(2\)。この翻訳コストが「微分形式は汚い」という印象の正体。
\end{itemize}
\end{quote}

\begin{center}\rule{0.5\linewidth}{0.5pt}\end{center}

\hypertarget{ux30b9ux30c8ux30fcux30afux30b9ux306eux5b9aux7406ux3092ux7ffbux8a33ux3059ux308b}{%
\subsection{§8.3
ストークスの定理を翻訳する}\label{ux30b9ux30c8ux30fcux30afux30b9ux306eux5b9aux7406ux3092ux7ffbux8a33ux3059ux308b}}

第5章 §5.6 で、我々はすでに \(\int_{\partial S}\omega = \int_S d\omega\)
という形のストークスの定理を導いた。第7章 §7.8.1 では、同じ定理を
\(\nabla\) の言葉で
\(\oint_C \mathbf{F}\cdot d\mathbf{r} = \iint_S (\nabla\times\mathbf{F})\cdot\mathbf{n}\,dS\)
と書いた。この二つが同じものであることを、§8.2 の辞書を使って確認する。

\hypertarget{nabla-ux304bux3089-d-ux3078}{%
\subsubsection{\texorpdfstring{8.3.1 \(\nabla\) から \(d\)
へ}{8.3.1 \textbackslash nabla から d へ}}\label{nabla-ux304bux3089-d-ux3078}}

\(\nabla\) 表記のストークスの定理を出発点とする。

\[\oint_C \mathbf{F} \cdot d\mathbf{r} = \iint_S (\nabla \times \mathbf{F}) \cdot \mathbf{n}\,dS\]

左辺の \(\oint_C \mathbf{F}\cdot d\mathbf{r}\) は、ベクトル場
\(\mathbf{F}\) を曲線 \(C\) に沿って線積分したものだ。\(\mathbf{F}\)
に対応する \(1\)-form を \(\omega = F_x\,dx + F_y\,dy + F_z\,dz\)
とすれば、第3章の定義によりこれは \(\int_C \omega\) に等しい。

\[\oint_C \mathbf{F} \cdot d\mathbf{r} = \int_{\partial S} \omega\]

右辺の \(\iint_S (\nabla\times\mathbf{F})\cdot\mathbf{n}\,dS\)
は、回転の法線成分の面積分である。ここでは、ベクトル場の法線成分つき面積分を、対応する
\(2\)-form の面積分として読んでいる。辞書
\(\nabla\times\mathbf{F} \leftrightarrow \ast d\omega\)
を使うと、\((\nabla\times\mathbf{F})\cdot\mathbf{n}\,dS\) は \(2\)-form
\(\ast(\ast d\omega) = d\omega\)
の面積分に対応する（\(\ast\ast = \mathrm{id}\) により \(\ast\)
が打ち消し合う）。

\[\iint_S (\nabla \times \mathbf{F}) \cdot \mathbf{n}\,dS = \int_S d\omega\]

したがって、

\[\int_{\partial S} \omega = \int_S d\omega\]

これは第5章 §5.6 で得た式そのものである。\(\nabla\)
の世界のストークスの定理は、辞書を通じて \(d\)
の言葉に翻訳すると、ただの \(\int_{\partial M}\omega = \int_M d\omega\)
になる。

\hypertarget{ux7ffbux8a33ux306eux5f80ux5fa9ux3069ux3061ux3089ux304bux3089ux3067ux3082ux884cux3051ux308b}{%
\subsubsection{8.3.2
翻訳の往復------どちらからでも行ける}\label{ux7ffbux8a33ux306eux5f80ux5fa9ux3069ux3061ux3089ux304bux3089ux3067ux3082ux884cux3051ux308b}}

逆方向もできる。\(d\) の言葉で書かれた式に辞書を逆に当てれば、\(\nabla\)
の言葉に戻る。この双方向の翻訳が自由にできることこそ、第6章で \(\ast\)
を導入した最大の成果である。

\begin{quote}
\textbf{注}
（どちらの言葉で考えるべきか）結論から言えば、問題による。境界の形状が単純で、場の対称性が高い問題では
\(\nabla\)
表記の方が見通しがよい。一方、座標系が歪んでいたり、場の「型」を明確に区別したい問題では
\(d,\ast\) 表記が有利である。重要なのは、両方の言語を使えることだ。
\end{quote}

\begin{quote}
\textbf{注}（筆者の本音）本文では「問題に応じて二つの言語を使い分ける」と書いた。これは読者のための建前である。筆者自身は、可能な限り
\(d,\ast\) の言葉で考える。なぜなら、\(\nabla\)
の記法は計算には便利だが、勾配・回転・発散の型の違いを同じ矢印やスカラーの中に押し込めてしまうからである。本書の目的は、ベクトル解析を禁止することではなく、その背後で暗黙に行われている型変換を見える場所に出すことにある。
\end{quote}

\begin{center}\rule{0.5\linewidth}{0.5pt}\end{center}

\hypertarget{ux30acux30a6ux30b9ux306eux5b9aux7406ux3092ux7ffbux8a33ux3059ux308b}{%
\subsection{§8.4
ガウスの定理を翻訳する}\label{ux30acux30a6ux30b9ux306eux5b9aux7406ux3092ux7ffbux8a33ux3059ux308b}}

第7章 §7.8.2 のガウスの発散定理も、同様に翻訳できる。

\[\oiint_{\partial V} \mathbf{F} \cdot \mathbf{n}\,dS = \iiint_V (\nabla \cdot \mathbf{F})\,dV\]

左辺の面積分は、\(\mathbf{F}\) に対応する \(1\)-form \(\omega\) を使って
\(\int_{\partial V} \ast\omega\) と書ける（\(\ast\omega\) は \(2\)-form
であり、閉曲面上の面積分の対象となる）。

右辺の発散の体積分は、辞書
\(\nabla\cdot\mathbf{F} \leftrightarrow \ast d\ast\omega\) により
\(\int_V \ast d\ast\omega\)
と読める。ただし、スカラー場そのものは、まだ積分される \(3\)-form
ではない。本書の流儀では、体積分の対象は \(3\)-form
である。したがって、\(\ast d\ast\omega\) を体積分するには、もう一度
\(\ast\) を作用させて \(d\ast\omega\) という \(3\)-form
として読む。すなわち \(\int_V d\ast\omega\) である。

結局、ガウスの定理は次の形に翻訳される。

\[\int_{\partial V} \ast\omega = \int_V d\ast\omega\]

ここで \(\eta = \ast\omega\)（\(2\)-form）と置けば、

\[\int_{\partial V} \eta = \int_V d\eta\]

これは第5章 §5.7 で得たガウスの定理を \(d\)
の言葉で書いた表現そのものだ。そしてこの式は、すでに §5.9.1 で
\(\int_{\partial M}\omega = \int_M d\omega\) の \(k=2\)
の場合として統合されている。

\hypertarget{ux4e09ux3064ux306eux5b9aux7406ux4e00ux3064ux306eux5f0f}{%
\subsubsection{8.4.1
三つの定理、一つの式}\label{ux4e09ux3064ux306eux5b9aux7406ux4e00ux3064ux306eux5f0f}}

第7章で別々の定理として並べた三つの積分定理------グリーン、ストークス、ガウス------は、\(\nabla\)
の言葉では確かに異なる顔をしていた。しかし \(d\)
の言葉に翻訳すれば、すべてが
\(\int_{\partial M}\omega = \int_M d\omega\) の \(k\)
の値が違うだけの同一の式になる。

\begin{longtable}[]{@{}llll@{}}
\toprule\noalign{}
\(k\) & \(\partial M\) & \(M\) & 定理 \\
\midrule\noalign{}
\endhead
\bottomrule\noalign{}
\endlastfoot
\(0\) & \(B-A\)（二点） & 曲線 & 微積分学の基本定理 \\
\(1\) & 閉曲線 & 曲面 & ストークス（グリーンを含む） \\
\(2\) & 閉曲面 & 立体 & ガウス \\
\end{longtable}

\(k\)
が違うだけだ。もはや三つの定理を別々に暗記する理由はない。\(\int_{\partial M}\omega = \int_M d\omega\)
一つ覚えれば、あとは \(M\) の次元に応じて適用するだけである。

\begin{quote}
\textbf{【ここまでのチェックポイント — §8.3–§8.4】}

\begin{itemize}
\item
  \(\nabla\) 表記のストークスの定理は、辞書を当てると
  \(\int_{\partial S}\omega = \int_S d\omega\) に還元される。
\item
  \(\nabla\) 表記のガウスの定理は、辞書を当てると
  \(\int_{\partial V}\eta = \int_V d\eta\) に還元される。
\item
  どちらも \(\int_{\partial M}\omega = \int_M d\omega\) の \(k\)
  が違うだけ。三定理は一つの式の異なる現れである。
\end{itemize}
\end{quote}

\begin{center}\rule{0.5\linewidth}{0.5pt}\end{center}

\hypertarget{ux306aux305cux4e8cux3064ux306eux8a00ux8a9eux306fux4e00ux81f4ux3059ux308bux306eux304b}{%
\subsection{§8.5
なぜ二つの言語は一致するのか}\label{ux306aux305cux4e8cux3064ux306eux8a00ux8a9eux306fux4e00ux81f4ux3059ux308bux306eux304b}}

第6章 §6.3.1 で、スカラーを得る二つの方法を導入した。§8.1 では、それを
\(d\) の世界と \(\nabla\)
の世界として対比した。ここでは同じ説明を繰り返さない。§8.2〜§8.4
で見た翻訳を踏まえて、なぜ二つの言語が積分で同じ値を与えるのかを整理する。

\hypertarget{ux4e8cux3064ux306eux65b9ux6cd5ux306eux9055ux3044}{%
\subsubsection{8.5.1
二つの方法の違い}\label{ux4e8cux3064ux306eux65b9ux6cd5ux306eux9055ux3044}}

§8.1 で見たように、\(d\) と \(\nabla\)
は同じことをしているわけではない。\(d\)
は測定器の次数を上げ、\(0\)-form、\(1\)-form、\(2\)-form、\(3\)-form
のあいだを進む。一方、\(\nabla\)
は同じ形式的な微分演算子から、スカラー場や縦ベクトル場を作る。

つまり、\(d\) の世界では「何を測る測定器なのか」が変わり、\(\nabla\)
の世界では「どのような場が生まれたのか」が変わる。この違いが、第7章で見た「すべてが同じ矢印に見える」問題の正体である。

\hypertarget{ux306aux305cux4e00ux81f4ux3059ux308bux306eux304b}{%
\subsubsection{8.5.2
なぜ一致するのか}\label{ux306aux305cux4e00ux81f4ux3059ux308bux306eux304b}}

方法①と方法②は、やっていることが\textbf{根本的に違う}。一方は測定器を作り替え、他方は新たな場を作る。操作方法も、出てくる対象の種類も異なる。

なのに積分すると一致する。たとえばスカラー場 \(f\)
の変化を線積分するとき、\(df\)（横ベクトルの測定器）を使っても
\(\nabla f\)（縦ベクトルの矢印）を使っても、結果はどちらも \(f(B)-f(A)\)
だ。

\[\int_C df = \int_C \nabla f \cdot d\mathbf{r} = f(B) - f(A)\]

なぜか。計量 \(g\) と \(\ast\) が両者を繫いでいるからだ。\(\nabla f\)
のような「矢印」と \(df\)
のような「測定器」は、型としては別物である。デカルト座標では \(g=I\)
なので同じ成分に見えるが、一般には \(g\)
による横・縦の変換を経由して対応づける。この変換を経由することで、\(\nabla f\)
と \(df\)、\(\nabla\times\mathbf{F}\) と
\(\ast d\omega\)、\(\nabla\cdot\mathbf{F}\) と \(\ast d\ast\omega\)
が、積分のレベルで同じ結果を与えることが保証される。

そして \(\ast\ast = \mathrm{id}\)
があるから、変換を何度往復しても情報は失われない。この往復の自由こそ、両方の言語を話せることの力である。

\begin{quote}
\textbf{注}（二つの方法の使い分け------筆者の立場）筆者は、問題に応じて両方の方法を行き来する。対称性の高い系では
\(\nabla\)
で見通しをつけ、座標系が歪んだり場の種類が入り組んだりする系では
\(d,\ast\)
に翻訳して型の整合性を確認する。どちらか一方に固執する必要はない。重要なのは、翻訳の辞書を頭に入れておくことだ。
\end{quote}

\begin{quote}
\textbf{【ここまでのチェックポイント — §8.5】}

\begin{longtable}[]{@{}
  >{\raggedright\arraybackslash}p{(\columnwidth - 4\tabcolsep) * \real{0.3333}}
  >{\raggedright\arraybackslash}p{(\columnwidth - 4\tabcolsep) * \real{0.3333}}
  >{\raggedright\arraybackslash}p{(\columnwidth - 4\tabcolsep) * \real{0.3333}}@{}}
\toprule\noalign{}
\begin{minipage}[b]{\linewidth}\raggedright
\end{minipage} & \begin{minipage}[b]{\linewidth}\raggedright
\textbf{$d$ の言語}
\end{minipage} & \begin{minipage}[b]{\linewidth}\raggedright
\textbf{$\nabla$ の言語}
\end{minipage} \\
\midrule\noalign{}
\endhead
\bottomrule\noalign{}
\endlastfoot
\textbf{操作} & 外微分 \(d\) & ナブラ \(\nabla\) \\
\textbf{変化} & 測定器の次数が変わる &
同じ形式的な微分演算子から新たな場が生まれる \\
\textbf{結果} & \(1\)-form、\(2\)-form、\(3\)-form などの測定器 &
スカラー場や縦ベクトル場 \\
\end{longtable}

積分での一致を支えるのが、計量 \(g\) と \(\ast\)
による翻訳辞書である。\(\ast\ast=\mathrm{id}\) により往復は自由。
\end{quote}

\begin{center}\rule{0.5\linewidth}{0.5pt}\end{center}

\hypertarget{ux66f2ux7ddaux5ea7ux6a19ux3068ux4e8cux3064ux306eux65b9ux6cd5}{%
\subsection{§8.6
曲線座標と二つの方法}\label{ux66f2ux7ddaux5ea7ux6a19ux3068ux4e8cux3064ux306eux65b9ux6cd5}}

ここからは第9章の実践編への入口である。本格的な計算練習は次章で行うが、その前に、辞書がデカルト座標を離れても機械的に作れることを一度だけ確認しておく。

§8.2 の辞書はデカルト座標 \((x,y,z)\) で \(\mathbf{g}=I\)
を前提としていた。一般の座標系では計量 \(g = J^T J\)
が単位行列ではなくなり、\(\ast\)
の辞書も変わる。ここでは円柱座標を例に、辞書の変化を確認したあと、§8.5
で整理した二つの言語の違いを目の前で回してみる。

\hypertarget{ux8f9eux66f8ux306fux8a08ux91cfux3067ux5909ux308fux308b}{%
\subsubsection{8.6.1
辞書は計量で変わる}\label{ux8f9eux66f8ux306fux8a08ux91cfux3067ux5909ux308fux308b}}

第6章 §6.1.3 で円柱座標の計量を導出した。

\[\mathbf{g} = \begin{pmatrix} 1 & 0 & 0 \\ 0 & r^2 & 0 \\ 0 & 0 & 1 \end{pmatrix}\]

この \(\mathbf{g} \neq I\) が \(\ast\)
の辞書にどう影響するか。デカルト座標では \(\ast(dx) = dy \wedge dz\)
だった。円柱座標の基底 \(1\)-form を \((dr, d\theta, dz)\)
とすると、\(\ast\) の辞書は次のように変わる。

\[\begin{aligned}
\ast(dr) &= r\,d\theta \wedge dz \\
\ast(d\theta) &= \frac{1}{r}\,dz \wedge dr \\
\ast(dz) &= r\,dr \wedge d\theta
\end{aligned}\]

\(r\) や \(1/r\) という係数は \(\mathbf{g} = J^T J\) の対角成分
\(1, r^2, 1\) から決まる。\(\ast\)
の辞書は、その場所の計量を反映して場所ごとに異なる係数を持つ。

\hypertarget{ux4e8cux3064ux306eux65b9ux6cd5ux3092ux56deux3057ux3066ux307fux308b}{%
\subsubsection{8.6.2
二つの方法を回してみる}\label{ux4e8cux3064ux306eux65b9ux6cd5ux3092ux56deux3057ux3066ux307fux308b}}

辞書が変わるということは、§8.1
で対比した二つの方法の違いが具体的な計算に現れるということだ。ここで、同じ問題を二つの方法で実際に解いてみよう。題材は「二次元平面上のある点での発散を求める」である。

\begin{quote}
\textbf{注} （二次元で考える）第1章以来、二次元の問題を扱うときは
\(z=0\)
と断ってきた。しかし本章まで読み進めた読者に、もはやそのような断りは不要だろう。以下では二次元平面
\((x,y)\) とその極座標 \((r,\theta)\) だけで話を進める。
\end{quote}

\textbf{方法①——場はそのまま、測定器を変える。}

パラメータ空間 \((r,\theta)\) と実空間
\((x,y)\)、二枚のデカルト座標を用意する。パラメータ空間の用紙には、どこも同じ大きさの正方形グリッドが敷き詰められている。この正方形グリッドの上で、我々は測定を行う。

成分が \(F_x,F_y\) であるベクトル場 \(\mathbf{F}\)
は実空間に住んでいる。その測定器 \(\omega = F_x\,dx + F_y\,dy\)
を、パラメータ空間の用紙に\textbf{引き戻す}。引き戻し（pullback）とは、実空間の測定器をパラメータ空間の言葉に焼き直す操作だ。写像
\(\phi: (r,\theta) \to (x,y)\)
で実空間に送るのと逆向きに、測定器を戻してくる。\(dx = \cos\theta\,dr - r\sin\theta\,d\theta\),
\(dy = \sin\theta\,dr + r\cos\theta\,d\theta\) を使って、

\[\tilde{\omega} = F_r\,dr + (r F_\theta)\,d\theta\]

が得られる。

\begin{quote}
\textbf{注} （正規直交成分と form 係数）ベクトル解析では \(\mathbf{F}\)
を「長さ1の矢印」に対する成分 \((F_r, F_\theta)\)
で表す。しかし微分形式の自然な基底は座標微分 \((dr, d\theta)\)
であり、\(d\theta\) 方向の目盛り間隔は場所によって \(r\)
倍異なる。以後この節で \(F_\theta\)
と書くときは、ベクトル解析で使う正規直交方向の成分を意味する。\(d\theta\)
の係数そのものではない。\(\tilde{\omega} = F_r\,dr + \tilde{F}_\theta\,d\theta\)
と書いたときの form 係数 \(\tilde{F}_\theta\) は \(r F_\theta\)
であり、ベクトル解析の成分 \(F_\theta\)
とは次元も値も異なる。この違いが、方法①と方法②の公式の見た目の差を生んでいる。
\end{quote}

ここから先は \(\ast d\ast\) の機械的な計算だ。極座標の計量
\(\mathbf{g} = \begin{pmatrix} 1 & 0 \\ 0 & r^2 \end{pmatrix}\) から
\(\ast\) の辞書は \(\ast(dr) = r\,d\theta\),
\(\ast(d\theta) = -\frac{1}{r}\,dr\) となる（§8.6.1 の二次元版）。\(d\)
は偏微分の係数をそのまま拾うだけ------途中に \(\frac{1}{r}\)
などを挟み込む必要は一切ない。

\[\begin{aligned}
\ast\tilde{\omega} &= F_r \cdot r\,d\theta + (r F_\theta) \cdot \left(-\frac{1}{r}\,dr\right) = r F_r\,d\theta - F_\theta\,dr \\[0.3em]
d\ast\tilde{\omega} &= \left(\frac{\partial}{\partial r}(r F_r) + \frac{\partial F_\theta}{\partial \theta}\right) dr \wedge d\theta \\[0.3em]
\ast d\ast\tilde{\omega} &= \frac{1}{r}\frac{\partial}{\partial r}(r F_r) + \frac{1}{r}\frac{\partial F_\theta}{\partial \theta}
\end{aligned}\]

\(\frac{1}{r}\) が現れたのは最後の
\(\ast(dr\wedge d\theta) = \frac{1}{r}\) の一箇所だけだ。途中の \(r\) と
\(\frac{1}{r}\) は、\(d\theta\) の係数に付いていた \(r\) と
\(\ast(d\theta)\) の \(\frac{1}{r}\)
が打ち消し合った結果である。\(\ast\) が計量の情報を一手に引き受け、\(d\)
は純粋に偏微分だけを担当する------この分業こそが方法①の最大の強みだ。

測るグリッドの形は均一なまま考えられる。場所ごとの換算係数の違いは
pullback が織り込み、\(\ast\)
が整理する。極限操作は最後にまとめて行えばよい。

\textbf{方法②——同じ演算子から新たな場を考える。}

今度は実空間に一枚のグラフ用紙を使う。極座標 \((r,\theta)\) の上に、点
\((r,\theta)\) を中心とする小さな面積要素を直接描く。\(r\) 方向に
\(\Delta r\)、\(\theta\) 方向に \(\Delta\theta\)
の幅を持った小片だ。ここでは\textbf{$\Delta r, \Delta\theta$ は微小でなければならない}------角度があるために、有限の
\(\Delta\theta\)
では動径方向の辺が平行にならず、単純な四角形にならないからだ。

この小片は扇形の一部である。\(\theta\) 方向の内側の辺の長さは
\(r\Delta\theta\)、外側の辺の長さは \((r+\Delta r)\Delta\theta\)
で、\(r\) が大きいほど辺が長い。成分が \(F_r,F_\theta\) であるベクトル場
\(\mathbf{F}\) がこの小片の四辺を通り抜ける流束を、一辺ずつ勘定する。

\(r\) 方向：内側の辺から \(-F_r(r,\theta) \cdot r\Delta\theta\)
が流入し、外側の辺から
\(+F_r(r+\Delta r,\theta) \cdot (r+\Delta r)\Delta\theta\)
が流出する。差を \(\Delta r\) で割り、さらに面積
\(r\Delta r\Delta\theta\) で割って極限をとれば
\(\frac{1}{r}\frac{\partial}{\partial r}(rF_r)\) が現れる。

\(\theta\) 方向：\(\theta\) 側の辺から
\(-F_\theta(r,\theta) \cdot \Delta r\) が流入し、\(\theta+\Delta\theta\)
側の辺から \(+F_\theta(r,\theta+\Delta\theta) \cdot \Delta r\)
が流出する。同様に処理すれば
\(\frac{1}{r}\frac{\partial F_\theta}{\partial\theta}\) が現れる。

\[\nabla \cdot \mathbf{F} = \frac{1}{r}\frac{\partial}{\partial r}(r F_r) + \frac{1}{r}\frac{\partial F_\theta}{\partial \theta}\]

実空間に直接小片を描き、壁を通る流れを数え上げる------絵は直感的だ。しかし座標系が変わるたびに小片の形が変わり、毎回異なる勘定を強いられる。\(\frac{1}{r}\frac{\partial}{\partial r}(rF_r)\)
の \(\frac{1}{r}\) は扇形の面積 \(r\Delta r\Delta\theta\)
で割るために出てくるが、デカルト座標で同じ計算をしたときには現れなかった項だ。

\textbf{どちらが優れているというわけではない}

方法①は均一な正方形グリッドで測り、方法②は扇形の小片を描いて流れを数える。やっていることはまったく違う。しかし出てくる答えは同じ
\(\frac{1}{r}\frac{\partial}{\partial r}(rF_r) + \frac{1}{r}\frac{\partial F_\theta}{\partial\theta}\)
だ。そして、この一致は発散に限らない。勾配も回転もラプラシアンも、すべての微分演算がこの二つの方法で構成できる。単純な座標変換なら方法②のほうが早いこともあるが、座標系が歪んだり場の種類が入り組んだりしたら方法①に切り替えればよい。

\hypertarget{ux66f2ux7ddaux5ea7ux6a19ux3067ux306eux767aux6563ux3068ux56deux8ee2}{%
\subsubsection{8.6.3
曲線座標での発散と回転}\label{ux66f2ux7ddaux5ea7ux6a19ux3067ux306eux767aux6563ux3068ux56deux8ee2}}

§8.6.1 の辞書と §8.6.2
の方法①を組み合わせれば、円柱座標での発散と回転の公式が機械的に導出できる。ここでの
\(F_r,F_\theta,F_z\) は、円柱座標の正規直交成分である。結果だけを示す。

\[\nabla \cdot \mathbf{F} = \frac{1}{r}\frac{\partial}{\partial r}(r F_r) + \frac{1}{r}\frac{\partial F_\theta}{\partial \theta} + \frac{\partial F_z}{\partial z}\]

\[\nabla \times \mathbf{F} = \begin{pmatrix}
\frac{1}{r}\frac{\partial F_z}{\partial \theta} - \frac{\partial F_\theta}{\partial z} \\[0.3em]
\frac{\partial F_r}{\partial z} - \frac{\partial F_z}{\partial r} \\[0.3em]
\frac{1}{r}\frac{\partial}{\partial r}(r F_\theta) - \frac{1}{r}\frac{\partial F_r}{\partial \theta}
\end{pmatrix}\]

これらの公式を暗記するのは骨が折れる。しかし \(\ast d\ast\omega\) という
\(d\) と \(\ast\) の組み合わせで考えれば、\(g = J^T J\) から \(\ast\)
の辞書をその場で導出し、機械的に計算できる。第9章では、この「辞書をその場で作る」技術を実践する。

\begin{quote}
\textbf{【ここまでのチェックポイント — 第8章全体】}

\begin{itemize}
\item
  \(d\) と \(\nabla\)
  は根本的に異なる世界観に立つが、積分では一致する。\(d\)
  は測定器の次数を変え、\(\nabla\)
  は同じ形式的な微分演算子からスカラー場や縦ベクトル場を作る。
\item
  \(\nabla f \leftrightarrow df\),
  \(\nabla\times\mathbf{F} \leftrightarrow \ast d\omega\),
  \(\nabla\cdot\mathbf{F} \leftrightarrow \ast d\ast\omega\)
  の辞書により、両者は自由に翻訳できる。
\item
  \(\int_{\partial M}\omega = \int_M d\omega\)
  一本で、微積分学の基本定理・ストークス・ガウスのすべてを統合できる。
\item
  \(\ast\) の辞書は計量 \(g = J^T J\) に依存し、曲線座標では \(r\) や
  \(1/r\) の係数が現れる。
\item
  重要なのは、一つの言語に固執せず、問題に応じて二つの言語を使い分けること。そのための辞書は、もう手に入れた。第9章では、この辞書をその場で作り、実際の問題に適用する。
\end{itemize}
\end{quote}

\newpage

\hypertarget{ux7b2c9ux7ae0ux5b9fux6226-ux8f9eux66f8ux3092ux4f5cux308aux96e3ux554fux3092ux89e3ux304f}{%
\chapter{第9章：実戦 ------
辞書を作り、難問を解く}\label{ux7b2c9ux7ae0ux5b9fux6226-ux8f9eux66f8ux3092ux4f5cux308aux96e3ux554fux3092ux89e3ux304f}}

\hypertarget{ux672cux66f8ux306eux4e2dux5fc3ux7684ux306aux9053ux5177ux306fux63c3ux3063ux305f}{%
\subsection{§9.0
本書の中心的な道具は揃った}\label{ux672cux66f8ux306eux4e2dux5fc3ux7684ux306aux9053ux5177ux306fux63c3ux3063ux305f}}

　第8章 §8.0
で「本章はこの本のハイライトである」と述べ、実際に二つの言語のあいだを翻訳し、積分定理を一本化し、二つの方法論を対比した。本書の中心的な道具立ては、ここまででほぼ揃っている。

しかし、これだけの道具を揃えて、使わずに終わるのはもったいない。本章では、ここまでに構築した辞書を実際の問題に適用し、ベクトル解析の難問を機械的に解いてみせる。理論は第8章までで十分だ------ここから先は\textbf{演習}である。

読者は、以下の計算を眺めるだけでよい。あるいは、自分の手で追ってみてもよい。いずれにせよ、ここで見るのは「公式を暗記しなくても、辞書さえあればすべて導出できる」という事実の実演だ。

\begin{quote}
\textbf{注}
（ベクトル解析を学んだことがない読者へ）以下の計算のうち、grad, div, rot
の導出はどれも数行で済んでいる。\(\ast d\ast\)
に辞書を代入し、偏微分を展開する------それだけだ。だから簡単そうに見える。しかし、これと同じ結果をベクトル解析の流儀で導出しようとすると、\(\nabla\)
に \(\frac{1}{r}\frac{\partial}{\partial r}(r\,\cdot\,)\) や
\(\frac{1}{\rho^2\sin\theta}\)
といった因子を手作業で管理しながら、数ページにわたる式変形を強いられる。簡単そうに見えるのは、次数の異なる測定器を
\(\ast\) が自動的に整理してくれるからだ。§9.4
のラプラシアンはさすがに手強いが、それでもやっていることは辞書を引いて偏微分を展開するだけ------要領は変わらない。この差こそ、本書が第1章から積み上げてきた「測定器」の枠組みの真価である。
\end{quote}

\begin{center}\rule{0.5\linewidth}{0.5pt}\end{center}

\hypertarget{ux8f9eux66f8ux3092ux305dux306eux5834ux3067ux4f5cux308bux6a5fux68b0ux7684ux624bux9806}{%
\subsection{§9.1
辞書をその場で作る------機械的手順}\label{ux8f9eux66f8ux3092ux305dux306eux5834ux3067ux4f5cux308bux6a5fux68b0ux7684ux624bux9806}}

直交曲線座標で \(\mathrm{grad}, \mathrm{div}, \mathrm{rot}\)
の公式を導出する手順は、以下の3ステップに集約される。

\begin{enumerate}
\def\labelenumi{\arabic{enumi}.}
\item
  \textbf{ヤコビ行列 $J$ を書く}。座標変換 \(x = x(q_1,q_2,q_3)\),
  \(y = y(q_1,q_2,q_3)\), \(z = z(q_1,q_2,q_3)\)
  から偏微分を並べる。\(J\)
  の各列はパラメータ空間の各軸方向の一歩が実空間でどう動くかを表す。
\item
  \textbf{計量 $\mathbf{g} = J^T J$ を計算する}。これは第6章 §6.1.3
  以来の手順だ。\(\mathbf{g}\)
  の対角成分が各軸方向の目盛りの二乗、非対角成分が軸どうしの直交度を表す。
\item
  \textbf{$\mathbf{g}$ から $\ast$ の辞書を導き、$d$ と組み合わせる}。\(\mathrm{grad} = d\)、\(\mathrm{rot}=\ast d\)、\(\mathrm{div}=\ast d\ast\)
  の各公式に、その座標系での \(\ast\)
  の辞書を代入すれば、求める公式が機械的に得られる。
\end{enumerate}

ただし、通常のベクトル解析の \(\mathrm{grad}\) や \(\mathrm{rot}\)
は矢印としてのベクトル場を返す。一方、本書の
\(\mathrm{grad}=d\)、\(\mathrm{rot}=\ast d\)、\(\mathrm{div}=\ast d\ast\)
は、まず測定器側の式として読む。微分形式側で得られた \(1\)-form
を縦ベクトルとして読むには、まず計量 \(\mathbf{g}\) を使って
\(\mathbf{g}^{-1}(\cdot)^T\)
の形に戻す。これは座標に対応する縦成分である。ベクトル解析で使う正規直交成分として読むには、さらに各方向のスケール因子を考慮する。

また、曲線座標では、ベクトル解析で使う正規直交成分と、座標に対応する
\(1\)-form を基準にした係数は一般に一致しない。計算は \(1\)-form
側で行うため、まずベクトル場を対応する \(1\)-form
に直す。そのとき、\(d\theta\) や \(d\phi\)
の係数にはスケール因子が入る。

\begin{quote}
\textbf{注}（ベクトル成分と \(1\)-form 係数）
曲線座標では、座標に対応する \(1\)-form と、実際の長さを1単位として測る
\(1\)-form を区別する必要がある。たとえば極座標では、角度の目盛り
\(d\theta\) そのものと、長さ1の目盛り \(r\,d\theta\)
は違う。ホッジ・スターの辞書は、このスケール因子の情報をすべて引き受けている。これにより、我々は「歪んだ座標系」の上で「真っ直ぐな計算」を行うことができるのである。
\end{quote}

以下、この手順を円柱座標と球座標で実演する。まず \(1\)-form や
\(2\)-form
として計算し、最後に正規直交ベクトル成分へ戻す流れを意識してほしい。

\begin{center}\rule{0.5\linewidth}{0.5pt}\end{center}

\hypertarget{ux5186ux67f1ux5ea7ux6a19-rthetaz}{%
\subsection{\texorpdfstring{§9.2 円柱座標
\((r,\theta,z)\)}{§9.2 円柱座標 (r,\textbackslash theta,z)}}\label{ux5186ux67f1ux5ea7ux6a19-rthetaz}}

\hypertarget{ux8f9eux66f8ux306eux5c0eux51fa}{%
\subsubsection{9.2.1 辞書の導出}\label{ux8f9eux66f8ux306eux5c0eux51fa}}

円柱座標への変換は \(x = r\cos\theta\), \(y = r\sin\theta\), \(z = z\)。

\begin{quote}
\textbf{注}（定義域と特異点） 以下では \(r > 0\)
の座標パッチで計算する。\(r = 0\) では \(\theta\)
が定まらず、円柱座標そのものが特異になる。また \(\theta\) は \(2\pi\)
周期を持つ角度座標なので、厳密には局所的な座標（例えば
\(0 < \theta < 2\pi\)）として扱う。本節の公式は、この特異軸を含まない領域で用いることを前提としている。
\end{quote}

\[J = \begin{pmatrix}
\cos\theta & -r\sin\theta & 0 \\
\sin\theta & r\cos\theta & 0 \\
0 & 0 & 1
\end{pmatrix}, \qquad
\mathbf{g} = J^T J = \begin{pmatrix}
1 & 0 & 0 \\
0 & r^2 & 0 \\
0 & 0 & 1
\end{pmatrix}\]

\(\mathbf{g}\) から \(\ast\) の辞書を導く。第6章で見たデカルト座標の
\(\ast\) 辞書と同じ原理で、\(\mathbf{g}\) の対角成分の平方根が \(\ast\)
の係数に現れる。

\[\begin{aligned}
\ast(dr) &= r\,d\theta \wedge dz \\
\ast(d\theta) &= \frac{1}{r}\,dz \wedge dr \\
\ast(dz) &= r\,dr \wedge d\theta
\end{aligned}\]

また \(\ast(1) = r\,dr\wedge d\theta\wedge dz\),
\(\ast(dr\wedge d\theta\wedge dz) = 1/r\)。

\begin{quote}
\textbf{注}（円柱座標で区別すべき4つのもの）
ベクトル場を正しく微分形式として扱うために、以下の 4
つのレイヤーを区別しよう。

\begin{longtable}[]{@{}ll@{}}
\toprule\noalign{}
名称 & 円柱座標での具体例 \\
\midrule\noalign{}
\endhead
\bottomrule\noalign{}
\endlastfoot
\textbf{座標に対応する $1$-form} & \(dr, d\theta, dz\) \\
\textbf{長さ1の目盛りに対応する $1$-form} & \(dr, r\,d\theta, dz\) \\
\textbf{ベクトル解析で使う成分} & \(F_r, F_\theta, F_z\) \\
\textbf{$1$-form としての係数} & \(F_r, rF_\theta, F_z\) \\
\end{longtable}

ベクトル解析で使う単位ベクトル方向の成分を \(F_r,F_\theta,F_z\)
と書くなら、対応する \(1\)-form は
\(F_r dr + F_\theta (r d\theta) + F_z dz\)、すなわち \(1\)-form
としての係数は \(F_r,rF_\theta,F_z\) となる。本章の \(F_\theta\)
はつねに「ベクトル解析で使う成分」を指す。
\end{quote}

\hypertarget{ux52feux914d-2}{%
\subsubsection{9.2.2 勾配}\label{ux52feux914d-2}}

\(\mathrm{grad}\,f = df\) は \(\ast\) を使わない。単に偏微分係数を基底
\(1\)-form に載せるだけだ。

\[df = \frac{\partial f}{\partial r}\,dr + \frac{\partial f}{\partial \theta}\,d\theta + \frac{\partial f}{\partial z}\,dz\]

ベクトル解析の表記（縦ベクトル）では、

\[
\nabla f
=
\begin{pmatrix}
\frac{\partial f}{\partial r}\\[0.4em]
\frac{1}{r}\frac{\partial f}{\partial \theta}\\[0.4em]
\frac{\partial f}{\partial z}
\end{pmatrix}
\]

\(\theta\) 成分に \(\frac{1}{r}\) がつくのは、\(\nabla f\)
が正規直交基底に対する成分だからだ。\(df\) の \(d\theta\) の係数
\(\frac{\partial f}{\partial\theta}\)
は「角度あたりの変化率」だが、\(\nabla f\) の \(\mathbf{e}_\theta\)
成分は「長さあたりの変化率」であり、\(d\theta\) が弧長 \(r\,d\theta\)
に対応するための換算 \(\frac{1}{r}\) が必要になる。

\hypertarget{ux767aux6563-2}{%
\subsubsection{9.2.3 発散}\label{ux767aux6563-2}}

正規直交成分が \(F_r,F_\theta,F_z\) であるベクトル場 \(\mathbf{F}\)
に対応する \(1\)-form は、以下に述べる「式に現れるスケール因子によって
\(1\)-form
としての係数を明示する」という約束に基づき、\(\omega = F_r\,dr + rF_\theta\,d\theta + F_z\,dz\)
である。

\begin{quote}
\textbf{注} （\(\tilde{\omega}\) のチルダについて） 第8章 §8.6.2
では、正規直交成分 \(F_\theta\) と \(1\)-form としての係数 \(rF_\theta\)
を区別するために \(\tilde{\omega}\)
という記法を使った。本章では記号を軽くするため、以下ではチルダを省いて
\(\omega\) と書く。ただし、\(F_r,F_\theta,F_z\)
は常に正規直交ベクトル成分を表すものとする。対応する \(1\)-form
は、その都度スケール因子を含めて明示する。円柱座標では \[
\omega = F_r\,dr + rF_\theta\,d\theta + F_z\,dz
\] であり、球座標では \[
\omega = F_\rho\,d\rho + \rho F_\theta\,d\theta + \rho\sin\theta\,F_\phi\,d\phi
\]
である。したがって、読者が文脈から推測するのではなく、式に現れるスケール因子によって
\(1\)-form としての係数を明示する、という約束で読む。
\end{quote}

\[\begin{aligned}
\ast\omega &= F_r(r\,d\theta\wedge dz) + rF_\theta\!\left(\frac{1}{r}\,dz\wedge dr\right) + F_z(r\,dr\wedge d\theta) \\
&= rF_r\,d\theta\wedge dz + F_\theta\,dz\wedge dr + rF_z\,dr\wedge d\theta \\[0.3em]
d\ast\omega &= \left(\frac{\partial}{\partial r}(rF_r) + \frac{\partial F_\theta}{\partial\theta} + \frac{\partial}{\partial z}(rF_z)\right) dr\wedge d\theta\wedge dz \\[0.3em]
\ast d\ast\omega &= \frac{1}{r}\frac{\partial}{\partial r}(rF_r) + \frac{1}{r}\frac{\partial F_\theta}{\partial\theta} + \frac{\partial F_z}{\partial z}
\end{aligned}\]

これが円柱座標での発散の公式である。

\hypertarget{ux56deux8ee2-2}{%
\subsubsection{9.2.4 回転}\label{ux56deux8ee2-2}}

回転は \(\mathrm{rot}\,\mathbf{F} \leftrightarrow \ast d\omega\)。

\[d\omega = \left(\frac{\partial}{\partial r}(rF_\theta) - \frac{\partial F_r}{\partial\theta}\right) dr\wedge d\theta + \left(\frac{\partial F_z}{\partial\theta} - \frac{\partial}{\partial z}(rF_\theta)\right) d\theta\wedge dz + \left(\frac{\partial F_r}{\partial z} - \frac{\partial F_z}{\partial r}\right) dz\wedge dr\]

\[\ast d\omega = \frac{1}{r}\!\left(\frac{\partial F_z}{\partial\theta} - \frac{\partial}{\partial z}(rF_\theta)\right)\! dr + r\!\left(\frac{\partial F_r}{\partial z} - \frac{\partial F_z}{\partial r}\right)\! d\theta + \frac{1}{r}\!\left(\frac{\partial}{\partial r}(rF_\theta) - \frac{\partial F_r}{\partial\theta}\right)\! dz\]

\(1\)-form の \(dr, d\theta, dz\)
の係数を読み取り、正規直交成分に焼き直すと \(d\theta\) 係数はさらに
\(\frac{1}{r}\) される。最終的に、

\[\nabla \times \mathbf{F} = \begin{pmatrix}
\frac{1}{r}\frac{\partial F_z}{\partial \theta} - \frac{\partial F_\theta}{\partial z} \\[0.3em]
\frac{\partial F_r}{\partial z} - \frac{\partial F_z}{\partial r} \\[0.3em]
\frac{1}{r}\frac{\partial}{\partial r}(r F_\theta) - \frac{1}{r}\frac{\partial F_r}{\partial \theta}
\end{pmatrix}\]

\(r\) が \(\frac{\partial}{\partial r}\)
の内側と外側に現れる構造は、前章で扇形の面積要素を数えたときの
\(\frac{1}{r}\) とまったく同じ起源を持つ。

\begin{center}\rule{0.5\linewidth}{0.5pt}\end{center}

\hypertarget{ux7403ux5ea7ux6a19-rhothetaphi}{%
\subsection{\texorpdfstring{§9.3 球座標
\((\rho,\theta,\phi)\)}{§9.3 球座標 (\textbackslash rho,\textbackslash theta,\textbackslash phi)}}\label{ux7403ux5ea7ux6a19-rhothetaphi}}

\hypertarget{ux8f9eux66f8ux306eux5c0eux51fa-1}{%
\subsubsection{9.3.1
辞書の導出}\label{ux8f9eux66f8ux306eux5c0eux51fa-1}}

球座標への変換は \(x = \rho\sin\theta\cos\phi\),
\(y = \rho\sin\theta\sin\phi\), \(z = \rho\cos\theta\)（\(\theta\) は
\(z\) 軸からの角度、\(\phi\) は \(xy\) 平面内の方位角）。

\[J = \begin{pmatrix}
\sin\theta\cos\phi & \rho\cos\theta\cos\phi & -\rho\sin\theta\sin\phi \\
\sin\theta\sin\phi & \rho\cos\theta\sin\phi & \rho\sin\theta\cos\phi \\
\cos\theta & -\rho\sin\theta & 0
\end{pmatrix},\qquad
\mathbf{g} = J^T J = \begin{pmatrix}
1 & 0 & 0 \\
0 & \rho^2 & 0 \\
0 & 0 & \rho^2\sin^2\theta
\end{pmatrix}\]

\(\ast\) の辞書は対角成分
\(\sqrt{1}, \sqrt{\rho^2}, \sqrt{\rho^2\sin^2\theta}\) すなわち
\(1, \rho, \rho\sin\theta\) から決まる。

\[\begin{aligned}
\ast(d\rho) &= \rho^2\sin\theta\; d\theta \wedge d\phi \\
\ast(d\theta) &= \sin\theta\; d\phi \wedge d\rho \\
\ast(d\phi) &= \frac{1}{\sin\theta}\; d\rho \wedge d\theta
\end{aligned}\]

\begin{quote}
\textbf{注}（球座標で区別すべき4つのもの）

\begin{longtable}[]{@{}
  >{\raggedright\arraybackslash}p{(\columnwidth - 2\tabcolsep) * \real{0.5000}}
  >{\raggedright\arraybackslash}p{(\columnwidth - 2\tabcolsep) * \real{0.5000}}@{}}
\toprule\noalign{}
\begin{minipage}[b]{\linewidth}\raggedright
名称
\end{minipage} & \begin{minipage}[b]{\linewidth}\raggedright
球座標での具体例
\end{minipage} \\
\midrule\noalign{}
\endhead
\bottomrule\noalign{}
\endlastfoot
\textbf{座標に対応する $1$-form} & \(d\rho, d\theta, d\phi\) \\
\textbf{長さ1の目盛りに対応する $1$-form} &
\(d\rho, \rho\,d\theta, \rho\sin\theta\,d\phi\) \\
\textbf{ベクトル解析で使う成分} & \(F_\rho, F_\theta, F_\phi\) \\
\textbf{$1$-form としての係数} &
\(F_\rho, \rho F_\theta, \rho\sin\theta F_\phi\) \\
\end{longtable}

正規直交成分 \(F_\rho,F_\theta,F_\phi\) に対する \(1\)-form は
\(\omega = F_\rho d\rho + \rho F_\theta d\theta + \rho\sin\theta F_\phi d\phi\)
となる。
\end{quote}

ここでは \(\rho>0\) かつ \(0<\theta<\pi\)
の座標パッチで計算している。原点 \(\rho=0\) と極軸 \(\sin\theta=0\)
では球座標そのものが特異になり、\(1/\rho\) や \(1/\sin\theta\)
を含む公式をそのまま使うことはできない。特異点をまたぐ問題では、別の座標パッチを使うか、特異点を除いた領域で計算してから境界条件として扱う。

\hypertarget{ux767aux6563-3}{%
\subsubsection{9.3.2 発散}\label{ux767aux6563-3}}

正規直交成分が \(F_\rho,F_\theta,F_\phi\) であるベクトル場
\(\mathbf{F}\) の \(1\)-form への変換では、\(d\rho\) の係数は
\(F_\rho\)（\(\rho\) はすでに長さ）、\(d\theta\) の係数は
\(\rho F_\theta\)、\(d\phi\) の係数は \(\rho\sin\theta\,F_\phi\)
となる（それぞれ弧長 \(\rho\,d\theta\), \(\rho\sin\theta\,d\phi\)
に対応）。

\[\omega = F_\rho\,d\rho + \rho F_\theta\,d\theta + \rho\sin\theta\,F_\phi\,d\phi\]

\(\ast d\ast\omega\) を計算する。

\[\ast\omega = \rho^2\sin\theta\,F_\rho\,d\theta\wedge d\phi + \rho\sin\theta\,F_\theta\,d\phi\wedge d\rho + \rho\,F_\phi\,d\rho\wedge d\theta\]

\[d\ast\omega = \left(\frac{\partial}{\partial\rho}(\rho^2\sin\theta\,F_\rho) + \frac{\partial}{\partial\theta}(\rho\sin\theta\,F_\theta) + \frac{\partial}{\partial\phi}(\rho\,F_\phi)\right) d\rho\wedge d\theta\wedge d\phi\]

\(\ast(d\rho\wedge d\theta\wedge d\phi) = 1/(\rho^2\sin\theta)\) より、

\[\nabla \cdot \mathbf{F} = \frac{1}{\rho^2}\frac{\partial}{\partial\rho}(\rho^2 F_\rho) + \frac{1}{\rho\sin\theta}\frac{\partial}{\partial\theta}(\sin\theta\,F_\theta) + \frac{1}{\rho\sin\theta}\frac{\partial F_\phi}{\partial\phi}\]

\hypertarget{ux56deux8ee2-3}{%
\subsubsection{9.3.3 回転}\label{ux56deux8ee2-3}}

同様に \(\ast d\omega\) から、

\[\nabla \times \mathbf{F} = \begin{pmatrix}
\frac{1}{\rho\sin\theta}\frac{\partial}{\partial\theta}(\sin\theta\,F_\phi) - \frac{1}{\rho\sin\theta}\frac{\partial F_\theta}{\partial\phi} \\[0.3em]
\frac{1}{\rho\sin\theta}\frac{\partial F_\rho}{\partial\phi} - \frac{1}{\rho}\frac{\partial}{\partial\rho}(\rho F_\phi) \\[0.3em]
\frac{1}{\rho}\frac{\partial}{\partial\rho}(\rho F_\theta) - \frac{1}{\rho}\frac{\partial F_\rho}{\partial\theta}
\end{pmatrix}\]

\begin{center}\rule{0.5\linewidth}{0.5pt}\end{center}

\hypertarget{ux5faeux7a4dux5206ux306eux96e3ux554fux7403ux5ea7ux6a19ux3067ux306eux30d9ux30afux30c8ux30ebux30e9ux30d7ux30e9ux30b7ux30a2ux30f3}{%
\subsection{§9.4
微積分の難問------球座標でのベクトルラプラシアン}\label{ux5faeux7a4dux5206ux306eux96e3ux554fux7403ux5ea7ux6a19ux3067ux306eux30d9ux30afux30c8ux30ebux30e9ux30d7ux30e9ux30b7ux30a2ux30f3}}

ベクトルラプラシアン
\(\nabla^2\mathbf{F} = \nabla(\nabla\cdot\mathbf{F}) - \nabla\times(\nabla\times\mathbf{F})\)
は、ベクトル解析の公式集で最も長い式の一つとして知られる。ここでは全成分が辞書から機械的に出ることを示す。

方針は第8章の辞書をそのまま使うだけだ。

\begin{itemize}
\item
  \(\nabla(\nabla\cdot\mathbf{F}) \leftrightarrow d(\ast d\ast\omega)\)
\item
  \(\nabla\times(\nabla\times\mathbf{F}) \leftrightarrow \ast d(\ast d\omega)\)
\end{itemize}

したがって、結果を \(1\)-form として見るなら、

\[\nabla^2\mathbf{F} \leftrightarrow d\ast d\ast\omega - \ast d\ast d\omega\]

球座標で
\(\omega = F_\rho\,d\rho + \rho F_\theta\,d\theta + \rho\sin\theta\,F_\phi\,d\phi\)
から出発し、\(d\ast d\ast\omega\) と \(\ast d\ast d\omega\)
を個別に計算して差をとる。

\hypertarget{ux7b2cux4e00ux9805ux52feux914dux306eux767aux6563}{%
\subsubsection{9.4.1
第一項------勾配の発散}\label{ux7b2cux4e00ux9805ux52feux914dux306eux767aux6563}}

§9.3.2 より、発散は

\[\ast d\ast\omega = \frac{1}{\rho^2}\frac{\partial}{\partial\rho}(\rho^2 F_\rho) + \frac{1}{\rho\sin\theta}\frac{\partial}{\partial\theta}(\sin\theta\,F_\theta) + \frac{1}{\rho\sin\theta}\frac{\partial F_\phi}{\partial\phi}\]

これを \(S\) とおく。\(S\) はスカラー場だから、その勾配は
\(dS = \frac{\partial S}{\partial\rho}\,d\rho + \frac{\partial S}{\partial\theta}\,d\theta + \frac{\partial S}{\partial\phi}\,d\phi\)
である。正規直交成分に焼き直すと、

\[\nabla(\nabla\cdot\mathbf{F}) = \begin{pmatrix}
\frac{\partial S}{\partial\rho} \\[0.5em]
\frac{1}{\rho}\frac{\partial S}{\partial\theta} \\[0.5em]
\frac{1}{\rho\sin\theta}\frac{\partial S}{\partial\phi}
\end{pmatrix}\]

\hypertarget{ux7b2cux4e8cux9805ux56deux8ee2ux306eux56deux8ee2}{%
\subsubsection{9.4.2
第二項------回転の回転}\label{ux7b2cux4e8cux9805ux56deux8ee2ux306eux56deux8ee2}}

§9.3.3 の回転の結果を \(C_\rho, C_\theta, C_\phi\) とおく。

\[\begin{aligned}
C_\rho &= \frac{1}{\rho\sin\theta}\frac{\partial}{\partial\theta}(\sin\theta\,F_\phi) - \frac{1}{\rho\sin\theta}\frac{\partial F_\theta}{\partial\phi} \\[0.3em]
C_\theta &= \frac{1}{\rho\sin\theta}\frac{\partial F_\rho}{\partial\phi} - \frac{1}{\rho}\frac{\partial}{\partial\rho}(\rho F_\phi) \\[0.3em]
C_\phi &= \frac{1}{\rho}\frac{\partial}{\partial\rho}(\rho F_\theta) - \frac{1}{\rho}\frac{\partial F_\rho}{\partial\theta}
\end{aligned}\]

\(\mathbf{C} = \nabla\times\mathbf{F}\) に再度 \(\nabla\times\)
を作用させるには、\(C_\rho, C_\theta, C_\phi\) を §9.3.3 の公式の
\(F_\rho, F_\theta, F_\phi\) の位置に代入すればよい。

\[\begin{aligned}
(\nabla\times(\nabla\times\mathbf{F}))_\rho &= \frac{1}{\rho\sin\theta}\frac{\partial}{\partial\theta}(\sin\theta\,C_\phi) - \frac{1}{\rho\sin\theta}\frac{\partial C_\theta}{\partial\phi} \\[0.5em]
(\nabla\times(\nabla\times\mathbf{F}))_\theta &= \frac{1}{\rho\sin\theta}\frac{\partial C_\rho}{\partial\phi} - \frac{1}{\rho}\frac{\partial}{\partial\rho}(\rho C_\phi) \\[0.5em]
(\nabla\times(\nabla\times\mathbf{F}))_\phi &= \frac{1}{\rho}\frac{\partial}{\partial\rho}(\rho C_\theta) - \frac{1}{\rho}\frac{\partial C_\rho}{\partial\theta}
\end{aligned}\]

\hypertarget{ux5deeux3092ux3068ux308bux30d9ux30afux30c8ux30ebux30e9ux30d7ux30e9ux30b7ux30a2ux30f3}{%
\subsubsection{9.4.3
差をとる------ベクトルラプラシアン}\label{ux5deeux3092ux3068ux308bux30d9ux30afux30c8ux30ebux30e9ux30d7ux30e9ux30b7ux30a2ux30f3}}

\(\nabla^2\mathbf{F} = \nabla(\nabla\cdot\mathbf{F}) - \nabla\times(\nabla\times\mathbf{F})\)
だから、各成分は第一項から第二項を引けばよい。\(\rho\) 成分を展開する。

\[(\nabla^2\mathbf{F})_\rho = \frac{\partial S}{\partial\rho} - \frac{1}{\rho\sin\theta}\frac{\partial}{\partial\theta}(\sin\theta\,C_\phi) + \frac{1}{\rho\sin\theta}\frac{\partial C_\theta}{\partial\phi}\]

\(S\) と \(C_\theta, C_\phi\)
の表式を代入し、偏微分を展開して整理すると以下を得る。展開は長いので、ここでは代入後に同類項を整理した結果だけを示す。

\[(\nabla^2\mathbf{F})_\rho = \nabla^2 F_\rho - \frac{2F_\rho}{\rho^2} - \frac{2}{\rho^2\sin\theta}\frac{\partial}{\partial\theta}(\sin\theta\,F_\theta) - \frac{2}{\rho^2\sin\theta}\frac{\partial F_\phi}{\partial\phi}\]

ここで
\(\nabla^2 F_\rho = \frac{1}{\rho^2}\frac{\partial}{\partial\rho}\!\left(\rho^2\frac{\partial F_\rho}{\partial\rho}\right) + \frac{1}{\rho^2\sin\theta}\frac{\partial}{\partial\theta}\!\left(\sin\theta\frac{\partial F_\rho}{\partial\theta}\right) + \frac{1}{\rho^2\sin^2\theta}\frac{\partial^2 F_\rho}{\partial\phi^2}\)
である。

同様に \(\theta\) 成分と \(\phi\) 成分も、

\[(\nabla^2\mathbf{F})_\theta = \nabla^2 F_\theta - \frac{F_\theta}{\rho^2\sin^2\theta} + \frac{2}{\rho^2}\frac{\partial F_\rho}{\partial\theta} - \frac{2\cos\theta}{\rho^2\sin^2\theta}\frac{\partial F_\phi}{\partial\phi}\]

\[(\nabla^2\mathbf{F})_\phi = \nabla^2 F_\phi - \frac{F_\phi}{\rho^2\sin^2\theta} + \frac{2}{\rho^2\sin\theta}\frac{\partial F_\rho}{\partial\phi} + \frac{2\cos\theta}{\rho^2\sin^2\theta}\frac{\partial F_\theta}{\partial\phi}\]

以上で全成分の導出が完了した。ベクトル解析の流儀で一から導出すれば1ページを優に超える計算量だが、辞書を経由すれば各ステップは偏微分と係数の掛け算の繰り返しに過ぎない。個別公式を丸暗記する必要はない------\(d\ast d\ast\omega - \ast d\ast d\omega\)
の一本と、§9.3 の \(\ast\) 辞書さえあれば、その場で再構成できる。

\hypertarget{ux96fbux78c1ux6c17ux5b66ux306eux96e3ux554fux70b9ux96fbux8377ux306eux96fbux5834ux3068ux767aux6563}{%
\subsection{§9.5
電磁気学の難問------点電荷の電場と発散}\label{ux96fbux78c1ux6c17ux5b66ux306eux96e3ux554fux70b9ux96fbux8377ux306eux96fbux5834ux3068ux767aux6563}}

最後に、電磁気学からの例題を一つ。第10章への橋渡しでもある。

\begin{quote}
\textbf{注}
（電磁気学が未習の読者へ）以下に出てくる「電場」や「点電荷」といった用語は、理解できなくても構わない。単に「\(\frac{1}{\rho^2}\)
に比例するベクトル場」の計算だと思って眺めてほしい。第10章でも、電磁気学そのものを深く掘り下げるわけではない。ただ、マクスウェル方程式を微分形式で書くと恐ろしく簡潔になること、そして「みんなここで美しいと言って終わるが、行列まで落とすとどうなるか」を実際に全部書き下す------それだけだ。
\end{quote}

点電荷 \(q\) が原点にあるとき、電場は
\(\mathbf{E} = \frac{q}{4\pi\varepsilon_0}\frac{\hat{\boldsymbol{\rho}}}{\rho^2}\)（\(\hat{\boldsymbol{\rho}}\)
は動径方向の単位ベクトル）。ベクトル解析では
\(\nabla\cdot\mathbf{E} = 0\)（原点以外）を示すために球座標の発散公式に代入するが、電場に対応する
\(1\)-form を直接扱う方が見通しがよい。

\(\mathbf{E}\) の \(\rho\) 成分は
\(E_\rho = \frac{k}{\rho^2}\)（\(k = q/4\pi\varepsilon_0\)）、\(\theta,\phi\)
成分はゼロ。対応する \(1\)-form は

\[\omega = \frac{k}{\rho^2}\,d\rho\]

\(d\ast\omega\) を計算する。球座標では
\(\ast(d\rho) = \rho^2\sin\theta\,d\theta\wedge d\phi\)（§9.3.1）だから、

\[\ast\omega = \frac{k}{\rho^2} \cdot \rho^2\sin\theta\,d\theta\wedge d\phi = k\sin\theta\,d\theta\wedge d\phi\]

\[d\ast\omega = \frac{\partial}{\partial\rho}(k\sin\theta)\,d\rho\wedge d\theta\wedge d\phi + \frac{\partial}{\partial\theta}(k\sin\theta)\,d\theta\wedge d\theta\wedge d\phi + \cdots\]

第二項は \(d\theta\wedge d\theta = 0\) で消える。第三項も同様。第一項は
\(\frac{\partial}{\partial\rho}(k\sin\theta) = 0\)（\(\rho\)
に依存しない）。よって

\[d\ast\omega = 0 \quad (\rho \neq 0)\]

\(\ast d\ast\omega = \nabla\cdot\mathbf{E} = 0\) である。\(\rho^2\) が
\(\ast(d\rho)\) の \(\rho^2\) と打ち消し合い、残った
\(\sin\theta\,d\theta\wedge d\phi\)
の外微分がゼロ------これだけの計算で、原点以外での発散ゼロが示された。

\begin{quote}
\textbf{注} （原点では？）\(\rho = 0\) では \(\omega\)
が定義できない。原点を含む全空間で扱うには、通常の関数としての発散ではなく、分布（\(\delta\)
関数）や積分形式で点電荷を表す必要がある。本書ではそこには立ち入らず、原点を除いた領域での計算に限ることにする。電荷密度を扱う第10章では、動径座標
\(\rho\) と混同しないように、電荷密度を \(\rho_{\mathrm e}\)
と書いて右辺に \(\rho_{\mathrm e}/\varepsilon_0\) が現れる形を扱う。
\end{quote}

\begin{center}\rule{0.5\linewidth}{0.5pt}\end{center}

\hypertarget{ux8f9eux66f8ux306fux7d42ux308fux308aux65c5ux306fux7d9aux304f}{%
\subsection{§9.6
辞書は終わり、旅は続く}\label{ux8f9eux66f8ux306fux7d42ux308fux308aux65c5ux306fux7d9aux304f}}

本章では、第8章までに構築した辞書を使って、ベクトル解析の公式を機械的に導出し、微積分と電磁気学の難問を解いた。重要なのは、どの計算も「公式を思い出す」ではなく「辞書を引いて手順をなぞる」だけで済んだことだ。

もはや、円柱座標や球座標の発散・回転の公式を暗記する必要はない。\(J \to g = J^T J \to \ast\)
辞書 \(\to\) \(d\) と \(\ast\)
の組み合わせ------この一本道を覚えておけば、直交曲線座標の公式をその場で再構成できる。非直交座標でも原理は同じだが、\(\ast\)
の辞書には非対角計量に由来する混合項が現れるため、本章のような単純なスケール因子だけの表にはならないことに注意されたい。

第10章では、この道具をマクスウェル方程式に適用する。
ベクトル解析では4本の基礎方程式が、微分形式では2本に集約される。多くの教科書はここで「なんと美しい」と止まってしまうが、我々はさらに進む。マクスウェル方程式を微分形式で書き、その全成分を実際に展開するとどうなるか------それを実際にやってみせる。本章の計算練習を経ていれば、その翻訳はもはや難しくないはずだ。

\begin{quote}
\textbf{【ここまでのチェックポイント — 第9章全体】}

\begin{itemize}
\item
  直交曲線座標での公式導出は \(J \to g = J^T J \to \ast\) 辞書 \(\to\)
  \(d\) と \(\ast\) の組み合わせ、という手順で進む。
\item
  円柱座標・球座標での \(\ast\) 辞書、発散・回転の公式を導出した。
\item
  ベクトルラプラシアン \(\nabla^2\mathbf{F}\)
  のような難問も、辞書と恒等式の組み合わせで機械的に解ける。
\item
  点電荷の電場の発散ゼロは、\(\ast(d\rho)\) と \(E_\rho\) の \(\rho^2\)
  が打ち消し合うだけの単純な計算で示せる。
\item
  公式は丸暗記しなくてもよい。辞書さえあれば、そしてその辞書すらも、すべてその場で再構成できる。
\end{itemize}
\end{quote}

\newpage

\hypertarget{ux7b2c10ux7ae0ux30deux30afux30b9ux30a6ux30a7ux30ebux65b9ux7a0bux5f0f-ux7f8eux3057ux3055ux306eux305dux306eux5148ux3078}{%
\chapter{第10章：マクスウェル方程式 ------
美しさのその先へ}\label{ux7b2c10ux7ae0ux30deux30afux30b9ux30a6ux30a7ux30ebux65b9ux7a0bux5f0f-ux7f8eux3057ux3055ux306eux305dux306eux5148ux3078}}

\hypertarget{ux304aux7d04ux675f}{%
\subsection{§10.0 お約束}\label{ux304aux7d04ux675f}}

さて、いよいよ本書も終わりが近づいてきた。微分形式の物理数学の教科書というのは、最後にマクスウェル方程式を微分形式で書き直して締めるというのが、どうやらお約束らしい。筆者もこの慣習に倣うことにしよう。

ただし、一言断っておく。本書は電磁気学の教科書ではない。電磁気学を知っていれば意味づけも楽しめるが、知らなくても成分計算として追えるように書く。ここでは、\(E_x, E_y, E_z, B_x, B_y, B_z\)
という6つの関数が登場し、それらがとある反対称行列に収まり、\(d\) と
\(\ast\)
を作用させるとベクトル解析で有名な式が現れる------その流れを、純粋に代数的な操作として眺めてほしい。

\begin{quote}
\textbf{注}
（電磁気学が未習の読者へ）本章に出てくる物理用語はすべて読み飛ばしてよい。必要なのは「6つの関数を4×4行列に並べ、\(d\)
と \(\ast\)
を作用させたら何が出るか」を見ることだけだ。電磁気学の知識がなくても、ここまでの章を読み通した読者なら、行列の成分を追うだけで十分に楽しめるはずである。
\end{quote}

\begin{quote}
\textbf{注}
（本章の立ち位置）本書の核心は第9章までで完結している。本章は、いわばおまけ、あるいは蛇足に近い。あまり真面目に身構えず、気楽に読んでほしい。
\end{quote}

\begin{center}\rule{0.5\linewidth}{0.5pt}\end{center}

\hypertarget{ux30deux30afux30b9ux30a6ux30a7ux30ebux65b9ux7a0bux5f0f2ux672cux3067ux66f8ux304f}{%
\subsection{§10.1
マクスウェル方程式------2本で書く}\label{ux30deux30afux30b9ux30a6ux30a7ux30ebux65b9ux7a0bux5f0f2ux672cux3067ux66f8ux304f}}

電磁気学の基本法則は、ベクトル解析記法で次の4本の式にまとまる。なお、SI単位系を採用する。

\[\begin{aligned}
\mathrm{div}\,\mathbf{E} &= \frac{\rho_{\mathrm e}}{\varepsilon_0} &
\mathrm{div}\,\mathbf{B} &= 0 \\[0.5em]
\mathrm{rot}\,\mathbf{B} &= \mu_0\mathbf{J} + \frac{1}{c^2}\frac{\partial \mathbf{E}}{\partial t} &
\mathrm{rot}\,\mathbf{E} &= -\frac{\partial \mathbf{B}}{\partial t}
\end{aligned}\]

このままでは、これから \(4\times4\)
行列に押し込むための「型」が揃っていない。電場 \(\mathbf{E}\)（次元
\(\mathrm{V/m}\)）と磁束密度 \(\mathbf{B}\)（次元
\(\mathrm{T}\)）は単位が違い、さらに空間の微分（\(\mathrm{rot}\) や
\(\mathrm{div}\)）と時間 \(t\) の微分も次元が異なる。

そこで、すべてを「長さ」の次元に変換する。

\begin{enumerate}
\def\labelenumi{\arabic{enumi}.}
\item
  \textbf{時間の変換}：時間 \(t\) に光速 \(c\)
  を掛けて、長さの次元を持つ座標 \(w = ct\) を導入する。時間微分は
  \(\frac{\partial}{\partial t} = c\frac{\partial}{\partial w}\)
  となる。
\item
  \textbf{磁場の変換}：磁束密度 \(\mathbf{B}\) に \(c\)
  を掛けて、電場と同じ \(\mathrm{V/m}\) の次元を持つ
  \(\mathbf{B}' = c\mathbf{B}\) を定義する。
\end{enumerate}

これらを上の4式に代入する。ファラデーの法則は、

\[\mathrm{rot}\,\mathbf{E} = -\frac{\partial}{\partial t}\!\left(\frac{1}{c}\mathbf{B}'\right) = -c\frac{\partial}{\partial w}\!\left(\frac{1}{c}\mathbf{B}'\right) = -\frac{\partial \mathbf{B}'}{\partial w}\]

アンペールの法則も同様に整理すると、

\[\begin{aligned}
\mathrm{div}\,\mathbf{E} &= \frac{\rho_{\mathrm e}}{\varepsilon_0} &
\mathrm{div}\,\mathbf{B}' &= 0 \\[0.5em]
\mathrm{rot}\,\mathbf{B}' &= c\mu_0\mathbf{J} + \frac{\partial \mathbf{E}}{\partial w} &
\mathrm{rot}\,\mathbf{E} &= -\frac{\partial \mathbf{B}'}{\partial w}
\end{aligned}\]

時間微分が \(\frac{\partial}{\partial w}\)
になり、空間微分（\(\mathrm{rot}, \mathrm{div}\)）と分母の次元（長さ）が揃った。

以降、\(\mathbf{B}'\) のことを改めて \(\mathbf{B}\) と呼び、\(w\)
を改めて \(t\) と書く。つまり、これ以降の \(\mathbf{B}\)
は本来の磁束密度ではなく \(c\) 倍された「次元を揃えた磁場」であり、\(t\)
は本来の時間ではなく \(c\) 倍された「長さの次元を持つ時間座標」である。

これで準備が整った。\(\mathbf{E}\) と \(\mathbf{B}\)
は同じ次元を持ち、同じ行列に並べられる。電磁場をまとめた\(2\)-form \(F\)
と、電流・電荷をまとめた\(1\)-form \(\mathcal{J}\)
を用意すれば、4本の式は本書の正規化・符号規約のもとで\textbf{2本}にまとまる。

\[dF = 0, \qquad d(\ast F) = \mu_0(\ast\mathcal{J})\]

本章の規約では、以下では \(\mu_0\ast\mathcal{J}\)
を成分で直接定義していく。

以上である。4本が2本に------この簡潔さは、微分形式の記述能力を象徴するものだ。

しかし、本書はここで終わらない。\(F\) の中身は何なのか。\(dF=0\)
を成分で展開すると、本当に
\(\mathrm{rot}\,\mathbf{E} = -\frac{\partial\mathbf{B}}{\partial t}\) と
\(\mathrm{div}\,\mathbf{B} = 0\)
が出てくるのか。それを、この章で全部書き下す。

\begin{center}\rule{0.5\linewidth}{0.5pt}\end{center}

\hypertarget{ux96fbux78c1ux5834-f-ux3068ux7b26ux53f7ux898fux7d04ux306eux56faux5b9a}{%
\subsection{\texorpdfstring{§10.2 電磁場 \(F\)
と符号規約の固定}{§10.2 電磁場 F と符号規約の固定}}\label{ux96fbux78c1ux5834-f-ux3068ux7b26ux53f7ux898fux7d04ux306eux56faux5b9a}}

電磁場 \(F\)
を具体的な成分で書き下す前に、本章で用いる時空の符号規約を厳密に固定しておこう。符号の選択には複数の流儀があるが、本書では以下のセットを採用する。

\begin{itemize}
\item
  \textbf{座標順序}： \((x^0, x^1, x^2, x^3) = (t, x, y, z)\)
\item
  \textbf{時間座標}： \(t = ct_{\text{SI}}\) （長さ次元に正規化済み）
\item
  \textbf{磁場}： \(\mathbf{B} = c\mathbf{B}_{\text{SI}}\)
  （電場と次元を揃えた量）
\item
  \textbf{向き（基準の $4$-form）}：
  \(dt \wedge dx \wedge dy \wedge dz\)
\item
  \textbf{計量シグネチャ}： \((-, +, +, +)\)
\item
  \textbf{2-form の行列表現}： 第2章の規約に従い、基底
  \(dx^\mu \wedge dx^\nu\) の係数を、行列の \((\mu, \nu)\)
  成分に配置する。
\end{itemize}

この規約のもとで、電磁場 \(2\)-form \(F\) を次のように定義する。

\[
F = -E_x\,dt\wedge dx - E_y\,dt\wedge dy - E_z\,dt\wedge dz + B_x\,dy\wedge dz + B_y\,dz\wedge dx + B_z\,dx\wedge dy
\]

電場項（\(dt\)
を含む項）に負号がついているのは、後述するファラデーの法則
\(\mathrm{rot}\,\mathbf{E} = -\frac{\partial\mathbf{B}}{\partial t}\)
およびポテンシャル構成 \(F = -d\mathcal{A}\)
との整合性を保つための意図的な選択である。

これを \(4\times4\) の反対称行列で表示すれば、次のようになる。

\[F = \begin{pmatrix}
0 & -E_x & -E_y & -E_z \\
E_x & 0 & B_z & -B_y \\
E_y & -B_z & 0 & B_x \\
E_z & B_y & -B_x & 0
\end{pmatrix}\]

\begin{quote}
\textbf{注} （電場項の符号と行列） 本書の行列規約では、\(dt \wedge dx\)
の係数 \(-E_x\) を行列の \((t,x)\) 成分に置く。したがって
\(F_{tx}=-E_x,\;F_{xt}=E_x\) である。多くの教科書では \(F_{t x}\) を正の
\(E_x\) と定義するが、その場合は \(F = E_x dx \wedge dt + \dots\)
（順序が逆）とするか、あるいはポテンシャルの定義 \(F=d\mathcal{A}\)
の符号を調整する必要がある。本書では \(F=-d\mathcal{A}\) と
\(dt \wedge dx\) の順序を優先し、この符号を採る。
\end{quote}

\begin{center}\rule{0.5\linewidth}{0.5pt}\end{center}

\hypertarget{ux30dfux30f3ux30b3ux30d5ux30b9ux30adux30fcux8a08ux91cfmathbbr3-ux304bux30894ux6b21ux5143ux3078}{%
\subsection{\texorpdfstring{§10.3
ミンコフスキー計量------\(\mathbb{R}^3\)
から4次元へ}{§10.3 ミンコフスキー計量------\textbackslash mathbb\{R\}\^{}3 から4次元へ}}\label{ux30dfux30f3ux30b3ux30d5ux30b9ux30adux30fcux8a08ux91cfmathbbr3-ux304bux30894ux6b21ux5143ux3078}}

本書は一貫して \(\mathbb{R}^3\) のデカルト座標 \((x,y,z)\)
を舞台にしてきた。しかし \(F\) を扱うには、時間 \(t\) を加えた4次元時空
\((t,x,y,z)\) が必要になる。

幸い、ここまでに積み上げた知識があれば、拡張は容易だ。第6章で
\(\mathbf{g} = J^T J\) を導出し、第9章で \(\mathbf{g}\) から \(\ast\)
の辞書を作る手順を練習した。時空でも同じことをすればよい。計量は次のようになる（時間成分の符号が空間成分と異なる点が唯一の新しい要素だ）。

\[\mathbf{g} = \begin{pmatrix}
-1 & 0 & 0 & 0 \\
0 & 1 & 0 & 0 \\
0 & 0 & 1 & 0 \\
0 & 0 & 0 & 1
\end{pmatrix}\]

\begin{quote}
\textbf{注}
（正定値性からの離脱）本書の計量はこれまでつねに正定値（\(\mathbf{v}^T\mathbf{g}\,\mathbf{v} > 0\)
for
\(\mathbf{v}\neq\mathbf{0}\)）だった。ミンコフスキー計量は正定値ではなく、これまでの
\(\mathbf{g}=J^T J\) 型の計量とは異なる。ここでは「計量行列から \(\ast\)
の辞書を作る」という手順だけを引き継いでいる。この違いは \(\ast\)
の辞書に符号の変化として現れるが、計算手順そのものは変わらない。この点は、次章で多様体や計量を眺めるときの伏線になる。
\end{quote}

\begin{quote}
\textbf{注} （4次元ホッジ・スターの規約） 本章では向き（基底の順序）を
\(dt\wedge dx\wedge dy\wedge dz\) とし、時空の計量符号を \((-,+,+,+)\)
とする。この規約のもとで、以下の \(\ast F\) や \(\ast\mathcal{J}\)
を用いる。異なるシグネチャや向きの規約では、いくつかの符号が反転することに注意されたい。
\end{quote}

\begin{center}\rule{0.5\linewidth}{0.5pt}\end{center}

\hypertarget{df0-ux3092ux5168ux90e8ux66f8ux304dux4e0bux3059}{%
\subsection{\texorpdfstring{§10.4 \(dF=0\)
を全部書き下す}{§10.4 dF=0 を全部書き下す}}\label{df0-ux3092ux5168ux90e8ux66f8ux304dux4e0bux3059}}

\(F\) は \(2\)-form だから、\(dF\) は \(3\)-form になる。4次元空間での
\(3\)-form の独立成分は4つ。§10.2 で再定義した
\(F\)（電場項に負号を含む）について \(d\) を作用させ、同じ基底
\(3\)-form の係数をまとめる。

\[F = -E_x\,dt\wedge dx - E_y\,dt\wedge dy - E_z\,dt\wedge dz + B_x\,dy\wedge dz + B_y\,dz\wedge dx + B_z\,dx\wedge dy\]

各項に \(d\) を作用させる。電場項は次のようになる（符号の変化に注意）：

\[\begin{aligned}
d\!\left(-E_x\,dt\wedge dx\right) &= -\frac{\partial E_x}{\partial y}\,dy\wedge dt\wedge dx - \frac{\partial E_x}{\partial z}\,dz\wedge dt\wedge dx \\
&= -\frac{\partial E_x}{\partial y}\,dt\wedge dx\wedge dy + \frac{\partial E_x}{\partial z}\,dt\wedge dz\wedge dx \\
d\!\left(-E_y\,dt\wedge dy\right) &= -\frac{\partial E_y}{\partial z}\,dt\wedge dy\wedge dz + \frac{\partial E_y}{\partial x}\,dt\wedge dx\wedge dy \\
d\!\left(-E_z\,dt\wedge dz\right) &= -\frac{\partial E_z}{\partial x}\,dt\wedge dz\wedge dx + \frac{\partial E_z}{\partial y}\,dt\wedge dy\wedge dz
\end{aligned}\]

磁場項は時間微分 \(t\) も含み、次のようになる：

\[\begin{aligned}
d(B_x\,dy\wedge dz) &= \frac{\partial B_x}{\partial t}\,dt\wedge dy\wedge dz + \frac{\partial B_x}{\partial x}\,dx\wedge dy\wedge dz \\
d(B_y\,dz\wedge dx) &= \frac{\partial B_y}{\partial t}\,dt\wedge dz\wedge dx + \frac{\partial B_y}{\partial y}\,dy\wedge dz\wedge dx \\
d(B_z\,dx\wedge dy) &= \frac{\partial B_z}{\partial t}\,dt\wedge dx\wedge dy + \frac{\partial B_z}{\partial z}\,dz\wedge dx\wedge dy
\end{aligned}\]

これらをすべて足し合わせ、基底 \(3\)-form ごとに整理する。たとえば
\(dt\wedge dy\wedge dz\) の係数は次のようになる：

\[
\left(\frac{\partial B_x}{\partial t} + \frac{\partial E_z}{\partial y} - \frac{\partial E_y}{\partial z}\right) dt\wedge dy\wedge dz = \left(\frac{\partial B_x}{\partial t} + (\mathrm{rot}\,\mathbf{E})_x\right) dt\wedge dy\wedge dz
\]

各基底係数をゼロとおくと、\(dF=0\) は次の 4 本の式と同値になる。

\[dF = 0 \Longleftrightarrow \begin{cases}
\displaystyle \frac{\partial B_x}{\partial x} + \frac{\partial B_y}{\partial y} + \frac{\partial B_z}{\partial z} = 0 & (\mathrm{div}\,\mathbf{B} = 0) \\[1em]
\displaystyle (\mathrm{rot}\,\mathbf{E})_x = -\frac{\partial B_x}{\partial t} & \\[0.5em]
\displaystyle (\mathrm{rot}\,\mathbf{E})_y = -\frac{\partial B_y}{\partial t} & (\mathrm{rot}\,\mathbf{E} = -\frac{\partial\mathbf{B}}{\partial t}) \\[0.5em]
\displaystyle (\mathrm{rot}\,\mathbf{E})_z = -\frac{\partial B_z}{\partial t} & 
\end{cases}\]

1行目は磁場に関するガウスの法則、2〜4行目はファラデーの電磁誘導の法則である。電場項に負号を置いたことにより、微分形式の一本の方程式
\(dF=0\) から、ベクトル解析でおなじみの正しい符号の公式が導き出された。

\begin{center}\rule{0.5\linewidth}{0.5pt}\end{center}

\hypertarget{ast-f-ux3068ux6b8bux308aux306e2ux672c}{%
\subsection{\texorpdfstring{§10.5 \(\ast F\)
と残りの2本}{§10.5 \textbackslash ast F と残りの2本}}\label{ast-f-ux3068ux6b8bux308aux306e2ux672c}}

もう一つの式 \(d(\ast F) = \mu_0(\ast\mathcal{J})\) も、§10.4
と同じ手順で全部書き下そう。

まず、§10.2 で再定義した \(F\) から \(\ast F\) を求める。§10.3
のミンコフスキー計量シグネチャ \((-,+,+,+)\) と向き
\(dt\wedge dx\wedge dy\wedge dz\)
のもとでホッジ・スターを作用させると、\(\ast F\) は次のようになる。

\[
\ast F = B_x\,dt\wedge dx + B_y\,dt\wedge dy + B_z\,dt\wedge dz + E_x\,dy\wedge dz + E_y\,dz\wedge dx + E_z\,dx\wedge dy
\]

行列表示では次の通りだ。

\[\ast F = \begin{pmatrix}
0 & B_x & B_y & B_z \\
-B_x & 0 & E_z & -E_y \\
-B_y & -E_z & 0 & E_x \\
-B_z & E_y & -E_x & 0
\end{pmatrix}\]

電場項の符号を反転させた \(F\) を出発点としたことで、\(\ast\) 作用後の
\(\ast F\) では逆に磁場項が \(dt\)
を含む項になり、電場項が空間のみの項（\(2\)-form）へ移動している。

次に、\(d(\ast F)\) を展開する。基底 \(3\)-form
ごとに整理すると、次のようになる。

\[\begin{aligned}
d(\ast F) = &\left(\frac{\partial E_x}{\partial x} + \frac{\partial E_y}{\partial y} + \frac{\partial E_z}{\partial z}\right) dx\wedge dy\wedge dz \\
{+} &\left(\frac{\partial E_x}{\partial t} - \left(\frac{\partial B_z}{\partial y} - \frac{\partial B_y}{\partial z}\right)\right) dt\wedge dy\wedge dz \\
{+} &\left(\frac{\partial E_y}{\partial t} - \left(\frac{\partial B_x}{\partial z} - \frac{\partial B_z}{\partial x}\right)\right) dt\wedge dz\wedge dx \\
{+} &\left(\frac{\partial E_z}{\partial t} - \left(\frac{\partial B_y}{\partial x} - \frac{\partial B_x}{\partial y}\right)\right) dt\wedge dx\wedge dy
\end{aligned}\]

1行目は、空間上の \(2\)-form \(\ast_3E\) に \(d_3\) を作用させた
\(d_3(\ast_3E)\) である。ベクトル解析の辞書で読めば、この係数が
\(\mathrm{div}\,\mathbf{E}\) になる。2〜4行目では、係数全体が
\(\partial_t(\ast_3E)-d_3(\ast_3B)\) に対応している。

右辺の \(\mu_0(\ast\mathcal{J})\) は、電荷密度 \(\rho_{\mathrm e}\)
と電流密度 \(\mathbf{J}\) を含む \(3\)-form である。§10.1
の正規化のもとで次のように書ける。

\[
\mu_0(\ast\mathcal{J}) = \ast_3\left(\frac{\rho_{\mathrm e}}{\varepsilon_0}\right) - \mu_0 c J_x\,dt\wedge dy\wedge dz - \mu_0 c J_y\,dt\wedge dz\wedge dx - \mu_0 c J_z\,dt\wedge dx\wedge dy
\]

第1項は、各時刻の空間上でスカラー場 \(\rho_{\mathrm e}/\varepsilon_0\)
を \(3\)-form にしたものである。デカルト座標では

\[
\ast_3\left(\frac{\rho_{\mathrm e}}{\varepsilon_0}\right)=\frac{\rho_{\mathrm e}}{\varepsilon_0}\,dx\wedge dy\wedge dz
\]

である。

両辺の係数を比較すれば、まず空間上の微分形式として

\[
d_3(\ast_3E)=\ast_3\left(\frac{\rho_{\mathrm e}}{\varepsilon_0}\right)
\]

および

\[
d_3(\ast_3B)=\mu_0 c\,J+\frac{\partial(\ast_3E)}{\partial t}
\]

が得られる。

特に1本目は、さらに両辺に \(\ast_3\) を作用させて

\[
\ast_3d_3(\ast_3E)=\frac{\rho_{\mathrm e}}{\varepsilon_0}
\]

と読む。左辺が、本書の辞書でいう \(\mathrm{div}\,\mathbf E\) である。

同様に、2本目をベクトル解析の記法に戻せば、

\[
\mathrm{rot}\,\mathbf{B} = \mu_0 c \mathbf{J} + \frac{\partial \mathbf{E}}{\partial t}
\]

これによって、4 本すべてのマクスウェル方程式が、微分形式の 2
本の式から正しい符号と係数で導かれることが完全に確認できた。

つまり、マクスウェル方程式の4本は、微分形式では

\[dF = 0, \qquad d(\ast F) = \mu_0(\ast\mathcal{J})\]

の2本に集約される。そして、この2本を実際に行列と偏微分で展開すれば、見慣れたベクトル解析の式がすべて再現される。美しさの裏に、ちゃんと泥臭い計算が息づいている------それを見届けたことが、本章の成果である。

\begin{quote}
\textbf{注}
（なぜ2本なのか）鋭い読者ならこう思うかもしれない------「結局2本か。\(dF=0\)
と \(d(\ast F)=\mu_0(\ast\mathcal{J})\)
を、\textbf{1行}にまとめられないのか」と。できる。第12章では
\(\mathbf{E}\) と \(\mathbf{B}\)
をひとまとめにした複素ベクトルと、パウリ行列から作られるディラック演算子によって、マクスウェル方程式をたったの一撃に統合する。楽しみにしていてほしい。
\end{quote}

\begin{quote}
\textbf{注} （なぜここで終わるのか）多くの教科書は \(dF=0\)
を示した時点で「かくしてマクスウェル方程式は幾何学である」と締める。しかし本書の流儀は違う。行列の成分を全部書き下し、\(d\)
と \(\ast\)
の辞書を通じてベクトル解析の式を再現する------その往復ができることこそ、第1章から積み上げてきた「測定器」の枠組みの到達点だ。
\end{quote}

\begin{center}\rule{0.5\linewidth}{0.5pt}\end{center}

\hypertarget{ux30ddux30c6ux30f3ux30b7ux30e3ux30ebux69cbux6210f-dmathcala-ux304bux3089ux59cbux3081ux308b}{%
\subsection{\texorpdfstring{§10.6
ポテンシャル構成------\(F=-d\mathcal{A}\)
から始める}{§10.6 ポテンシャル構成------F=-d\textbackslash mathcal\{A\} から始める}}\label{ux30ddux30c6ux30f3ux30b7ux30e3ux30ebux69cbux6210f-dmathcala-ux304bux3089ux59cbux3081ux308b}}

§10.2 では、電磁場 \(F\) を \(\mathbf{E}\) と \(\mathbf{B}\)
の寄せ集めとして定義した。

しかし、ここにはもう一段深い見方がある。

\(F\) は、ある \(1\)-form から外微分によって作ることができる。

その \(1\)-form を \(\mathcal{A}\) と書く。もし

\[
F=-d\mathcal{A}
\]

と置けるなら、

\[
dF
=
d(-d\mathcal{A})
=
-d(d\mathcal{A})
=
0
\]

である。

つまり、§10.4 で成分を全部書き下して得た

\[
\mathrm{div}\,\mathbf{B}=0,
\qquad
\mathrm{rot}\,\mathbf{E}
=
-\frac{\partial\mathbf{B}}{\partial t}
\]

は、外微分の恒等式

\[
dd=0
\]

の中にすでに入っていたことになる。

これがポテンシャル表示の核心である。

§10.4 では、\(dF=0\) を4つの基底 \(3\)-form
の係数として展開した。しかし、\(F\) を \(-d\mathcal{A}\)
として作れば、その4本は一行で自動的に従う。

では、そのような \(\mathcal{A}\) を実際に書いてみよう。

\hypertarget{ux5143ux30ddux30c6ux30f3ux30b7ux30e3ux30ebux3092ux7f6eux304f}{%
\subsubsection{4元ポテンシャルを置く}\label{ux5143ux30ddux30c6ux30f3ux30b7ux30e3ux30ebux3092ux7f6eux304f}}

4元ポテンシャル \(\mathcal{A}\) は \(1\)-form
である。電磁気学では、スカラーポテンシャル \(\phi\)
とベクトルポテンシャル \(\mathbf{A}\) が組になる。

本章では、§10.2 の \(F\) と、標準的な

\[
\mathbf{E}
=
-\nabla\phi
-
\frac{\partial\mathbf{A}}{\partial t},
\qquad
\mathbf{B}
=
\nabla\times\mathbf{A}
\]

を同時に満たすため、ポテンシャル \(1\)-form を次の符号で置く。

\[
\mathcal{A}
=
\phi\,dt
-
A_x\,dx
-
A_y\,dy
-
A_z\,dz
\]

ここで \(\phi,A_x,A_y,A_z\) は、いずれも \((t,x,y,z)\) の関数である。

\begin{quote}
\textbf{注}（ポテンシャルの正規化） ここでの \(\mathbf{A}\)
も、本章の正規化に合わせた量として読む。第10章の \(\mathbf{B}\)
が本来の磁束密度ではなく \(c\mathbf{B}_{\mathrm{SI}}\)
であるのと同様に、\(\mathbf{A}\) も必要なら
\(c\mathbf{A}_{\mathrm{SI}}\)
として正規化されたベクトルポテンシャルだと思えばよい。この約束のもとで、\(\mathbf{B}=\nabla\times\mathbf{A}\)
と \(\mathbf{E}=-\nabla\phi-\partial\mathbf{A}/\partial t\)
が同時に成り立つ。
\end{quote}

ここから \(d\mathcal{A}\) を計算し、最後に \(F=-d\mathcal{A}\) と置く。

この計算の目的は、§10.2 で定義した \(F\)
の成分と、通常のポテンシャル公式が同時に整合することを確認することである。

\hypertarget{dmathcala-ux3092ux6210ux5206ux3067ux8a08ux7b97ux3059ux308b}{%
\subsubsection{\texorpdfstring{\(d\mathcal{A}\)
を成分で計算する}{d\textbackslash mathcal\{A\} を成分で計算する}}\label{dmathcala-ux3092ux6210ux5206ux3067ux8a08ux7b97ux3059ux308b}}

\(d\mathcal{A}\)
の計算は、第5章以来の手順そのものである。\(\mathcal{A}\) の4項それぞれに
\(d\) を作用させる。

\[
\begin{aligned}
d(\phi\,dt)
&=
\frac{\partial\phi}{\partial x}\,dx\wedge dt
+
\frac{\partial\phi}{\partial y}\,dy\wedge dt
+
\frac{\partial\phi}{\partial z}\,dz\wedge dt,
\\[0.5em]
d(-A_x\,dx)
&=
-\frac{\partial A_x}{\partial t}\,dt\wedge dx
-
\frac{\partial A_x}{\partial y}\,dy\wedge dx
-
\frac{\partial A_x}{\partial z}\,dz\wedge dx,
\\[0.5em]
d(-A_y\,dy)
&=
-\frac{\partial A_y}{\partial t}\,dt\wedge dy
-
\frac{\partial A_y}{\partial x}\,dx\wedge dy
-
\frac{\partial A_y}{\partial z}\,dz\wedge dy,
\\[0.5em]
d(-A_z\,dz)
&=
-\frac{\partial A_z}{\partial t}\,dt\wedge dz
-
\frac{\partial A_z}{\partial x}\,dx\wedge dz
-
\frac{\partial A_z}{\partial y}\,dy\wedge dz.
\end{aligned}
\]

\(\phi\) の時間微分は \(dt\wedge dt=0\) で消える。

同様に、\(A_x\) の \(x\) 微分、\(A_y\) の \(y\) 微分、\(A_z\) の \(z\)
微分も、それぞれ \(dx\wedge dx=0\)、\(dy\wedge dy=0\)、\(dz\wedge dz=0\)
によって消える。

これらを足し合わせ、ウェッジ積の反対称性で基底 \(2\)-form
ごとに整理する。

まず、\(dt\wedge dx\) の係数を見よう。

\(d(\phi\,dt)\) から

\[
\frac{\partial\phi}{\partial x}\,dx\wedge dt
=
-\frac{\partial\phi}{\partial x}\,dt\wedge dx
\]

が出る。また、\(d(-A_x\,dx)\) から

\[
-\frac{\partial A_x}{\partial t}\,dt\wedge dx
\]

が出る。したがって、\(d\mathcal{A}\) の \(dt\wedge dx\) 係数は

\[
-\frac{\partial\phi}{\partial x}
-
\frac{\partial A_x}{\partial t}
\]

である。

次に、\(dy\wedge dz\) の係数を見よう。

\(d(-A_y\,dy)\) から

\[
-\frac{\partial A_y}{\partial z}\,dz\wedge dy
=
\frac{\partial A_y}{\partial z}\,dy\wedge dz
\]

が出る。また、\(d(-A_z\,dz)\) から

\[
-\frac{\partial A_z}{\partial y}\,dy\wedge dz
\]

が出る。したがって、\(d\mathcal{A}\) の \(dy\wedge dz\) 係数は

\[
-\frac{\partial A_z}{\partial y}
+
\frac{\partial A_y}{\partial z}
\]

である。

全6基底について同様に整理すると、

\[
\begin{aligned}
d\mathcal{A}
&=
\left(
-\frac{\partial\phi}{\partial x}
-
\frac{\partial A_x}{\partial t}
\right)
dt\wedge dx
\\[0.5em]
&\quad+
\left(
-\frac{\partial\phi}{\partial y}
-
\frac{\partial A_y}{\partial t}
\right)
dt\wedge dy
\\[0.5em]
&\quad+
\left(
-\frac{\partial\phi}{\partial z}
-
\frac{\partial A_z}{\partial t}
\right)
dt\wedge dz
\\[0.5em]
&\quad+
\left(
-\frac{\partial A_z}{\partial y}
+
\frac{\partial A_y}{\partial z}
\right)
dy\wedge dz
\\[0.5em]
&\quad+
\left(
-\frac{\partial A_x}{\partial z}
+
\frac{\partial A_z}{\partial x}
\right)
dz\wedge dx
\\[0.5em]
&\quad+
\left(
-\frac{\partial A_y}{\partial x}
+
\frac{\partial A_x}{\partial y}
\right)
dx\wedge dy.
\end{aligned}
\]

\hypertarget{ux7b26ux53f7ux306eux7406ux7531}{%
\subsubsection{符号の理由}\label{ux7b26ux53f7ux306eux7406ux7531}}

§10.2 で定義した \(F\) は

\[
F
=
-E_x\,dt\wedge dx
-
E_y\,dt\wedge dy
-
E_z\,dt\wedge dz
+
B_x\,dy\wedge dz
+
B_y\,dz\wedge dx
+
B_z\,dx\wedge dy
\]

であった。

ここで、もし \(F=d\mathcal{A}\) と置くと、\(dt\wedge dx\) 係数の比較から

\[
-E_x
=
-\frac{\partial\phi}{\partial x}
-
\frac{\partial A_x}{\partial t}
\]

となり、

\[
E_x
=
\frac{\partial\phi}{\partial x}
+
\frac{\partial A_x}{\partial t}
\]

を得てしまう。これは標準的な

\[
E_x
=
-\frac{\partial\phi}{\partial x}
-
\frac{\partial A_x}{\partial t}
\]

とは逆符号である。

そこで本章では

\[
F=-d\mathcal{A}
\]

と置く。

このとき、\(dt\wedge dx\) の係数は

\[
-E_x
=
\frac{\partial\phi}{\partial x}
+
\frac{\partial A_x}{\partial t}
\]

となるので、

\[
E_x
=
-\frac{\partial\phi}{\partial x}
-
\frac{\partial A_x}{\partial t}
\]

が得られる。

同様に、\(dy\wedge dz\) の係数を見ると、\(-d\mathcal{A}\) の
\(dy\wedge dz\) 係数は

\[
\frac{\partial A_z}{\partial y}
-
\frac{\partial A_y}{\partial z}
\]

である。したがって、

\[
B_x
=
\frac{\partial A_z}{\partial y}
-
\frac{\partial A_y}{\partial z}
\]

が得られる。

全成分をまとめると、

\[
\mathbf{E}
=
-\nabla\phi
-
\frac{\partial\mathbf{A}}{\partial t},
\qquad
\mathbf{B}
=
\nabla\times\mathbf{A}.
\]

つまり、

\[
\mathcal{A}
=
\phi\,dt
-
A_x\,dx
-
A_y\,dy
-
A_z\,dz,
\qquad
F=-d\mathcal{A}
\]

と置けば、§10.2 の \(F\)
の係数表示と、標準的なポテンシャル公式が同時に整合する。

\begin{quote}
\textbf{注}（\(F\) とポテンシャルの符号規約） \(F=-d\mathcal{A}\)
という定義は、§10.2 の \(F\) の \(dt\wedge dx\) 係数が \(-E_x\)
であることと、\(\mathbf{E}=-\nabla\phi-\partial\mathbf{A}/\partial t\)、\(\mathbf{B}=\nabla\times\mathbf{A}\)
を同時に満たすために採用している。\(F=+d\mathcal{A}\)、\(\mathcal{A}=-\phi\,dt+\cdots\)
など、異なる符号の組み合わせを使う流儀も存在する。
\end{quote}

\hypertarget{dd0-ux304c-df0-ux3092ux8a3cux660eux3059ux308b}{%
\subsubsection{\texorpdfstring{\(dd=0\) が \(dF=0\)
を証明する}{dd=0 が dF=0 を証明する}}\label{dd0-ux304c-df0-ux3092ux8a3cux660eux3059ux308b}}

ここで、冒頭の話に戻ろう。

\(F=-d\mathcal{A}\) だから、

\[
dF
=
d(-d\mathcal{A})
=
-d(d\mathcal{A})
=
0.
\]

これは、第5章 §5.8 で見た外微分の基本性質

\[
dd=0
\]

そのものである。

したがって、§10.4 で成分展開して得た

\[
\mathrm{div}\,\mathbf{B}=0,
\qquad
\mathrm{rot}\,\mathbf{E}
=
-\frac{\partial\mathbf{B}}{\partial t}
\]

は、\(F=-d\mathcal{A}\) と置いた瞬間に自動的に従う。

ここで起きていることは、単なる記法の圧縮ではない。

第5章で見た \(dd=0\)
が、ここではマクスウェル方程式の半分として現れている。

\(F\) を \(-d\mathcal{A}\) として作った瞬間、\(dF=0\)
は証明すべき法則ではなく、外微分の構造から自動的に従う恒等式になる。

\begin{quote}
\textbf{注}（\(\mathbf{B}\) と \(\mathrm{div}\,\mathbf{B}=0\)）
ベクトル解析でも「\(\mathbf{B}\) が何かの回転で書けるなら
\(\mathrm{div}\,\mathbf{B}=0\)」は公式として知られている。\(\mathrm{div}\,\mathrm{rot}\equiv0\)
である。これが \(dd=0\) の3次元ベクトル解析版だ。
\end{quote}

\hypertarget{ux6b8bux308bux306eux306fux3082ux30461ux672cux306eux65b9ux7a0bux5f0fux3067ux3042ux308b}{%
\subsubsection{残るのは、もう1本の方程式である}\label{ux6b8bux308bux306eux306fux3082ux30461ux672cux306eux65b9ux7a0bux5f0fux3067ux3042ux308b}}

\(dF=0\) は、\(F=-d\mathcal{A}\) から自動的に従う。

では、もう1本

\[
d(\ast F)
=
\mu_0(\ast\mathcal{J})
\]

はどうなるだろうか。

こちらは、自動的にゼロになるわけではない。\(F=-d\mathcal{A}\)
を代入すると、

\[
d(\ast(-d\mathcal{A}))
=
\mu_0(\ast\mathcal{J})
\]

すなわち

\[
-d(\ast d\mathcal{A})
=
\mu_0(\ast\mathcal{J})
\]

となる。

これが、ポテンシャル \(\mathcal{A}\) が満たす方程式である。

\begin{quote}
\textbf{注}（ゲージ自由度）\(\mathcal{A}'=\mathcal{A}+d\chi\)（\(\chi\)
は任意の
\(0\)-form）と変換する。このとき\(F'=-d\mathcal{A}'=-d\mathcal{A}-d(d\chi)=-d\mathcal{A}=F\)となり、物理的な電磁場
\(F\) は不変である。これがゲージ変換であり、\(dd=0\)
のもう一つの現れだ。電磁気学では、スカラーポテンシャル \(\phi\)
とベクトルポテンシャル \(\mathbf{A}\) を同時に調整することで、\(F\)
を変えずに計算に都合のよいゲージを選ぶことができる。
\end{quote}

\begin{quote}
\textbf{【ここまでのチェックポイント — 第10章全体】}

\begin{itemize}
\item
  電磁場 \(F\) は \(4\times4\) 反対称行列。6つの独立成分に
  \(E_x,E_y,E_z,B_x,B_y,B_z\) が収まる。
\item
  \(dF=0\) を成分展開すると \(\mathrm{div}\,\mathbf{B}=0\) と
  \(\mathrm{rot}\,\mathbf{E}=-\partial\mathbf{B}/\partial t\) が現れる。
\item
  この章のシグネチャ・向きの規約のもとでは、\(\ast F\) は \(E\) と \(B\)
  を入れ替える形に見える。\(d(\ast F)=\mu_0(\ast\mathcal{J})\)
  から残りの2本が出る。
\item
  \textbf{ポテンシャル構成}：\(\mathcal{A}=\phi\,dt-A_x\,dx-A_y\,dy-A_z\,dz\)
  から \(F=-d\mathcal{A}\) が得られる。すると \(dF=0\) は \(dd=0\)
  より自動的に成り立つ。これは \(\mathrm{div}\,\mathbf{B}=0\) と
  \(\mathrm{rot}\,\mathbf{E}=-\partial\mathbf{B}/\partial t\)
  が外微分の構造に含まれている、ということである。
\item
  もう1本の方程式 \(d(\ast F)=\mu_0(\ast\mathcal{J})\)
  は、\(F=-d\mathcal{A}\) を代入すると、ポテンシャル \(\mathcal{A}\)
  が満たす方程式になる。
\end{itemize}
\end{quote}

\begin{center}\rule{0.5\linewidth}{0.5pt}\end{center}

\hypertarget{ux4ed8ux9332edf-ux3068-dast-f-ux306eux30b9ux30e9ux30a4ux30b9ux884cux5217ux8868ux793a-4times4times4-ux914dux5217ux3067ux898bux308bux30deux30afux30b9ux30a6ux30a7ux30ebux65b9ux7a0bux5f0f}{%
\section{\texorpdfstring{付録E：\(dF\) と \(d(\ast F)\)
のスライス行列表示 ------ \(4\times4\times4\)
配列で見るマクスウェル方程式}{付録E：dF と d(\textbackslash ast F) のスライス行列表示 ------ 4\textbackslash times4\textbackslash times4 配列で見るマクスウェル方程式}}\label{ux4ed8ux9332edf-ux3068-dast-f-ux306eux30b9ux30e9ux30a4ux30b9ux884cux5217ux8868ux793a-4times4times4-ux914dux5217ux3067ux898bux308bux30deux30afux30b9ux30a6ux30a7ux30ebux65b9ux7a0bux5f0f}}

§10.4・§10.5 では \(dF\) と \(d(\ast F)\) を基底 \(3\)-form
の係数として展開した。付録Eでは、同じ計算を
\textbf{$4\times4$ スライス行列の束}------実質的には \(4\times4\times4\)
の3階テンソル------で可視化する。付録Aの延長線上にあり、本書の「全部行列に載せる」流儀の到達点でもある。

\hypertarget{e.1-ux57faux5e95-3-form-ux3068ux305dux306eux30b9ux30e9ux30a4ux30b9ux884cux5217-ux516816ux679a}{%
\subsection{\texorpdfstring{E.1 基底 \(3\)-form とそのスライス行列
------
全16枚}{E.1 基底 3-form とそのスライス行列 ------ 全16枚}}\label{e.1-ux57faux5e95-3-form-ux3068ux305dux306eux30b9ux30e9ux30a4ux30b9ux884cux5217-ux516816ux679a}}

4次元の基底 \(3\)-form は次の4つである（§10.4
の並びと比べると、\(\omega_2\) と \(\omega_3\)
の符号・順序が異なるが、内容は同じ4つである）。

\[
\omega_1 = dt\wedge dx\wedge dy,\qquad
\omega_2 = dt\wedge dx\wedge dz,\qquad
\omega_3 = dt\wedge dy\wedge dz,\qquad
\omega_4 = dx\wedge dy\wedge dz
\]

各 \(\omega_i\) について、座標方向 \(t,x,y,z\) の\textbf{スライス行列}
\(\mathbf{S}_{t}^{(\omega_i)}, \mathbf{S}_{x}^{(\omega_i)}, \mathbf{S}_{y}^{(\omega_i)}, \mathbf{S}_{z}^{(\omega_i)}\)
を定義する。スライス行列とは「その方向の座標を除いた残りの \(3\)
方向が作る \(3\)-form」を \(4\times4\)
反対称行列に焼き直したものである。行・列の順は
\((t,x,y,z)\)。非ゼロ成分には \(\pm1\) が入る。\(4\) 基底 \(\times 4\)
スライス \(=\) 全 \(16\) 枚。

\textbf{（1）} \(\omega_1 = dt\wedge dx\wedge dy\) のスライス：

\[
\mathbf{S}_{t}^{(\omega_1)}=\begin{pmatrix}0&0&0&0\\[2pt]0&0&1&0\\[2pt]0&-1&0&0\\[2pt]0&0&0&0\end{pmatrix},\qquad
\mathbf{S}_{x}^{(\omega_1)}=\begin{pmatrix}0&0&-1&0\\[2pt]0&0&0&0\\[2pt]1&0&0&0\\[2pt]0&0&0&0\end{pmatrix}
\]

\[
\mathbf{S}_{y}^{(\omega_1)}=\begin{pmatrix}0&1&0&0\\[2pt]-1&0&0&0\\[2pt]0&0&0&0\\[2pt]0&0&0&0\end{pmatrix},\qquad
\mathbf{S}_{z}^{(\omega_1)}=\begin{pmatrix}0&0&0&0\\[2pt]0&0&0&0\\[2pt]0&0&0&0\\[2pt]0&0&0&0\end{pmatrix}
\]

\textbf{（2）} \(\omega_2 = dt\wedge dx\wedge dz\) のスライス：

\[
\mathbf{S}_{t}^{(\omega_2)}=\begin{pmatrix}0&0&0&0\\[2pt]0&0&0&1\\[2pt]0&0&0&0\\[2pt]0&-1&0&0\end{pmatrix},\qquad
\mathbf{S}_{x}^{(\omega_2)}=\begin{pmatrix}0&0&0&-1\\[2pt]0&0&0&0\\[2pt]0&0&0&0\\[2pt]1&0&0&0\end{pmatrix}
\]

\[
\mathbf{S}_{y}^{(\omega_2)}=\begin{pmatrix}0&0&0&0\\[2pt]0&0&0&0\\[2pt]0&0&0&0\\[2pt]0&0&0&0\end{pmatrix},\qquad
\mathbf{S}_{z}^{(\omega_2)}=\begin{pmatrix}0&1&0&0\\[2pt]-1&0&0&0\\[2pt]0&0&0&0\\[2pt]0&0&0&0\end{pmatrix}
\]

\textbf{（3）} \(\omega_3 = dt\wedge dy\wedge dz\) のスライス：

\[
\mathbf{S}_{t}^{(\omega_3)}=\begin{pmatrix}0&0&0&0\\[2pt]0&0&0&0\\[2pt]0&0&0&1\\[2pt]0&0&-1&0\end{pmatrix},\qquad
\mathbf{S}_{x}^{(\omega_3)}=\begin{pmatrix}0&0&0&0\\[2pt]0&0&0&0\\[2pt]0&0&0&0\\[2pt]0&0&0&0\end{pmatrix}
\]

\[
\mathbf{S}_{y}^{(\omega_3)}=\begin{pmatrix}0&0&0&-1\\[2pt]0&0&0&0\\[2pt]0&0&0&0\\[2pt]1&0&0&0\end{pmatrix},\qquad
\mathbf{S}_{z}^{(\omega_3)}=\begin{pmatrix}0&0&-1&0\\[2pt]0&0&0&0\\[2pt]1&0&0&0\\[2pt]0&0&0&0\end{pmatrix}
\]

\textbf{（4）} \(\omega_4 = dx\wedge dy\wedge dz\) のスライス：

\[
\mathbf{S}_{t}^{(\omega_4)}=\begin{pmatrix}0&0&0&0\\[2pt]0&0&0&0\\[2pt]0&0&0&0\\[2pt]0&0&0&0\end{pmatrix},\qquad
\mathbf{S}_{x}^{(\omega_4)}=\begin{pmatrix}0&0&0&0\\[2pt]0&0&0&0\\[2pt]0&0&0&1\\[2pt]0&0&-1&0\end{pmatrix}
\]

\[
\mathbf{S}_{y}^{(\omega_4)}=\begin{pmatrix}0&0&0&0\\[2pt]0&0&0&-1\\[2pt]0&0&0&0\\[2pt]0&1&0&0\end{pmatrix},\qquad
\mathbf{S}_{z}^{(\omega_4)}=\begin{pmatrix}0&0&0&0\\[2pt]0&0&-1&0\\[2pt]0&1&0&0\\[2pt]0&0&0&0\end{pmatrix}
\]

16枚のうち、\(\mathbf{S}_{z}^{(\omega_1)}, \mathbf{S}_{y}^{(\omega_2)}, \mathbf{S}_{x}^{(\omega_3)}, \mathbf{S}_{t}^{(\omega_4)}\)
の4枚はゼロ行列である。残る12枚に \(\pm1\)
が一つずつ（と反対称のペアが一つずつ）入っている。これが
\(4\times4\times4\) テンソルの正体だ。

\hypertarget{e.2-df-ux3092ux30b9ux30e9ux30a4ux30b9ux884cux5217ux3067ux66f8ux304f}{%
\subsection{\texorpdfstring{E.2 \(dF\)
をスライス行列で書く}{E.2 dF をスライス行列で書く}}\label{e.2-df-ux3092ux30b9ux30e9ux30a4ux30b9ux884cux5217ux3067ux66f8ux304f}}

§10.4 の展開により、\(dF\) の各基底係数は次で与えられる。

\[
\begin{aligned}
A_{txy} &= \frac{\partial E_y}{\partial x} - \frac{\partial E_x}{\partial y} + \frac{\partial B_z}{\partial t} = (\mathrm{rot}\,\mathbf{E})_z + \frac{\partial B_z}{\partial t}
\quad (\omega_1 = dt\wedge dx\wedge dy\text{ の係数}) \\[6pt]
A_{txz} &= \frac{\partial E_z}{\partial x} - \frac{\partial E_x}{\partial z} - \frac{\partial B_y}{\partial t} = -(\mathrm{rot}\,\mathbf{E})_y - \frac{\partial B_y}{\partial t}
\quad (\omega_2 = dt\wedge dx\wedge dz\text{ の係数}) \\[6pt]
A_{tyz} &= \frac{\partial E_z}{\partial y} - \frac{\partial E_y}{\partial z} + \frac{\partial B_x}{\partial t} = (\mathrm{rot}\,\mathbf{E})_x + \frac{\partial B_x}{\partial t}
\quad (\omega_3 = dt\wedge dy\wedge dz\text{ の係数}) \\[6pt]
A_{xyz} &= \frac{\partial B_x}{\partial x} + \frac{\partial B_y}{\partial y} + \frac{\partial B_z}{\partial z} = \mathrm{div}\,\mathbf{B}
\quad (\omega_4 = dx\wedge dy\wedge dz\text{ の係数})
\end{aligned}
\]

\(dF\) の各スライスは、4つの基底スライス行列を係数 \(A_{\cdots}\)
で線形結合したものである。たとえば \(t\)-スライス
\(\mathbf{S}_{t}^{(dF)}\) は、全 \(16\) 成分が明示された単一の
\(4\times4\) 反対称行列にまとまる。

\[
\mathbf{S}_{t}^{(dF)}
{=}
\left(\begin{array}{c|cccc}
 & t & x & y & z \\\hline
t & 0 & 0 & 0 & 0 \\[6pt]
x & 0 & 0 &
\displaystyle \frac{\partial E_y}{\partial x} - \frac{\partial E_x}{\partial y} + \frac{\partial B_z}{\partial t} &
\displaystyle \frac{\partial E_z}{\partial x} - \frac{\partial E_x}{\partial z} - \frac{\partial B_y}{\partial t} \\[14pt]
y & 0 &
\displaystyle \frac{\partial E_x}{\partial y} - \frac{\partial E_y}{\partial x} - \frac{\partial B_z}{\partial t} &
0 &
\displaystyle \frac{\partial E_z}{\partial y} - \frac{\partial E_y}{\partial z} + \frac{\partial B_x}{\partial t} \\[14pt]
z & 0 &
\displaystyle \frac{\partial E_x}{\partial z} - \frac{\partial E_z}{\partial x} + \frac{\partial B_y}{\partial t} &
\displaystyle \frac{\partial E_y}{\partial z} - \frac{\partial E_z}{\partial y} - \frac{\partial B_x}{\partial t} &
0
\end{array}\right)
\]

\(x,y,z\) のスライス
\(\mathbf{S}_{x}^{(dF)}, \mathbf{S}_{y}^{(dF)}, \mathbf{S}_{z}^{(dF)}\)
も同様に4項の線形結合から得られる。\(dF\)
全体は、これら4枚のスライス行列の束である。

\hypertarget{e.3-ux30d5ux30edux30d9ux30cbux30a6ux30b9ux7a4dux3067ux4fc2ux6570ux3092ux629cux304dux51faux3059}{%
\subsection{E.3
フロベニウス積で係数を抜き出す}\label{e.3-ux30d5ux30edux30d9ux30cbux30a6ux30b9ux7a4dux3067ux4fc2ux6570ux3092ux629cux304dux51faux3059}}

付録Dで導入したフロベニウス積は、\(4\times4\)
行列にもそのまま拡張できる。

\begin{quote}
\textbf{注}
（計量由来の内積ではない）ここでのフロベニウス積は、ミンコフスキー計量に由来する時空の内積ではない。反対称行列表示から係数を抜き出すための、配列上のペアリングである。基底スライス行列との係数抽出演算として使っており、\(\frac{1}{2}\operatorname{tr}(A^T B)\)
はそのための計算手段だ。
\end{quote}

\[
A \cdot B = \frac{1}{2}\operatorname{tr}(A^T B) = \frac{1}{2}\sum_{i=1}^{4}\sum_{j=1}^{4} A_{ij}B_{ij}
\]

この内積を使うと、各スライスから係数を\textbf{一行で抽出}できる。たとえば
\(A_{txy}\) は

\[
A_{txy} = \mathbf{S}_{t}^{(\omega_1)} \cdot \mathbf{S}_{t}^{(dF)}
\]

である。\(\mathbf{S}_{t}^{(\omega_1)}\) の非ゼロ成分は
\((x,y)=+1, (y,x)=-1\) だけなので、フロベニウス積は
\((1 \cdot S_{xy} + (-1) \cdot S_{yx})/2 = (S_{xy} - S_{yx})/2 = S_{xy}\)（\(\mathbf{S}_{t}^{(dF)}\)
が反対称だから \(S_{yx}=-S_{xy}\)）。一発で係数が抜ける。

同様に、\textbf{すべての係数は対応する基底スライスとのフロベニウス積で得られる}。この構造は、付録Dで
\(\ast_{2\to1}(\mathbf{M}) = (E_1\!\cdot\!\mathbf{M},\;E_2\!\cdot\!\mathbf{M},\;E_3\!\cdot\!\mathbf{M})\)
が \(2\)-form \(\mathbf{M}\)
の各成分を抽出したのと完全に相似である。次数が \(1\to2\) から \(2\to3\)
に上がり、内積をとる相手が縦ベクトルから \(4\times4\)
行列の束に変わっただけだ。

\hypertarget{e.4-df0-ux3092ux30b9ux30e9ux30a4ux30b9ux3067ux8aadux3080}{%
\subsection{\texorpdfstring{E.4 \(dF=0\)
をスライスで読む}{E.4 dF=0 をスライスで読む}}\label{e.4-df0-ux3092ux30b9ux30e9ux30a4ux30b9ux3067ux8aadux3080}}

\(dF=0\) とは、4枚のスライス行列すべてがゼロ行列になることである。

\[
\mathbf{S}_{t}^{(dF)} = \mathbf{0},\quad
\mathbf{S}_{x}^{(dF)} = \mathbf{0},\quad
\mathbf{S}_{y}^{(dF)} = \mathbf{0},\quad
\mathbf{S}_{z}^{(dF)} = \mathbf{0}
\]

\(t\)-スライス（上記）の非ゼロ成分を読めば、

\begin{itemize}
\item
  \((x,y)\) 成分 \(=0\) から
  \(\mathrm{rot}_z\,\mathbf{E} + \partial B_z/\partial t = 0\)
\item
  \((x,z)\) 成分 \(=0\) から
  \(-\mathrm{rot}_y\,\mathbf{E} - \partial B_y/\partial t = 0\)
\item
  \((y,z)\) 成分 \(=0\) から
  \(\mathrm{rot}_x\,\mathbf{E} + \partial B_x/\partial t = 0\)
\end{itemize}

これら3つは
\(\mathrm{rot}\,\mathbf{E} = -\partial\mathbf{B}/\partial t\)
にほかならない。\(z\)-スライス（\(\mathbf{S}_{z}^{(\omega_4)}\)
由来）の非ゼロ成分からは \(\mathrm{div}\,\mathbf{B} = 0\)
が現れる。見かけは巨大だが、中身は §10.4
の4本の係数比較式の繰り返しにすぎない。

\hypertarget{e.5-dast-f-mu_0astmathcalj-ux306eux30b9ux30e9ux30a4ux30b9ux8868ux793a}{%
\subsection{\texorpdfstring{E.5 \(d(\ast F) = \mu_0(\ast\mathcal{J})\)
のスライス表示}{E.5 d(\textbackslash ast F) = \textbackslash mu\_0(\textbackslash ast\textbackslash mathcal\{J\}) のスライス表示}}\label{e.5-dast-f-mu_0astmathcalj-ux306eux30b9ux30e9ux30a4ux30b9ux8868ux793a}}

有源側も同型の構造を持つ。\(\ast F\) と \(\ast\mathcal{J}\) は §10.5
で展開済みであり、\(d(\ast F)\) は \(dF\)
と同型のスライス行列になる------\(\mathbf{E}\) と \(\mathbf{B}\)
の係数位置が入れ替わるだけだ。

\(d(\ast F)\) の \(\omega_1\sim\omega_4\) 係数を
\(B_{txy}, B_{txz}, B_{tyz}, B_{xyz}\) と書く。

\[
\begin{aligned}
B_{txy} &= \frac{\partial B_x}{\partial y} - \frac{\partial B_y}{\partial x} + \frac{\partial E_z}{\partial t}
\quad (\text{§10.5 の } dt\wedge dx\wedge dy \text{ 係数}) \\[6pt]
B_{txz} &= \frac{\partial B_x}{\partial z} - \frac{\partial B_z}{\partial x} - \frac{\partial E_y}{\partial t} \\[6pt]
B_{tyz} &= \frac{\partial B_y}{\partial z} - \frac{\partial B_z}{\partial y} + \frac{\partial E_x}{\partial t} \\[6pt]
B_{xyz} &= \frac{\partial E_x}{\partial x} + \frac{\partial E_y}{\partial y} + \frac{\partial E_z}{\partial z}
\quad (= \mathrm{div}\,\mathbf{E})
\end{aligned}
\]

\(d(\ast F)\) の各スライスも \(dF\) と同様に、4つの係数 \(B_{\cdots}\)
で基底スライス行列を線形結合すればよい。\(t\)-スライスを項別に書く。

\[
\begin{aligned}
\mathbf{S}_{t}^{(d(\ast F))}
&= \left(\frac{\partial B_x}{\partial y} - \frac{\partial B_y}{\partial x} + \frac{\partial E_z}{\partial t}\right)
\begin{pmatrix}0&0&0&0\\[2pt]0&0&1&0\\[2pt]0&-1&0&0\\[2pt]0&0&0&0\end{pmatrix} \\[10pt]
&\quad+ \left(\frac{\partial B_x}{\partial z} - \frac{\partial B_z}{\partial x} - \frac{\partial E_y}{\partial t}\right)
\begin{pmatrix}0&0&0&0\\[2pt]0&0&0&1\\[2pt]0&0&0&0\\[2pt]0&-1&0&0\end{pmatrix} \\[10pt]
&\quad+ \left(\frac{\partial B_y}{\partial z} - \frac{\partial B_z}{\partial y} + \frac{\partial E_x}{\partial t}\right)
\begin{pmatrix}0&0&0&0\\[2pt]0&0&0&0\\[2pt]0&0&0&1\\[2pt]0&0&-1&0\end{pmatrix}
{+} B_{xyz}
\begin{pmatrix}0&0&0&0\\[2pt]0&0&0&0\\[2pt]0&0&0&0\\[2pt]0&0&0&0\end{pmatrix}
\end{aligned}
\]

同じ位置の成分を一枚の行列に足し合わせる。\(dF\) の
\(t\)-スライスと見比べてほしい------\(\mathbf{E}\) と
\(\mathbf{B}\)（および添字の巡回）がきれいに入れ替わっている。

\[
\mathbf{S}_{t}^{(d(\ast F))}
{=}
\left(\begin{array}{c|cccc}
 & t & x & y & z \\\hline
t & 0 & 0 & 0 & 0 \\[6pt]
x & 0 & 0 &
\displaystyle\frac{\partial B_x}{\partial y} {-} \frac{\partial B_y}{\partial x} {+} \frac{\partial E_z}{\partial t} &
\displaystyle\frac{\partial B_x}{\partial z} {-} \frac{\partial B_z}{\partial x} {-} \frac{\partial E_y}{\partial t} \\[14pt]
y & 0 &
\displaystyle{-}\frac{\partial B_x}{\partial y} {+} \frac{\partial B_y}{\partial x} {-} \frac{\partial E_z}{\partial t} &
0 &
\displaystyle\frac{\partial B_y}{\partial z} {-} \frac{\partial B_z}{\partial y} {+} \frac{\partial E_x}{\partial t} \\[14pt]
z & 0 &
\displaystyle{-}\frac{\partial B_x}{\partial z} {+} \frac{\partial B_z}{\partial x} {+} \frac{\partial E_y}{\partial t} &
\displaystyle{-}\frac{\partial B_y}{\partial z} {+} \frac{\partial B_z}{\partial y} {-} \frac{\partial E_x}{\partial t} &
0
\end{array}\right)
\]

右辺 \(\mu_0(\ast\mathcal{J})\)
も同じ4枚のスライス行列の線形結合で書ける。\(\ast\mathcal{J}\)
を本付録の基底順に展開し直すと（§10.5
参照）、\(\omega_1\!\sim\!\omega_4\) の係数は順に
\(-\mu_0 c J_z,\; +\mu_0 c J_y,\; -\mu_0 c J_x,\; \rho_{\mathrm e}/\varepsilon_0\)
となる。すなわち

\[
\mathbf{S}_{t}^{(d(\ast F))} = \mathbf{S}_{t}^{(\mu_0\ast\mathcal{J})},\quad
\mathbf{S}_{x}^{(d(\ast F))} = \mathbf{S}_{x}^{(\mu_0\ast\mathcal{J})},\quad
\mathbf{S}_{y}^{(d(\ast F))} = \mathbf{S}_{y}^{(\mu_0\ast\mathcal{J})},\quad
\mathbf{S}_{z}^{(d(\ast F))} = \mathbf{S}_{z}^{(\mu_0\ast\mathcal{J})}
\]

\(t\)-スライスの非ゼロ成分を読めば
\(-\mathrm{rot}\,\mathbf{B} + \partial\mathbf{E}/\partial t = -c\mu_0\mathbf{J}\)、すなわち
\(\mathrm{rot}\,\mathbf{B} = c\mu_0\mathbf{J} + \partial\mathbf{E}/\partial t\)
の各成分。\(z\)-スライスの \((x,y)\) 成分からは
\(\mathrm{div}\,\mathbf{E} = \rho_{\mathrm e}/\varepsilon_0\) が現れる。

\begin{center}\rule{0.5\linewidth}{0.5pt}\end{center}

以上で、マクスウェル方程式の全容が \(4\times4\times4\)
のスライス行列の束として可視化された。この「巨大配列」の各マス目に書かれているのは偏微分の組み合わせにすぎず、外微分
\(d\) とホッジ・スター \(\ast\)
という二つの代数操作が、いかに整然とした行列の文法で物理法則を記述しているか------それを見届けたことが、本付録の成果である。

\begin{center}\rule{0.5\linewidth}{0.5pt}\end{center}

\hypertarget{ux4ed8ux9332fux7b2c5ux7ae0ux306e4ux672cux306eux5f0fux3068ux7b2c10ux7ae0ux306e2ux672cux306eux5f0f}{%
\section{付録F：第5章の4本の式と、第10章の2本の式}\label{ux4ed8ux9332fux7b2c5ux7ae0ux306e4ux672cux306eux5f0fux3068ux7b2c10ux7ae0ux306e2ux672cux306eux5f0f}}

第10章本文では、マクスウェル方程式を4次元時空上の2本の式

\[
dF=0,
\qquad
d(\ast F)=\mu_0(\ast\mathcal{J})
\]

として書いた。

しかし第5章では、いきなり4次元時空上の \(2\)-form
にまとめるのではなく、各時刻の空間上で、電場・磁場・電流・電荷をそれぞれ微分形式として扱う見方をしていた。

この付録では、第5章で見た「空間上の4本の式」と、第10章本文で見た「時空上の2本の式」が、同じマクスウェル方程式を別の切り方で書いているだけであることを確認する。

\hypertarget{f.1-ux7a7aux9593ux4e0aux306e-ebjrho_mathrm-e}{%
\subsection{\texorpdfstring{F.1 空間上の
\((E,B,J,\rho_{\mathrm e})\)}{F.1 空間上の (E,B,J,\textbackslash rho\_\{\textbackslash mathrm e\})}}\label{f.1-ux7a7aux9593ux4e0aux306e-ebjrho_mathrm-e}}

各時刻 \(t\) ごとの空間 \((x,y,z)\) を考える。

電場を空間上の \(1\)-form として

\[
E = E_x\,dx + E_y\,dy + E_z\,dz
\]

と置く。

磁場を空間上の \(2\)-form として

\[
B = B_x\,dy\wedge dz + B_y\,dz\wedge dx + B_z\,dx\wedge dy
\]

と置く。

電流密度も、空間を貫く量として \(2\)-form で表す。

\[
J = J_x\,dy\wedge dz + J_y\,dz\wedge dx + J_z\,dx\wedge dy
\]

電荷密度は、空間上の \(3\)-form として

\[
\rho_{\mathrm e}\,dx\wedge dy\wedge dz
\]

と書く。

ここで使う \(d\) と \(\ast\) は、空間3次元の \(d\) と \(\ast\)
である。第10章本文で使った4次元時空上の \(d\) と \(\ast\)
と区別したいときは、

\[
d_3,\quad \ast_3,
\qquad
d_4,\quad \ast_4
\]

と書く。

\hypertarget{f.2-ux7a7aux9593ux4e0aux306e4ux672cux306eux5f0f}{%
\subsection{F.2
空間上の4本の式}\label{f.2-ux7a7aux9593ux4e0aux306e4ux672cux306eux5f0f}}

この記法では、マクスウェル方程式は次の4本になる。

\[
d_3B=0
\]

\[
d_3E+\frac{\partial B}{\partial t}=0
\]

\[
d_3(\ast_3E)=\ast_3\left(\frac{\rho_{\mathrm e}}{\varepsilon_0}\right)
\]

\[
d_3(\ast_3B)=\mu_0 c\,J+\frac{\partial(\ast_3E)}{\partial t}
\]

これが、第5章で見た微分形式スタイルの4方程式である。

ここで右辺は、スカラー場 \(\rho_{\mathrm e}/\varepsilon_0\) を空間上の
\(3\)-form として読むために \(\ast_3\)
を作用させたものである。デカルト座標では、

\[
\ast_3\left(\frac{\rho_{\mathrm e}}{\varepsilon_0}\right)=\frac{\rho_{\mathrm e}}{\varepsilon_0}\,dx\wedge dy\wedge dz
\]

である。

本書では、ベクトル解析の \(\mathrm{div}\)
を最初から基本演算として置くのではなく、\(d\) と \(\ast\) の合成

\[
\mathrm{div}=\ast_3 d_3\ast_3
\]

として読む。そのため、電荷密度のようなスカラー場も、積分される式の右辺では
\(\ast_3\) によって \(3\)-form にしてから比較する。特に3本目は、両辺に
\(\ast_3\) を作用させて

\[
\ast_3d_3(\ast_3E)=\frac{\rho_{\mathrm e}}{\varepsilon_0}
\]

と読む。左辺が、本書の辞書でいう \(\mathrm{div}\,\mathbf E\) である。

ベクトル解析の記法に戻せば、これらはそれぞれ

\[
\mathrm{div}\,\mathbf B=0,
\qquad
\mathrm{rot}\,\mathbf E=-\frac{\partial\mathbf B}{\partial t},
\]

\[
\mathrm{div}\,\mathbf E=\frac{\rho_{\mathrm e}}{\varepsilon_0},
\qquad
\mathrm{rot}\,\mathbf B=\mu_0 c\,\mathbf J+\frac{\partial\mathbf E}{\partial t}
\]

である。

ここでの \(t\) と \(\mathbf B\) は、第10章本文と同じく、すでに

\[
t=ct_{\mathrm{SI}},
\qquad
\mathbf B=c\mathbf B_{\mathrm{SI}}
\]

と正規化された量である。

\begin{quote}
\textbf{注} （計量を使わない書き方もある） 本書では、ベクトル解析の
\(\mathrm{div}\)
を基本演算として出発点にするのではなく、これまで作ってきた辞書に合わせて
\(\mathrm{div}=\ast d\ast\) と読む。ガウスの法則も
\(d_3(\ast_3E) = \ast_3\left(\frac{\rho_{\mathrm e}}{\varepsilon_0}\right)\)
と書いてから、必要に応じて
\(\mathrm{div}\,\mathbf E=\rho_{\mathrm e}/\varepsilon_0\)
と読む。つまり、計量から作った \(\ast\)
を使って、スカラー場・ベクトル場・form
を行き来している。しかし、マクスウェル方程式には「計量を使わない」立場もある。Hehl--Obukhov
の \emph{Foundations of Classical Electrodynamics: Charge, Flux, and
Metric} で強調されるような見方では、まず \(F\) と \(H\) という二つの
form を分け、\(dF=0, dH=J\)
と書く。ここでは、方程式そのものにはホッジ・スターを入れない。計量や媒質の性質は、あとから
\(H\) と \(F\)
を結びつける関係として入る。本書では、この方向には深入りしない。目的はあくまで、これまで作ってきた
\(d\) と \(\ast\)
の辞書を使い、通常のベクトル解析の式がどう成分として現れるかを見届けることである。
\end{quote}

\hypertarget{f.3-d_4f0-ux304bux30892ux672cux304cux51faux308b}{%
\subsection{\texorpdfstring{F.3 \(d_4F=0\)
から2本が出る}{F.3 d\_4F=0 から2本が出る}}\label{f.3-d_4f0-ux304bux30892ux672cux304cux51faux308b}}

第10章本文では、電磁場 \(F\) を

\[
F=-E_x\,dt\wedge dx-E_y\,dt\wedge dy-E_z\,dt\wedge dz+B_x\,dy\wedge dz+B_y\,dz\wedge dx+B_z\,dx\wedge dy
\]

と定義した。

空間上の \(E\) と \(B\) を使えば、これは

\[
F=-dt\wedge E+B
\]

と書ける。

空間上の form \(\alpha(t)\) に対しては、

\[
d_4\alpha = dt\wedge\frac{\partial\alpha}{\partial t}+d_3\alpha
\]

である。

したがって、

\[
d_4F=d_4(-dt\wedge E+B)
\]

を計算すると、

\[
d_4(-dt\wedge E)=dt\wedge d_3E
\]

であり、

\[
d_4B=dt\wedge\frac{\partial B}{\partial t}+d_3B
\]

である。

よって

\[
d_4F=dt\wedge\left(d_3E+\frac{\partial B}{\partial t}\right)+d_3B.
\]

したがって、

\[
d_4F=0
\]

は、\(dt\)
を含む部分と含まない部分をそれぞれゼロにすることと同じであり、

\[
d_3E+\frac{\partial B}{\partial t}=0,
\qquad
d_3B=0
\]

を与える。

つまり、第10章本文の1本目

\[
d_4F=0
\]

は、第5章スタイルでは

\[
d_3B=0,
\qquad
d_3E+\frac{\partial B}{\partial t}=0
\]

という2本の式に分かれる。

\hypertarget{f.4-d_4ast_4fmu_0ast_4mathcal-j-ux304bux3089ux6b8bux308a2ux672cux304cux51faux308b}{%
\subsection{\texorpdfstring{F.4 \(d_4(\ast_4F)=\mu_0(\ast_4\mathcal J)\)
から残り2本が出る}{F.4 d\_4(\textbackslash ast\_4F)=\textbackslash mu\_0(\textbackslash ast\_4\textbackslash mathcal J) から残り2本が出る}}\label{f.4-d_4ast_4fmu_0ast_4mathcal-j-ux304bux3089ux6b8bux308a2ux672cux304cux51faux308b}}

次に、有源側を見る。

第10章本文の規約では、

\[
\ast_4F= B_x\,dt\wedge dx+B_y\,dt\wedge dy+B_z\,dt\wedge dz+E_x\,dy\wedge dz+E_y\,dz\wedge dx+E_z\,dx\wedge dy.
\]

空間上の記法で書けば、これは

\[
\ast_4F=dt\wedge(\ast_3B)+\ast_3E
\]

である。

実際、

\[
\ast_3B=B_x\,dx+B_y\,dy+B_z\,dz
\]

なので、

\[
dt\wedge(\ast_3B)=B_x\,dt\wedge dx+B_y\,dt\wedge dy+B_z\,dt\wedge dz
\]

となる。

また、

\[
\ast_3E=E_x\,dy\wedge dz+E_y\,dz\wedge dx+E_z\,dx\wedge dy
\]

である。

これに \(d_4\) を作用させる。

まず、

\[
d_4\bigl(dt\wedge(\ast_3B)\bigr)=-dt\wedge d_3(\ast_3B)
\]

である。一方、

\[
d_4(\ast_3E)=dt\wedge\frac{\partial(\ast_3E)}{\partial t}+d_3(\ast_3E).
\]

したがって、

\[
d_4(\ast_4F)=dt\wedge\left(\frac{\partial(\ast_3E)}{\partial t}-d_3(\ast_3B)\right)+d_3(\ast_3E).
\]

右辺は、第10章本文の規約では

\[
\mu_0(\ast_4\mathcal J)=\ast_3\left(\frac{\rho_{\mathrm e}}{\varepsilon_0}\right)-\mu_0 c\,dt\wedge J
\]

である。

ここで

\[
J=J_x\,dy\wedge dz+J_y\,dz\wedge dx+J_z\,dx\wedge dy
\]

である。

したがって、\(dt\) を含まない部分を比較すると、

\[
d_3(\ast_3E)=\ast_3\left(\frac{\rho_{\mathrm e}}{\varepsilon_0}\right).
\]

また、\(dt\) を含む部分を比較すると、

\[
\frac{\partial(\ast_3E)}{\partial t}-d_3(\ast_3B)=-\mu_0 c\,J.
\]

これを移項して、

\[
d_3(\ast_3B)=\mu_0 c\,J+\frac{\partial(\ast_3E)}{\partial t}.
\]

したがって、第10章本文の2本目

\[
d_4(\ast_4F)=\mu_0(\ast_4\mathcal J)
\]

は、第5章スタイルでは

\[
d_3(\ast_3E)=\ast_3\left(\frac{\rho_{\mathrm e}}{\varepsilon_0}\right)
\]

と

\[
d_3(\ast_3B)=\mu_0 c\,J+\frac{\partial(\ast_3E)}{\partial t}
\]

の2本に分かれる。

\hypertarget{f.5-ux307eux3068ux3081}{%
\subsection{F.5 まとめ}\label{f.5-ux307eux3068ux3081}}

第5章の微分形式スタイルでは、マクスウェル方程式は空間上の4本の式として現れる。

\[
d_3B=0
\]

\[
d_3E+\frac{\partial B}{\partial t}=0
\]

\[
d_3(\ast_3E)=\ast_3\left(\frac{\rho_{\mathrm e}}{\varepsilon_0}\right)
\]

\[
d_3(\ast_3B)=\mu_0 c\,J+\frac{\partial(\ast_3E)}{\partial t}
\]

一方、第10章本文では、時間を空間と同じく座標の一つとして扱い、電場と磁場を
\(F\) にまとめた。

\[
F=-dt\wedge E+B
\]

すると、上の4本は

\[
d_4F=0
\]

と

\[
d_4(\ast_4F)=\mu_0(\ast_4\mathcal J)
\]

の2本にまとまる。

つまり、4本が2本になるというのは、式を圧縮したというより、空間と時間を分けて見ていたものを、時空上の
\(2\)-form という一つの測定器にまとめ直したということである。

第5章の4本は、空間と時間を分けて眺める書き方。 第10章の2本は、時空の
\(2\)-form として眺める書き方。

どちらも同じマクスウェル方程式を見ている。違うのは、何をひとまとまりの測定器として見るかである。

\newpage

\hypertarget{ux7b2c11ux7ae0ux66f2ux304cux3063ux305fux7a7aux9593ux3078-ux672cux66f8ux306eux5148ux306bux3042ux308bux3082ux306e}{%
\chapter{第11章：曲がった空間へ ------
本書の先にあるもの}\label{ux7b2c11ux7ae0ux66f2ux304cux3063ux305fux7a7aux9593ux3078-ux672cux66f8ux306eux5148ux306bux3042ux308bux3082ux306e}}

\hypertarget{ux3053ux306eux7ae0ux306eux7acbux3061ux4f4dux7f6e-ux3055ux3089ux306bux5148ux3092ux898bux305fux3044ux8aadux8005ux306eux305fux3081ux306eux9053ux6a19}{%
\subsection{§11.0 この章の立ち位置 ------
さらに先を見たい読者のための道標}\label{ux3053ux306eux7ae0ux306eux7acbux3061ux4f4dux7f6e-ux3055ux3089ux306bux5148ux3092ux898bux305fux3044ux8aadux8005ux306eux305fux3081ux306eux9053ux6a19}}

　本書『ナブラ解体新書』の目的は、\(\nabla\) を \(d\) と \(\ast\)
に解体し、ベクトル解析の公式群を見通しよく再構築することだった。第8章・第9章でその目的は達成されている。第10章は「おまけ」として、この枠組みがマクスウェル方程式をいかに簡潔に記述するかを見た。

本章は、そのさらに先------本書を読み終えて「この \(d\) と \(\ast\)
という道具は、もっと広い世界でどう使えるのか」と興味を持った読者のための\textbf{道標}である。完結した説明はしない。専門書への橋渡しとして、何がどう変わり、何が変わらないかの見取り図だけを示す。それからもう一つ------本章では、これまで筆者が意図的に避けてきた数学書的な言い回しを、特に読者に配慮せず、定義もせずに遠慮なく使う。まだまだ世界は広いのだ、ということを感じてほしい。

\begin{quote}
\textbf{注} （建前と本音）
「さらに先を見たい読者のための道標」というのは、一応の建前である。
本音を言えば、詳しい読者に「こいつは抽象数学を分かっていないな」と思われてしまうことへの、著者なりの先回りした回答でもある。
だが同時に、厳密な抽象論をそのまま最初から並べても、初心者の手には負えないだろうという確信もある。
本章はそのような「理解のステップ」としての妥協と、標準的な数学への敬意が混ざった場所である。
\end{quote}

\begin{quote}
\textbf{注} （筆者にとっての楽な章）
本章では、比喩の吟味や、厳密さと分かりやすさのバランス、情報量の配分を考える必要がない。
読者の認知負荷も、ここではあまり気にしなくてよい。
そのうえ論理の流れが明快なので、非常に書きやすかった。
------ブルバキスタイルの論理的堅牢さに、まったくもって敬意を表する。
\end{quote}

\begin{quote}
\textbf{注} （本書の意外な帰結）
本書の直接的な目的は、ベクトル解析を行列形式で書き直し、\(d\) と
\(\ast\) の辞書によって見通しよく再構築することだった。
しかし振り返ってみると、本書は実質的に「成分を直接扱う」という点で、副次的にテンソル解析を行っていたのかもしれない。
テンソル解析の弱点は、計算の過程で \(k\)-form
としての幾何学的な「型」の直感が失われがちなことにある。
本書が行列表示を採用したのは、逆にその型------次数、反対称性、横ベクトルと縦ベクトルの区別------を視覚的に保ちながら、成分計算を可能にするためだった。
つまり本書は、「微分形式の幾何学的直感」と「テンソル解析の計算力」の両方を手に入れるための、中間的なアプローチだった、と言えるかもしれない。
------ただし、どちらの伝統から見ても中途半端な側面は否めないことも認めざるを得ない。
\end{quote}

\hypertarget{ux591aux69d8ux4f53-mathbfgx-ux304bux3089ux59cbux3081ux308b}{%
\subsection{\texorpdfstring{§11.1 多様体 ------ \(\mathbf{g}(x)\)
から始める}{§11.1 多様体 ------ \textbackslash mathbf\{g\}(x) から始める}}\label{ux591aux69d8ux4f53-mathbfgx-ux304bux3089ux59cbux3081ux308b}}

本書では第6章で計量 \(\mathbf{g}\) を導入し、第9章で \(\mathbf{g}\) から
\(\ast\) の辞書を作る手順を練習した。第6章のデカルト座標では
\(\mathbf{g}=I\) という定数行列だったが、第9章の円柱座標や球座標では
\(\mathbf{g}\) の成分はすでに \(r,\rho,\theta\)
に依存していた。では、これを座標表示の都合ではなく、空間そのものの各点に与えられた計量
\(\mathbf{g}(p)\) として見ると何が起きるか。

\(\ast\) の辞書は \(\mathbf{g}\) から作られる。\(\mathbf{g}\)
が場所ごとに変われば、\(\ast\)
の辞書も場所ごとに変わる。\(dx \mapsto \ast(dx) = \cdots\,dy\wedge dz\)
という対応表が、空間の各点で異なる係数を持つようになるのだ。

この「各点で計量 \(\mathbf{g}(p)\) を持ち、その計量から各点の \(\ast\)
を作る」という見方は、リーマン多様体へ進むための最初の模型である。平たい
\(\mathbb{R}^3\) では \(\mathbf{g}\) が単位行列だったから \(\ast\)
の辞書は至って単純だった。一般の座標表示では、\(\mathbf{g}\)
はもはや単位行列とは限らず、場所の関数として現れる。それでも
\(\mathbf{g}\) から \(\ast\)
の辞書を書き下す手順そのものは、第9章でやったのと同じだ------ただ毎回、各点で辞書を引き直す必要があるだけである。

\begin{quote}
\textbf{注} （計量の場所依存と曲率の関係）計量成分 \(g_{ij}(x)\)
が場所に依存すること自体は、ただちに空間が曲がっていることを意味しない。ユークリッド空間を曲線座標（極座標・球座標など）で書けば、平坦な空間でも
\(g_{ij}(x)\)
は場所に依存する。曲がっているかどうかは、計量成分そのものや特定の座標での接続係数
\(\Gamma^i_{jk}\) の見た目ではなく、そこから構成される曲率テンソル
\(R^i{}_{mjk}\) が消えるかどうかで判定される。
\end{quote}

この「場所ごとに座標や計量を扱う」世界を数学的に整理する枠組みが\textbf{多様体（manifold）}である。多様体そのものは、まず局所的に
\(\mathbb{R}^n\)
と見なせる滑らかな空間として定義される。その上に、各点の接空間に非退化な対称双線形形式
\(g_p\) を滑らかに与えたものが\textbf{計量付き多様体}であり、特に
\(g_p\)
が正定値なら\textbf{リーマン多様体}、ローレンツ符号なら\textbf{ローレンツ多様体}になる。

\begin{quote}
\textbf{注}
（「曲がった時空」という言葉について）一般相対性理論の解説では「時空が曲がる」と表現される。個人的なこだわりではあるが、筆者はこの言い回しがあまり好きではない。というのも、任意の一点では局所慣性系を選び、その点で計量をミンコフスキー計量にし、接続係数を消すことができるが、曲率は一般には消えず、有限の近傍全体が平坦になるわけではないからである。曲率は、まさにその局所慣性系を二階以上で貼り合わせるときのずれとして現れる。つまり、単体の「曲がったゴム膜」のようなお絵かきだけでは不十分で、むしろ「計量と接続を通じて、異なる点の接空間をどう比べるかが問題になる」と捉えたほうが正確だと感じている。とはいえこれは筆者が偏屈なのであり、素直に「曲がっている」と考えたほうが直感的だろう。
\end{quote}

\hypertarget{ux5b9aux7fa9-ux30c1ux30e3ux30fcux30c8ux3068ux30a2ux30c8ux30e9ux30b9}{%
\subsubsection{11.1.1 定義 ------
チャートとアトラス}\label{ux5b9aux7fa9-ux30c1ux30e3ux30fcux30c8ux3068ux30a2ux30c8ux30e9ux30b9}}

\(n\) 次元\textbf{位相多様体} \(M\)
とは、ハウスドルフかつ第二可算公理を満たす位相空間であって、任意の点
\(p \in M\) に対し、\(p\) を含む開集合 \(U \subset M\) と、\(U\) から
\(\mathbb{R}^n\) の開集合への同相写像
\(\varphi: U \to \varphi(U) \subset \mathbb{R}^n\)
が存在するものである。この組 \((U, \varphi)\)
を\textbf{チャート}（座標近傍）と呼ぶ。\(\varphi(p) = (x^1(p), \dots, x^n(p))\)
と書いたとき、\(x^i\) を\textbf{局所座標}という。

\(M\) を覆うチャートの族 \(\{(U_\alpha, \varphi_\alpha)\}\)
であって、任意の重なるチャート対に対して\textbf{座標変換}

\[\varphi_\beta \circ \varphi_\alpha^{-1}: \varphi_\alpha(U_\alpha \cap U_\beta) \to \varphi_\beta(U_\alpha \cap U_\beta)\]

が \(C^\infty\)
級であるとき、これを\textbf{滑らかなアトラス}と呼び、\(M\)
は\textbf{滑らかな多様体（$C^\infty$ 多様体）}
になる。座標変換が滑らかであること------この一点が、多様体上で微積分を可能にする根幹である。

本書がこれまで舞台としてきた \(\mathbb{R}^3\) は、恒等写像
\(\mathrm{id}: \mathbb{R}^3 \to \mathbb{R}^3\)
という一枚のチャートで全体を覆える、最も自明な多様体の例だ。

よく知られた別の例として、\(n\) 次元球面
\(S^n = \{x \in \mathbb{R}^{n+1} \mid \|x\| = 1\}\) がある。\(S^n\)
は一枚のチャートでは覆えないが、北極と南極を除いた二枚のチャート（立体射影）でアトラスが構成できる。一般相対性理論の時空はローレンツ多様体------計量の符号が
\((-,+,+,+)\) である4次元多様体------として定式化される。

\hypertarget{ux63a5ux7a7aux9593-ux66f2ux304cux3063ux305fux7a7aux9593ux3067ux306eux5faeux5206ux306eux5b9aux7fa9}{%
\subsubsection{11.1.2 接空間 ------
「曲がった空間での微分」の定義}\label{ux63a5ux7a7aux9593-ux66f2ux304cux3063ux305fux7a7aux9593ux3067ux306eux5faeux5206ux306eux5b9aux7fa9}}

\(\mathbb{R}^n\)
では「ベクトル」を「原点から伸びる矢印」として素朴に定義できた。多様体上ではそうはいかない。球面の表面に矢印を描こうとしても、矢印は曲面からはみ出してしまう。

接空間の構成には二つの同値な方法がある。

\textbf{曲線の同値類による定義。} \(p\) を通る滑らかな曲線
\(\gamma: (-\varepsilon, \varepsilon) \to M\)（\(\gamma(0)=p\)）に対し、チャート
\((U, \varphi)\) を一つ選び、\(\mathbb{R}^n\) における速度ベクトル
\((\varphi \circ \gamma)'(0)\) を計算する。別のチャート \((V, \psi)\)
を選んでも、\((\psi \circ \gamma)'(0) = D(\psi \circ \varphi^{-1})_{\varphi(p)} \cdot (\varphi \circ \gamma)'(0)\)
とヤコビ行列で変換されるから、ある一つのチャートで速度ベクトルが一致すれば、座標変換の微分により任意のチャートでも一致する。したがって、この一致関係はチャートの選び方に依存せず、接ベクトルの同値類を定める。この同値類
\([\gamma]\) が \(p\) における\textbf{接ベクトル}である。

\textbf{方向微分（derivation）による定義。} 接ベクトル \(v \in T_p M\)
を、\(p\) の近傍の滑らかな関数 \(f\) に実数 \(v(f)\)
を対応させる線形写像であって、ライプニッツ則
\(v(fg) = v(f)\,g(p) + f(p)\,v(g)\)
を満たすものとして定義する。この定義はチャートに依存しない内在的なものであり、代数幾何などへの一般化にも耐える。チャート
\((U, \varphi)\) と局所座標 \(x^i\) が与えられれば、

\[\left.\frac{\partial}{\partial x^i}\right|_p (f) = \left.\frac{\partial (f \circ \varphi^{-1})}{\partial x^i}\right|_{\varphi(p)}\]

が \(T_p M\) の基底をなす。\(\dim T_p M = \dim M = n\) である。

二つのベクトル場 \(X, Y\)
が与えられたとき、その\textbf{リー括弧（Lie bracket）}
\([X, Y] = XY - YX\)
は再びベクトル場になる。\([X,Y]^i = \sum_j (X^j \partial_j Y^i - Y^j \partial_j X^i)\)
であり、本書では陽に扱わなかったが、フロベニウスの定理（分布の積分可能性）やリー群論の根幹をなす演算である。

すべての点における接空間の非交和が\textbf{接束}
\(TM = \bigsqcup_{p \in M} T_p M\) である。\(TM\) 自身も \(2n\)
次元の多様体になる。\(M\) 上の\textbf{ベクトル場} \(X\) とは、各点
\(p \in M\) に接ベクトル \(X_p \in T_p M\)
を滑らかに割り当てる写像、すなわち接束の\textbf{切断} \(X: M \to TM\)
にほかならない。

座標変換と基底の変換則にも触れておく。二つのチャート \((U, x^i)\) と
\((V, y^j)\) が重なっているとき、接空間の基底は

\[\frac{\partial}{\partial y^j} = \sum_i \frac{\partial x^i}{\partial y^j}\,\frac{\partial}{\partial x^i}\]

と変換される。この変換行列 \(\partial x^i/\partial y^j\)
は、本書第4章で座標変換のヤコビ行列として繰り返し登場したものだ。接ベクトルの成分はこのヤコビ行列の逆行列で変換される（反変ベクトル）。これに対し、1-form
の成分はヤコビ行列そのもので変換される（共変ベクトル）。この「反変」と「共変」の区別------本書ではベクトルが「縦」、1-form
が「横」として視覚化されてきた------が、多様体論におけるテンソル解析の出発点である。

\hypertarget{ux5faeux5206ux5f62ux5f0f-ux4f59ux63a5ux675fux3068ux30c6ux30f3ux30bdux30ebux5834}{%
\subsubsection{11.1.3 微分形式 ------
余接束とテンソル場}\label{ux5faeux5206ux5f62ux5f0f-ux4f59ux63a5ux675fux3068ux30c6ux30f3ux30bdux30ebux5834}}

接空間 \(T_p M\) の双対空間を\textbf{余接空間} \(T_p^* M\)
と呼ぶ。\(\omega_p \in T_p^* M\) は接ベクトルを食べて実数を返す線形写像
\(\omega_p: T_p M \to \mathbb{R}\)、すなわち\textbf{1-form}である。本書第1章で
\(dx\) を「変位ベクトルから \(x\)
成分を抽出する横ベクトル」と定義したのは、まさにこの双対性の具現化だった。

\begin{quote}
\textbf{注}（本編との対応） これは、第1章で
\(dx=\begin{pmatrix}1&0&0\end{pmatrix}\)
と書いた約束を、標準デカルト座標における余接基底 \(dx^1|_p\)
の成分表示として言い直したものである。第1章では点 \(p\)
への依存を抑えて書いたが、多様体論では各点の余接空間 \(T_p^* M\)
に属する対象として扱う。
\end{quote}

局所座標 \(x^i\) に対し、\(dx^i|_p \in T_p^* M\) を
\(dx^i|_p(\partial/\partial x^j|_p) = \delta^i_j\)
で定めると、\(\{dx^i|_p\}\) は \(T_p^* M\) の
\(\{\partial/\partial x^i|_p\}\) に対する双対基底になる。任意の1-form
\(\omega\) は局所的に \(\omega = \sum_i \omega_i\,dx^i\) と展開できる。

これを \(k\) 重に拡張したものが \textbf{$k$-form} である。多様体上の
\(k\)-form とは、各点 \(p\) に \(\Lambda^k T_p^* M\)
の元を滑らかに割り当てる切断である。つまり、各点では \(k\)
個の接ベクトルを引数にとる交代多重線形写像

\[\omega_p: \underbrace{T_p M \times \cdots \times T_p M}_{k\ \text{個}} \to \mathbb{R}\]

であり、その成分が局所座標で滑らかに変化するものをいう。ウェッジ積
\(\wedge\)
は、二つの形式の引数を「すべての順列にわたって符号付きで和をとる」操作として定義される。本書第2章で反対称行列の成分計算として導入したウェッジ積が、この抽象的な定義のもとで自然に再現される。

\textbf{引き戻し（pullback）}も多様体間の写像へと一般化される。滑らかな写像
\(f: M \to N\) に対し、\(N\) 上の \(k\)-form \(\omega\) の引き戻し
\(f^*\omega\) は \(M\) 上の \(k\)-form であり、

\[(f^*\omega)_p(v_1, \dots, v_k) = \omega_{f(p)}(df_p(v_1), \dots, df_p(v_k))\]

で定義される。\(df_p: T_p M \to T_{f(p)} N\) は \(f\)
の微分（接写像）である。本書第4章で座標変換のヤコビアンとして計算した引き戻しは、この一般公式において
\(f\) をチャート間の座標変換とし、局所座標で成分表示したものに一致する。

本書が一貫して使ってきた「\(dx\)
は横ベクトル」という記法は、実は多様体論の標準言語では \(dx^i\)
たちが余接空間の基底であるという事実にほかならない。\(\omega = \sum_i \omega_i\,dx^i\)
と書いたとき、係数 \(\omega_i\) は本書の言葉では「横ベクトルの第 \(i\)
成分」であり、多様体論の言葉では「1-form
の局所座標に関する成分」である。\(\omega(v) = \sum_i \omega_i\,dx^i(v) = \sum_i \omega_i v^i\)
という計算は、本書第1章の行列の積
\(\begin{pmatrix}\omega_1&\omega_2&\omega_3\end{pmatrix}\begin{pmatrix}v^1\\v^2\\v^3\end{pmatrix}\)
に完全に対応する。

この双対性は高次の形式にも拡張される。本書第2章で \(k\)-form
を行列（\(k=1\) では横ベクトル、\(k=2\) では反対称行列、\(k=3\)
では反対称3階テンソル）で表示したのは、交代多重線形写像としての
\(k\)-form の成分表示にほかならない。付録Eで \(4\times4\times4\)
のスライス行列を扱ったとき、我々はすでに多様体上の3-form
の成分計算を実践していたのだ。

\hypertarget{ux5916ux5faeux5206-d-ux591aux69d8ux4f53ux4e0aux3067ux5909ux308fux3089ux306aux3044ux3082ux306e}{%
\subsubsection{\texorpdfstring{11.1.4 外微分 \(d\) ------
多様体上で変わらないもの}{11.1.4 外微分 d ------ 多様体上で変わらないもの}}\label{ux5916ux5faeux5206-d-ux591aux69d8ux4f53ux4e0aux3067ux5909ux308fux3089ux306aux3044ux3082ux306e}}

外微分 \(d\)
は、多様体上で\textbf{計量を一切使わずに}定義できる数少ない微分演算子の一つである。\(k\)-form
\(\omega\) に対し、\(d\omega\) は \((k+1)\)-form
であり、その定義は本書第5章で偏微分とウェッジ積の組み合わせとして学んだものと完全に一致する。

外微分 \(d\) は、\(\Omega^k(M)\) から \(\Omega^{k+1}(M)\)
への線形作用素であり、関数に対して通常の全微分を与え、次数付きライプニッツ則を満たし、局所座標では偏微分とウェッジ積によって与えられる。この作用素について
\(d^2 = 0\)
が成り立つ。計量はどこにも登場しない。曲がっていようがいまいが、\(d\)
は本書で慣れ親しんだ計算規則をそのまま使える。

この普遍性こそが微分形式の最大の武器であり、\textbf{ストークスの定理}

\[\int_M d\omega = \int_{\partial M} \omega\]

が向き付けられた \(n\) 次元多様体 \(M\)
と、適切な滑らかさ・台条件を満たす \((n-1)\)-form
に対して成立する根拠でもある。本書第8章で「統一ストークスの定理」として出会ったものは、この一般公式の特別な場合（\(M \subset \mathbb{R}^3\),
\(n \le 3\)）だった。

\hypertarget{ux8a08ux91cfux3068ux30dbux30c3ux30b8ux30b9ux30bfux30fc-ux5834ux6240ux3054ux3068ux306bux5909ux308fux308bux3082ux306e}{%
\subsubsection{11.1.5 計量とホッジ・スター ------
場所ごとに変わるもの}\label{ux8a08ux91cfux3068ux30dbux30c3ux30b8ux30b9ux30bfux30fc-ux5834ux6240ux3054ux3068ux306bux5909ux308fux308bux3082ux306e}}

\(d\) が普遍的であるのに対し、\(\ast\)
はそうではない。\textbf{リーマン計量} \(g\) とは、各点 \(p \in M\)
において接空間 \(T_p M\) 上の正定値内積
\(g_p: T_p M \times T_p M \to \mathbb{R}\)
を滑らかに与える2階の共変テンソル場である。局所座標では

\[g = \sum_i\sum_j g_{ij}(x)\,dx^i \otimes dx^j\]

と書け、\(g_{ij}(x)\) は場所の関数である。本書では \(g_{ij}\)
を成分とする行列を \(\mathbf{g}\) と書いてきた。\(\mathbb{R}^3\)
のデカルト座標では \(\mathbf{g}\)
は単位行列、ミンコフスキー時空では対角成分が \((-1,1,1,1)\)
の定数行列だが、一般の多様体では \(\mathbf{g}(x)\) は場所ごとに異なる。

向き付けられたリーマン多様体では、計量から\textbf{体積形式}
\(\mathrm{vol}_g = \sqrt{\det g}\;dx^1 \wedge \cdots \wedge dx^n\)（擬リーマン計量では規約に応じて
\(\sqrt{|\det g|}\)
が現れる）が定まり、これを用いて\textbf{ホッジ・スター作用素}

\[\ast: \Omega^k(M) \to \Omega^{n-k}(M)\]

が定義される。\(\ast\) の定義式
\(\omega \wedge \ast\eta = \langle\omega,\eta\rangle_g\,\mathrm{vol}_g\)
には明示的に \(g\) が入っている。\(\mathbf{g}(x)\)
が場所の関数なら、\(\ast\)
の辞書も場所ごとに異なる------本書第9章で練習した「\(\mathbf{g} \to \ast\)
辞書」の手順を、各点で繰り返す必要がある。擬リーマン計量の場合も同様の構成は可能だが、符号規約に注意が必要である。

この \(d\) と \(\ast\) の非対称性------\(d\) は普遍的、\(\ast\)
は計量依存------こそが、Burke
の「計量遅延」が最終的に照射する構造である。一般相対性理論では、計量
\(g\)
そのものがアインシュタイン方程式の解として動的に決まる。電磁場の方程式
\(dF = 0\), \(d(\ast F) = \mu_0(\ast\mathcal{J})\)
は、時空計量から定まる \(\ast\)
と適切な電流形式を用いれば、曲がった時空上でも同じ形で書ける。

\hypertarget{ux7a4dux5206-1ux306eux5206ux5272ux3068ux5411ux304dux4ed8ux3051ux53efux80fdux6027}{%
\subsubsection{11.1.6 積分 ------
1の分割と向き付け可能性}\label{ux7a4dux5206-1ux306eux5206ux5272ux3068ux5411ux304dux4ed8ux3051ux53efux80fdux6027}}

多様体上の積分もまた、本書のリーマン和の延長線上にある。\(\mathbb{R}^n\)
では、\(n\)-form \(\omega = f\,dx^1 \wedge \cdots \wedge dx^n\)
の積分を素朴に \(\int f\,dx^1 \cdots dx^n\) と定義できた。多様体上で
\(n\)-form
を積分するには、チャートごとに局所積分を計算し、それらを貼り合わせる必要がある。この貼り合わせを可能にするのが\textbf{1の分割（partition of unity）}------多様体を被覆するチャート族に付随する、足して1になる滑らかな関数の族------である。

また、多様体上で \(n\)-form
の積分が座標の取り方に依存しないためには、多様体が\textbf{向き付け可能（orientable）}でなければならない。本書で一貫して右手系
\((x,y,z)\) を仮定してきたのは、\(\mathbb{R}^3\)
に向きを固定していたことに相当する。メビウスの帯のように向き付け不能な多様体では、大域的な体積形式が定義できない。

こうして多様体上の積分論が整備されると、適切な滑らかさ・台条件のもとで、ストークスの定理
\(\int_M d\omega = \int_{\partial M} \omega\) は向き付けられた \(n\)
次元多様体とその上の \((n-1)\)-form
に対して成立する。本書の旅が到達した統一ストークスの定理は、この一般的枠組みの
\(n \le 3\) への適用だったのだ。

\hypertarget{ux63a5ux7d9aux3068ux66f2ux7387-ux4e00ux8a00ux3060ux3051}{%
\subsubsection{11.1.7 接続と曲率 ------
一言だけ}\label{ux63a5ux7d9aux3068ux66f2ux7387-ux4e00ux8a00ux3060ux3051}}

多様体上でベクトル場を「平行移動」するための追加構造が\textbf{接続（connection）}、あるいは\textbf{共変微分}
\(\nabla\)
である。接続が与えられると、曲がり具合を測る\textbf{曲率テンソル} \(R\)
が定義される。一般相対性理論では、レヴィ＝チヴィタ接続（計量と整合的で捩れのない唯一の接続）から得られる曲率が、時空の重力を記述する。

微分形式の言葉では、曲率は\textbf{曲率2-form} \(\Omega\)
として表され、ビアンキ恒等式 \(D\Omega = 0\)（外共変微分で \(\Omega\)
が閉じている）といった構造が現れる。この先は、Flanders や Burke
の専門書に譲る。

\hypertarget{ux672cux66f8ux3068ux306eux5bfeux5fdc}{%
\subsubsection{11.1.8
本書との対応}\label{ux672cux66f8ux3068ux306eux5bfeux5fdc}}

本書の概念が多様体論のどこに対応するかを一覧にしておく。

\begin{longtable}[]{@{}
  >{\raggedright\arraybackslash}p{(\columnwidth - 2\tabcolsep) * \real{0.4000}}
  >{\raggedright\arraybackslash}p{(\columnwidth - 2\tabcolsep) * \real{0.6000}}@{}}
\toprule\noalign{}
\begin{minipage}[b]{\linewidth}\raggedright
本書
\end{minipage} & \begin{minipage}[b]{\linewidth}\raggedright
多様体論
\end{minipage} \\
\midrule\noalign{}
\endhead
\bottomrule\noalign{}
\endlastfoot
\(dx\) は横ベクトル（第1章） & \(dx^i\) は余接空間 \(T_p^*M\) の基底 \\
\(dx(v) = v_x\)（第1章） & 1-form と接ベクトルの双対ペアリング
\(\langle\omega, v\rangle\) \\
ウェッジ積の反対称行列表示（第2章） & 交代多重線形写像としての
\(k\)-form の成分表示 \\
引き戻し＝測定器の焼き直し（第4章） & 滑らかな写像による \(k\)-form
の引き戻し \(f^*\omega\) \\
\(d\) の偏微分＋ウェッジ積（第5章） & 外微分
\(d\)（多様体上で計量不要） \\
\(\mathbf{g} = J^T J\)（第6章） &
ユークリッド空間内の座標変換・埋め込みから誘導される計量の具体例 \\
\(\ast\) の辞書（第6章、第9章） & ホッジ・スター
\(\ast: \Omega^k(M) \to \Omega^{n-k}(M)\)（計量依存） \\
統一ストークスの定理（第8章） &
\(\int_M d\omega = \int_{\partial M} \omega\)（一般次元） \\
\(dF=0\), \(d(\ast F)=\mu_0(\ast\mathcal{J})\)（第10章） &
曲がった時空上でも、時空計量から定まる \(\ast\)
と適切な電流形式を用いれば同じ形で書ける \\
付録Eの \(4\times4\times4\) スライス（第10章） & 4次元多様体上の 3-form
の成分計算 \\
\end{longtable}

\hypertarget{ux30eaux30fcux30deux30f3ux5e7eux4f55ux5b66-ux30c6ux30f3ux30bdux30ebux89e3ux6790ux306eux30b9ux30bfux30a4ux30eb}{%
\subsection{§11.2 リーマン幾何学 ------
テンソル解析のスタイル}\label{ux30eaux30fcux30deux30f3ux5e7eux4f55ux5b66-ux30c6ux30f3ux30bdux30ebux89e3ux6790ux306eux30b9ux30bfux30a4ux30eb}}

多様体論が「幾何学的対象を座標に依存せずに定義する」ことを目指すのに対し、古典的なリーマン幾何学------アインシュタインの一般相対性理論で使われるテンソル解析------は「座標を選び、成分の変換則を正面から扱う」スタイルをとる。本書が行列の成分を明示してきた流儀に近いのは、むしろこちらである。

\hypertarget{ux30c6ux30f3ux30bdux30eb-ux5909ux63dbux5247ux306bux3088ux308bux5b9aux7fa9}{%
\subsubsection{11.2.1 テンソル ------
変換則による定義}\label{ux30c6ux30f3ux30bdux30eb-ux5909ux63dbux5247ux306bux3088ux308bux5b9aux7fa9}}

古典的なテンソル解析では、テンソルを「座標変換に対する変換規則」で定義する。座標
\(x^i\) から \(\bar{x}^i\) への変換に対して、

\textbf{反変ベクトル} \(V^i\) は
\(\bar{V}^i = \sum_j \frac{\partial \bar{x}^i}{\partial x^j} V^j\)
と変換され、 \textbf{共変ベクトル} \(\omega_i\) は
\(\bar{\omega}_i = \sum_j \frac{\partial x^j}{\partial \bar{x}^i} \omega_j\)
と変換される。

この区別は本書第4章で既に出会っている。変位ベクトル \(\mathbf{v}\)
の成分 \(v^i\) は反変（ヤコビ行列の逆行列で変換）、1-form \(\omega\)
の係数 \(\omega_i\) は共変（ヤコビ行列で変換）だった。一般の
\((r,s)\)-テンソル \(T^{i_1\cdots i_r}_{j_1\cdots j_s}\) は、\(r\)
個の反変添字と \(s\)
個の共変添字を持ち、添字ごとに上記の変換則を適用する。

本書で計量を表す行列 \(\mathbf{g}\) の成分 \(g_{ij}\) は
\((0,2)\)-テンソル（二階共変テンソル）である。座標変換に対して
\(g_{ij}\) は

\[\bar{g}_{ij} = \sum_{p,q} \frac{\partial x^p}{\partial \bar{x}^i} \frac{\partial x^q}{\partial \bar{x}^j}\,g_{pq}\]

と変換される。本書第6章で導出した \(\mathbf{g} = J^T J\)
という関係式は、デカルト座標で \(g_{ij} = \delta_{ij}\)
であるときに曲線座標での \(\bar{g}_{ij}\)
を計算する公式にほかならない。一般のリーマン計量は、必ずしも一つの座標変換のヤコビ行列
\(J\) から \(J^T J\) として得られるわけではない。第6章の \(J^T J\)
は、あくまでユークリッド空間内で誘導計量を計算する最も具体的な例である。

\hypertarget{ux30afux30eaux30b9ux30c8ux30c3ux30d5ux30a7ux30ebux8a18ux53f7-ux5faeux5206ux304cux30c6ux30f3ux30bdux30ebux3067ux306aux304fux306aux308bux554fux984c}{%
\subsubsection{11.2.2 クリストッフェル記号 ------
微分がテンソルでなくなる問題}\label{ux30afux30eaux30b9ux30c8ux30c3ux30d5ux30a7ux30ebux8a18ux53f7-ux5faeux5206ux304cux30c6ux30f3ux30bdux30ebux3067ux306aux304fux306aux308bux554fux984c}}

ベクトル場 \(V^i\) の偏微分 \(\partial_j V^i\)
は、そのままではテンソルとして振る舞わない。座標変換すると余分な項が現れてしまう。計量
\(g\)
と整合的で捩れのないレヴィ＝チヴィタ接続を選ぶと、その接続係数、すなわち\textbf{クリストッフェル記号（Christoffel symbols）}は

\[\Gamma^i_{jk} = \frac{1}{2}\sum_m g^{im}(\partial_j g_{km} + \partial_k g_{jm} - \partial_m g_{jk})\]

で与えられる。ここで \(g^{im}\) は \(g_{ij}\)
の逆行列の成分である。本書の用語では、\(\Gamma^i_{jk}\)
は「その座標表示において、基底の変化や計量の変化を補正する係数」であり、一般には
\(g_{ij}\) の一階微分を含む。

クリストッフェル記号を用いて\textbf{共変微分}

\[\nabla_j V^i = \partial_j V^i + \sum_k \Gamma^i_{jk} V^k\]

を定義すると、\(\nabla_j V^i\) は
\((1,1)\)-テンソルとして正しく変換される。共変微分こそが、曲がった空間における「まっすぐ」を定義する道具である。接続の概念（§11.1.7）の成分表示が
\(\Gamma^i_{jk}\) にほかならない。

本書の経験で言い換えよう。第9章では、円柱座標や球座標の \(\mathbf{g}\)
から \(\mathrm{div}\) や \(\mathrm{rot}\) の公式を導出した。第9章の
\(\ast\)
辞書による導出は、クリストッフェル記号を明示的に計算する方法とは異なる。むしろ、\(d\)
と \(\ast\) を使うことで、\(\Gamma^i_{jk}\)
を表に出さずに同じ座標公式へ到達していた、と言うべきである。テンソル解析で同じ公式を導けば、そこには共変微分とクリストッフェル記号が現れる。リーマン幾何学を学ぶとは、同じ座標公式を、接続係数
\(\Gamma^i_{jk}\) と曲率テンソル \(R^i_{\,mjk}\)
の側から見直すことである。

\hypertarget{ux30eaux30fcux30deux30f3ux66f2ux7387ux30c6ux30f3ux30bdux30eb}{%
\subsubsection{11.2.3
リーマン曲率テンソル}\label{ux30eaux30fcux30deux30f3ux66f2ux7387ux30c6ux30f3ux30bdux30eb}}

共変微分の非可換性が、空間の曲がり具合を定量化する。ベクトル場 \(V^i\)
に対して
\[(\nabla_j \nabla_k - \nabla_k \nabla_j) V^i = \sum_m R^i_{\,mjk} V^m\]
として定義される。クリストッフェル記号で具体的に書き下せば、\(R^i_{\,mjk}\)
は \(\Gamma\) の一階微分と \(\Gamma\Gamma\)
の組み合わせで表される。具体的な添字順序と符号は文献の規約に依存するが、いずれにせよ
\(\mathbf{g}\) が定数行列なら \(\Gamma\)
はすべてゼロであり、\(R^i_{\,mjk} = 0\) ------空間は平坦である。

\begin{quote}
\textbf{注} （座標依存計量と曲率）\(\mathbf{g}(x)\)
が場所に依存すること自体は、曲率が非ゼロであることを意味しない。ユークリッド空間を曲線座標（極座標・球座標など）で書けば、\(g_{ij}(x)\)
は場所に依存しても曲率はゼロである。曲率が非ゼロになるのは、\(\Gamma\)
とその微分から構成される \(R^i_{\,mjk}\) が消えない場合である。
\end{quote}

曲率テンソルは \(n^4\)
個の成分を持つが、多数の対称性により独立な成分ははるかに少ない。全共変曲率テンソルを
\(R_{abcd}\) と書くと、代表的な規約では
\[R_{abcd}=-R_{bacd},\qquad R_{abcd}=-R_{abdc},\qquad R_{abcd}=R_{cdab}\]
といった対称性を持つ。さらに第一ビアンキ恒等式も成り立つ。これにより、4次元では独立成分は20個になる。

\begin{quote}
\textbf{注}（曲率テンソルの符号と添字順序）
曲率テンソルの符号と添字順序には複数の規約がある。本項では一つの代表的な規約を示したが、他の文献では添字の順序や符号が異なる場合がある。
\[
\nabla_i R_{mjkl}+\nabla_j R_{mkil}+\nabla_k R_{mijl}=0
\] この恒等式は、外微分の \(d^2=0\)
と同じ「二度境界を取ると消える」型の構造を思わせるが、厳密には接続を含む外共変微分
\(D\) に関する恒等式であり、単なる \(d^2=0\) ではない。
\end{quote}

リッチテンソルは曲率テンソルの一組の添字を縮約して得られる。縮約の具体的な位置や符号は曲率テンソルの規約に依存する。ここでは詳細な規約の固定には立ち入らない。これと\textbf{スカラー曲率}
\(R = \sum_i\sum_j g^{ij} \mathrm{Ric}_{ij}\)
が、アインシュタイン方程式の左辺を構成する。方程式は

\[R_{ij} - \frac{1}{2}R\,g_{ij} = \frac{8\pi G}{c^4}\,T_{ij}\]

である。左辺が時空の幾何（曲率）、右辺が物質のエネルギー運動量テンソル
\(T_{ij}\) を表す。

\begin{quote}
\textbf{注}
（アインシュタイン方程式の符号規約と宇宙項）この式は一つの符号規約と宇宙項
\(\Lambda = 0\)
を選んだ結果である。符号規約（計量の符号、座標系、曲率テンソルの定義）によって式の形は変わり、宇宙項を含む拡張もある。本書では導出も証明もしないが、読者が他の文献を読むときの注意喚起として記す。
\end{quote}

本書はこの式を導出しないし、読者に理解を求めもしない。ただ、この式の背後にある土台------多様体、計量、共変微分、曲率------は、本書で扱った
\(d\) と \(\ast\) だけでは尽くせない。むしろ、第11章で初めて現れた接続
\(\nabla\)
が不可欠である。それでも、微分形式による見方は、これらの構造へ進むための有力な入口になる。

\hypertarget{ux4e8cux3064ux306eux30b9ux30bfux30a4ux30eb-ux672cux66f8ux3068ux306eux95a2ux4fc2}{%
\subsubsection{11.2.4 二つのスタイル ------
本書との関係}\label{ux4e8cux3064ux306eux30b9ux30bfux30a4ux30eb-ux672cux66f8ux3068ux306eux95a2ux4fc2}}

微分形式による記述と、添字を用いたテンソル解析による記述は、どちらもテンソル場を扱うための言語である。前者は反対称テンソルと外微分を前面に出し、後者は成分と変換則を前面に出す。

本書は一貫して成分を明示するスタイル------行列、添字なしの偏微分、ウェッジ積の展開------をとってきた。この意味で、本書の自然な延長は微分形式の抽象論よりもむしろ、テンソル解析の具体的計算にある。すでに第9章で曲線座標の
\(\ast\) 辞書を引いたとき、我々は \(\Gamma^i_{jk}\) を明示せずに
\(\mathrm{rot}\) や \(\mathrm{div}\)
の座標公式へ到達した。テンソル解析で同じ公式を導くなら、その背後には
\(g_{ij}\)、接続係数
\(\Gamma^i_{jk}\)、共変微分の世界が現れる。その先にリーマン幾何学の全容がある。

\hypertarget{ux305dux306eux5148ux3078}{%
\subsection{§11.3 その先へ}\label{ux305dux306eux5148ux3078}}

本書はここで終わりである。第1章から積み上げてきた \(d\) と \(\ast\)
の枠組みは、ベクトル解析の解体という当初の目的を果たした。

ただ、ここまで来ると趣味の領域だが------別の流派の姿を少しだけ覗いてみよう。

現代数理物理には、我々が苦労して分離した外積（\(\wedge\)）と内積（計量）を、最初から「幾何積」として一つの代数に統合する\textbf{幾何代数（Geometric Algebra）}
という世界がある。\(d\) と \(\ast\)
という二つの道具も、そこではより根源的な演算子に吸収される。分離の苦しみを知った者だけが、統合の凄みを味わえる------それだけを伝えて、本書を閉じることにする。

\newpage

\hypertarget{ux7b2c12ux7ae0ux771fux306eux30caux30d6ux30e9-ux30afux30eaux30d5ux30a9ux30fcux30c9ux30d1ux30a6ux30eaux30c7ux30a3ux30e9ux30c3ux30afux30cfux30dfux30ebux30c8ux30f3}{%
\chapter{第12章：真のナブラ ------
クリフォード・パウリ・ディラック・ハミルトン}\label{ux7b2c12ux7ae0ux771fux306eux30caux30d6ux30e9-ux30afux30eaux30d5ux30a9ux30fcux30c9ux30d1ux30a6ux30eaux30c7ux30a3ux30e9ux30c3ux30afux30cfux30dfux30ebux30c8ux30f3}}

\hypertarget{ux672cux30921ux672cux306b}{%
\subsection{§12.0 2本を1本に}\label{ux672cux30921ux672cux306b}}

　第10章の注釈内で、実はこんなことを書いていた------「なぜ2本なのか。\(dF=0\)
と \(d(\ast F)=\mu_0(\ast\mathcal{J})\) を1本にまとめられないのか」と。

できるのだ。それも、あっさりと。虚数単位 \(i\)
を持ち出せば、一行の複素方程式にまとめられる。本章ではまずその「トリック」を見てから、その先にあるもっと統合的な視点------パウリ行列とディラック演算子------へと進む。

\begin{quote}
\textbf{注}
（本章の立ち位置）本章は発展的な補足であり、本書の本線からは外れている。第9章でベクトル解析の書き換えは完了し、第10章でマクスウェル方程式の行列展開も終えた。本章では、厳密な幾何代数（クリフォード代数）の教科書を書くつもりはない。目的は、本書で解体してきた
grad・rot・div
が、パウリ行列とディラック演算子の中で再び一つの演算へ統合される様子を、計算の入口として見ることである。
\end{quote}

\begin{center}\rule{0.5\linewidth}{0.5pt}\end{center}

\hypertarget{ux865aux6570ux3067ux4e00ux672cux306b-ux30deux30afux30b9ux30a6ux30a7ux30ebux65b9ux7a0bux5f0fux306eux7d71ux5408}{%
\subsection{§12.1 虚数で一本に ------
マクスウェル方程式の統合}\label{ux865aux6570ux3067ux4e00ux672cux306b-ux30deux30afux30b9ux30a6ux30a7ux30ebux65b9ux7a0bux5f0fux306eux7d71ux5408}}

第10章で我々は、マクスウェル方程式を

\[dF = 0, \qquad d(\ast F) = \mu_0(\ast\mathcal{J})\]

の2本に集約した。ここで \(F\) も \(\ast F\) も \(2\)-form
であることに注目しよう。4次元時空では、ホッジ・スター \(\ast\) は
\(2\)-form を \(2\)-form に写す。だから \(F\) と \(\ast F\)
は同じ次数の形式だ。

同じ次数なら、足し算に支障はない。虚数単位 \(i\) を使って一つの複素
\(2\)-form を作る（本章では第10章の符号規約に従う。ローレンツ計量上の
Hodge star では \(\ast^2\)
の符号が次元やシグネチャの規約に依存するため、ここでは成分展開との整合性を優先して記述する）。

\[G = F + i\ast F\]

\(G\) に外微分 \(d\) を作用させる。\(d\)
は実数の微分演算子だから、\(d(i\ast F) = i\,d(\ast F)\) である。

\[dG = dF + i\,d(\ast F) = 0 + i\mu_0(\ast\mathcal{J})\]

すなわち

\[d(F + i\ast F) = i\mu_0(\ast\mathcal{J})\]

これだけだ。\(dF=0\) と \(d(\ast F)=\mu_0(\ast\mathcal{J})\)
の2本が、たった一本の複素方程式に統合された。\(F\) と \(\ast F\)
という実部と虚部が、それぞれ無源側と有源側のマクスウェル方程式を分担している。

しかし、ここで立ち止まって考えてほしい。このトリックが成り立つのは、\textbf{たまたま}
4次元で \(\ast\) が \(2\)-form を \(2\)-form に写すからだ。一般の \(n\)
次元では \(\ast\) は \(k\)-form を \((n-k)\)-form に写す。\(2\)-form と
\(2\)-form を足せるのは \(n=4\) の特権であり、ましてや \(1\)-form と
\(2\)-form を直接足すことは、通常の微分形式の枠組みでは許されていない。

\(d\) と \(\ast\)
を組み合わせたとき、grad（\(0\to1\)）、rot（\(1\to1\)）、div（\(1\to0\)）といった異なる次数の間の変換が現れる。これらの元を一つの代数の中でまとめて扱うことができれば、\(\nabla\)
の解体ではなく、\(\nabla\) の\textbf{統合}ができるはずだ。

異なる次数の形式を足し合わせる------そんな「魔法の箱」はないものだろうか。

\begin{center}\rule{0.5\linewidth}{0.5pt}\end{center}

\hypertarget{ux30d1ux30a6ux30eaux884cux5217-ux7570ux306aux308bux6b21ux6570ux3092ux8db3ux3059ux9b54ux6cd5}{%
\subsection{§12.2 パウリ行列 ------
異なる次数を足す魔法}\label{ux30d1ux30a6ux30eaux884cux5217-ux7570ux306aux308bux6b21ux6570ux3092ux8db3ux3059ux9b54ux6cd5}}

「魔法の箱」が、実はある。物理学者たちが量子力学のために発見した道具が、そのままこの魔法の箱として使える。パウリ行列だ。より正確には、ここで見るのは
3次元ユークリッド空間のクリフォード代数を、パウリ行列で表現したものである。微分形式そのものではなく、各次数（0〜3次）の基底をこの代数の中に（本章の規約に沿って）対応させている。

\hypertarget{ux30d1ux30a6ux30eaux884cux5217ux306eux5b9aux7fa9ux3068ux639bux3051ux7b97}{%
\subsubsection{12.2.1
パウリ行列の定義と掛け算}\label{ux30d1ux30a6ux30eaux884cux5217ux306eux5b9aux7fa9ux3068ux639bux3051ux7b97}}

パウリ行列 \(\sigma_1, \sigma_2, \sigma_3\) は次の \(2\times2\)
行列である。

\[
\sigma_1 = \begin{pmatrix}0 & 1 \\ 1 & 0\end{pmatrix},\quad
\sigma_2 = \begin{pmatrix}0 & -i \\ i & 0\end{pmatrix},\quad
\sigma_3 = \begin{pmatrix}1 & 0 \\ 0 & -1\end{pmatrix}
\]

これらの掛け算規則を調べてみよう。単独では
\(\sigma_1^2 = \sigma_2^2 = \sigma_3^2 = I\)（単位行列）であり、異なるもの同士では
\(\sigma_1\sigma_2 = -\sigma_2\sigma_1\)
という\textbf{反可換性}が成り立つ。積は全部で9通りある。

\[
\begin{aligned}
\sigma_1\sigma_1 = I,\quad &\sigma_1\sigma_2 = i\sigma_3,\quad &\sigma_1\sigma_3 = -i\sigma_2 \\
\sigma_2\sigma_1 = -i\sigma_3,\quad &\sigma_2\sigma_2 = I,\quad &\sigma_2\sigma_3 = i\sigma_1 \\
\sigma_3\sigma_1 = i\sigma_2,\quad &\sigma_3\sigma_2 = -i\sigma_1,\quad &\sigma_3\sigma_3 = I
\end{aligned}
\]

この9通りの積を、対称部分と反対称部分に分解してみよう。第2章でウェッジ積
\(\wedge\) を反対称化によって導入したのと同じ発想である。

同じ添字どうし（\(\sigma_1\sigma_1\), \(\sigma_2\sigma_2\),
\(\sigma_3\sigma_3\)）の対称部分は \(I\)、反対称部分は \(0\)
だ。異なる添字の場合は、たとえば \(\sigma_1\sigma_2 = i\sigma_3\)
について、

\[
\frac{1}{2}(\sigma_1\sigma_2 + \sigma_2\sigma_1) = \frac{i\sigma_3 + (-i\sigma_3)}{2} = 0,\qquad
\frac{1}{2}(\sigma_1\sigma_2 - \sigma_2\sigma_1) = \frac{i\sigma_3 - (-i\sigma_3)}{2} = i\sigma_3
\]

という具合に、異なる添字の積は純粋に反対称である。対称部分は添字が同じときだけ
\(I\)、それ以外は
\(0\)------これは第6章で導入した計量（デカルト座標では単位行列）が、パウリ行列の積の対称部分として姿を現したことにほかならない。反対称部分は添字を入れ替えると符号が反転する------第2章のウェッジ積
\(dx\wedge dy = -dy\wedge dx\) とまったく同じ構造だ。

パウリ行列の積は、つねに「内積（対称部分）\(+\)
ウェッジ積（反対称部分）」に分解される。この一点が、パウリ行列のすべての魔法の源である。

\hypertarget{ux30d9ux30afux30c8ux30ebux3092ux30d1ux30a6ux30eaux884cux5217ux3067ux66f8ux304f}{%
\subsubsection{12.2.2
ベクトルをパウリ行列で書く}\label{ux30d9ux30afux30c8ux30ebux3092ux30d1ux30a6ux30eaux884cux5217ux3067ux66f8ux304f}}

3次元ベクトル \(\mathbf{v} = (v_1, v_2, v_3)\)
をパウリ行列で表してみよう。

\[V = v_1\sigma_1 + v_2\sigma_2 + v_3\sigma_3 = \mathbf{v}\!\cdot\!\bm{\sigma}\]

二つのベクトル \(\mathbf{v}, \mathbf{w}\) に対応するパウリ行列
\(V = \mathbf{v}\!\cdot\!\bm{\sigma}\),
\(W = \mathbf{w}\!\cdot\!\bm{\sigma}\) の積を計算する。

\[
\begin{aligned}
VW &= (v_1\sigma_1 + v_2\sigma_2 + v_3\sigma_3)(w_1\sigma_1 + w_2\sigma_2 + w_3\sigma_3) \\[4pt]
&= v_1w_1\sigma_1\sigma_1 + v_1w_2\sigma_1\sigma_2 + v_1w_3\sigma_1\sigma_3 \\
&\quad + v_2w_1\sigma_2\sigma_1 + v_2w_2\sigma_2\sigma_2 + v_2w_3\sigma_2\sigma_3 \\
&\quad + v_3w_1\sigma_3\sigma_1 + v_3w_2\sigma_3\sigma_2 + v_3w_3\sigma_3\sigma_3
\end{aligned}
\]

上記の9通りの積を各項に適用する。非ゼロの置き換えだけ書けば

\[
\sigma_1\sigma_1 = I,\quad \sigma_2\sigma_2 = I,\quad \sigma_3\sigma_3 = I,\quad
\sigma_1\sigma_2 = i\sigma_3,\quad \sigma_2\sigma_3 = i\sigma_1,\quad \sigma_3\sigma_1 = i\sigma_2
\]

と、これらの添字を逆順にした \(\sigma_2\sigma_1 = -i\sigma_3\),
\(\sigma_3\sigma_2 = -i\sigma_1\), \(\sigma_1\sigma_3 = -i\sigma_2\)
の計9通りである。\(I\) の係数（対称部分＝内積）と
\(\sigma_1,\sigma_2,\sigma_3\)
の係数（反対称部分＝ウェッジ積）をそれぞれ集める。

\[
\begin{aligned}
I\text{ の係数（内積）} &: v_1w_1 + v_2w_2 + v_3w_3 \;=\; \mathbf{v}\!\cdot\!\mathbf{w} \\
\sigma_1\text{ の係数（$\wedge$）} &: i\,(v_2w_3 - v_3w_2) \\
\sigma_2\text{ の係数（$\wedge$）} &: i\,(v_3w_1 - v_1w_3) \\
\sigma_3\text{ の係数（$\wedge$）} &: i\,(v_1w_2 - v_2w_1)
\end{aligned}
\]

\(\sigma_1,\sigma_2,\sigma_3\)
の係数カッコ内は、第2章で定義したウェッジ積
\(\mathbf{v}\wedge\mathbf{w}\) の各成分そのものである。したがって

\[VW = (\mathbf{v}\!\cdot\!\mathbf{w})\,I \;+\; i\,(\mathbf{v}\times\mathbf{w})\!\cdot\!\bm{\sigma}\]

と、内積 \(+\)
ウェッジ積の形にまとまる。内積と外積が、\textbf{一つの積 $VW$ から同時に}現れた。

これが「魔法の箱」の正体だ。より正確には、ここで見ているのは3次元ユークリッド空間のクリフォード代数を、パウリ行列で表現したものである。その積は、内積（\(0\)-form
に相当するスカラー）と外積（\(2\)-form
に相当するウェッジ積）を、\textbf{一つの演算で同時に}生み出す。通常の微分形式では別々の次数に属していたものが、この行列表現の中では一つの
\(2\times2\) 行列として共存している。

\begin{quote}
\textbf{注} （なぜ \(2\times2\) なのか）ベクトルを \(2\times2\)
行列で表すのは奇妙に思えるかもしれない。しかし、内積とウェッジ積の両方を一つの積で表現するには、少なくとも非可換な代数が必要であり、その最小の実現が
\(2\times2\)
複素行列なのである。本書の流儀でいえば、\(1\)-form（横ベクトル
\(1\times3\)）と \(2\)-form（反対称 \(3\times3\)
行列）という次数の異なる二つの対象を、一つの \(2\times2\)
行列の中に押し込めた------そう考えればよい。
\end{quote}

\hypertarget{ux30b9ux30abux30e9ux30fcux3082ux6df7ux305cux3066ux8db3ux3059}{%
\subsubsection{12.2.3
スカラーも混ぜて足す}\label{ux30b9ux30abux30e9ux30fcux3082ux6df7ux305cux3066ux8db3ux3059}}

ここまでで、ベクトル（\(1\)-form）とウェッジ積（\(2\)-form）が \(VW\)
という一通りの計算で同時に出てくることがわかった。さらに一歩進めて、スカラー（\(0\)-form）そのものも同じ代数に混ぜよう。

パウリ行列の線形結合で表される最も一般的な \(2\times2\) 行列は

\[\psi = \varphi\,I + A_1\sigma_1 + A_2\sigma_2 + A_3\sigma_3\]

の形をしている。\(\varphi\) はスカラー（\(0\)-form）、\(A_1,A_2,A_3\)
はベクトルの成分（\(1\)-form）である。さらに \(\psi\) に左から
\(\sigma_i\) を掛ければ、\(2\)-form に相当する \(i\sigma_k\)
の項まで現れる。つまり、本章の規約では、\(0,1,2,3\)-form
の各次数の基底をこの代数の中に対応させられる。

\begin{quote}
\textbf{注}
（何ができるようになったのか）微分形式でも異なる次数の形式を直和として形式的に並べることはできる。しかし、内積（縮約）と外積（反対称化）を\textbf{一つの積として同時に}扱うには、通常のウェッジ積だけでは足りない。パウリ行列の代数では、それらが一つの積
\(VW\) の中で自然に分離・共存する。これが幾何代数の核心である。
\end{quote}

\begin{quote}
\textbf{注}
（全次数の対応）本章の規約における次数対応を整理しておく。\(I\) が
\(0\)-form（スカラー）、\(\sigma_1,\sigma_2,\sigma_3\) が
\(1\)-form（ベクトル）、\(i\sigma_1,i\sigma_2,i\sigma_3\) が
\(2\)-form（バイベクトル）、\(iI\) が
\(3\)-form（擬スカラー）に対応する。これは微分形式そのものとの同一視ではなく、3次元ユークリッドのクリフォード代数の行列表現の中に、各次数の基底を対応させる約束である。最も一般的な元は
\(\psi = aI + \mathbf{v}\!\cdot\!\bm{\sigma} + i\,\mathbf{w}\!\cdot\!\bm{\sigma} + ibI\)
の形で、4つの次数すべてを複素係数で含んでいる。
\end{quote}

\begin{quote}
\textbf{注}
（四元数------すでにあった魔法の箱）実はこの「スカラーとベクトルを一つの数に混ぜる」という発想は、パウリ行列よりもずっと古い。行列代数よりも古い。1843年にハミルトンが発見した\textbf{四元数（quaternion）}
\(q = a + bi + cj + dk\)
がまさにそれだ。マクスウェルは自身の電磁気学を四元数で定式化していた。しかしその後、ギブズとヘヴィサイドが四元数からスカラー部分とベクトル部分を\textbf{分離}し、現代のベクトル解析（内積とクロス積）が誕生した。本書の視点からは、パウリ行列は、いったん内積・外積へ分かれた構造を、量子力学の文脈で再び一つの非可換積として眺め直す道具にも見える。
\end{quote}

\begin{quote}
\textbf{注}
（行列なき時代のベクトル解析）ギブズとヘヴィサイドがベクトル解析を整備した1880年代には、まだ行列代数が一般的ではなかった。行列式はあったが、行列そのものの代数------積の非可換性、固有値、線形変換としての統一的理解------は、ケイリーの先駆的研究を経て、20世紀初頭のヒルベルトらによってようやく整備される。ギブズは線形変換を表すために\textbf{ダイアド（dyad）}
\(\mathbf{a}\mathbf{b}\) という独自の記法を発明したが、これは現代の行列
\(\mathbf{a}\mathbf{b}^T\)
に相当する。本書の視点から見ると、ベクトル解析は四元数的な統合から、内積・外積・ダイアド的な道具へと分化していったものとして眺めることができる。本書が「ナブラ解体新書」と題して
\(\nabla\) を \(d\) と \(\ast\)
に分解し、最終章でパウリ行列による再統合を見せたことには、この分離と統合を繰り返してきた道具立てを、別の角度から見直す意味もある。
\end{quote}

\begin{quote}
\textbf{注}
（行列代数が物理学者の標準語でなかったこと）なお、行列代数が物理学者の標準語でなかったことは、1925年の行列力学にもよく現れている。ハイゼンベルク自身は当初、自分の計算が行列の非可換積であることを知らず、それを行列として読んだのはボルンだった。
\end{quote}

\begin{center}\rule{0.5\linewidth}{0.5pt}\end{center}

\hypertarget{ux30c7ux30a3ux30e9ux30c3ux30afux6f14ux7b97ux5b50ux3068ux771fux306eux30caux30d6ux30e9}{%
\subsection{§12.3
ディラック演算子と「真のナブラ」}\label{ux30c7ux30a3ux30e9ux30c3ux30afux6f14ux7b97ux5b50ux3068ux771fux306eux30caux30d6ux30e9}}

\hypertarget{bmsigmacdotnabla-ux7d71ux5408ux7684ux306aux5faeux5206ux6f14ux7b97ux5b50}{%
\subsubsection{\texorpdfstring{12.3.1 \(\bm{\sigma}\!\cdot\!\nabla\)
------
統合的な微分演算子}{12.3.1 \textbackslash bm\{\textbackslash sigma\}\textbackslash!\textbackslash cdot\textbackslash!\textbackslash nabla ------ 統合的な微分演算子}}\label{bmsigmacdotnabla-ux7d71ux5408ux7684ux306aux5faeux5206ux6f14ux7b97ux5b50}}

パウリ行列とナブラ
\(\nabla = (\frac{\partial}{\partial x}, \frac{\partial}{\partial y}, \frac{\partial}{\partial z})\)
を組み合わせて、次の演算子を定義する。

\[D = \sigma_1\frac{\partial}{\partial x} + \sigma_2\frac{\partial}{\partial y} + \sigma_3\frac{\partial}{\partial z} = \bm{\sigma}\!\cdot\!\nabla\]

\(D\) は \(2\times2\) 行列の形をした微分演算子である。これをスカラー場
\(\varphi\) に作用させてみよう。

\[D\varphi = \sigma_1\frac{\partial\varphi}{\partial x} + \sigma_2\frac{\partial\varphi}{\partial y} + \sigma_3\frac{\partial\varphi}{\partial z} = (\nabla\varphi)\!\cdot\!\bm{\sigma}\]

これは \(\mathrm{grad}\,\varphi\) をパウリ行列で表したものだ。

次に、ベクトル場 \(\mathbf{A} = (A_1, A_2, A_3)\) をパウリ行列で
\(A = \mathbf{A}\!\cdot\!\bm{\sigma}\) と書いて \(D\)
を作用させる。§12.2.2 の積の公式を思い出そう。

\[
\begin{aligned}
D A &= (\sigma_1\tfrac{\partial}{\partial x} + \sigma_2\tfrac{\partial}{\partial y} + \sigma_3\tfrac{\partial}{\partial z})(A_1\sigma_1 + A_2\sigma_2 + A_3\sigma_3) \\[4pt]
&= \tfrac{\partial A_1}{\partial x}\,\sigma_1\sigma_1 + \tfrac{\partial A_2}{\partial x}\,\sigma_1\sigma_2 + \tfrac{\partial A_3}{\partial x}\,\sigma_1\sigma_3 \\
&\quad + \tfrac{\partial A_1}{\partial y}\,\sigma_2\sigma_1 + \tfrac{\partial A_2}{\partial y}\,\sigma_2\sigma_2 + \tfrac{\partial A_3}{\partial y}\,\sigma_2\sigma_3 \\
&\quad + \tfrac{\partial A_1}{\partial z}\,\sigma_3\sigma_1 + \tfrac{\partial A_2}{\partial z}\,\sigma_3\sigma_2 + \tfrac{\partial A_3}{\partial z}\,\sigma_3\sigma_3
\end{aligned}
\]

§12.2.2 と同じ \(\sigma_i\sigma_j\) の置き換えを9項に適用し、\(I\)
の係数（内積＝発散）と \(\sigma_1,\sigma_2,\sigma_3\)
の係数（ウェッジ積＝回転）を集める。

\[
\begin{aligned}
I\text{ の係数} &: \frac{\partial A_1}{\partial x} + \frac{\partial A_2}{\partial y} + \frac{\partial A_3}{\partial z} = \mathrm{div}\,\mathbf{A} \\
\sigma_1\text{ の係数（$\wedge$）} &: i\,\left(\frac{\partial A_3}{\partial y} - \frac{\partial A_2}{\partial z}\right) \\
\sigma_2\text{ の係数（$\wedge$）} &: i\,\left(\frac{\partial A_1}{\partial z} - \frac{\partial A_3}{\partial x}\right) \\
\sigma_3\text{ の係数（$\wedge$）} &: i\,\left(\frac{\partial A_2}{\partial x} - \frac{\partial A_1}{\partial y}\right)
\end{aligned}
\]

\(\sigma_k\) の係数カッコ内は \(\mathrm{rot}\,\mathbf{A}\)
の各成分------外微分 \(d\) が \(1\)-form \(\mathbf{A}\) から生み出す
\(2\)-form \(d\mathbf{A}\) の成分------にほかならない。したがって

\[D(\mathbf{A}\!\cdot\!\bm{\sigma}) = (\mathrm{div}\,\mathbf{A})\,I \;+\; i\,(\mathrm{rot}\,\mathbf{A})\!\cdot\!\bm{\sigma}\]

\textbf{一つの演算子 $D$ が、$\mathrm{grad}$（$D\varphi$）、$\mathrm{div}$（$DA$ の実部）、$\mathrm{rot}$（$DA$ の虚部）を同時に導き出す。}
これこそが、\(\nabla\) を \(d\) と \(\ast\)
に分解して理解を深めてきた旅の、一つの終着点------統合------にほかならない。

\begin{quote}
\textbf{注} （\(d\) と \(\ast\) で書くと）\(D\)
は本書の言葉でも表現できる。符号規約を別にすれば、\(D\) は
\(d\)（次数を上げる）と
\(\ast d\ast\)（次数を下げる）を合わせた演算子として理解できる。\(\ast d\ast\)
は\textbf{余微分（codifferential）} \(\delta\)
と呼ばれ、\(\delta = \pm\ast d\ast\)（符号は次元と次数に依存）と定義される。ここでは厳密な符号には立ち入らず、\(D \sim d + \delta\)
という構造だけを押さえておく。\(d\) が次数を上げ、\(\delta\)
が次数を下げる------この二つを足した \(D\) が、grad・rot・div
を一つの演算子に束ねているのだ。
\end{quote}

\hypertarget{d2-ux30e9ux30d7ux30e9ux30b7ux30a2ux30f3}{%
\subsubsection{\texorpdfstring{12.3.2 \(D^2\) ------
ラプラシアン}{12.3.2 D\^{}2 ------ ラプラシアン}}\label{d2-ux30e9ux30d7ux30e9ux30b7ux30a2ux30f3}}

さらに \(D\) を二度作用させるとどうなるか。

\[D^2 = (\bm{\sigma}\!\cdot\!\nabla)(\bm{\sigma}\!\cdot\!\nabla) = (\sigma_1\frac{\partial}{\partial x} + \sigma_2\frac{\partial}{\partial y} + \sigma_3\frac{\partial}{\partial z})(\sigma_1\frac{\partial}{\partial x} + \sigma_2\frac{\partial}{\partial y} + \sigma_3\frac{\partial}{\partial z})\]

9項に展開する。

\[
\begin{aligned}
D^2 &= \sigma_1\sigma_1\,\frac{\partial^2}{\partial x^2} + \sigma_1\sigma_2\,\frac{\partial^2}{\partial x\partial y} + \sigma_1\sigma_3\,\frac{\partial^2}{\partial x\partial z} \\
&\quad + \sigma_2\sigma_1\,\frac{\partial^2}{\partial y\partial x} + \sigma_2\sigma_2\,\frac{\partial^2}{\partial y^2} + \sigma_2\sigma_3\,\frac{\partial^2}{\partial y\partial z} \\
&\quad + \sigma_3\sigma_1\,\frac{\partial^2}{\partial z\partial x} + \sigma_3\sigma_2\,\frac{\partial^2}{\partial z\partial y} + \sigma_3\sigma_3\,\frac{\partial^2}{\partial z^2}
\end{aligned}
\]

\(\sigma_i\sigma_j\) の反対称部分（\(i\neq j\) の項、すなわち
\(\sigma_1\sigma_2 = i\sigma_3\) など）を含む項は、偏微分の可換性
\(\frac{\partial^2}{\partial x\partial y} = \frac{\partial^2}{\partial y\partial x}\)
により、符号が逆のペア同士で打ち消し合う。たとえば
\(\sigma_1\sigma_2\,\frac{\partial^2}{\partial x\partial y}\) と
\(\sigma_2\sigma_1\,\frac{\partial^2}{\partial y\partial x}\)
は、\(\sigma_2\sigma_1 = -\sigma_1\sigma_2\) かつ
\(\frac{\partial^2}{\partial y\partial x} = \frac{\partial^2}{\partial x\partial y}\)
だから、足してゼロ。残るのは \(\sigma_i\sigma_i = I\)
の対角項だけである。

\[D^2 = \left(\frac{\partial^2}{\partial x^2} + \frac{\partial^2}{\partial y^2} + \frac{\partial^2}{\partial z^2}\right)I = \nabla^2 I\]

\(D^2\) はラプラシアンそのものである。\(D\)
はラプラシアンの「平方根」なのだ。この事実------微分演算子の平方根をとる------が、ディラックが相対論的量子力学で電子の方程式を導いたときの核心的な洞察だった。

\hypertarget{ux30deux30afux30b9ux30a6ux30a7ux30ebux3092ux4e00ux672cux306b-ux518dux8a2a}{%
\subsubsection{12.3.3 マクスウェルを一本に ------
再訪}\label{ux30deux30afux30b9ux30a6ux30a7ux30ebux3092ux4e00ux672cux306b-ux518dux8a2a}}

§12.1 では \(F + i\ast F\)
によって4次元でマクスウェル方程式を一本化した。\(D\)
を使えば、同じことが3次元の言葉だけで書ける。

電場 \(\mathbf{E}\) と磁場 \(\mathbf{B}\)
を一つの複素ベクトルに束ねる。これは\textbf{リーマン－ジルバーシュタイン・ベクトル}と呼ばれる。

\[\mathbf{F} = \mathbf{E} + i\mathbf{B}\]

パウリ行列で書けば
\(F = \mathbf{E}\!\cdot\!\bm{\sigma} + i\,\mathbf{B}\!\cdot\!\bm{\sigma}\)
である。\(F\)
はスカラー部を持たない純粋な複素パウリ・ベクトルだ。時間微分を含めた演算子
\(D + \frac{\partial}{\partial t}\) を作用させる。

\[(D + \frac{\partial}{\partial t})F = (\bm{\sigma}\!\cdot\!\nabla + \frac{\partial}{\partial t})(\mathbf{E}\!\cdot\!\bm{\sigma} + i\,\mathbf{B}\!\cdot\!\bm{\sigma})\]

§12.3.1 の \(D(\mathbf{A}\!\cdot\!\bm{\sigma})\) の公式を \(\mathbf{E}\)
と \(\mathbf{B}\) の両方に適用する。

\[
\begin{aligned}
(D + \frac{\partial}{\partial t})F &= (\mathrm{div}\,\mathbf{E})\,I + i\,(\mathrm{rot}\,\mathbf{E})\!\cdot\!\bm{\sigma} + \frac{\partial}{\partial t}\mathbf{E}\!\cdot\!\bm{\sigma} \\
&\quad + i\,(\mathrm{div}\,\mathbf{B})\,I - (\mathrm{rot}\,\mathbf{B})\!\cdot\!\bm{\sigma} + i\,\frac{\partial}{\partial t}\mathbf{B}\!\cdot\!\bm{\sigma}
\end{aligned}
\]

実部と虚部、そして \(I\)（スカラー）と
\(\bm{\sigma}\)（ベクトル）の各成分に整理する。

\[
\begin{aligned}
\text{実部・}I &: \mathrm{div}\,\mathbf{E} - 0 \\
\text{実部・}\bm{\sigma} &: \frac{\partial}{\partial t}\mathbf{E} - \mathrm{rot}\,\mathbf{B} \\
\text{虚部・}I &: 0 + \mathrm{div}\,\mathbf{B} \\
\text{虚部・}\bm{\sigma} &: \mathrm{rot}\,\mathbf{E} + \frac{\partial}{\partial t}\mathbf{B}
\end{aligned}
\]

これらをゼロとおけば、真空のマクスウェル方程式4本がすべて現れる。電流
\(\mathbf{J}\) と電荷 \(\rho\)
がある場合は、第10章と同じ正規化のもとで右辺を
\(\rho/\epsilon_0 - \mu_0 c \mathbf{J}\!\cdot\!\bm{\sigma}\)
と置けば、源のあるマクスウェル方程式が一行にまとまる。

\[(D + \frac{\partial}{\partial t})F = \frac{\rho}{\varepsilon_0} - \mu_0 c \mathbf{J}\!\cdot\!\bm{\sigma}\]

\begin{quote}
\textbf{注} （第12章での正規化と符号） ここでの
\(t, \mathbf{B}, \rho, \mathbf{J}\) は第10章 §10.1--10.5
と同様に正規化された量である。右辺の源の係数と符号は、第10章の
\(d(\ast F) = \mu_0(\ast\mathcal{J})\)
の展開結果と整合するように選ばれている。
\end{quote}

\(dF=0\) と \(d(\ast F)=\mu_0(\ast\mathcal{J})\) の2本が、\(D\)
の前ではたった一本の複素方程式に統合された。§12.1
の4次元トリックが特殊な次元に依存していたのに対し、こちらは \(D\)
という統合された演算子そのものが統合を実現している。

\hypertarget{cancelpartial-f-j-4ux6b21ux5143ux6642ux7a7aux306eux30c7ux30a3ux30e9ux30c3ux30afux6f14ux7b97ux5b50}{%
\subsubsection{\texorpdfstring{12.3.4 \(\cancel{\partial} F = J\) ------
4次元時空のディラック演算子}{12.3.4 \textbackslash cancel\{\textbackslash partial\} F = J ------ 4次元時空のディラック演算子}}\label{cancelpartial-f-j-4ux6b21ux5143ux6642ux7a7aux306eux30c7ux30a3ux30e9ux30c3ux30afux6f14ux7b97ux5b50}}

§12.3.3 の
\((D + \frac{\partial}{\partial t})F = \rho + \mathbf{J}\!\cdot\!\bm{\sigma}\)
には、時間と空間が分離して現れている。これは3次元のパウリ行列
\(\sigma_1,\sigma_2,\sigma_3\)
をベースに組み立てたからだ。4次元時空を最初から扱えば、この分離すら消える。

パウリ行列を \(4\times4\) に拡張したものが\textbf{ガンマ行列}
\(\gamma^0,\gamma^1,\gamma^2,\gamma^3\) である。これらは反可換性
\(\gamma^\mu\gamma^\nu + \gamma^\nu\gamma^\mu = 2g^{\mu\nu}I\)
を満たす（\(g^{\mu\nu}\)
はミンコフスキー計量）。この式に現れる実数係数の符号は、選んだ計量署名や添字の上下の規約に依存する。本章では、第10章と同じく時空計量の署名を
\((-,+,+,+)\)
とした規約を採用する。このガンマ行列を用いて、4次元のディラック演算子を

\[\cancel{\partial} = \gamma^0\frac{\partial}{\partial t} + \gamma^1\frac{\partial}{\partial x} + \gamma^2\frac{\partial}{\partial y} + \gamma^3\frac{\partial}{\partial z}\]

と定義する。第10章の電磁場 \(F\)（\(4\times4\)
反対称行列）をガンマ行列のウェッジ積 \(\gamma^\mu\!\wedge\!\gamma^\nu\)
で展開したものと同一視すれば、マクスウェル方程式は

\[\cancel{\partial} F = J\]

の一行にまとめられる、というのが幾何代数側の標準的な見取り図である。\(F\)
は電磁場、\(J\) は4元電流（\(\gamma^\mu\) で展開されたもの）。真空では
\(\cancel{\partial} F = 0\)。ここでは詳細な符号規約や単位系の換算までは扱わず、第10章までの演算が一つの演算子へ束ねられることだけを見る。

\begin{quote}
\textbf{注} （\(\cancel{\partial} F\)
の積について）ここでの積は通常の行列積ではなく、ガンマ行列が生成する幾何代数の積である。第10章の反対称行列
\(F\) は、ここでは \(\gamma^\mu\!\wedge\!\gamma^\nu\)
で展開された2-vectorとして再解釈される。細部は幾何代数の教科書に譲るが、\(\sigma_i\sigma_j = \delta_{ij}I + i\varepsilon_{ijk}\sigma_k\)（§12.2.1）の4次元版がここでも成り立っていると考えればよい。
\end{quote}

§12.1 の \(F + i\ast F\) も、§12.3.3 の
\((D + \frac{\partial}{\partial t})F\)
も、すべてこの一行に吸収される。これが、\(\nabla\)
の解体から始まった旅の、最も遠くにある風景である。ディラックが1928年に電子の方程式を導くために発見したこの演算子は、時空の幾何と物理法則を一つの代数で記述する\textbf{幾何代数（Geometric Algebra）}の核であり、Doran
\& Lasenby の Geometric Algebra for Physicists に詳しい。

\hypertarget{ux6f14ux7b97ux5b50nabla}{%
\subsection{\texorpdfstring{§12.4
演算子\(\nabla\)}{§12.4 演算子\textbackslash nabla}}\label{ux6f14ux7b97ux5b50nabla}}

さあ、ここで最初の問いに戻ろう。\(\nabla\) とは何だったのか。

ベクトル解析では、\(\nabla\)
は「偏微分を並べた形式的なベクトル」として導入され、grad・rot・div
の三つの顔を持つ謎めいた記号だった。本書はそれを \(d\) と \(\ast\)
に\textbf{解体}することで、各演算の正体を明らかにしてきた。

そして今、これを逆方向に辿り、\(d\) と \(\ast\)
というバラバラだった部品が、幾何代数の中で3次元のディラック演算子\(D = \bm{\sigma}\!\cdot\!\nabla\)
として\textbf{統合}され、更にその4次元版のディラック演算子\(\cancel{\partial}\)が、マクスウェル方程式のより統合的な表現
\(\cancel{\partial} F = J\) を与えた。

ところで、幾何代数の文脈では、この演算子を必ずしも\(\cancel{\partial}\)
とは書かない。むしろ、この種のベクトル微分演算子を

\[
\nabla
\]

と書く。

\newpage

\hypertarget{ux304aux308fux308aux306bux30caux30d6ux30e9ux89e3ux4f53ux65b0ux66f8ux306fux3044ux304bux306bux3057ux3066ux751fux307eux308cux305fux304b}{%
\chapter{おわりに：『ナブラ解体新書』はいかにして生まれたか}\label{ux304aux308fux308aux306bux30caux30d6ux30e9ux89e3ux4f53ux65b0ux66f8ux306fux3044ux304bux306bux3057ux3066ux751fux307eux308cux305fux304b}}

長い旅路を終えた読者の皆さんに、まずは心からの敬意を表したいと思います。

本書を閉じるにあたり、この奇妙で、いささか強引な構成を持つ『ナブラ解体新書』がいかにして生まれたのか、少しだけ舞台裏の話をさせてください。

ことの発端は、私自身の苦い記憶にあります。多くの物理学徒と同じように、私もまた「ベクトル解析」の泥沼で溺れかけた一人でした。円柱座標や球座標の
\(\mathrm{rot}\) や \(\mathrm{div}\)
を前に、「なぜこんな形になるのか」という問いは「とにかく公式を暗記せよ」という圧力に封殺されました。微分形式がその特効薬になりうることは耳にしていましたが、数学書の公理的アプローチは当時の私には抽象度が高すぎました。

「この消化不良を一発で解決する、なにか小ずるい、『銀の弾丸』はないものか」------そう考えていた私に、第一のブレイクスルーが訪れました。個人的な話になりますが、当時の私は量子力学の勉強になかなかのめり込んでおり、Dirac
の「ブラ・ケット記法」の扱いに慣れ親しんでいました。状態ベクトル（縦ベクトル）を双対ベクトル（横ベクトル）として寝かせ、行列の積として自然に計算していく、あの感覚です。

ただし、私にとって重要だったのは、ブラ・ケット記法が「内積を計算する便利な記号」だったことではありません。むしろ逆です。ブラは本来、ケットを食べる双対ベクトルです。横ベクトルが縦ベクトルを食べるという双対ペアリングと、計量によって二つのベクトルを比べる内積は、本来は別の操作です。ところが物理数学の記法では、この二つがしばしば同じような見た目で現れ、いつのまにか区別が潰れてしまう。Dirac
自身も、そのあまりにも有名な著作で、``set up a second set of vectors,
which mathematicians call the \textbf{dual} vectors''
と述べています。にもかかわらず、教育の現場では、ブラはしばしば「内積記号の左半分」のように扱われてしまう。私はここに、当時は憤りを覚えていたほどです。

そんな折、ふと、この感覚を古典的な3次元ユークリッド空間の \(dx,dy,dz\)
にそのまま応用してみたらどうだろうか、と考えました。そう考えたとき、本書の最初の核が定まりました。積分記号の尻尾の
\(dx\)
を「変位ベクトルから成分を抽出する横ベクトル（行列）」として定義し直せば、「ベクトルを食べる関数（1-form）」という抽象概念を、誰もが知る「行列の掛け算」へと瞬時に変換できるのです。少なくとも、学生時代の筆者にとっては、これがベクトル解析を理解するための決定的な足場になりました。

これは内積ではありません。変位ベクトルを食べる測定器です。計量による同一視は、そのあとに来るべきものです。これは私個人の理解にとっては、まさに『銀の弾丸』でした。ベクトル解析の公式が、行列の足し引きや掛け算のルールから、すべて手を動かして導き出せるようになったのです。

しかし、これだけでは本書を書き始める動機にはまだ足りませんでした。そこで第二のブレイクスルーをもたらしてくれたのが、W.
L. Burke が提唱した「計量導入の遅延（put the metric later and later into
the course）」という教育的哲学でした。

Burke
の意図は、おそらく、長さ・角度・内積を入れる前から成立する構造を、計量に依存する構造から切り離して見せることにあったのだと思います。外微分
\(d\)、ウェッジ積、境界、ストークス型の構造は、計量を入れる前からかなりのところまで動きます。一般相対論やベクトル解析では、最初から計量を前面に出すと、外微分、境界、保存則、座標変換、そしてホッジ・スターによる計量依存の変換が一つに混ざって見えます。Burke
はそこに、かなり早い段階で問題意識を抱いていたのだろうと思います。

もちろん、これは筆者なりの受け取り方です。Burke
自身の問題意識が、私の「双対と内積を分けたい」という要求と同じだったとは思いません。それでも、計量を急いで導入すると、本来別々に見えるはずの構造が一つに潰れてしまう、という点で、私は
Burke の問題意識に強く共感しました。

Dirac から受け取ったのは、「双対を内積に潰すな」という感覚でした。Burke
から受け取ったのは、「計量を急いで入れると、計量なしで見える構造まで見えなくなる」という感覚でした。本書は、この二つの感覚が私の中で結びつき、\(dx\)
を測定器として置き、引き戻しを計量から分離し、その後で計量とホッジ・スターを導入する、という形に落ち着いたものです。

私にとって、とりわけ重要だったのは、まず座標変換と引き戻しを、長さや角度の話から切り離して扱えると示すことでした。座標変換こそ、ベクトル解析の鬼門です。そこに単位ベクトル、スケール因子、面積要素、計量を一度に流し込むから、何が何に由来するのかが見えなくなる。本書では、まず引き戻しをヤコビアンの行列計算として処理し、その後で計量とホッジ・スターを導入する、という順序を選びました。そして私は、この順序そのものを、読者を巻き込む「目次構成の教育的エンターテインメント」にしてやろう、と決意しました。

第1章から第5章までは、計量という概念を形式的に定義せず、実空間
\((x,y,z)\)
における素朴な長さ・面積の直観に基づいた代数で押し通しました。より進んだ立場から見れば「面積には計量が必要だ」と言いたくなる場面でしょうが、本書の序盤では
\(xyz\)
のデカルト座標に固定し、ピタゴラスの定理が成り立つ範囲で、素朴な面積・体積の直観をあえて使いました。計量
\(g = J^T J\)
が真に必要になるのは、パラメータ空間のような目盛りの歪んだ座標系へ一般化するときである------その事実を、第6章で行列計算を通じて発見的に導き出しました。

そして迎えた第6章の冒頭。「筆者もそろそろ飽きてきた」と冗談めかして、内積を正式に定義しました。同時に、計量
\(g = J^T J\) の意味を明らかにし、ホッジ・スター \(\ast\)
を「微分形式とベクトル解析を繫ぐ型変換アダプタ」として位置づけました。この発見的構成を通じて、「\(df\)
は計量に依存しないが、通常の勾配ベクトル・回転・発散として読むには計量やホッジ・スターが必要である」という構造を、読者に納得してもらうための仕掛けでした。

これらのアイデアのうち、Burke の「計量遅延」哲学や Flanders
の微分形式の物理応用は、いずれも既存の英語文献に書かれています。本書の内容に数学的な新規性はありません。ただ、\(dx\)
を行列として定義し全編を通じて行列計算で押し切るこの構成は、筆者の知る限り洋書にも見当たりません。東アジアの学生が記号計算に妙に訓練されているだけなのだろうか------などと茶化しつつ、いずれにせよ筆者が自分自身のために書いたものを整理したに過ぎません。本書が立っている巨人たちの肩については、巻末の参考文献を見ていただければと思います。

本書は、ベクトル解析に挫折したかつての私自身に向けて書いた「サバイバルガイド」です。もはやベクトル解析を飛び越え、「双対」「外微分」「引き戻し」「計量」「ホッジ・スター
\(\ast\)」といった強力なツールと、デカルト座標の限界を知ったあなたなら、すでに次の世界への扉を開いています。

ここから先は蛇足です。けれども、本書を閉じる前に、どうしても表明しておきたい筆者の信念があります。

本書は、微分形式という一見すると抽象的な道具を、できるかぎり行列計算へ落とし込むことで進んできました。\(dx\)
を横ベクトルとして読むことも、ウェッジ積を反対称な行列の計算として扱うことも、ホッジ・スターを計量行列から作る辞書として読むことも、すべてそのための仕掛けでした。

振り返ってみると、これは単なる筆者個人の便法ではなく、少なくとも入試数学をかなり本気で通過してきた理系学生がすでに身につけている代数的訓練を、大学数学の入口で「悪用」する試みでもあったのかもしれません。日本の難関入試は、良くも悪くも、式を変形し、文字を置き換え、パラメータを追い、条件を整理する能力を学生に叩き込みます。微分積分学ですら例外ではありません。典型的には、置換積分や媒介変数表示の処理は、ほとんど代数的な記号操作として訓練されます。

そのような訓練はしばしば過剰で、無意味な技巧にも見えるでしょう。けれども、これを無意味だ、抽象論への妨げだ、と嘆くのではなく、身体化された強力な武器として活用する余地があるのです。

その表れとして、本書は読者にかなり多くの手計算を求めます。しかもその多くは、極限を評価したり、関数の収束を調べたりする解析的な計算というより、記号を並べ替え、成分を追い、行列を掛け、外積を展開する代数的な計算です。これは、読者を低く見積もっていないということでもあります。読者はこの程度の記号操作には耐えられるはずだ、と筆者は期待しています。あるいは、少し買いかぶっているのかもしれません。

大学数学は、あまりにも早く集合論的「定義」へ飛ぶことがあります。もちろん、抽象的な定義には強さがあります。ブルバキ的な構造の整理には、今でも学ぶべきものがあります。しかし、すでに手に入れている代数的な身体を使って、抽象概念を一度、具体的な操作として掴む道もあるはずです。本書は、その道をかなり極端な形で試したものでした。\(dx=(1\ 0\ 0)\)
という断言は、その象徴です。

本書の構成は、決して唯一の正解ではありません。おそらく不格好で、ところどころ強引で、標準的な教科書から見れば遠回りに見えるでしょう。それでも、もし読者がこの本を通じて、抽象的な記号を自分の手で動かせるものとして感じられたなら、本書の役目は果たされたことになります。

『ナブラ解体新書』が、あなたの物理学と数学の旅において、すでに持っている代数的な手を捨てず、より高く、より遠くへ進むための強靭な足場となることを、著者として心から願っています。

\newpage

\hypertarget{ux53c2ux8003ux6587ux732eux3068ux8457ux8005ux304bux3089ux306eux30b3ux30e1ux30f3ux30c8}{%
\chapter{参考文献（と著者からのコメント）}\label{ux53c2ux8003ux6587ux732eux3068ux8457ux8005ux304bux3089ux306eux30b3ux30e1ux30f3ux30c8}}

本書の特異な代数的アプローチは、以下の先導者たちの知見を組み合わせて構築されました。

\begin{center}\rule{0.5\linewidth}{0.5pt}\end{center}

\textbf{1. Daniel Fleisch, *A Student's Guide to Vectors and Tensors*, Cambridge University Press (2011)}

【コメント：伝統的アプローチの丁寧なガイド】伝統的なベクトル解析とテンソル解析の成分計算を極めて丁寧に扱う、定評ある入門書。本書が別ルート（微分形式）での回避を図った、クリストッフェル記号や添字操作のジャングルを正面から突破する力を身につけるには最適の一冊。

\begin{center}\rule{0.5\linewidth}{0.5pt}\end{center}

\textbf{2. David Bachman, *A Geometric Approach to Differential Forms*, Birkhäuser (2nd ed. 2011)}

【コメント：微分形式への幾何学的導入】微分形式を公理からではなく「影を測る」といった幾何学的な直観から説明する良書。本書の「
\(dx\)
は測定器である」という直観の多くは本書から得ていますが、本書では計算としての完結性を優先し、その直観を行列代数へと落とし込んで採用しました。

\begin{center}\rule{0.5\linewidth}{0.5pt}\end{center}

\textbf{3. Harley Flanders, *Differential Forms with Applications to the Physical Sciences*, Dover Publications (1989/1963)}

【コメント：物理数学への応用の古典】微分形式を物理学へ応用する際のバイブル的な名著。外微分
\(d\) やホッジ・スター \(\ast\)
の厳密な操作ルールが凝縮されています。公理的な導入から始まるため初学者の独学にはやや硬派ですが、本書を終えた読者ならその強力な計算を十分に使いこなせるはずです。

\begin{center}\rule{0.5\linewidth}{0.5pt}\end{center}

\textbf{4. William L. Burke, *Applied Differential Geometry*, Cambridge University Press (1985)}
\textbf{William L. Burke, *Div, Grad, Curl are Dead* (Unfinished Manuscript)}

【コメント：本書の構成上の支柱------「計量遅延」の哲学】本書の骨格を決定づけた最重要文献。Burke
は「計量をコースのどんどん後の方に遅らせる（put the metric later and
later into the
course）」という教育的哲学を提唱しました。計量なしでも外微分
\(d\)（トポロジー）で語れる部分があまりに多いことを明示したこの哲学は、本書の第5章までの構成を支えています。

\begin{center}\rule{0.5\linewidth}{0.5pt}\end{center}

\textbf{5. Leonard Susskind \& Art Friedman, *Quantum Mechanics: The Theoretical Minimum*, Basic Books (2014)}

【コメント：行列表現の直観】抽象的な
\(1\)-form（双対空間）を「横ベクトルとして寝かせて行列の積をとる」という物理学者の母語に翻訳する直観は、量子力学のブラ・ケット記法の操作感から強く影響を受けています。

\begin{center}\rule{0.5\linewidth}{0.5pt}\end{center}

\textbf{6. Chris Doran \& Anthony Lasenby, *Geometric Algebra for Physicists*, Cambridge University Press (2003)}

【コメント：本書の先にある一つの統合】現代物理学の幾何代数（クリフォード代数）への入口。本書では徹底的に分離した「内積」と「外積」を、最初から一つの非可換積として統合した風景が見られます。分離の苦しみを知った読者こそ、この統合がいかに強力な解放感をもたらすかを感じ取れるでしょう。

\newpage

\hypertarget{ux4ed8ux9332ux672cux66f8ux3067ux8a9eux3089ux306aux304bux3063ux305fux3082ux306e}{%
\chapter{付録：本書で語らなかったもの}\label{ux4ed8ux9332ux672cux66f8ux3067ux8a9eux3089ux306aux304bux3063ux305fux3082ux306e}}

本付録は、本書の記述に対して想定される流儀上の相違、理論的限界、または読者からの典型的な指摘を整理したものです。

本書は、一般多様体上の微分形式論を公理的に展開する本ではありません。3次元デカルト空間におけるベクトル解析を、行列表示の微分形式を通して再構成するための教育的な本です。したがって、本書で採用した記法や説明順序は、標準的な数学書の記述と一致しない箇所があります。それらは原則として意図的な選択であり、必要に応じて第11章で標準的理論への接続を示しました。

\begin{quote}
\textbf{注}（本音）
このセクションのタイトルは「本書で語らなかったもの」と穏やかなものにしてみましたが、それは建前です。本音を言えば、これは本書の射程を読まずに投げられる典型的な批判への対応集です。

ただし、そう書いてしまえば、本当に応援したい読者が怖がって逃げ出してしまうでしょう。だから、表向きは穏やかな名前をつけました。
\end{quote}

\begin{quote}
\textbf{注}（「私」は怒っている）
筆者は若いころ、本書と同じ方向を向いた一連の SNS
投稿を公開し、立場のある方も含む「数学に詳しい」人々から手ひどく批判されたことがあります。もちろん、当時の草稿には未熟な点が多かったと思います。あれから10年経っても未だにこんな調子の文章を書く筆者の、若かりし時代です。それはもう、「勢いのある」投稿だったに違いありません。だが、そのとき筆者が受け取ったものは、定義を読み、射程を確認し、どう直せばよいかを示す批評というより、専門的な立場から未完成な試みを一方的に裁く態度でした。

それは、SNSのエンゲージメント構造と、人間の群集心理が生み出す、ある種の必然でもあったのでしょう。無料で文章を公開することには、想像以上の消耗が伴うと身をもって知りました。

その構造を理解したうえで、なお筆者はいまでも怒っています。厳密性を掲げながら、相手が何を定義し、どの射程で語っているかを読まない態度は、少なくとも筆者の考える数学的態度ではありません。

それでも、本書に限っては無料公開することに決めました。ある意味では、私的な代償行為でもあります。その上で、本来の読者を守るために------いえ、ごめんなさい、これも少し建前です。私自身の自尊心を守るためにも、精一杯の予防線を張ることにしました。

本書の定義が数学者向けの教科書と異なるように見える箇所があるなら、まず本書内の定義を読んでほしいと思います。そのうえで、本当に矛盾しているのか、単に異なる文脈の記号を持ち込んでいるだけなのかを峻別してほしいのです。その後には遠慮なく、私を今一度、手ひどく批判してほしい。定義を読んだうえで、筆者の間違いを指摘してほしいのです。
\end{quote}

\begin{quote}
\textbf{注}（高級な「きはじ」として）
数学書を読み進めるうちに、筆者は抽象的な定義順序の美しさも、その強さも、多少は内面化してしまいました。その結果、本書でやっていることが、ふと高級な「きはじ」にすぎないように見えることがあります。

これについて言葉を尽くせば言い訳に見えるし、黙れば誤解される。何より、自分自身がこのような批判を半ば信じているので、うまく語れません。

それでもなお、便法は必ずしも悪ではありません。便法は専門的な体系を置き換えませんし、それらに従属するだけのものでもありません。三次元の物理数学を、行列表示という手触りのある言葉で扱うための強力な言語にはなりえます。少なくとも具体的なシンボルに対する操作の体系にはなる。

改めて、本書の第一義はやはり、筆者自身の代償行為です。とはいえ、読者が手で動かせる言葉を一つ持って帰ってくれるなら、それは望外の副産物です。
\end{quote}

\begin{quote}
\textbf{注}（惑わぬ日は来るか）
三十にして立つなどといいますが、実際にその節目を迎えると、立つどころか未だに式と腹が立つばかりです。オヤジギャグは言うようになってしまいました。
\end{quote}

\begin{center}\rule{0.5\linewidth}{0.5pt}\end{center}

\hypertarget{i.-ux7406ux8ad6ux7684ux67a0ux7d44ux307fux3068ux53b3ux5bc6ux6027ux306bux3064ux3044ux3066}{%
\section{I.
理論的枠組みと厳密性について}\label{i.-ux7406ux8ad6ux7684ux67a0ux7d44ux307fux3068ux53b3ux5bc6ux6027ux306bux3064ux3044ux3066}}

\begin{itemize}
\item
  \textbf{多様体、アトラス、接空間などの数学的定義なしに外微分 $d$ を導入するのは不正確だ}
  →
  本書は一般多様体上の微分形式の教科書ではなく、3次元ユークリッド空間上でベクトル解析へ到達するための導入書です。一般多様体への移行については第11章で道標を示します。
\item
  \textbf{反変ベクトルと共変ベクトル（双対空間）の区別を序盤で明確にすべきだ}
  → 第1章 §1.1.6
  で「縦ベクトル（変位）」と「横ベクトル（測定器）」の区別を導入しており、これは双対空間の概念の具体的な表現です。より抽象的な双対空間の定式化については第6章および第11章で接続します。
\item
  \textbf{$dx$ を横ベクトルとして扱う記法は数学書の標準的な表記ではない}
  → 意図的な選択です。第1章 §1.1
  で定義し、全章で一貫して使用しています。この記法の狙いは「1-form =
  ベクトルを食べる関数」という抽象概念を、読者が既知の「行列の掛け算」に翻訳することにあります。標準的な微分形式の記法との対応は第11章で示します。
\item
  \textbf{外微分 $d$ は公理（ライプニッツ則や $d^2=0$）から一般的に定義すべきだ}
  → 本書では公理から始めるのではなく、第5章 §5.3
  で「微小ループを一周したときのズレ」という物理的直感から発見的に \(d\)
  を導出しています。\(d^2=0\) は第5章 §5.8
  で混合偏微分の対称性に由来する性質として理解します。公理的定義は第11章で補足します。
\item
  \textbf{一般の $n$ 次元多様体におけるストークスの定理の完全な証明が記載されていない}
  →
  本書では3次元ユークリッド空間での具体的なケース（ガウスの定理・ストークスの定理・グリーンの定理）を積分の定義から直接導出します（第5章、第8章）。一般次元のストークスの定理の証明は本書の射程外であり、第11章で道標のみ示します。
\item
  \textbf{ウェッジ積の定義において外積代数の普遍性質を用いた公理的アプローチを経由していない}
  → 本書では §2.4
  でテンソル積の反対称化としてウェッジ積を構築しています。これは公理的アプローチの最も具体的な実践であり、読者が行列の引き算として手を動かせるようにするための意図的な選択です。普遍性質による定義は第11章で紹介します。
\item
  \textbf{幾何学的な対象そのものに絶対的な座標表示はないのに成分表示に依存しすぎている}
  →
  意図的な選択です。本書の主題は「微分形式をどう使うか」であり、「座標に依存しない幾何学的主体としての微分形式」の哲学は目的としていません。本書は標準的なベクトル解析の教科書と同様に、実用的な計算手段として成分表示を採用しています。座標不変な視点は第11章で接続します。
\item
  \textbf{標準的な記法を避けているため、独自流派を主張しているように見える}
  →
  本書の記法は「初学者が行列計算で手を動かせる」ことを最優先に選ばれています。標準的な記法（添字による縮約、\(\partial_\mu\)
  など）を敢えて使わない理由は第11章で説明しています。この方針自体は数値計算の分野などでは広く見られるもので、著者が特に新しいことを主張しているつもりはない。意外と和書がなかったので、筆者が書いたまでです。
\end{itemize}

\begin{center}\rule{0.5\linewidth}{0.5pt}\end{center}

\hypertarget{ii.-ux624bux6cd5ux306eux9078ux629eux3068ux9069ux7528ux9650ux754cux306bux3064ux3044ux3066}{%
\section{II.
手法の選択と適用限界について}\label{ii.-ux624bux6cd5ux306eux9078ux629eux3068ux9069ux7528ux9650ux754cux306bux3064ux3044ux3066}}

\begin{itemize}
\item
  \textbf{アインシュタインの縮約記法（テンソル解析）を用いれば行列の束を使わずとも計算できる}
  →
  その通りです。本書では、行列の成分を毎回明示的に書き下す方針をとっています。これにより、読者は「ブラックボックス化された添字操作」に惑わず、各成分がどこから来て、どこへ作用するのかを追うことができます。むしろ、本書を読みこなした読者なら、上付き添字・下付き添字の規則とアインシュタイン縮約を練習するだけで、テンソル解析の成分計算へ比較的すぐ移行できるはずです。添字記法との対応は第11章で示します。
\item
  \textbf{最初からクリフォード代数（幾何代数）を教えれば外積も内積も統合できる}
  →
  本書はあえて「分離」の道を選びました。第1章〜第5章で外積（\(\wedge\)）を、第6章で内積（計量）を構成し、その後に
  第6章 §6.3 でホッジ・スター \(\ast\)
  を導入することで「なぜ統合が必要になるのか」を読者が実感できるようにしています。幾何代数による統合は第12章で紹介します。
\item
  \textbf{定数行列の計量 $\mathbf{g}$ を前提としているため、空間が曲がった場合に破綻するのではないかという懸念がある}
  → 第6章で \(xyz\) 実空間の計量 \(\mathbf{g}=I\)
  から出発し、パラメータ空間では \(\mathbf{g}=J^T J\)
  が場所の関数になることを示しています。曲がった空間（リーマン多様体）への一般化は第11章で道標を示します。本書の
  \(g=J^T J\)
  という表式は、ユークリッド空間内での座標変換や埋め込みから誘導される最も具体的な例です。
\item
  \textbf{第6章の正定値計量の導入から、第10章の擬リーマン計量への適用は論理的に飛躍している}
  → 第10章 §10.3
  の注で明記しているように、第10章では正定値性を放棄する代わりに「計量行列から
  \(\ast\)
  の辞書を作る」という手順だけを第6章から引き継いでいます。正定値から擬リーマンへの一般化の理論的根拠は第11章で補足します。
\item
  \textbf{捩れ（トーション）のある空間や非対称な接続を持つ多様体が考慮されていない}
  →
  本書の射程外です。本書が扱うのは3次元デカルト空間（およびパラメータ空間）であり、接続は常に対称・計量
  compatible な Levi-Civita
  接続を暗黙に前提としています。一般の接続の理論は第11章で道標を示します。
\item
  \textbf{スピンやスピノルといった対象を微分形式のみで扱うには代数的な限界がある}
  →
  その通りです。微分形式は反対称テンソルを扱う枠組みであり、スピノル（二価表現）を自然に扱うには、クリフォード代数やスピン構造など追加の道具が必要になります。本書ではこの方向には立ち入りませんが、第12章で幾何代数への道標を示します。
\item
  \textbf{行列の成分のみで計算を進めると、背景にある空間の幾何学的性質（計量等）が隠蔽される}
  →
  本書では幾何を隠蔽しているのではなく、アトミックな行列操作と幾何学的直感を常に対応させて与えています。たとえば第6章
  §6.1 では、パラメータ空間の計量 \(\mathbf{g}=J^T J\)
  の各成分が「軸方向の目盛りの伸び縮み」を直接表現していることを、行列の対角成分を見ながら確認できます。
\item
  \textbf{微分形式を行列表示した際、成分がどの基底に関するものかが判別できない}
  →
  記法上の限界であることは認めます。熱力学の偏微分で固定する変数を下付きにするように、座標系や基底を記号に埋め込む独自記法も考えましたが、一般記法との乖離をこれ以上広げないため採用しませんでした。本書では第1章
  §1.1.6 と §1.4
  以降、「どの座標系の基底で書いているか」を毎回本文中で明示する方針をとっています。標準的な
  \(dx^i\) 記法との対応は第11章を参照してください。
\item
  \textbf{外微分 $d$ と共変微分 $\nabla$ は別物であり、本書の $d$ の説明だけでは曲がった空間での微分が扱えないのではないか}
  → その通りであり、意図的な分離です。本書の \(d\)
  は反対称微分（外微分）であり、計量に依存しないトポロジカルな操作です。一方、共変微分
  \(\nabla\)
  は計量に依存する対称微分であり、曲がった空間で必要になります。この分離（\(d\)
  と \(\nabla\)）は Burke
  の「計量遅延」から強い影響を受けた筆者の解釈であり、第6章と第11章で扱います。
\item
  \textbf{曲線座標（円柱座標・極座標）において、本書の $dr$ や $d\theta$ がデカルト座標の $dx$ と同じ「横ベクトル」として扱えるのか不明瞭}
  → 第1章 §1.4 と第4章で明示的に扱っています。パラメータ空間の
  \(dr,d\theta,dz\) もデカルト座標の \(dx,dy,dz\)
  と同様に、パラメータ空間内での横ベクトル（一次形式）である。両者の違いは「引き戻し」によって吸収されます。
\end{itemize}

\begin{center}\rule{0.5\linewidth}{0.5pt}\end{center}

\hypertarget{iii.-ux7269ux7406ux7684ux5b9aux7fa9ux3068ux6559ux80b2ux7684ux914dux616eux306bux3064ux3044ux3066}{%
\section{III.
物理的定義と教育的配慮について}\label{iii.-ux7269ux7406ux7684ux5b9aux7fa9ux3068ux6559ux80b2ux7684ux914dux616eux306bux3064ux3044ux3066}}

\begin{itemize}
\item
  \textbf{直感的な図解（右ねじの法則等）がなく、代数計算ばかりで物理的イメージが湧かない}
  →
  本書は「図解」ではなく「行列計算」と「矢印ベクトル」によって物理的イメージを形成することを目的としています。たとえば第2章
  §2.4
  では、面積測定器の内部構造を行列として視覚化することで、ウェッジ積の交代性が「引き算」として現れることを示しています。これが本書における「イメージ」です。
\item
  \textbf{「ベクトル解析は間違いだ」とするのは物理学の歴史に対する誤解だ}
  →
  本書は「ベクトル解析は間違いだ」とは一言も述べていません。第8章では、ベクトル解析の定理（ガウス・ストークス・グリーン）を微分形式の積分定理の特殊ケースとして再導出しており、両者が等価であること（翻訳可能であること）を示しています。第12章では、この翻訳を代数的に完全自動化します。
\item
  \textbf{内容は既存の数学書に書かれており、著者の独自研究ではない} →
  その通りです。この流儀自体は Burke や Flanders
  の先行例があり、数値計算の分野でも見られるアプローチです。意外と和書がなかったため筆者が書いたに過ぎず、新規性を主張する気はありません。本書が立っている巨人たちの肩については「参考文献」を参照してください。
\item
  \textbf{ホッジ・スター $\ast$ は計量と空間の向きに依存するのに、単なる「次数反転」として扱いすぎではないか}
  → 第6章 §6.3 で \(\ast\) の定義は計量 \(\mathbf{g}\)
  に依存することを明示しています。直交デカルト座標（\(\mathbf{g}=I\)）では
  \(\ast(dx)=dy\wedge dz\) という簡潔な辞書になるが、直交曲線座標では
  \(\ast\) の係数が場所の関数になることを第8章 §8.6
  と第9章で具体的に示しています。
\item
  \textbf{マクスウェル方程式の微分形式表示において、単位系の正規化・符号規約・計量の符号が曖昧ではないか}
  → 第10章 §10.1 で \(w=ct\) と \(\mathbf{B}'=c\mathbf{B}\)
  による次元の正規化を明示し、§10.3
  でミンコフスキー計量の符号規約を定義し、§10.2 の注で \(F_{\mu\nu}\)
  の符号規約を明記し、§10.6 で \(F=-d\mathcal{A}\)
  の規約を宣言しています。§10.5 の注では \(\mathcal{J}\)
  の正規化が標準的な相対論的記法と係数が異なる場合があることも注記しています。
\item
  \textbf{$dx$ の行列表現 $\begin{pmatrix}1\&0\&0\end{pmatrix}$ は座標変換後に同じ行列として残るわけではない}
  → 第1章 §1.4 で明示的に扱っています。デカルト座標での
  \(dx = \begin{pmatrix}1&0&0\end{pmatrix}\) は、円柱座標の目盛りでは
  \(\begin{pmatrix}\cos\theta & -r\sin\theta & 0\end{pmatrix}\)
  に変わる。これが「引き戻し」の具体化であり、座標変換による行列の成分変化を隠さずに正面から見せています。
\end{itemize}

\begin{center}\rule{0.5\linewidth}{0.5pt}\end{center}

\begin{quote}
\textbf{注}（イースター・エッグ）
かつて私は、「ブルバキ」という名前に、聞きかじりの憧れを抱いていました。数学に大きな転換をもたらした存在だということです。ならばそこには、既存の数学書とは違う、何か鋭く新しい言葉があるのだろうと思っていました。

ところが実際に読んでみると、そこにあったのは、私にとってはむしろ見慣れた「大学数学の教科書」の顔でした。もちろん、それはブルバキがつまらなかったということではありません。むしろ逆です。かつて新しかった公理的で構造的な書き方は、ブルバキだけによってではないにせよ、現代数学の標準的な文体の一部として深く定着しています。だからこそ、いまの読者にはそれが「普通」に見えるのでしょう。

そして本書は、その「普通」の顔の前で迷子になる読者のために書かれています。
\end{quote}
